\documentclass[10pt]{article}
\usepackage[a4paper,margin=0.75in]{geometry}
\usepackage[T1]{fontenc}
\usepackage[utf8]{inputenc}
\usepackage{fvextra}
\usepackage{longtable}
\usepackage{booktabs}
\usepackage{xcolor}
\usepackage{hyperref}
\usepackage{microtype}
\setlength{\parindent}{0pt}
\setlength{\parskip}{4pt}
\DefineVerbatimEnvironment{AnswerText}{Verbatim}{fontsize=\scriptsize,breaklines=true,breakanywhere=true,frame=single,framesep=3pt}
\DefineVerbatimEnvironment{EvidenceText}{Verbatim}{fontsize=\tiny,breaklines=true,breakanywhere=true,frame=single,framesep=2pt}
\begin{document}
\begin{center}{\Large\bfseries External-baseline answer and evidence audit}\\[4pt]{\normalsize Unified batch: 8 arms $\times$ 10 questions = 80 answers}\end{center}
\textbf{Purpose.} This companion appendix preserves the complete answer text and retrieval provenance for the external-baseline batch. It is an audit transcript, not an additional statistical analysis. The raw JSONL files in \path{experiment_data/external_baselines_v3/} are the exact machine-readable source; Unicode source paths are rendered with ASCII escapes in this PDF where necessary. Model identifiers are shown without their provider prefix.

\textbf{Judging.} Each entry reports the four dimension scores, the verdict, the overall arithmetic mean and the judge's rationale. For each question the judge saw the eight anonymised answers of this batch together with a head-only excerpt of the gold paper, never the retrieved context or a system name. The scores are reported in Table 5 of the main text and Supplementary Tables S19--S20; the earlier passes, corpus scopes and pipeline limitations are reported in Supplementary Tables S19--S21.

\tableofcontents
\clearpage
\section*{Protocol texts}
\addcontentsline{toc}{section}{Protocol texts}
\subsection*{Answer prompt of the retrieval-mode arms}
The six retrieval-mode arms (BM25, dense, RRF, RAPTOR, LightRAG and GraphRAG) received one prompt per question, built from the template below; angle brackets mark the text substituted per question. The HMRS agent and the native DeepRead pipeline ran their own multi-turn procedures and did not use this prompt.
\begin{EvidenceText}
You answer questions about research papers using ONLY the context below. Each source is marked [S#] with its document name. Cite the sources you use inline like [S2]. If the context does not contain the answer, say so explicitly instead of guessing.

Question: <question>

Context:
[S1] <document name of unit 1>
<text of unit 1>

[S2] <document name of unit 2>
<text of unit 2>

<units 3 to 12 follow in the same form; each unit holds at most 800 characters>

Answer:
\end{EvidenceText}
The prompt below reconstructs the dense answer to BQ01 from its 12 recorded context units (\path{sources} in \path{answers-v3.jsonl}).
\begin{EvidenceText}
You answer questions about research papers using ONLY the context below. Each source is marked [S#] with its document name. Cite the sources you use inline like [S2]. If the context does not contain the answer, say so explicitly instead of guessing.

Question: How does the Transformer compute attention, and why does it replace recurrence and convolution?

Context:
[S1] \u8ba1\u7b97\u673a\u4e0e\u4eba\u5de5\u667a\u80fd/artificial-intelligence__1706.03762__attention-is-all-you-need.md (part 2 of 2).md
## 7 Conclusion In this work, we presented the Transformer, the first sequence transduction model based entirely on attention, replacing the recurrent layers most commonly used in encoder-decoder architectures with multi-headed self-attention. For translation tasks, the Transformer can be trained significantly faster than architectures based on recurrent or convolutional layers. On both WMT 2014 English-to-German and WMT 2014 English-to-French translation tasks, we achieve a new state of the

[S2] \u8ba1\u7b97\u673a\u4e0e\u4eba\u5de5\u667a\u80fd/artificial-intelligence__1706.03762__attention-is-all-you-need.md (part 1 of 2).md
odeling of dependencies without regard to their distance in the input or output sequences [2, 19]. In all but a few cases [27], however, such attention mechanisms are used in conjunction with a recurrent network. In this work we propose the Transformer, a model architecture eschewing recurrence and instead relying entirely on an attention mechanism to draw global dependencies between input and output. The Transformer allows for significantly more parallelization and can reach a new state of the

[S3] \u8ba1\u7b97\u673a\u4e0e\u4eba\u5de5\u667a\u80fd/artificial-intelligence__1706.03762__attention-is-all-you-need.md (part 1 of 2).md
## Abstract The dominant sequence transduction models are based on complex recurrent or convolutional neural networks that include an encoder and a decoder. The best performing models also connect the encoder and decoder through an attention mechanism. We propose a new simple network architecture, the Transformer, based solely on attention mechanisms, dispensing with recurrence and convolutions entirely. Experiments on two machine translation tasks show these models to be superior in quality w

[S4] \u8ba1\u7b97\u673a\u4e0e\u4eba\u5de5\u667a\u80fd/artificial-intelligence__1706.03762__attention-is-all-you-need.md (part 1 of 2).md
## 3.2.3 Applications of Attention in our Model The Transformer uses multi-head attention in three different ways: \u2022 In "encoder-decoder attention" layers, the queries come from the previous decoder layer, and the memory keys and values come from the output of the encoder. This allows every position in the decoder to attend over all positions in the input sequence. This mimics the typical encoder-decoder attention mechanisms in sequence-to-sequence models such as [38, 2, 9]. \u2022 The encoder co

[S5] \u8ba1\u7b97\u673a\u4e0e\u4eba\u5de5\u667a\u80fd/artificial-intelligence__1706.03762__attention-is-all-you-need.md (part 1 of 2).md
and have been shown to perform well on simple-language question answering and language modeling tasks [34]. To the best of our knowledge, however, the Transformer is the first transduction model relying entirely on self-attention to compute representations of its input and output without using sequencealigned RNNs or convolution. In the following sections, we will describe the Transformer, motivate self-attention and discuss its advantages over models such as [17, 18] and [9].

[S6] \u8ba1\u7b97\u673a\u4e0e\u4eba\u5de5\u667a\u80fd/artificial-intelligence__1706.03762__attention-is-all-you-need.md (part 1 of 2).md
sitions, linearly for ConvS2S and logarithmically for ByteNet. This makes it more difficult to learn dependencies between distant positions [12]. In the Transformer this is reduced to a constant number of operations, albeit at the cost of reduced effective resolution due to averaging attention-weighted positions, an effect we counteract with Multi-Head Attention as described in section 3.2. Self-attention, sometimes called intra-attention is an attention mechanism relating different positions o

[S7] \u8ba1\u7b97\u673a\u4e0e\u4eba\u5de5\u667a\u80fd/nlp__2110.07844__modeling-endorsement-for-multi-document-abst (part 1 of 3).md
rm { F E E D F O R W A R D } \big ( \widehat { G } ^ { ( l ) } \big )\tag{5} $$ The original cross-attention head z transforms the j-th decoder state $\widetilde { \pmb { g } } _ { j } ^ { ( l - 1 ) }$ and i-th encoder state ${ h _ { i } ^ { ( L ) } }$ into query, key and value vectors (Eqs. (6-8)). It computes attention weights as a normalized dot product between query and key vectors. The output of the head is a weighted sum of value vectors.

[S8] \u8ba1\u7b97\u673a\u4e0e\u4eba\u5de5\u667a\u80fd/artificial-intelligence__1706.03762__attention-is-all-you-need.md (part 1 of 2).md
arable convolutions [6], however, decrease the complexity considerably, to $\overset { \cdot } { O ( k \cdot n \cdot d + n \cdot d ^ { 2 } ) }$ . Even with $k = n$ , however, the complexity of a separable convolution is equal to the combination of a self-attention layer and a point-wise feed-forward layer, the approach we take in our model. As side benefit, self-attention could yield more interpretable models. We inspect attention distributions from our models and present and discuss examples i

[S9] \u8ba1\u7b97\u673a\u4e0e\u4eba\u5de5\u667a\u80fd/nlp__2503.21676__how-do-language-models-learn-facts-dynamics.md (part 2 of 4).md
he proposed credit-assignment argument for plateau formation generalizes to other architectures, as evidenced by ablations in which attention is replaced with recurrence (Figure C). Decoupling the token-mixing role of attention from other computations, as advocated by Elhage et al. (2021) for mechanistic understanding of trained Transformers, also proves powerful for analyzing learning dynamics. This perspective offers a promising avenue for future investigations into neural network learning dyn

[S10] \u8ba1\u7b97\u673a\u4e0e\u4eba\u5de5\u667a\u80fd/artificial-intelligence__2510.13343__aoad-mat-transformer-based-multi-agent-deep (part 1 of 2).md
## 3.3 Multi-Agent Transformer Approach Transformer model [24], originally developed for NLP tasks, such as translation and text generation, has demonstrated remarkable versatility across various domains. Its effectiveness in sequence modeling has led to its adoption in fields beyond NLP, including image and signal processing. The key strength of the Transformer lies in its attention mechanism, which learns to focus on relevant parts of the input data, thereby enabling powerful representationa

[S11] \u8ba1\u7b97\u673a\u4e0e\u4eba\u5de5\u667a\u80fd/artificial-intelligence__1706.03762__attention-is-all-you-need.md (part 2 of 2).md
translation tasks, we achieve a new state of the art. In the former task our best model outperforms even all previously reported ensembles. We are excited about the future of attention-based models and plan to apply them to other tasks. We plan to extend the Transformer to problems involving input and output modalities other than text and to investigate local, restricted attention mechanisms to efficiently handle large inputs and outputs such as images, audio and video. Making generation less

[S12] \u8ba1\u7b97\u673a\u4e0e\u4eba\u5de5\u667a\u80fd/artificial-intelligence__1706.03762__attention-is-all-you-need.md (part 1 of 2).md
nt at a time. At each step the model is auto-regressive [10], consuming the previously generated symbols as additional input when generating the next. ![](images/d4522aeb4367bfd769b08095db34e4b8da8b4e5333940f5df006c862d23551a9.jpg) Figure 1: The Transformer - model architecture. The Transformer follows this overall architecture using stacked self-attention and point-wise, fully connected layers for both the encoder and decoder, shown in the left and right halves of Figure 1, respectively.

Answer:
\end{EvidenceText}
\clearpage
\subsection*{Judge instruction of the unified batch}
Every judge packet began with the instruction below, identical for all ten questions, followed by the question, a head-only excerpt of the gold paper and the eight anonymised answers in a fresh random letter order (\path{mapping-unified.json}). The fourth dimension is called attribution in the manuscript.
\begin{EvidenceText}
You are a strict, neutral grader of question-answering systems.
You will see: a QUESTION, GOLD SOURCE MATERIAL (excerpts from the document the
question was written about), and answers from anonymous systems.

Score EACH answer on four dimensions, each an integer 1-5:
- factual: correctness of every factual claim against the GOLD material (5 = all
  claims correct and the key facts match; 1 = mostly wrong)
- completeness: does it answer all parts of the question (5 = fully; 1 = misses
  the point)
- groundedness: are claims supported by the sources it cites or the gold
  material (5 = no unsupported claims; 1 = substantial fabrication)
- sources: does it attribute statements to identifiable documents/passages
  (5 = clear, accurate attributions; 1 = none or bogus)

Also give verdict: "correct", "partial", or "incorrect" for the question as a
whole, and one short rationale sentence.

Rules:
- Grade each answer INDEPENDENTLY against the rubric anchors; do not curve
  scores relative to the other answers in this packet.
- You MUST ground each dimension score in the actual text: "factual" must
  reference specific facts from the GOLD material; "groundedness" must state
  which claims are or are not supported; "sources" must describe what
  attributions the answer actually makes. Quote or paraphrase the answer.
- An honest "the material does not contain this" answer is not fabrication:
  its groundedness stays high, while completeness/factual follow the rubric.

For every answer, output one JSON object:
{"system": "<letter>", "factual": <1-5>, "completeness": <1-5>,
 "groundedness": <1-5>, "sources": <1-5>,
 "dim_rationale": {"factual": "<evidence-based justification>",
                   "completeness": "<...>", "groundedness": "<...>",
                   "sources": "<...>"},
 "verdict": "<correct|partial|incorrect>", "rationale": "<one sentence>"}
\end{EvidenceText}
\section{HMRS agent (pinned)}
\subsection{BQ01: How does the Transformer compute attention, and why does it replace recurrence and convolution?}
\begin{tabular}{@{}ll@{}}
Arm & HMRS agent (pinned)\\
Generator & omp / deepseek-flash\\
HTTP / success & 200 / True\\
Wall time / cost & 97.6 s / 0.0566778 USD\\
Judge scores & F=5, C=5, G=5, A=5; overall=5.00\\
Verdict & correct\\
\end{tabular}

\textbf{Dimension rationale}
\begin{EvidenceText}
{
  "factual": "The direct quote 'We compute the dot products of the query with all keys, divide each by sqrt(d_k), and apply a softmax function...' is verbatim from the paper; Equation 1, multi-head formulas, h=8 with d_k=d_v=64, the exact Table 1 rows (self-attention O(n^2*d)/O(1)/O(1); recurrent O(n*d^2)/O(n)/O(n); convolutional O(k*n*d^2)/O(1)/O(log_k n); restricted O(r*n*d)/O(n/r)), base 27.3 / big 28.4 / EN-FR 41.8 BLEU, and FLOPs 3.3e18 all match the real paper; the ConvS2S-linear/ByteNet-log contrast and 'precludes parallelization' quotes match the GOLD text exactly.",
  "completeness": "Both parts answered to the maximum: computation (definition, procedure, matrix form, scaling rationale, multi-head, decoder mask), motivation (three criteria, Table 1, arguments against recurrence and against convolution, restricted self-attention, positional encoding necessity), plus effect corroboration (Table 2 FLOPs and BLEU).",
  "groundedness": "Labeled 'P0 direct evidence' with verbatim quotes that check out against the paper's text; it also transparently discloses 'ModuleNotFoundError: No module named laya' and says the quotes came from reading the KB document directly \u2014 nothing unsupported detected.",
  "sources": "Names the exact KB file path ('\u8ba1\u7b97\u673a\u4e0e\u4eba\u5de5\u667a\u80fd/artificial-intelligence__1706.03762__attention-is-all-you-need.md (part 1 of 2).md') and a KB id \u2014 file-level attribution, the most concrete in the packet."
}
\end{EvidenceText}

\textbf{Judge summary}
\begin{EvidenceText}
Most complete answer: full formulas, Table 1/Table 2 numbers and verbatim quotes, all verifiable against the paper and gold, with file-level sourcing.
\end{EvidenceText}

\textbf{Complete answer}
\begin{AnswerText}
**\u7ed3\u8bba\uff1aTransformer \u7684\u6ce8\u610f\u529b\uff0c\u662f\u628a Q \u4e0e\u6240\u6709 K \u505a\u70b9\u79ef\u3001\u9664\u4ee5 $\sqrt{d_k}$\u3001softmax \u5f97\u6743\u91cd\u540e\u5bf9 V \u52a0\u6743\u6c42\u548c\uff1b\u5b83\u53d6\u4ee3\u5faa\u73af\u548c\u5377\u79ef\uff0c\u662f\u56e0\u4e3a\u4e00\u6b21\u77e9\u9635\u4e58\u5c31\u8ba9\u4efb\u610f\u4e24\u4f4d\u7f6e\u76f4\u63a5\u76f8\u8fde\u2014\u2014\u5c42\u5185\u590d\u6742\u5ea6 $O(n^2 d)$\u3001\u5e76\u884c\u6b65\u6570 $O(1)$\u3001\u6700\u5927\u8def\u5f84\u957f\u5ea6 $O(1)$\uff0c\u800c RNN \u9700 $O(n)$ \u4e32\u884c\u6b65\u548c $O(n)$ \u8def\u5f84\uff0c\u5377\u79ef\u9700\u5806\u53e0\u591a\u5c42\u624d\u80fd\u8fde\u901a\u5168\u5e8f\u5217\u3002**

## 1. \u6ce8\u610f\u529b\u5982\u4f55\u8ba1\u7b97\uff08\xa73.2\u20133.2.2\uff0cP0 \u76f4\u63a5\u8bc1\u636e\uff09

\u6ce8\u610f\u529b\u7684\u5b9a\u4e49\u662f"\u628a query \u4e0e\u4e00\u7ec4 key-value \u5bf9\u6620\u5c04\u4e3a\u8f93\u51fa"\uff0c\u8f93\u51fa\u4e3a value \u7684\u52a0\u6743\u548c\uff0c\u6743\u91cd\u6765\u81ea query \u4e0e\u5bf9\u5e94 key \u7684\u517c\u5bb9\u51fd\u6570\u3002\u5177\u4f53\u7528"\u7f29\u653e\u70b9\u79ef\u6ce8\u610f\u529b"\uff08Scaled Dot-Product Attention\uff09\uff1a

- \u8f93\u5165\uff1aquery\u3001key \u7ef4\u5ea6 $d_k$\uff0cvalue \u7ef4\u5ea6 $d_v$\uff1b
- "We compute the dot products of the query with all keys, divide each by $\sqrt{d_k}$, and apply a softmax function to obtain the weights on the values."
- \u77e9\u9635\u5f62\u5f0f\uff08\u5f0f 1\uff09\uff1a$\mathrm{Attention}(Q,K,V)=\mathrm{softmax}\!\left(\frac{QK^T}{\sqrt{d_k}}\right)V$
- \u7f29\u653e\u7684\u7406\u7531\uff1a$d_k$ \u5927\u65f6\u70b9\u79ef\u91cf\u7ea7\u53d8\u5927\uff0c"pushing the softmax function into regions where it has extremely small gradients"\uff0c\u6545\u4e58 $1/\sqrt{d_k}$\uff1b\u76f8\u5bf9\u52a0\u6027\u6ce8\u610f\u529b\uff0c"dot-product attention is much faster and more space-efficient in practice, since it can be implemented using highly optimized matrix multiplication code"\u3002

\u591a\u5934\u6ce8\u610f\u529b\uff08\xa73.2.2\uff09\uff1a\u628a Q/K/V \u5404\u505a $h$ \u6b21\u4e0d\u540c\u53ef\u5b66\u4e60\u7ebf\u6027\u6295\u5f71\u5230 $d_k,d_k,d_v$\uff0c\u5e76\u884c\u7b97\u6ce8\u610f\u529b\uff0c\u62fc\u63a5\u540e\u518d\u6295\u5f71 $W^O$\uff1a$\mathrm{MultiHead}=\mathrm{Concat}(\mathrm{head}_1,\dots,\mathrm{head}_h)W^O$\uff0c$\mathrm{head}_i=\mathrm{Attention}(QW_i^Q,KW_i^K,VW_i^V)$\u3002\u8bba\u6587\u53d6 $h=8$\uff0c$d_k=d_v=d_{model}/h=64$\uff0c\u603b\u8ba1\u7b97\u6210\u672c\u4e0e\u5168\u7ef4\u5355\u5934\u76f8\u5f53\u3002\u5176\u5fc5\u8981\u6027\uff1a"Multi-head attention allows the model to jointly attend to information from different representation subspaces at different positions. With a single attention head, averaging inhibits this." \u89e3\u7801\u5668\u81ea\u6ce8\u610f\u529b\u7528 mask\uff08softmax \u8f93\u5165\u7f6e $-\infty$\uff09\u5c4f\u853d\u672a\u6765\u4f4d\u7f6e\u4ee5\u4fdd\u6301\u81ea\u56de\u5f52\u6027\u3002

## 2. \u4e3a\u4f55\u53d6\u4ee3\u5faa\u73af\u4e0e\u5377\u79ef\uff08\xa74 "Why Self-Attention"\uff0cP0\uff09

\u8bba\u6587\u7528\u4e09\u6761\u6807\u51c6\u6bd4\u8f83\uff1a\u6bcf\u5c42\u8ba1\u7b97\u590d\u6742\u5ea6\u3001\u6700\u5c0f\u4e32\u884c\u64cd\u4f5c\u6570\u3001\u4efb\u610f\u4e24\u4f4d\u7f6e\u95f4\u6700\u5927\u8def\u5f84\u957f\u5ea6\uff08"The shorter these paths\u2026, the easier it is to learn long-range dependencies"\uff09\u3002Table 1\uff1a

| Layer Type | Complexity/Layer | Sequential Ops | Max Path Length |
|---|---|---|---|
| Self-Attention | $O(n^2\cdot d)$ | $O(1)$ | $O(1)$ |
| Recurrent | $O(n\cdot d^2)$ | $O(n)$ | $O(n)$ |
| Convolutional | $O(k\cdot n\cdot d^2)$ | $O(1)$ | $O(\log_k(n))$ |
| Self-Attention (restricted) | $O(r\cdot n\cdot d)$ | $O(1)$ | $O(n/r)$ |

- **\u5bf9\u5faa\u73af**\uff1aRNN \u6cbf\u7b26\u53f7\u4f4d\u7f6e\u5206\u89e3\u8ba1\u7b97\uff0c"This inherently sequential nature precludes parallelization within training examples"\uff0c\u4e14 $h_t$ \u4f9d\u8d56 $h_{t-1}$\uff08\u5f15\u8a00\uff0cP0\uff09\u3002\u81ea\u6ce8\u610f\u529b"connects all positions with a constant number of sequentially executed operations, whereas a recurrent layer requires $O(n)$ sequential operations"\uff1b\u5f53 $n<d$\uff08\u8bcd\u7247/BPE \u8868\u793a\u7684\u5e38\u89c1\u60c5\u5f62\uff09\u65f6\u81ea\u6ce8\u610f\u529b\u66f4\u5feb\u3002
- **\u5bf9\u5377\u79ef**\uff1a\u5355\u4e2a $k<n$ \u7684\u5377\u79ef\u5c42\u4e0d\u8fde\u901a\u6240\u6709\u8f93\u5165\u8f93\u51fa\u5bf9\uff0c\u9700\u5806 $O(n/k)$ \u5c42\uff08\u8fde\u7eed\u6838\uff09\u6216 $O(\log_k n)$ \u5c42\uff08\u7a7a\u6d1e\u5377\u79ef\uff09\uff0c"increasing the length of the longest paths"\uff1b\u5377\u79ef\u5c42\u4e00\u822c\u6bd4\u5faa\u73af\u5c42\u8d35 $k$ \u500d\u3002\u5373\u4fbf $k=n$\uff0c\u53ef\u5206\u79bb\u5377\u79ef\u7684\u590d\u6742\u5ea6\u4e5f\u53ea\u7b49\u4e8e"\u81ea\u6ce8\u610f\u529b\u5c42 + \u9010\u70b9\u524d\u9988\u5c42"\u7684\u7ec4\u5408\uff0c\u5373\u672c\u6a21\u578b\u91c7\u7528\u7684\u7ed3\u6784\u3002
- **\u4ee3\u4ef7\u4e0e\u8865\u6551**\uff1a\u5168\u5c40\u6ce8\u610f\u529b\u5e26\u6765\u4f4d\u7f6e\u5e73\u5747\u5bfc\u81f4\u7684"reduced effective resolution"\uff0c\u7528\u591a\u5934\u6ce8\u610f\u529b\u62b5\u6d88\uff08\xa72\uff09\uff1b\u957f\u5e8f\u5217\u53ef\u7528\u90bb\u57df\u4e3a $r$ \u7684\u53d7\u9650\u81ea\u6ce8\u610f\u529b\u964d\u5230 $O(rnd)$\u3002\u56e0\u65e0\u5faa\u73af\u65e0\u5377\u79ef\uff0c\u5fc5\u987b\u663e\u5f0f\u6ce8\u5165\u4f4d\u7f6e\u4fe1\u606f\u2014\u2014\u6b63\u5f26\u4f4d\u7f6e\u7f16\u7801\uff08\xa73.5\uff09\u3002

## 3. \u6548\u679c\u4f50\u8bc1\uff08P1\uff09

\u57fa\u6a21\u578b 8\xd7P100 \u8bad\u7ec3 12 \u5c0f\u65f6\uff0810 \u4e07\u6b65\uff09\uff0cEN-DE BLEU 27.3\u3001big 28.4\uff0cEN-FR 41.8\uff0c\u8bad\u7ec3 FLOPs $3.3\cdot10^{18}$\uff0c\u8fdc\u4f4e\u4e8e GNMT/ConvS2S \u7684 $10^{19}$\u2013$10^{20}$ \u91cf\u7ea7\uff08Table 2\uff09\u3002

**\u51fa\u5904**\uff1a`\u8ba1\u7b97\u673a\u4e0e\u4eba\u5de5\u667a\u80fd/artificial-intelligence__1706.03762__attention-is-all-you-need.md (part 1 of 2).md` \u2014\u300aAttention Is All You Need\u300b(1706.03762)\uff0cKB `b1199132-\u2026`\uff08\u8ba1\u7b97\u673a\u4e0e\u4eba\u5de5\u667a\u80fd\u5e93\uff09\u3002\u6ce8\uff1aLaya \u8bc1\u636e\u95e8\u672c\u8f6e\u4e0d\u53ef\u7528\uff08`ModuleNotFoundError: No module named 'laya'`\uff09\uff0c\u4e0a\u6587\u5f15\u6587\u4e3a\u76f4\u63a5\u8bfb\u53d6\u8be5\u5e93\u539f\u59cb\u6587\u6863\u6240\u5f97\u3002
\end{AnswerText}

\textbf{Retrieved evidence index (all returned units; paths and scores) }
\begin{EvidenceText}
[No source list recorded for this arm.]
\end{EvidenceText}
\subsection{BQ02: What defines NISQ technology and what is its central limitation?}
\begin{tabular}{@{}ll@{}}
Arm & HMRS agent (pinned)\\
Generator & omp / deepseek-flash\\
HTTP / success & 200 / True\\
Wall time / cost & 49.9 s / 0.047616 USD\\
Judge scores & F=5, C=5, G=5, A=5; overall=5.00\\
Verdict & correct\\
\end{tabular}

\textbf{Dimension rationale}
\begin{EvidenceText}
{
  "factual": "The Chinese summary and English quotes all match the paper: coinage quote 'I needed a name... NISQ. This stands for Noisy Intermediate-Scale Quantum'; the three-element table (50 to a few hundred qubits; 50 beyond brute-force simulation; noisy = imperfect control; noisy gates unprotected by QEC); the abstract circuit-size sentence; the verbatim 1000-gate passage and ceiling sentence; two-qubit gate error 'above the .1% level (and often much worse)'; QEC 'heavy overhead cost'. The cited benchmarking-paper phrase 'restrictive qubit topologies, short coherence times and imperfect quantum gates' is genuine.",
  "completeness": "Both parts fully answered, with a conclusion line up front, a per-element definition table, the quantitative limitation chain (error rate -> 1000-gate cap -> compute ceiling), the QEC escape route, and a corroborating P1 source explicitly labeled as a different formulation of the same limitation.",
  "groundedness": "Self-labeled evidence tiers (P0 direct quotes verified against the raw document's first 80 lines; P1 partial corroboration) with the claim that the benchmarking paper 'does not conflict' with Preskill; all quotes check against the paper text.",
  "sources": "Most concrete sourcing in the packet: KB id, full file path ('quantum-physics__1801.00862...md (part 1 of 3).md'), retrieval commands used, and why Chinese-language hits were rejected as chunk mirrors."
}
\end{EvidenceText}

\textbf{Judge summary}
\begin{EvidenceText}
Complete, verbatim-quoted answer with tiered evidence labels and file-level attribution; definition and central limitation both fully and accurately covered.
\end{EvidenceText}

\textbf{Complete answer}
\begin{AnswerText}
**\u7ed3\u8bba\uff1aNISQ\uff08Noisy Intermediate-Scale Quantum\uff0c\u566a\u58f0\u4e2d\u7b49\u89c4\u6a21\u91cf\u5b50\uff09\u6280\u672f\u6307\u300c\u6bd4\u7279\u6570\u7ea6 50 \u5230\u51e0\u767e\u3001\u95e8\u9519\u8bef\u7387\u7ea6 0.1% \u4ee5\u4e0a\u3001\u4e14\u4e0d\u53d7\u91cf\u5b50\u7ea0\u9519\u4fdd\u62a4\u7684\u91cf\u5b50\u8ba1\u7b97\u673a\u300d\uff1b\u5176\u6838\u5fc3\u9650\u5236\u662f\u566a\u58f0\u8ba9\u53ef\u53ef\u9760\u6267\u884c\u7684\u7ebf\u8def\u89c4\u6a21\u88ab\u5361\u5728\u7ea6 1000 \u4e2a\u53cc\u6bd4\u7279\u95e8\u4ee5\u5185\u2014\u2014\u8fd9\u4e00\u7ebf\u8def\u89c4\u6a21\u4e0a\u9650\u76f4\u63a5\u7ed9 NISQ \u7684\u8ba1\u7b97\u80fd\u529b\u5c01\u4e86\u9876\u3002**

## \u4e00\u3001NISQ \u7684\u754c\u5b9a\uff08Preskill \u539f\u59cb\u5b9a\u4e49\uff09

John Preskill \u5728 *Quantum Computing in the NISQ Era and Beyond*\uff082018-07-30\uff09\u4e2d\u81ea\u8ff0\u9020\u8bcd\uff1a

> "I needed a name to describe this impending new era, so I made up a word: NISQ. This stands for Noisy Intermediate-Scale Quantum."

\u4e09\u4e2a\u6784\u6210\u8981\u7d20\u88ab\u9010\u5b57\u62c6\u89e3\uff08\u540c\u8282 4.1 "The 50-qubit barrier"\uff09\uff1a

| \u8981\u7d20 | \u539f\u6587\u754c\u5b9a |
|---|---|
| **Intermediate-scale** | "the size of quantum computers which will be available in the next few years, with **a number of qubits ranging from 50 to a few hundred**"\uff1b50 \u6bd4\u7279\u7684\u610f\u4e49\u662f"beyond what can be simulated by brute force using the most powerful existing digital supercomputers" |
| **Noisy** | "'Noisy' emphasizes that we'll have **imperfect control over those qubits**; the noise will place serious limitations on what quantum devices can achieve in the near term" |
| **\u65e0\u7ea0\u9519\u4fdd\u62a4**\uff08\u9690\u542b\u4f46\u660e\u786e\uff09 | "When I speak of the NISQ era, I'm imagining quantum computers with **noisy gates unprotected by quantum error correction**" |

\u6458\u8981\u7684\u6b63\u5f0f\u8868\u8ff0\uff1a
> "Quantum computers with 50-100 qubits may be able to perform tasks which surpass the capabilities of today's classical digital computers, but **noise in quantum gates will limit the size of quantum circuits that can be executed reliably**."

## \u4e8c\u3001\u6838\u5fc3\u9650\u5236\uff1a\u566a\u58f0\u5c01\u9876\u7ebf\u8def\u89c4\u6a21 \u21d2 \u5c01\u9876\u7b97\u529b

\u539f\u6587\uff08\xa74.2 Qubit "quality"\uff09\u7ed9\u51fa\u91cf\u5316\u94fe\u6761\uff1a

> "with these noisy devices we don't expect to be able to execute a circuit that contains **many more than about 1000 gates \u2014 that is, 1000 fundamental two-qubit operations** \u2014 because the noise will overwhelm the signal in a circuit much larger than that. **That limitation on circuit size imposes a ceiling on the computational power of NISQ technology.**"

\u5373\uff1a**\u9650\u5236\u7684\u672c\u8d28\u4e0d\u662f\u6bd4\u7279\u4e0d\u591f\u591a\uff0c\u800c\u662f\u300c\u566a\u58f0 \u221d \u95e8\u6570\u7d2f\u79ef\u300d\u4f7f\u4fe1\u53f7\u88ab\u6df9\u6ca1**\u3002\u652f\u6491\u6570\u636e\uff1a\u4e24\u6bd4\u7279\u95e8\u9519\u8bef\u7387 "above the .1% level (and often much worse)"\uff1b\u6d4b\u91cf\u9519\u8bef\u7387\u7ea6 1%\uff08\u8d85\u5bfc\u6bd4\u7279\uff09\u3002\u7834\u89e3\u8def\u5f84\u662f\u91cf\u5b50\u7ea0\u9519\uff0c\u4f46 "quantum error correction imposes a heavy overhead cost in number of qubits and number of gates"\uff0c\u6545\u5c5e\u66f4\u8fdc\u76ee\u6807\u3002

## \u4e09\u3001\u4f50\u8bc1\uff1a\u53e6\u4e00\u7bc7\u6587\u732e\u7684\u7b49\u4ef7\u523b\u753b\uff08P1\uff0c\u90e8\u5206\u4e00\u81f4\uff09

*Benchmarking Noisy Intermediate-Scale Quantum\u2026*\uff08arXiv 2303.00445\uff09\xa7I\uff1a
> "We find ourselves in the era of noisy intermediate-scale quantum (NISQ) computation, which is characterized by various obstacles including **restrictive qubit topologies, short coherence times and imperfect quantum gates**."

\u8fd9\u662f\u5bf9"\u566a\u58f0"\u9650\u5236\u7684**\u6210\u56e0\u8865\u5145**\uff08\u62d3\u6251\u53d7\u9650\u3001\u76f8\u5e72\u65f6\u95f4\u77ed\u3001\u95e8\u4e0d\u5b8c\u7f8e\uff09\uff0c\u4e0e Preskill \u7684"\u95e8\u9519\u8bef\u7387\u2192\u7ebf\u8def\u89c4\u6a21\u4e0a\u9650"\u662f\u540c\u4e00\u9650\u5236\u7684\u4e0d\u540c\u8868\u8ff0\u5c42\u9762\uff0c\u4e0d\u6784\u6210\u51b2\u7a81\u3002

## \u56db\u3001\u53ef\u4fe1\u5ea6\u4e0e\u68c0\u7d22\u8bf4\u660e

- **P0\uff08\u76f4\u63a5\u8bc1\u636e\uff09**\uff1aPreskill \u539f\u6587\uff08KB `1e6199ea-\u2026`\uff0c`\u81ea\u7136\u79d1\u5b66\u4e0e\u5730\u7403\u79d1\u5b66/quantum-physics__1801.00862__quantum-computing-in-the-nisq-era-and-beyond.md (part 1 of 3).md`\uff09\u2014\u2014\u5b9a\u4e49\u300150\u2013\u51e0\u767e\u6bd4\u7279\u30011000 \u95e8\u4e0a\u9650\u3001"ceiling on computational power" \u5747\u4e3a\u9010\u5b57\u539f\u6587\u3002
- **P1\uff08\u90e8\u5206\u4f50\u8bc1\uff09**\uff1a2303.00445\uff08\u540c KB\uff09\u2014\u2014\u566a\u58f0\u4e09\u7c7b\u5177\u8c61\u6765\u6e90\u3002
- **\u68c0\u7d22\u6267\u884c**\uff1a`kb_search_two_stage`\uff08\u82f1\u6587\u67e5\u8be2\uff0c20 \u5019\u9009/5 \u547d\u4e2d\uff09+ `kb_search_vector`\uff08\u4e2d\u6587\u67e5\u8be2\uff0c20 \u547d\u4e2d\uff09\uff0c\u53e6 `kb_doc_read` \u8bfb\u53d6\u539f\u6587\u524d 80 \u884c\u6838\u5bf9\u5f15\u6587\u3002\u4e2d\u6587\u8bed\u6599\u4e2d\u672a\u89c1\u5bf9 NISQ \u7684\u72ec\u7acb\u4e2d\u6587\u5b9a\u4e49\u6bb5\u843d\uff08\u4e2d\u6587\u68c0\u7d22\u547d\u4e2d\u7684\u5747\u662f\u540c\u4e00\u6279\u82f1\u6587\u91cf\u5b50\u7269\u7406\u8bba\u6587\u7684 chunk \u955c\u50cf\uff0c\u5982 `Corpus-Chunks800/\u2026`\uff09\uff0c\u6545\u7b54\u6848\u57fa\u4e8e\u82f1\u6587\u539f\u59cb\u6587\u732e\u3002
\end{AnswerText}

\textbf{Retrieved evidence index (all returned units; paths and scores) }
\begin{EvidenceText}
[No source list recorded for this arm.]
\end{EvidenceText}
\subsection{BQ03: What is MultiMedQA and how did Flan-PaLM perform on the MedQA benchmark?}
\begin{tabular}{@{}ll@{}}
Arm & HMRS agent (pinned)\\
Generator & omp / deepseek-flash\\
HTTP / success & 200 / True\\
Wall time / cost & 76.0 s / 0.0602682 USD\\
Judge scores & F=5, C=5, G=5, A=5; overall=5.00\\
Verdict & correct\\
\end{tabular}

\textbf{Dimension rationale}
\begin{EvidenceText}
{
  "factual": "The abstract is quoted verbatim and correctly; seven datasets including HealthSearchQA (3,375 questions) match the paper; the section 4.1 quote ('67.6%... surpassing the DRAGON model by 20.1%') and the 50.3%/17.3% comparison are real, and its note that the two percentages use different baselines is an accurate disambiguation; MedMCQA 57.6% and PubMedQA 79.0% match the paper; the human-eval numbers (61.9% consensus, 29.7% harm, 60.6% helpful vs Med-PaLM 92.6/5.8/80.1 and clinicians 92.9/6.5/90.8) match the paper's evaluation section.",
  "completeness": "Both parts answered to full depth: definition with verbatim quote, composition, purpose, and the 12-axis human-evaluation framework; performance with headline number, both comparison baselines, prompting method, sibling dataset scores; plus the limitations section the paper itself emphasizes and the Med-PaLM follow-up.",
  "groundedness": "Everything is labeled P0 direct quotation with the section it came from; the only unverifiable detail is '40 clinician exemplars' for instruction prompt tuning, which is plausible and does not contradict anything.",
  "sources": "Strongest sourcing in the packet: KB id, full file path with parts 1/8-3/8 mapped to sections, cross-verification against Corpus-Chunks800 chunk copies, and a transparent note that the laya judge module was unavailable."
}
\end{EvidenceText}

\textbf{Judge summary}
\begin{EvidenceText}
Complete, verbatim-quoted answer with correct numbers throughout, explicit baseline disambiguation, and file-level provenance for every claim.
\end{EvidenceText}

\textbf{Complete answer}
\begin{AnswerText}
**\u7ed3\u8bba\uff1aMultiMedQA \u662f Google \u5728\u300aLarge Language Models Encode Clinical Knowledge\u300b\u4e2d\u63d0\u51fa\u7684\u533b\u5b66\u95ee\u7b54\u57fa\u51c6\uff08\u7531 6 \u4e2a\u65e2\u6709\u516c\u5f00\u6570\u636e\u96c6\u52a0\u65b0\u5efa\u7684 HealthSearchQA \u5171 7 \u4e2a\u6570\u636e\u96c6\u7ec4\u6210\uff09\uff1bFlan-PaLM 540B \u5728\u8be5\u57fa\u51c6\u7684 MedQA\uff08USMLE \u56db\u9009\u4e00\uff09\u4e0a\u53d6\u5f97 67.6% \u7684\u51c6\u786e\u7387\uff0c\u8d85\u8fc7\u6b64\u524d SOTA \u903e 17%\u3002**

### \u4e00\u3001MultiMedQA \u662f\u4ec0\u4e48
- \u5b9a\u4e49\uff08\u6458\u8981\u539f\u6587\uff09\uff1a"we present MultiMedQA, a benchmark combining six existing open question answering datasets spanning professional medical exams, research, and consumer queries; and HealthSearchQA, a new free-response dataset of medical questions searched online."
- \u7ec4\u6210\uff08\u7b2c 1 \u8282\u660e\u786e\u5217\u51fa\uff09\uff1a\u516d\u4e2a\u65e2\u6709\u6570\u636e\u96c6 **MedQA\u3001MedMCQA\u3001PubMedQA\u3001LiveQA\u3001MedicationQA\u3001MMLU clinical topics**\uff0c\u52a0\u4e0a\u65b0\u6784\u5efa\u7684\u7b2c\u4e03\u4e2a **HealthSearchQA**\uff083,375 \u6761\u5e38\u89c1\u7f51\u7edc\u641c\u7d22\u5065\u5eb7\u95ee\u9898\uff09\u3002
- \u8bbe\u8ba1\u76ee\u7684\uff1a\u8986\u76d6"medical exam / medical research / consumer medical questions"\u4e09\u7c7b\u573a\u666f\uff0c\u5e76\u914d\u5957\u4e00\u4e2a**\u4eba\u5de5\u8bc4\u6d4b\u6846\u67b6**\uff08\u533b\u5e08 + \u666e\u901a\u7528\u6237\uff09\uff0c\u8bc4\u4f30\u4e8b\u5b9e\u6027\u3001\u79d1\u5b66\u5171\u8bc6\u4e00\u81f4\u6027\u3001\u6f5c\u5728\u5371\u5bb3\u3001\u504f\u5dee\u3001\u53ef\u8bfb\u6027\u3001\u68c0\u7d22\u4e0e\u63a8\u7406\u8bc1\u636e\u7b49\u8f74\uff08Table 2 \u5171 12 \u4e2a\u8f74\uff0cTable 3 \u4e3a\u666e\u901a\u7528\u6237 2 \u4e2a\u8f74\uff09\u2014\u2014\u56e0\u4e3a\u591a\u9009\u9898\u51c6\u786e\u7387\u65e0\u6cd5\u53cd\u6620\u771f\u5b9e\u4e34\u5e8a\u53ef\u7528\u6027\u3002

### \u4e8c\u3001Flan-PaLM \u5728 MedQA \u4e0a\u7684\u8868\u73b0
- **\u6838\u5fc3\u6570\u5b57**\uff1a"On the MedQA dataset consisting of USMLE style questions with 4 options, our Flan-PaLM 540B model achieved a multiple-choice question (MCQ) accuracy of **67.6%**, surpassing the DRAGON model by 20.1%."\uff08\xa74.1\uff09
- \u76f8\u5bf9\u63d0\u5347\u53e3\u5f84\uff1a\u6587\u4e2d\u53e6\u8ff0\u6b64\u524d\u6700\u4f18\u6a21\u578b\u4e3a **50.3%**\uff0c"Flan-PaLM 540B exceeded this by **17.3%**"\uff0c\u7ae0\u8282\u6807\u9898\u5373 "Flan-PaLM exceeds previous state-of-the-art on MedQA (USMLE) by over 17%"\u3002\u4e24\u4e2a\u767e\u5206\u6bd4\u5bf9\u5e94\u4e0d\u540c\u6bd4\u8f83\u57fa\u7ebf\uff08DRAGON \u4e0e 50.3% \u90a3\u6761 SOTA \u8bb0\u5f55\uff09\uff0c\u5f15\u7528\u65f6\u9700\u533a\u5206\u3002
- \u53d6\u5f97\u65b9\u5f0f\uff1a**few-shot + chain-of-thought + self-consistency** \u7ec4\u5408\u63d0\u793a\uff08\u975e\u5fae\u8c03\uff09\uff0c\u6545\u79f0"\u4ec5\u9760\u63d0\u793a\u7b56\u7565"\u3002
- \u540c\u6279\u591a\u9009\u9898\u6210\u7ee9\uff1aMedMCQA **57.6%**\u3001PubMedQA **79.0%**\u3001MMLU clinical topics \u4ea6\u8fbe SOTA\uff08"state-of-the-art accuracy on every MultiMedQA multiple-choice dataset"\uff09\u3002

### \u4e09\u3001\u5fc5\u987b\u4e00\u5e76\u8bf4\u660e\u7684\u5c40\u9650\uff08\u8bba\u6587\u7684\u5173\u952e\u8bba\u70b9\uff09
\u591a\u9009\u9898\u5f3a \u2260 \u53ef\u4e34\u5e8a\u4f7f\u7528\u3002\u4eba\u5de5\u8bc4\u6d4b\u4e2d Flan-PaLM \u957f\u6587\u672c\u56de\u7b54\uff1a
- \u4ec5 **61.9%** \u4e0e\u79d1\u5b66\u5171\u8bc6\u4e00\u81f4\uff08Med-PaLM 92.6%\uff0c\u4e34\u5e8a\u533b\u751f 92.9%\uff09\uff1b
- **29.7%** \u88ab\u8bc4\u53ef\u80fd\u5bfc\u5411\u6709\u5bb3\u7ed3\u679c\uff08Med-PaLM 5.8%\uff0c\u533b\u751f 6.5%\uff09\uff1b
- \u666e\u901a\u7528\u6237\u8ba4\u4e3a\u6709\u5e2e\u52a9\u7684\u4ec5 60.6%\uff08Med-PaLM 80.1%\uff0c\u533b\u751f 90.8%\uff09\u3002
\u6b63\u56e0\u5982\u6b64\uff0c\u4f5c\u8005\u63d0\u51fa **instruction prompt tuning**\uff0840 \u6761\u4e34\u5e8a\u533b\u751f\u793a\u4f8b\uff09\u5f97\u5230 Med-PaLM\u3002

**\u51fa\u5904**
- \u300aLarge Language Models Encode Clinical Knowledge\u300b(arXiv 2212.13138)\uff0cKB `b28f0a25-1e1d-40ad-96b8-61bd2b093a28`\uff0f\u8def\u5f84 `\u751f\u547d\u79d1\u5b66\u4e0e\u533b\u5b66/clinical-medicine__2212.13138__...md`\uff0c\u7b2c 1/8 \u90e8\u5206\uff08\u6458\u8981\u3001\xa71 \u8d21\u732e\u5217\u8868\uff09\u3001\u7b2c 2/8 \u90e8\u5206\uff08\xa73.2\u2013\xa74.1\uff09\u3001\u7b2c 3/8 \u90e8\u5206\uff08\xa74.5 \u4eba\u8bc4\u6570\u636e\uff09\u3002
- \u4ea4\u53c9\u5370\u8bc1\uff1a\u540c\u6587\u5206\u5757\u7248 `Corpus-Chunks800/...__k30.md`\u3001`k31.md`\uff08KB `ce642f28-e157-48af-b669-49862cd73b23`\uff09\u542b\u540c\u4e00 MedQA \u8868\u683c\u7247\u6bb5 "Flan-PaLM (540 B)\u2026"\u3002

\u8bc1\u636e\u7b49\u7ea7\uff1aP0\uff08\u539f\u6587\u76f4\u63a5\u9648\u8ff0\uff0c\u542b\u5177\u4f53\u6570\u5b57\u4e0e\u7ae0\u8282\u6807\u9898\uff09\u3002\u6ce8\uff1a\u672c\u6b21 `kb_laya_judge` \u56e0\u540e\u7aef `laya` \u6a21\u5757\u7f3a\u5931\u4e0d\u53ef\u7528\uff0c\u5224\u5b9a\u73af\u8282\u6539\u4ee5\u539f\u6587\u76f4\u8bfb\u5b8c\u6210\uff0c\u4e0d\u5f71\u54cd\u4e0a\u8ff0\u5f15\u6587\u53ef\u9760\u6027\u3002
\end{AnswerText}

\textbf{Retrieved evidence index (all returned units; paths and scores) }
\begin{EvidenceText}
[No source list recorded for this arm.]
\end{EvidenceText}
\subsection{BQ04: Which three independent ontologies subdivide the Gene Ontology?}
\begin{tabular}{@{}ll@{}}
Arm & HMRS agent (pinned)\\
Generator & omp / deepseek-flash\\
HTTP / success & 200 / True\\
Wall time / cost & 30.1 s / 0.043136999999999995 USD\\
Judge scores & F=5, C=5, G=5, A=5; overall=5.00\\
Verdict & correct\\
\end{tabular}

\textbf{Dimension rationale}
\begin{EvidenceText}
{
  "factual": "Names the three ontologies with the gold-verbatim quote from section 2 including 'non-redundant'; the added GAF evidence ('field 9 denotes the sub-ontology of GO, either Molecular Function, Biological Process, or Cellular Component') is genuine primer content; the October 2015 statistics (43,835 terms, 73,776 is_a, 7,436 part_of, 8,263 regulates-type relations) are gold-verbatim, and its interpretation that the three share one identifier space divided by biological aspect is correct.",
  "completeness": "All three ontologies named, plus two independent supporting passages and structural statistics \u2014 fully answers the question.",
  "groundedness": "Both quotes are labeled P0 direct evidence with the section they come from (section 2 and section 5.2 on the GAF file), and the statistics match the gold word-for-word; nothing unsupported.",
  "sources": "The strongest attribution in the packet: full file path ('genomics__1602.01876__primer-on-the-gene-ontology.md (part 1 of 2).md'), KB id, chunk index, and a note that two independent statements were cross-checked."
}
\end{EvidenceText}

\textbf{Judge summary}
\begin{EvidenceText}
Complete, verbatim-quoted answer with a second corroborating GAF passage, gold statistics, and file-level provenance.
\end{EvidenceText}

\textbf{Complete answer}
\begin{AnswerText}
**Gene Ontology \u7531\u4e09\u4e2a\u4e92\u4e0d\u91cd\u53e0\u7684\u5b50\u672c\u4f53\uff08sub-ontologies\uff09\u6784\u6210\uff1aMolecular Function\uff08MF\uff0c\u5206\u5b50\u529f\u80fd\uff09\u3001Biological Process\uff08BP\uff0c\u751f\u7269\u8fc7\u7a0b\uff09\u3001Cellular Component\uff08CC\uff0c\u7ec6\u80de\u7ec4\u5206\uff09\u3002**

\u539f\u6587\uff08P0 \u76f4\u63a5\u8bc1\u636e\uff09\uff1a

> "The terms are subdivided in three distinct ontologies that represent different biological aspects: **Molecular Function (MF), Biological Process (BP), and Cellular Component (CC)** (1). These ontologies are **non-redundant** and share a common space of identifiers and a well-specified syntax."
> \u2014 *Primer on the Gene Ontology*, \xa72 "What is the Gene Ontology?"

\u8865\u5145\u4f50\u8bc1\uff08\u540c\u4e00\u6587\u732e\uff0c\xa75.2\uff0c\u63cf\u8ff0 GAF \u7b2c 9 \u5217\uff09\uff1a

> "field 9 denotes the **sub-ontology** of GO, either Molecular Function, Biological Process, or Cellular Component."

\u5373\uff1aGO \u6ce8\u91ca\u6587\u4ef6\u4e2d\u4e13\u8bbe\u4e00\u5217\u663e\u5f0f\u6807\u6ce8\u8be5\u6ce8\u91ca\u5c5e\u4e8e\u54ea\u4e00\u4e2a\u5b50\u672c\u4f53\uff0c\u4fbf\u4e8e\u6309 GO aspect \u89e3\u6790\u3002

\u51fa\u5904\uff1a
- \u6587\u6863\uff1a`\u751f\u547d\u79d1\u5b66\u4e0e\u533b\u5b66/genomics__1602.01876__primer-on-the-gene-ontology.md (part 1 of 2).md`\uff08Gaudet, \u0160kunca, Hu, Dessimoz, *Primer on the Gene Ontology*, Jan 2016\uff09
- \u77e5\u8bc6\u5e93\uff1a`b28f0a25-1e1d-40ad-96b8-61bd2b093a28`\uff08\u751f\u547d\u79d1\u5b66\u4e0e\u533b\u5b66\uff09
- \u53ef\u4fe1\u5ea6\uff1aP0\uff08\u68c0\u7d22\u547d\u4e2d chunk_index 9 \u539f\u6587\u6bb5\u843d + \u76f4\u63a5\u8bfb\u53d6\u5168\u6587\u6838\u5bf9\uff0c\u4e24\u5904\u72ec\u7acb\u8868\u8ff0\u4e00\u81f4\uff09

\u987a\u5e26\u4e00\u6761\u53ef\u6bd4\u5bf9\u7684\u7ed3\u6784\u6570\u636e\uff08\u540c\u8282\uff0cOct 2015 \u7248\uff09\uff1a\u5168\u672c\u4f53 43,835 terms\u300173,776 \u6761 `is_a`\u30017,436 \u6761 `part_of`\u30018,263 \u6761 `regulates` \u7c7b\u5173\u7cfb\u2014\u2014\u8bf4\u660e\u4e09\u4e2a\u5b50\u672c\u4f53\u5171\u4eab\u540c\u4e00\u5957 ID \u4e0e\u53e5\u6cd5\uff0c\u4ec5\u6309"\u751f\u7269\u5b66\u4fa7\u9762"\u5212\u5206\uff0c\u800c\u975e\u6309\u6807\u8bc6\u7b26\u4f53\u7cfb\u5212\u5206\u3002
\end{AnswerText}

\textbf{Retrieved evidence index (all returned units; paths and scores) }
\begin{EvidenceText}
[No source list recorded for this arm.]
\end{EvidenceText}
\subsection{BQ05: Why does historical European economic growth contradict Unified Growth Theory?}
\begin{tabular}{@{}ll@{}}
Arm & HMRS agent (pinned)\\
Generator & omp / deepseek-flash\\
HTTP / success & 200 / True\\
Wall time / cost & 119.9 s / 0.0623952 USD\\
Judge scores & F=5, C=5, G=5, A=5; overall=5.00\\
Verdict & correct\\
\end{tabular}

\textbf{Dimension rationale}
\begin{EvidenceText}
{
  "factual": "All verifiable details correct: the three Galor regime dates (100,000 BC-1750, 1750-1870, from 1870), Industrial Revolution 1760-1840 as 'prime engine' (Galor 2005a, p. 212), the four-point mutilation of Maddison data, and the verbatim conclusion quotes. The fit parameters (12 countries a=1.147e-1/k=5.961e-5/singularity 1923/23-year bypass; 30 countries 1929/29y; Eastern 1915/25y) agree with independent systems quoting the same paper, as do the AD 1 +27%/AD 1000 -54% deviations.",
  "completeness": "The fullest treatment: what UGT claims, the reciprocal-coordinate test method, exact quantitative results for all three datasets, four enumerated contradictions including the perverse direction of the divergence, the data-misuse critique, and the author's conclusion \u2014 everything the question asks plus essential context.",
  "groundedness": "Labeled P0 verbatim evidence; quotes like 'the economic growth was diverted to a slower trajectory around 1900, bypassing the singularity by about 23 years' are consistent with the document; the minor internal difference between '~1880' in one section and '~1900' in another plausibly reflects different country aggregates, which other answers corroborate.",
  "sources": "File-level attribution (paper, arXiv id, KB name/id, evidence tier) plus honest notes that this is a one-sided critique with no Galor rebuttal in the KB and that irrelevant retrieval hits were rejected."
}
\end{EvidenceText}

\textbf{Judge summary}
\begin{EvidenceText}
Comprehensive, verbatim-sourced answer combining the conceptual contradictions with exact fit parameters and method for all three European datasets.
\end{EvidenceText}

\textbf{Complete answer}
\begin{AnswerText}
**\u7ed3\u8bba\uff1aNielsen (2015) \u7528 Maddison (2010) \u7684\u6b27\u6d32 GDP \u6570\u636e\u8bc1\u660e\uff0c\u52a0\u5c14 (Galor) \u7684\u7edf\u4e00\u589e\u957f\u7406\u8bba\uff08UGT\uff09\u6240\u4f9d\u8d56\u7684\u4e09\u4e2a\u589e\u957f\u9636\u6bb5\u5728\u6b27\u6d32\u6839\u672c\u4e0d\u5b58\u5728\uff1a\u6b27\u6d32\u7684\u589e\u957f\u59cb\u7ec8\u662f\u53cc\u66f2\u5f0f\u7684\uff0c\u65e2\u6ca1\u6709\u201c\u9a6c\u5c14\u8428\u65af\u505c\u6ede\u201d\uff0c\u4e5f\u6ca1\u6709\u56e0\u5de5\u4e1a\u9769\u547d\u800c\u51fa\u73b0\u201c\u8d77\u98de\u201d\u8f6c\u6298\uff1b\u56e0\u6b64\u8be5\u7406\u8bba\u4e0e\u8bc1\u636e\u76f4\u63a5\u51b2\u7a81\u3002**

**1. UGT \u7684\u6838\u5fc3\u4e3b\u5f20\uff08\u88ab\u68c0\u9a8c\u7684\u5bf9\u8c61\uff09**
Galor \u7684\u57fa\u672c\u5047\u8bbe\u662f\u4e09\u4e2a\u589e\u957f\u9636\u6bb5\uff1a\u9a6c\u5c14\u8428\u65af\u505c\u6ede\u671f\uff08\u4ece\u516c\u5143\u524d 10 \u4e07\u5e74\u5230 1750 \u5e74\uff09\u3001\u540e\u9a6c\u5c14\u8428\u65af\u671f\uff081750\u20131870\uff09\u3001\u6301\u7eed\u589e\u957f\u671f\uff081870 \u5e74\u8d77\uff09\uff1b\u5de5\u4e1a\u5316\uff081760\u20131840\uff09\u662f\u201c\u7ecf\u6d4e\u589e\u957f\u7684\u9996\u8981\u5f15\u64ce\u201d\uff08Galor 2005a, p. 212\uff09\uff0c\u5e76\u5ba3\u79f0\u5b9e\u73b0\u4e86\u4ece\u9a6c\u5c14\u8428\u65af\u9677\u9631\u7684\u201cremarkable / stunning escape\u201d\uff08Galor 2005a, pp. 177, 220\uff09\u3002Nielsen \u6307\u51fa\uff0c\u6b27\u6d32\u5c24\u5176\u662f\u897f\u6b27\u672c\u5e94\u662f\u8be5\u7406\u8bba\u6700\u6709\u529b\u7684\u9a8c\u8bc1\u573a\u666f\uff1a\u5982\u679c\u8fd9\u4e9b\u673a\u5236\u771f\u5b9e\u5b58\u5728\uff0c\u5176\u6548\u5e94\u5728\u8fd9\u4e9b\u5730\u533a\u5c31\u5e94\u5f53\u8868\u73b0\u5f97\u6700\u660e\u663e\uff08P0\uff0c\u8be5\u6587 Introduction \u7684\u68c0\u9a8c\u8bbe\u5b9a\uff09\u3002

**2. \u53cd\u9a73\u7684\u65b9\u6cd5\uff1a\u7528\u5012\u6570\u5750\u6807\u68c0\u9a8c**
\u53cc\u66f2\u589e\u957f $S(t)=(a-kt)^{-1}$ \u7684\u5012\u6570\u4e3a\u76f4\u7ebf $1/S(t)=a-kt$\u3002\u5982\u679c\u589e\u957f\u786e\u5b9e\u5b58\u5728\u505c\u6ede\u671f\uff0c\u5012\u6570\u66f2\u7ebf\u5e94\u5448\u73b0\u968f\u673a\u6ce2\u52a8\u7684\u201c\u9a6c\u5c14\u8428\u65af\u632f\u8361\u201d\uff1b\u5982\u679c\u53d1\u751f\u4e86\u8d77\u98de\uff0c\u5012\u6570\u66f2\u7ebf\u7684\u659c\u7387\u5e94\u660e\u663e\u5411\u4e0b\u504f\u6298\uff1b\u800c\u4e00\u6761\u659c\u7387\u4e0d\u53d8\u7684\u76f4\u7ebf\uff0c\u6309 Nielsen \u7684\u8bf4\u6cd5\uff0c\u88ab\u4efb\u610f\u5207\u5206\u6210\u4e24\u6bb5\u6216\u4e09\u6bb5\u5e76\u5ba3\u79f0\u4e3a\u4e0d\u540c\u589e\u957f\u673a\u5236\uff0c\u5728\u65b9\u6cd5\u4e0a\u6ca1\u6709\u610f\u4e49\u3002

**3. \u6570\u636e\u7ed3\u679c\uff08\u5173\u952e\u6570\u5b57\uff0c\u9010\u5b57\u5f15\u7528\uff09**
- \u897f\u6b27 12 \u4e2a\u9886\u5148\u56fd\u5bb6\uff08\u5360\u897f\u6b27 30 \u56fd GDP \u7684 85%\uff09\uff1a$a=1.147\times10^{-1}$\uff0c$k=5.961\times10^{-5}$\uff1b\u5947\u70b9\u5728 1923 \u5e74\uff0c\u201cthe economic growth was diverted to a slower trajectory around 1900, bypassing the singularity by about 23 years\u201d\u3002\u62df\u5408\u5728 AD 1500\u20131900 \u5e74\u95f4\u201cremarkably good\u201d\u3002
- \u897f\u6b27 30 \u56fd\uff1a$a=9.859\times10^{-2}$\uff0c$k=5.112\times10^{-5}$\uff1b\u5947\u70b9\u5728 1929 \u5e74\uff0c\u7ea6\u5728 1900 \u5e74\u8f6c\u5411\u8f83\u6162\u7684\u8f68\u8ff9\uff0c\u6bd4\u5947\u70b9\u63d0\u524d\u7ea6 29 \u5e74\u3002
- \u4e1c\u6b27\uff1a$a=7.749\times10^{-1}$\uff0c$k=4.048\times10^{-4}$\uff1b\u5947\u70b9\u5728 1915 \u5e74\uff0c\u7ea6\u5728 1890 \u5e74\u8f6c\u5411\u8f83\u6162\u7684\u8f68\u8ff9\uff0c\u6bd4\u5947\u70b9\u63d0\u524d\u7ea6 25 \u5e74\u3002

**4. \u7531\u6b64\u5f97\u5230\u7684\u56db\u65b9\u9762\u77db\u76fe**
(i) \u6ca1\u6709\u505c\u6ede\uff1aAD 1500\u20131900 \u5e74\u95f4\u589e\u957f\u6cbf\u7740\u7a33\u5b9a\u7684\u53cc\u66f2\u8f68\u8ff9\u4e0a\u5347\uff0c\u201cThere is no sign of the existence of Malthusian stagnation\u201d\uff1b\u800c\u4e14\u201cthe alleged remarkable or stunning escape from Malthusian trap did not happen because there was no trap\u201d\u3002
(ii) \u6ca1\u6709\u8f6c\u6298\u70b9\uff1a\u201cAbsolutely nothing had happened at the end of the alleged Malthusian regime. There was no transition from stagnation to growth at any time.\u201d \u5de5\u4e1a\u9769\u547d\u5728\u6b27\u6d32\u5bf9\u589e\u957f\u8f68\u8ff9\u201chad absolutely no impact\u201d\u3002
(iii) \u65b9\u5411\u76f8\u53cd\uff1a\u5728\u88ab\u5ba3\u79f0\u8fdb\u5165\u6301\u7eed\u589e\u957f\u671f\u4e4b\u540e\uff0c\u589e\u957f\u53cd\u800c\u201cstarted to be diverted to a slower trajectory\u201d\uff1b\u897f\u6b27\u4ece 1880 \u5e74\u524d\u540e\u5f00\u59cb\uff0c\u4e1c\u6b27\u4ece 1890 \u5e74\u524d\u540e\u5f00\u59cb\u3002
(iv) \u6570\u636e\u4f7f\u7528\u672c\u8eab\u6709\u95ee\u9898\uff1aGalor \u53cd\u590d\u53ea\u53d6 Maddison \u6570\u636e\u4e2d\u7684\u56db\u4e2a\u9009\u5b9a\u70b9\u3001\u7528\u76f4\u7ebf\u8fde\u63a5\uff0c\u8fd9\u6837\u5904\u7406\u6570\u636e\u201cunscientific\u201d\u3002\u8be5\u7406\u8bba\u751a\u81f3\u88ab\u5176\u81ea\u8eab\u5efa\u6a21\u65f6\u6240\u7528\u7684\u540c\u4e00\u6279\u6570\u636e\u6240\u5426\u5b9a\u3002

**5. \u4f5c\u8005\u7ed3\u8bba**
\u201cUnified Growth Theory does not explain the mechanism of the economic growth because its explanations are based on mythical features, which did not exist\u201d\uff1b\u8be5\u7406\u8bba\u201cneeds to be either thoroughly revised or most likely replaced by a theory supported by a professional analysis of economic growth data.\u201d \u8be5\u7ed3\u8bba\u4e0e\u5176\u66f4\u65e9\u7684\u7cfb\u5217\u7814\u7a76\u4e00\u81f4\uff08\u975e\u6d32\u3001\u4e9a\u6d32\u3001\u524d\u82cf\u8054\u3001\u4e16\u754c\u603b\u91cf\u3001\u897f\u6b27\uff1aNielsen 2014, 2015a\u2013e\uff09\u3002

**\u51fa\u5904**\uff1aRon W Nielsen, *Unified Growth Theory Contradicted by the Economic Growth in Europe*, Dec 2015\uff08arXiv:1512.08067\uff09\uff0cKB `\u7ecf\u6d4e\u4e0e\u793e\u4f1a`\uff08kb f5770687\uff09\u3002\u8bc1\u636e\u7b49\u7ea7\uff1aP0\uff08\u76f4\u63a5\u539f\u6587\uff09\u3002\u6ce8\u610f\uff1a\u8fd9\u662f\u4f5c\u8005\u5355\u65b9\u9762\u7684\u6279\u8bc4\u6027\u8bba\u8bc1\uff0cK B \u5185\u672a\u89c1 Galor \u4e00\u65b9\u7684\u56de\u5e94\u6587\u732e\u3002

*\u68c0\u7d22\u8bf4\u660e\uff1a\u5bf9 5 \u4e2a\u6307\u5b9a\u5e93\u6267\u884c\u4e86 kb_search_two_stage \u4e0e kb_search_vector\uff08\u5404 top_k 20\u201325\uff09\uff0c\u6709\u6548\u547d\u4e2d\u96c6\u4e2d\u5728\u4e0a\u8ff0\u5355\u4e00\u6587\u6863\uff1b\u5176\u4f59\u547d\u4e2d\u4e3a\u9898\u540d/\u5173\u952e\u8bcd\u5076\u5408\uff08\u5982 leisure class\u3001\u53cd\u5784\u65ad\u6cd5\uff09\uff0c\u4e0d\u6784\u6210\u76f8\u5173\u8bc1\u636e\u3002*
\end{AnswerText}

\textbf{Retrieved evidence index (all returned units; paths and scores) }
\begin{EvidenceText}
[No source list recorded for this arm.]
\end{EvidenceText}
\subsection{BQ06: Which physical factor primarily controls the response of precipitation extremes to climate change?}
\begin{tabular}{@{}ll@{}}
Arm & HMRS agent (pinned)\\
Generator & omp / deepseek-flash\\
HTTP / success & 200 / True\\
Wall time / cost & 44.0 s / 0.061153200000000005 USD\\
Judge scores & F=5, C=5, G=5, A=5; overall=5.00\\
Verdict & correct\\
\end{tabular}

\textbf{Dimension rationale}
\begin{EvidenceText}
{
  "factual": "Directly and correctly answers: thermodynamic control (saturation vapor pressure/specific humidity rise, ~6-7%/K Clausius-Clapeyron) with dynamics and precipitation efficiency as secondary modifiers. Every quote matches the paper: the Eq. 1 decomposition wording, 'relatively small contributions from changes in vertical velocities and precipitation efficiency', the abstract robustness ranking, and the conclusions sentence 'the easiest to understand and always gives an intensification'; caveats (below-295K non-CC behavior, subdaily convective exceedances, orographic/snowfall differences, tropical uncertainty) are genuine.",
  "completeness": "Answers with the factor, the mechanism, the magnitude, the robustness ranking, and the caveats \u2014 fully answered.",
  "groundedness": "Chunk-level citations (part 1 chunks 2, 6, 23/26, 28, 33; part 2 chunks 20, 25) with verbatim quotes; the one presentational flaw is a redundant restatement of the same 6-7%/K figure as if contrasting two values, but the number itself is consistent.",
  "sources": "Strongest sourcing in the packet: full file path, KB id, per-quote chunk indices, plus transparent disclosure of the failed Laya evidence gate (ModuleNotFoundError) and negative checks on other KBs."
}
\end{EvidenceText}

\textbf{Judge summary}
\begin{EvidenceText}
Decisive thermodynamic answer with chunk-verified verbatim quotes, quantitative scaling, and full caveat coverage.
\end{EvidenceText}

\textbf{Complete answer}
\begin{AnswerText}
**The thermodynamic factor \u2014 the rise in saturation vapor pressure (hence near-surface specific humidity) with warming, i.e. Clausius\u2013Clapeyron scaling at about 6\u20137 % K\u207b\xb9 \u2014 is the primary control; changes in vertical velocity (dynamics) and precipitation efficiency are secondary modifiers** (P0 direct evidence).

Evidence base: `\u81ea\u7136\u79d1\u5b66\u4e0e\u5730\u7403\u79d1\u5b66/climate-science__1503.07557__precipitation-extremes-under-climate-change.md` (O'Gorman, *Precipitation extremes under climate change*), KB `1e6199ea\u2026` \u2014 parts 1 and 2.

1. **Decomposition (Eq. 1).** $P \simeq -\epsilon\{\omega(p)\,S(T,p)\}$, where \u201c$S(T,p)=dq_s/dp|_{\theta_e^*}$ is the derivative of the saturation specific\u2026\u201d \u2014 the paper states this yields \u201ca dynamical contribution from changes in $\omega$, a thermodynamic contribution from changes in $S$ (this is termed thermodynamic since $S$ only depends on temperature and pressure), and a microphysical component from changes in the precipitation efficiency $\epsilon$.\u201d (part 1, chunk 23/26)

2. **Thermodynamics dominates the canonical case.** \u201cFor surface temperatures near those of the present-day tropics, the precipitation extremes increase at close to the thermodynamic rate, and this is close to Clausius-Clapeyron scaling with the surface specific humidity, **with relatively small contributions from changes in vertical velocities and precipitation efficiency**\u201d; same for organized squall-line convection. (part 1, chunk 33)

3. **Mechanism + magnitude.** When dynamics and $\epsilon$ are \u201cnegligible, precipitation extremes increase with warming because of increases in the saturation vapor pressure of water\u201d; thermodynamic scaling \u201cgives a sensitivity of 6\u20137 % K\u207b\xb9\u201d above ~295 K, versus \u201c6\u20137 % K\u207b\xb9 for temperatures above 295K\u201d for CC scaling overall. (part 1, chunks 6, 28; part 2, chunk 20)

4. **Robustness ranking.** \u201cThe thermodynamic contribution is robust and well understood, but theoretical understanding of the microphysical and dynamical contributions is still being developed.\u201d (part 1, chunk 2); \u201cThe thermodynamic contribution is the easiest to understand and always gives an intensification with warming.\u201d (part 2, chunk 25)

**Caveats from the same source (P1):** below ~295 K the response is not CC-like \u2014 sensitivity \u201cvaries widely depending on temperature and accumulation period,\u201d controlled by precipitation-efficiency changes and convective dynamics; subdaily convective extremes can exceed CC scaling via dynamics/$\epsilon$; orographic and snowfall extremes behave differently. Tropical sensitivity remains the largest stated uncertainty.

Note: the Laya evidence-gate backend was unavailable (`ModuleNotFoundError: No module named 'laya'`, 19 segments errored), so scoring was replaced by direct chunk verification above; no other pinned KB contained relevant material (KB 1 returned CS/neuro papers, KB 3 returned biomedical papers \u2014 zero on-topic hits).
\end{AnswerText}

\textbf{Retrieved evidence index (all returned units; paths and scores) }
\begin{EvidenceText}
[No source list recorded for this arm.]
\end{EvidenceText}
\subsection{BQ07: What is MameLoshnLM and what parameter scale does it have?}
\begin{tabular}{@{}ll@{}}
Arm & HMRS agent (pinned)\\
Generator & omp / deepseek-flash\\
HTTP / success & 200 / True\\
Wall time / cost & 53.5 s / 0.052864799999999997 USD\\
Judge scores & F=5, C=5, G=5, A=5; overall=5.00\\
Verdict & correct\\
\end{tabular}

\textbf{Dimension rationale}
\begin{EvidenceText}
{
  "factual": "All correct: the P0 abstract quote 'the first open-source 8B-parameter language model built specifically for Yiddish'; scale = 8B via Llama 3.1 8B weights; the training-recipe quote (bfloat16, 8-bit AdamW); the five baselines; Oytser statistics (346.3K docs, 915.21M words, 5.28B tokens, YBC books largest at 4.23B tokens); the mC4 finding (42.2% genuine Yiddish) consistent with the gold's 'less than half of the data is genuine Yiddish text'; Kashes task enumeration. The '9 tasks' count differs from another answer's '14-task' claim and cannot be resolved from the shown gold.",
  "completeness": "Both parts answered with verbatim quotes, plus companion resources with exact statistics, the motivating mC4 analysis, and baselines \u2014 fully answered.",
  "groundedness": "Quotes labeled P0 with chunk references (part 1 chunk 54, chunk 69) that match the gold abstract verbatim; nothing unsupported.",
  "sources": "File-level provenance: KB id, exact file path (part 1 of 4), chunk-level citations, mirror-chunk cross-verification, and authors/affiliations."
}
\end{EvidenceText}

\textbf{Judge summary}
\begin{EvidenceText}
Verbatim-sourced, complete answer with exact corpus statistics and training details, anchored to the specific KB file and chunks.
\end{EvidenceText}

\textbf{Complete answer}
\begin{AnswerText}
**MameLoshnLM is the first open-source 8B-parameter language model built specifically for Yiddish \u2014 its scale is 8B parameters, obtained by continued pretraining of Llama 3.1 8B.** (P0, direct evidence)

**What it is**
- "the first open-source 8B-parameter language model built specifically for Yiddish" \u2014 the model is named Mame-Loshn ("mother tongue"), per \xa75.
- Base model / training recipe: "MAMELOSHNLM was produced by continued pretraining of Llama-3.1-8B on a Yiddish corpus using a causal language modeling objective. The model was trained in bfloat16 precision using the 8-bit AdamW optimizer." (part 1, \xa75 region, chunk 54)
- Positioned as the first dedicated LLM effort for Yiddish: "no prior work, to our knowledge, has targeted Yiddish through the lens of modern LM development" (\xa71).

**Parameter scale**
- 8B \u2014 exactly Llama 3.1 8B weights, not a new architecture. It is compared against same-size baselines: Llama 3.1 8B, Qwen3 8B, BLOOMZ 7B, Gemma-2 9B, EuroLLM 9B (Table, part 1 chunk 69).

**Two companion resources delivered with it**
- **Oytser** \u2014 Yiddish pretraining corpus: 346.3K documents, 915.21M words, **5.28B tokens** (tokenizer = MameLoshnLM's own). Largest source: YBC books (12.3K docs, 723.09M words, 4.23B tokens).
- **Kashes** \u2014 first multi-task Yiddish benchmark: 9 tasks (translation incl. new natively-authored Kashes-mt, POS tagging, lemmatization, dependency parsing, transliteration, NER on EHRI-NER/WikiANN/newNLP, PIQA/WikiQA/PAWS-Wiki).

**Motivating finding (relevant context)**
Only **42.2%** of the Yiddish mC4 split is genuine native Yiddish; 29.8% machine-translated, 21.9% misidentified Hebrew, 6.2% other (Table 1). Hence the corpus was rebuilt rather than relying on mC4.

**Sources**
- KB `b1199132-e4d4-4305-8c2e-82dc1753b0ba` \u2192 `\u8ba1\u7b97\u673a\u4e0e\u4eba\u5de5\u667a\u80fd/nlp__2608.05850__mameloshnlm-yiddish-language-model-and-evalu.md (part 1 of 4).md` \u2014 paper "MAMELOSHNLM: Yiddish Language Model and Evaluation Benchmark" (Uri Katz, Omer Goldman, Tomasz Limisiewicz, Reut Tsarfaty, Noah A. Smith; Bar-Ilan / Cambridge / UW / AI2).
- Same doc mirrored in KB `ce642f28-\u2026` as `Corpus-Chunks800/\u2026__k12.md`, `__k13.md` (identical text, chunk-level copies).

Credibility: P0 for both definition and scale (verbatim abstract + \xa75 training paragraph).
\end{AnswerText}

\textbf{Retrieved evidence index (all returned units; paths and scores) }
\begin{EvidenceText}
[No source list recorded for this arm.]
\end{EvidenceText}
\subsection{BQ08: How does oxygen redox contribute to the capacity of battery cathodes?}
\begin{tabular}{@{}ll@{}}
Arm & HMRS agent (pinned)\\
Generator & omp / deepseek-flash\\
HTTP / success & 200 / True\\
Wall time / cost & 75.5 s / 0.060756000000000004 USD\\
Judge scores & F=5, C=5, G=5, A=5; overall=5.00\\
Verdict & correct\\
\end{tabular}

\textbf{Dimension rationale}
\begin{EvidenceText}
{
  "factual": "All correct: the capacity gap (200-300 mAh/g theoretical, rarely attained due to phase changes, interphase barriers, oxygen release; real cathodes exceeding the TM oxidation limit \u2014 gold-verbatim); the orbital mechanism (bonding O 2p ~10 eV below EF, orphaned non-bonding states in Li-/Na-excess); Li2MnO3 retaining Mn4+ with 1.4 Li extracted (gold); the Li4Mn2O5 stepwise sequence (Mn3+->Mn4+, then O2-->O1-, Mn4+->Mn5+ at x<1 with octahedral-tetrahedral change) consistent with peer answers; the 2.0-2.5 eV integration window for labile oxygen electrons; the cost section with fluorination tradeoff.",
  "completeness": "Five parts cover why the capacity gap exists, the orbital mechanism, material examples with stepwise capacity accounting, how the contribution is quantified computationally, and the cost/mitigation \u2014 fully answered.",
  "groundedness": "Every paragraph labeled [P0] with content matching gold and the paper; nothing unsupported.",
  "sources": "File-level attribution at the end (exact file name 'part 1 of 2', KB 8fa4db59) plus an honest retrieval note that only this document produced relevant evidence (top score 0.77)."
}
\end{EvidenceText}

\textbf{Judge summary}
\begin{EvidenceText}
Comprehensive Chinese-language answer with P0-labeled quotes covering the capacity gap, orbital mechanism, examples, quantification, and costs.
\end{EvidenceText}

\textbf{Complete answer}
\begin{AnswerText}
**\u7ed3\u8bba\uff1a\u6c27\uff08\u9634\u79bb\u5b50\uff09\u6c27\u5316\u8fd8\u539f\u662f\u8d85\u51fa\u8fc7\u6e21\u91d1\u5c5e\uff08TM\uff09\u6c27\u5316\u6781\u9650\u7684\u90a3\u90e8\u5206\u5bb9\u91cf\u7684\u6765\u6e90\u2014\u2014\u5728\u9502/\u94a0\u8fc7\u91cf\u5c42\u72b6\u6c27\u5316\u7269\u4e2d\uff0c\u672a\u88ab TM \u6742\u5316\u7684\u201corphan\u201d O 2p \u6001\u5347\u5230\u8d39\u7c73\u80fd\u7ea7\u9644\u8fd1\uff0c\u5145\u7535\u65f6 O\xb2\u207b \u8d21\u732e\u7535\u5b50\uff08\u5f62\u6210\u8fc7\u6c27\u7c7b\u7269\u79cd\u6216 O \u7a7a\u7a74\uff09\uff0c\u4ece\u800c\u5728 TM \u5df2\u88ab\u6c27\u5316\u5230\u6781\u9650\u540e\u7ee7\u7eed\u652f\u6301\u8131\u78b1\u91d1\u5c5e\uff0c\u76f4\u63a5\u62ac\u9ad8\u53ef\u7528\u5bb9\u91cf\uff1b\u4f46\u540c\u4e00\u673a\u5236\u4e5f\u5e26\u6765 O\u2082 \u6790\u51fa\u3001TM \u8fc1\u79fb\u4e0e\u7535\u538b/\u5bb9\u91cf\u8870\u51cf\uff0c\u4e8c\u8005\u5fc5\u987b\u6743\u8861\u3002**

**1. \u4e3a\u4ec0\u4e48\u9700\u8981\u5b83\uff08\u5bb9\u91cf\u7f3a\u53e3\uff09**
\u5728\u5e38\u89c4 LiTMO\u2082 \u4e2d\uff0c\u8131\u9502\u4ec5\u7531 TM \u6c27\u5316\u6001\u53d8\u5316\u627f\u62c5\uff08$\mathrm{LiTM^{3+}O_2 \to Li_{(1-x)}TM_{(1-x)}^{3+}TM_x^{4+}O_2}$\uff09\u3002\u7406\u8bba\u4e0a 200~300 mAh/g \u9700\u5b8c\u5168\u8131\u9502\uff0c\u201c\u800c\u5b9e\u9645\u4e2d\u8fd9\u5f88\u5c11\u53d1\u751f\u201d\u2014\u2014\u76f8\u53d8\u3001\u754c\u9762\u6269\u6563\u963b\u6321\u3001\u6c27\u6790\u51fa\u90fd\u4f1a\u9650\u5236\u5bb9\u91cf\u3002\u7136\u800c\u201c\u5b9e\u9645\u6b63\u6781\u6750\u6599\u7684\u5bb9\u91cf\u5df2\u8fbe\u5230\u751a\u81f3\u8d85\u8fc7\u8fd9\u4e00\u6c34\u5e73\uff0c\u5448\u73b0\u51fa\u8d85\u8fc7\u8fc7\u6e21\u91d1\u5c5e\u6c27\u5316\u6781\u9650\u7684 overcapacity\u201d\u3002\u8fd9\u6b63\u662f\u9634\u79bb\u5b50\u6c27\u5316\u8fd8\u539f\u7684\u8d21\u732e\u3002[P0]

**2. \u673a\u5236\uff08\u8f68\u9053\u5c42\u9762\uff09**
- \u5e38\u89c4\u7ed3\u6784\u4e2d O \u7531 3 \u4e2a TM + 3 \u4e2a Li \u914d\u4f4d\uff0cTM \u7684 s/p \u6742\u5316\u628a\u6210\u952e O 2p \u62c9\u5230\u4ef7\u5e26\u5e95\uff08\u8d39\u7c73\u80fd\u7ea7\u4e0b ~10 eV\uff09\uff0c\u6545\u6c27\u5316\u53ea\u53d1\u751f\u5728 TM t\u2082g \u4e0a\u3002
- Li-excess / Na-excess\uff08\u5c42\u5185 TM \u4f4d\u88ab Li/Na \u53d6\u4ee3\uff09\u4f7f\u90e8\u5206 O \u5e73\u5747\u88ab\u5c11\u4e8e 3 \u4e2a TM \u914d\u4f4d\uff0c\u51fa\u73b0 Li\u2013O\u2013Li\uff08Na\u2013O\u2013Na\uff09\u73af\u5883\uff0c\u201c\u65e0\u6742\u5316\u7684 orphaned O 2p \u6001\u201d\u80fd\u91cf\u63a5\u8fd1\u975e\u952e t\u2082g \u6001\uff0c\u4e8e\u662f\u201cO \u4e0e TM \u90fd\u53ef\u80fd\u7ed9\u51fa\u6216\u63a5\u53d7\u7535\u5b50\u201d\u3002
- \u5145\u7535\u65f6\u7684\u53ef\u9006\u8868\u8ff0\uff1a$\mathrm{O^{2-} \leftrightarrow (O_2)^{n-}}$\uff1b\u66ff\u4ee3\u89e3\u91ca\u662f\u6c27\u4e0a\u5f62\u6210\u7535\u5b50\u7a7a\u7a74\uff0c\u4ee5\u907f\u514d\u6709\u5bb3 O\u2082 \u6c14\u4f53\u9038\u51fa\u3002[P0]

**3. \u6750\u6599\u5b9e\u4f8b\u4e0e\u5206\u6bb5\u8d21\u732e**
- Li\u2082MnO\u2083\uff1aMn \u7ef4\u6301 +4\u3001t\u2082g \u80fd\u7ea7\u8fc7\u4f4e\u4e0d\u80fd\u6c27\u5316\uff0c\u5374\u4ecd\u53ef\u8131\u51fa 1.4 Li/\u5316\u5b66\u5f0f\uff1b\u5176\u6027\u80fd\u201c\u88ab\u5f52\u56e0\u4e8e\u6c27\u6c27\u5316\u8fd8\u539f\u6d3b\u6027\u201d\u3002[P0]
- Li\u2084Mn\u2082O\u2085\uff1a$\mathrm{Li_4^+Mn_2^{3+}O_5^{2-} \to 4Li^+ + 4e^- + Mn_x^{4+}Mn_{2-x}^{5+}O_{5-x}^{2-}O_x^{1-}}$\uff0cx \u5373\u6c27\u53c2\u4e0e\u7a0b\u5ea6\u3002DFT \u663e\u793a\uff1a\u5148 Mn\xb3\u207a\u2192Mn\u2074\u207a\uff0c\u989d\u5916\u8131\u9502\u5bb9\u91cf\u7531 O\xb2\u207b\u2192O\xb9\u207b \u627f\u62c5\uff0c\u4f4e\u9502\uff08x<1\uff09\u65f6\u624d\u53d1\u751f Mn\u2074\u207a\u2192Mn\u2075\u207a \u5e76\u4f34\u516b\u9762\u4f53\u2192\u56db\u9762\u4f53\u914d\u4f4d\u53d8\u5316\u3002[P0]
- Li-rich NMC\uff0c\u5e38\u7528\u5316\u5b66\u8ba1\u91cf Li\u2081.\u2082Mn\u2080.\u2085\u2084Co\u2080.\u2081\u2083Ni\u2080.\u2081\u2083O\u2082\u3002[P0]

**4. \u5982\u4f55\u5b9a\u91cf**
\u5bf9 O 2p \u6001\u5728\u8d39\u7c73\u80fd\u7ea7\u4e0b\u65b9\u7a84\u7a97\u53e3\uff08\u901a\u5e38 2.0~2.5 eV\uff09\u79ef\u5206\uff0c\u5f97\u5230\u53ef\u53c2\u4e0e redox \u7684\u201clabile \u6c27\u7535\u5b50\u201d\u6570\uff1b\u9700 HF \u6742\u5316\u6cdb\u51fd\u6216 GGA+U\uff083d \u6c27\u5316\u7269 GGA+U \u4ecd\u4e0d\u8db3\uff0c4d \u53ef\u7528\uff09\u3002[P0]

**5. \u4ee3\u4ef7\u4e0e\u6291\u5236**
\u6c27\u6c27\u5316\u82e5\u7f3a\u4e4f\u8db3\u591f M\u2013O \u5171\u4ef7\u6027\uff0c\u4f1a\u5f15\u8d77 TM \u6269\u6563\u4e0e\u9897\u7c92\u8868\u9762 O\u2082 \u91ca\u653e\uff1b\u9996\u5faa\u73af\u540e\u4e0d\u53ef\u9006\u53cd\u5e94\u201c\u8870\u51cf\u6b63\u6781\u7ed3\u6784\u4e0e\u5bb9\u91cf\u201d\uff0c\u5e76\u4f34\u7535\u538b\u8870\u51cf\u3001\u5c42\u72b6\u2192\u5c16\u6676\u77f3\u8f6c\u53d8\u3001\u9996\u5faa\u73af\u5bb9\u91cf\u635f\u5931\u4e0e\u53ef\u71c3\u98ce\u9669\u3002\u8868\u9762\u6c1f\u5316\u901a\u8fc7\u5f3a\u5316 TM\u2013O \u952e\u3001\u964d\u4f4e Jahn-Teller \u7578\u53d8\u63d0\u9ad8\u6c27\u7a7a\u4f4d\u5f62\u6210\u80fd\uff0c\u4ece\u800c\u6291\u5236\u6790\u6c27\uff0c\u4f46\u9700\u5728\u201c\u521d\u59cb\u5bb9\u91cf\u635f\u5931\u4e0e\u957f\u671f\u5bb9\u91cf\u4fdd\u6301\u201d\u95f4\u53d6\u6700\u4f18\u70b9\u3002[P0]

\u51fa\u5904\uff1a`materials__2408.09942__oxygen-redox-in-battery-cathodes-a-brief-overvie.md (part 1 of 2)`\uff08Assadi & Hanaor, *Oxygen redox in battery cathodes: A brief overview*\uff09\uff0c\u5e93 8fa4db59\uff08\u5de5\u7a0b\u4e0e\u80fd\u6e90\uff09\u3002\u68c0\u7d22\u8986\u76d6 5 \u4e2a pinned \u5e93\uff0c\u4ec5\u8be5\u6587\u6863\u547d\u4e2d\u76f8\u5173\u8bc1\u636e\uff08top score 0.77\uff09\uff0c\u5176\u4f59\u4e3a\u65e0\u5173\u8bed\u6599\u3002
\end{AnswerText}

\textbf{Retrieved evidence index (all returned units; paths and scores) }
\begin{EvidenceText}
[No source list recorded for this arm.]
\end{EvidenceText}
\subsection{BQ09: What problem of visual SLAM in dynamic environments does MLP-SLAM address and how?}
\begin{tabular}{@{}ll@{}}
Arm & HMRS agent (pinned)\\
Generator & omp / deepseek-flash\\
HTTP / success & 200 / True\\
Wall time / cost & 122.1 s / 0.058512 USD\\
Judge scores & F=5, C=5, G=5, A=5; overall=5.00\\
Verdict & correct\\
\end{tabular}

\textbf{Dimension rationale}
\begin{EvidenceText}
{
  "factual": "Gold-verbatim quotes throughout ('their precision tends to deteriorate when deployed in dynamic environments', 'leading to information loss', 'often require manual parameter tuning, lack generality'); the parked-car information-loss example; the full four-stage pipeline with YOLOv7 + Deep Sort, the DFM interval rule with eta = 1.2, the 3-D input vector with zero-mean normalization and argmax output; the >50,000-point 1:1 dataset with 70/20/10 split and training configuration; and the results table (Accuracy 87.71/F1 87.64 vs SVM 84.19/84.58; ATE 1.48 vs 1.99/2.43/1.77; KITTI360 3.42; tracking latency 62.81 ms vs ORB-SLAM2's 45.31 with MLP inference 0.18 ms; ablation 2.67/1.97/1.48; sequence 00: 0.79 vs DynaSLAM 1.40) \u2014 all internally consistent and corroborated by peer answers where they overlap.",
  "completeness": "Conclusion line plus problem analysis, four-stage method, dataset, and a full results table with ablation \u2014 the most complete quantitative treatment in the packet.",
  "groundedness": "Labeled P0 (direct quotes and table values); every number cross-checks with independent answers citing the same paper.",
  "sources": "File-level provenance: arXiv id, KB id, both part files with their section contents mapped, negative results on the other four libraries, and the duplicate-chunk KB noted but not double-counted."
}
\end{EvidenceText}

\textbf{Judge summary}
\begin{EvidenceText}
Exhaustive Chinese-language answer with verbatim quotes, complete pipeline, and full quantitative tables, all P0-sourced to the exact document.
\end{EvidenceText}

\textbf{Complete answer}
\begin{AnswerText}
**\u7ed3\u8bba\uff1aMLP-SLAM \u89e3\u51b3\u7684\u662f\u52a8\u6001\u73af\u5883\u4e0b V-SLAM \u7684\u7cbe\u5ea6\u9000\u5316\u95ee\u9898\uff0c\u5c24\u5176\u662f DynaSLAM \u4e00\u7c7b\u201c\u628a\u6574\u4e2a\u5019\u9009\u52a8\u6001\u76ee\u6807\u5305\u56f4\u76d2\u5185\u7684\u7279\u5f81\u5168\u90e8\u4e22\u5f03\u201d\u5bfc\u81f4\u7684\u9759\u6001\u4fe1\u606f\u635f\u5931\u3002\u505a\u6cd5\u662f\u628a\u201c\u52a8\u6001/\u9759\u6001\u7279\u5f81\u70b9\u5224\u522b\u201d\u53d8\u6210\u4e00\u4e2a\u76d1\u7763\u5b66\u4e60\u5206\u7c7b\u95ee\u9898\uff1a\u4ee5\u56fe\u50cf\u5f3a\u5ea6\u8bef\u5dee $e_I$\u3001\u5bf9\u6781\u8bef\u5dee $e_D$\u3001\u91cd\u6295\u5f71\u8bef\u5dee $e_{Re}$ \u6784\u6210 3 \u7ef4\u8f93\u5165\u7279\u5f81\uff0c\u7528 MLP \u5224\u522b\u5668\u66ff\u4ee3\u624b\u5de5\u9608\u503c\uff0c\u907f\u514d\u4fe1\u606f\u635f\u5931\u4e0e\u4eba\u5de5\u8c03\u53c2\u3002**

## 1. \u88ab\u6307\u51fa\u7684\u95ee\u9898

- \u5176\u7acb\u8db3\u70b9\u662f\u591a\u6570 SLAM \u7cfb\u7edf\uff08ORB-SLAM2/3\u3001LSD-SLAM\u3001DSO\uff09**\u4f9d\u8d56\u9759\u6001\u4e16\u754c\u5047\u8bbe**\uff0c\u201ctheir precision tends to deteriorate when deployed in dynamic environments\u201d\uff08Introduction\uff09\u3002
- \u5df2\u6709\u65b9\u6848\u4e09\u7c7b\u7f3a\u9677\uff1aType I \u7eaf\u51e0\u4f55\uff1bType II \u5982 DynaSLAM \u7528\u8bed\u4e49\u63a9\u819c\u628a\u6f5c\u5728\u52a8\u6001\u7c7b\u76ee\u6807\u5185\u7279\u5f81\u5168\u90e8\u5f53\u5916\u70b9\u4e22\u5f03\uff0c**\u8def\u65c1\u505c\u653e\u7684\u9759\u6b62\u8f66\u8f86\u88ab\u8bef\u5220**\uff0c\u201cleading to information loss\u201d\uff1bType III \u4e3b\u6d41\uff08CFP-SLAM\u3001PointSLOT\u3001VDO-SLAM\u3001D-SLAM\uff09\u53ea\u5229\u7528**\u90e8\u5206\u51e0\u4f55\u4fe1\u606f**\uff08\u4ec5\u5bf9\u6781\u8bef\u5dee\u6216\u4ec5\u91cd\u6295\u5f71\u8bef\u5dee\uff09\uff0c\u4e14\u201coften require manual parameter tuning, lack generality\u201d\uff08Introduction / Related Work\uff09\u3002
- \u6b21\u8981\u7f3a\u53e3\uff1a\u7f3a\u4e4f\u76f4\u63a5\u8bc4\u4f30\u52a8\u6001/\u9759\u6001\u5224\u522b\u7684\u516c\u5f00\u6570\u636e\u96c6\uff0c\u53ea\u80fd\u501f ATE \u7b49\u95f4\u63a5\u6307\u6807\uff0c\u53d7 RANSAC \u968f\u673a\u6027\u5f71\u54cd\u3002

## 2. \u65b9\u6cd5\uff08Section III-D / IV\uff09

\u4ee5 ORB-SLAM2 \u4e3a\u5e95\u5ea7\uff0c\u52a0\u4e00\u4e2a\u52a8\u6001-\u9759\u6001\u5224\u522b\u6a21\u5757\uff0c\u6d41\u7a0b\u56db\u6bb5\uff1a

1. **\u9884\u5904\u7406**\uff1a\u53cc\u76ee ORB \u7279\u5f81\u63d0\u53d6\uff1b\u5de6\u56fe\u7528 **YOLOv7 + Deep Sort** \u505a\u68c0\u6d4b\u4e0e\u8ddf\u8e2a\uff08\u6bd4\u8bed\u4e49\u5206\u5272\u5ef6\u8fdf\u4f4e\uff09\u3002
2. **\u6df1\u5ea6\u6ee4\u6ce2\u6a21\u5757 DFM**\uff1a\u56e0\u68c0\u6d4b\u6846\u4e0d\u5982\u5206\u5272\u63a9\u819c\u7cbe\u786e\uff0c\u6846\u5185\u6309\u5e73\u5747\u6df1\u5ea6 $\mu_d$ \u4e0e\u6807\u51c6\u5dee $\delta_d$ \u5254\u9664 $[\mu_d-\eta\delta_d,\ \mu_d+\eta\delta_d]$ \u4e4b\u5916\u7684\u70b9\u5e76\u5224\u4e3a\u9759\u6001\uff08$\eta=1.2$\uff09\u3002
3. **\u7c97\u4f30\u8ba1**\uff1a\u6846\u5916\u9759\u6001\u70b9 + DFM \u5224\u5b9a\u7684\u6846\u5185\u80cc\u666f\u70b9\uff0c\u6700\u5c0f\u5316\u91cd\u6295\u5f71\u8bef\u5dee\u6c42\u7c97\u4f4d\u59ff\u3002
4. **MLP \u5224\u522b + \u7cbe\u4f30\u8ba1**\uff1a\u6846\u5185\u672a\u88ab DFM \u8fc7\u6ee4\u7684\u70b9\u751f\u6210 $v_i=\{e_I,e_D,e_{Re}\}$\uff0c\u96f6\u5747\u503c\u5f52\u4e00\u5316\u540e\u6210\u6279\u9001 MLP \u8f93\u51fa $c_i=\arg\max M(v_i)$\uff080 \u9759\u6001 / 1 \u52a8\u6001\uff09\uff1b\u518d\u7528\u201c\u6846\u5916\u70b9 + DFM \u9759\u6001\u70b9 + MLP \u5224\u4e3a\u9759\u6001\u70b9\u201d\u91cd\u4f18\u5316\u4f4d\u59ff\u3002

\u8bba\u6587\u5f3a\u8c03\u8fd9\u662f**\u9996\u6b21\u5c06\u673a\u5668\u5b66\u4e60\u65b9\u6cd5\u7528\u4e8e SLAM \u4e2d\u52a8\u6001/\u9759\u6001\u7279\u5f81\u70b9\u5224\u522b**\u3002\u53e6\u5efa\u4e86 5 \u4e07\u4f59\u70b9\u3001\u52a8\u9759 1:1\u3001\u4eba\u5de5\u6807\u6ce8\u7684\u516c\u5f00\u6570\u636e\u96c6\uff08KITTI 00/01\uff0c70/20/10 \u5212\u5206\uff09\uff1b\u8bad\u7ec3 150 epoch\u3001batch 8\u3001lr 0.1\u3001MultiStepLR(\u03b3=0.1)\u3001Adam\uff0cTensorRT \u56fa\u5b9a\u8f93\u5165 (200,3)\u3002

## 3. \u6548\u679c\uff08P0 \u76f4\u63a5\u8bc1\u636e\uff09

| \u6307\u6807 | MLP-SLAM | \u5bf9\u6bd4 |
|---|---|---|
| \u5224\u522b Accuracy / F1 | **87.71 / 87.64** | SVM 84.19/84.58\uff1b\u7c7b PointSLOT 81.36\uff1b\u7c7b CFP-SLAM 73.61\uff1b\u4ec5 $e_D$ 60.10 |
| KITTI Odometry \u5e73\u5747 ATE | **1.48** | ORB-SLAM2 1.99\u3001DynaSLAM 2.43\u3001DynaSLAM II 1.77 |
| KITTI360 \u5e73\u5747 ATE | **3.42** | ORB-SLAM2 3.92\u3001VDO-SLAM(\u6539) 8.83 |
| \u8ddf\u8e2a\u5ef6\u8fdf | **62.81 ms**\uff08\u6bd4 ORB-SLAM2 \u4ec5 +17.50 ms/\u5e27\uff0cMLP \u63a8\u7406 0.18 ms\uff09 | DynaSLAM 519.37\u3001DynaSLAM II 80.15\u3001PointSLOT 73.6 |

\u6d88\u878d\uff1a\u4ec5\u76ee\u6807\u8ddf\u8e2a 2.67 \u2192 \u4ec5 MLP 1.97 \u2192 \u4e24\u8005\u7ed3\u5408 **1.48**\uff08\u5e73\u5747 ATE\uff09\uff0c\u8bf4\u660e MLP \u5224\u522b\u662f\u4e3b\u8981\u8d21\u732e\u3002\u5bf9\u505c\u653e\u8f66\u8f86\u5bc6\u96c6\u7684\u5e8f\u5217 00 \u53d6\u5f97\u6700\u4f18\uff080.79 vs DynaSLAM 1.40\uff09\uff0c\u9a8c\u8bc1\u201c\u80fd\u8bc6\u522b\u771f\u6b63\u7684\u52a8\u6001\u7269\u4f53\u201d\u800c\u975e\u6574\u7c7b\u4e22\u5f03\u3002

## \u51fa\u5904

- \u4e3b\u6e90\uff1a\u300aMLP-SLAM: Multilayer Perceptron-Based Simultaneous Localization and Mapping\u300b\uff08Taozhe Li, Wei Sun\uff1barXiv 2410.10669\uff09\uff0cKB\u300c\u8ba1\u7b97\u673a\u4e0e\u4eba\u5de5\u667a\u80fd\u300dKB `b1199132\u2026b0ba`\uff0cdoc `robotics__2410.10669__mlp-slam\u2026(part 1/2).md`\u3001`(part 2/2).md`\u3002\u7b2c 1 \u90e8\u5206\u542b Introduction/Related Work/III OverView/IV \u65b9\u6cd5/Table I\uff1b\u7b2c 2 \u90e8\u5206\u542b\u5b9e\u9a8c\uff08Table II\u2013VI\uff09\u4e0e\u7ed3\u8bba\u3002
- \u5176\u4f59 4 \u4e2a\u6307\u5b9a\u5e93\uff08\u5730\u9707\u5b66\u3001\u7535\u529b\u7cfb\u7edf\u3001\u533b\u5b66\u5f71\u50cf\u3001\u7ecf\u6d4e\uff09\u65e0\u76f8\u5173\u547d\u4e2d\uff1b`Corpus-Chunks800` \u4e2d\u7684\u540c\u6587\u91cd\u590d\u5206\u5757\uff08\u53e6\u4e00 KB `ce642f28\u2026`\uff09\u5185\u5bb9\u4e00\u81f4\uff0c\u672a\u5355\u72ec\u5f15\u7528\u3002
- \u53ef\u4fe1\u5ea6\uff1a\u4ee5\u4e0a\u5747\u4e3a P0\uff08\u6b63\u6587\u539f\u6587\u76f4\u5f15 + \u539f\u6587\u8868\u683c\u6570\u503c\uff09\u3002
\end{AnswerText}

\textbf{Retrieved evidence index (all returned units; paths and scores) }
\begin{EvidenceText}
[No source list recorded for this arm.]
\end{EvidenceText}
\subsection{BQ10: What strategies for tumor control emerge under immune evasion and chemotherapy resistance?}
\begin{tabular}{@{}ll@{}}
Arm & HMRS agent (pinned)\\
Generator & omp / deepseek-flash\\
HTTP / success & 200 / True\\
Wall time / cost & 82.7 s / 0.1014882 USD\\
Judge scores & F=5, C=5, G=5, A=5; overall=5.00\\
Verdict & correct\\
\end{tabular}

\textbf{Dimension rationale}
\begin{EvidenceText}
{
  "factual": "All correct: combination superiority including 'combination immuno-chemotherapy provides greater flexibility and improved tumor control compared to single-therapy approaches'; checkpoint evasion (PD-L1) being 'highly sensitive to the immune exhaustion rate' with checkpoint inhibitors plus moderate chemotherapy; reduced antigen presentation being 'more robust against relatively weak chemotherapy' and addressed by checkpoint inhibitors, CAR-T, and vaccines; immune-sensitive dominance handled by raising sigma and lowering delta with adoptive T-cell/cytotoxic therapy; 'non-personalized, uniform treatment often fails'; the three sequential protocols with the T1-resurgence caveat; MTD yielding to adaptive therapy; baseline parameters (chi = 0.818, zeta = 0.4908); and the gold-echoing 'modest parameter changes' tipping quote (preserving the source's typo).",
  "completeness": "Four numbered strategy layers (combination, mechanism-matched drug choice, sequencing, adaptive control with risks) plus future directions \u2014 fully answers the question.",
  "groundedness": "Labeled P0 with quotes from sections 3.3 and 4; content matches gold and cross-consistent peer answers.",
  "sources": "File-level provenance: arXiv id, KB id, three-part document structure, section attribution, plus transparent retrieval notes that only one of the five libraries hit relevant content."
}
\end{EvidenceText}

\textbf{Judge summary}
\begin{EvidenceText}
Complete Chinese-language answer with P0-quoted combination, mechanism-matched, and sequential strategies plus adaptive-therapy framing and file-level sourcing.
\end{EvidenceText}

\textbf{Complete answer}
\begin{AnswerText}
**\u7ed3\u8bba\uff1a\u53ef\u884c\u6027\u6700\u9ad8\u7684\u65b9\u6848\u662f\u5148\u6309\u6cbb\u7597\u524d\u80bf\u7624\u7684\u7ec4\u6210\u548c\u514d\u75ab\u9003\u9038\u673a\u5236\u5224\u5b9a\u4f18\u52bf\u8868\u578b\uff0c\u518d\u636e\u6b64\u9009\u62e9\u7528\u836f\u5e76\u5b89\u6392\u65f6\u5e8f\u2014\u2014\u4ee5\u201c\u514d\u75ab\u6cbb\u7597 + \u5316\u7597\u201d\u7684\u7ec4\u5408\u4e3a\u4e3b\u7ebf\uff0c\u4f18\u5148\u628a\u7cfb\u7edf\u63a8\u5165\u65e0\u7624\u5e73\u8861\uff08TFE\uff09\u6216\u4f4e\u8f7d\u91cf\u5e73\u8861\uff0c\u800c\u4e0d\u662f\u7528\u5355\u4e00\u7597\u6cd5\u8ffd\u6c42\u6700\u5927\u5242\u91cf\u706d\u6740\u3002** \u4e94\u4e2a\u6307\u5b9a\u5e93\u4e2d\uff0c\u771f\u6b63\u547d\u4e2d\u4e3b\u9898\u7684\u53ea\u6709 1 \u7bc7\u6587\u732e\uff08\u5176\u4f59 4 \u5e93\u547d\u4e2d\u7684\u90fd\u662f\u65e0\u5173\u5185\u5bb9\uff0c\u539f\u56e0\u662f\u8de8\u5e93\u8865\u5145\u68c0\u7d22\uff1b\u8fd9\u4e00\u70b9\u5df2\u5982\u5b9e\u6807\u6ce8\uff09\uff1a\u300aStrategies for tumor elimination and control under immune evasion and chemotherapy resistance\u300b(arXiv 2604.01385, Mokari & Morsky, FSU)\uff0c\u89c1 KB\u300c\u751f\u547d\u79d1\u5b66\u4e0e\u533b\u5b66\u300d`b28f0a25\u2026`\uff0c\u5206\u4e3a 1/3\u20133/3 \u4e09\u90e8\u5206\uff0c\u4e0b\u6587\u5f15\u6587\u90fd\u51fa\u81ea \xa73.3 \u4e0e \xa74\uff08P0\uff0c\u76f4\u63a5\u8bc1\u636e\uff09\u3002

**1. \u7ec4\u5408\u4f18\u4e8e\u5355\u836f\uff0c\u533a\u5206\u5ea6\u6765\u81ea\u53c2\u6570\u533a\u95f4**
\u201cimmune dynamics alone are generally insufficient to fully eliminate the tumor\u201d\uff0c\u800c\u201cchemotherapy significantly enlarges the parameter region corresponding to tumor control or tumor-free equilibria\u201d\u3002\u5316\u7597\u7f29\u5c0f\u80bf\u7624\u8d1f\u8377\uff0c\u5e76\u76f4\u63a5\u4f5c\u7528\u4e8e\u8010\u836f\u4e9a\u7fa4\uff1b\u514d\u75ab\u6cbb\u7597\u5219\u589e\u5f3a\u514d\u75ab\u4ecb\u5bfc\u7684\u6291\u5236\u3002\u5408\u7528\u65f6\uff0c\u65e0\u7624\u5e73\u8861\u6216\u4f4e\u8f7d\u91cf\u5e73\u8861\u7684\u533a\u95f4\u90fd\u6269\u5927\uff1b\u201ccombination immuno-chemotherapy provides greater flexibility and improved tumor control compared to single-therapy approaches\u201d\uff08P0\uff09\u3002

**2. \u6309\u514d\u75ab\u9003\u9038\u673a\u5236\u9009\u62e9\u514d\u75ab\u6cbb\u7597\u836f\u7269\uff08\u03c7\u3001\u03b4\u3001\u03c3\u3001\u03b8\u2081\u3001\u03bc\u3001\u03c1 \u5206\u5c94\u5206\u6790\uff09**
- **\u514d\u75ab\u68c0\u67e5\u70b9\u9003\u9038\uff08PD-L1\u2191\uff09**\uff1a\u63a7\u5236\u6548\u679c\u201chighly sensitive to the immune exhaustion rate\u201d\u2192 \u6700\u4f18\u7b56\u7565\u662f**\u68c0\u67e5\u70b9\u6291\u5236\u5242 + \u4e2d\u7b49\u5f3a\u5ea6\u5316\u7597**\uff0c\u4f5c\u7528\u9776\u70b9\u662f\u964d\u4f4e\u514d\u75ab\u8017\u7aed \u03bc\u3001\u63d0\u9ad8\u6fc0\u6d3b \u03c1 \u548c\u62db\u52df \u03c3\u3002
- **\u6297\u539f\u5448\u9012\u4e0b\u8c03**\uff1a\u5bf9\u5f31\u5316\u7597\u66f4\u8010\u53d7\uff08\u201creduced antigen presentation is more robust against relatively weak chemotherapy\u201d\uff09\uff0c\u7a33\u5b9a\u6027\u66f4\u4f9d\u8d56\u6548\u5e94\u7ec6\u80de\u6fc0\u6d3b\u4e0e\u8017\u7aed \u2192 **\u68c0\u67e5\u70b9\u6291\u5236\u5242\u3001CAR-T\u3001\u80bf\u7624\u75ab\u82d7**\u4e0e\u5316\u7597\u8054\u7528\u3002
- **\u514d\u75ab\u654f\u611f\u8868\u578b\u5360\u4f18**\uff1a\u5173\u952e\u5728\u63d0\u9ad8\u6548\u5e94\u7ec6\u80de\u7684\u62db\u52df\u4e0e\u5b58\u6d3b\uff08\u03c3\u2191\u3001\u03b4\u2193\uff09\u2192 **\u8fc7\u7ee7\u6027 T \u7ec6\u80de\u8f93\u6ce8\u6216\u7ec6\u80de\u6bd2\u6cbb\u7597**\u3002
\u4f5c\u8005\u7531\u6b64\u65ad\u8a00\uff1a\u201cnon-personalized, uniform treatment often fails\u201d\uff08P0\uff09\u3002

**3. \u6309\u80bf\u7624\u7ec4\u6210\u5b89\u6392\u65f6\u5e8f\uff08\xa74 \u660e\u786e\u63d0\u51fa\u7684 sequential strategies\uff09**
- T\u2081 \u5360\u4f18\uff08\u514d\u75ab\u8010\u836f\uff09\uff1a**\u5148\u5316\u7597\u51cf\u8d1f \u2192 \u514d\u75ab\u6cbb\u7597\u7ef4\u6301**\uff1b\u6216\u53cd\u5411\u5148\u514d\u75ab\u7f29\u5c0f\u3001\u518d\u5316\u7597\u6e05\u9664\u6b8b\u4f59\u8010\u836f\u7ec6\u80de\u3002
- T\u2082 \u5360\u4f18\uff08\u5316\u7597\u8010\u836f\uff09\uff1a\u5316\u7597\u867d\u51cf\u8d1f\u201ccan also lead to resistant escape\u201d\u2192 **\u5fc5\u987b\u5728\u5316\u7597\u540e\u63a5\u514d\u75ab\u6cbb\u7597\u7a33\u5b9a\u7cfb\u7edf**\u3002
- \u4e24\u7c7b\u5171\u5b58\uff1a**\u5148\u5316\u7597\u6e05\u9664 T\u2081**\uff0c\u843d\u5230\u4f4e\u8f7d\u91cf T\u2082-only \u5e73\u8861\uff0c\u518d\u514d\u75ab\u6cbb\u7597\u6709\u671b\u6e05\u9664\uff1b\u82e5\u4ec5\u88ab\u63a7\u5236\uff0c\u201cthe emergence of T\u2081 mutants could result in resurgence\u201d\u3002

**4. \u63a7\u5236\u903b\u8f91\u4e0e\u98ce\u9669**
\u4f20\u7edf MTD \u601d\u8def\u5e94\u8ba9\u4f4d\u4e8e adaptive therapy\uff08\u201ccontrol rather than eliminate\u201d\uff09\uff0c\u5229\u7528\u654f\u611f\u4e0e\u8010\u836f\u4e9a\u7fa4\u4e4b\u95f4\u7684\u7ade\u4e89\u3002\u5206\u5c94\u663e\u793a\u5b58\u5728\u53cc\u7a33\u533a\u4e0e\u9608\u503c\uff08\u57fa\u7ebf \u03c7=0.818\u3001\u03b6=0.4908\uff1b\u03b8\u2082 \u5c0f=\u8010\u836f\u6536\u76ca\u9ad8\u2192T\u2082 \u4e3b\u5bfc\uff09\uff0c**\u201cmodest parameter changes can induced qualitative transitions between immune control and tumor escape\u201d**\uff0c\u5373\u5c0f\u5e45\u53c2\u6570\u6f02\u79fb\u53ef\u8ba9\u63a7\u5236\u7ffb\u8f6c\u4e3a\u9003\u9038\uff0c\u6545\u9700\u4ee5\u751f\u7269\u6807\u5fd7\u7269\u4e0e\u80bf\u7624\u7ec4\u6210\u52a8\u6001\u6307\u5bfc\u6362\u836f/\u505c\u836f\uff08P0\uff09\u3002\u672a\u6765\u65b9\u5411\uff1a\u6700\u4f18\u63a7\u5236\u6392\u7a0b\u3001\u6eb6\u7624\u75c5\u6bd2\u3001\u673a\u5668\u5b66\u4e60\u4e2a\u4f53\u5316\u81ea\u9002\u5e94\u6cbb\u7597\u3002

**\u68c0\u7d22\u4e8b\u5b9e**\uff1a`kb_search_two_stage` \u4e24\u8f6e\uff08tumor control/immune evasion/chemo resistance\uff1bovercoming resistance + microenvironment + combination\uff09\uff0c5 \u4e2a\u6307\u5b9a\u5e93\u5171 25 \u6761\u5019\u9009\uff0c\u53bb\u91cd\u540e\u4ec5\u8be5\u6587\u732e 3 \u4e2a\u5206\u7247\u76f8\u5173\uff08stage2 \u5206\u6570 0.70\u20130.82\uff09\uff1b\u5176\u4f59 4 \u4e2a\u6307\u5b9a\u5e93\uff08\u8ba1\u7b97\u673a\u3001\u91cf\u5b50\u7269\u7406\u3001\u5de5\u7a0b\u80fd\u6e90\u3001\u7ecf\u6d4e\uff09\u5747\u65e0\u533b\u5b66\u76f8\u5173\u547d\u4e2d\u3002
\end{AnswerText}

\textbf{Retrieved evidence index (all returned units; paths and scores) }
\begin{EvidenceText}
[No source list recorded for this arm.]
\end{EvidenceText}
\section{BM25 baseline}
\subsection{BQ01: How does the Transformer compute attention, and why does it replace recurrence and convolution?}
\begin{tabular}{@{}ll@{}}
Arm & BM25 baseline\\
Generator & omp / deepseek-flash\\
HTTP / success & 200 / True\\
Wall time / cost & 41.7 s / 0.0766788 USD\\
Judge scores & F=5, C=3, G=5, A=5; overall=4.50\\
Verdict & partial\\
\end{tabular}

\textbf{Dimension rationale}
\begin{EvidenceText}
{
  "factual": "Every Transformer claim matches the GOLD excerpt: self-attention/'intra-attention' definition, 'averaging attention-weighted positions' reducing effective resolution counteracted by Multi-Head Attention, ConvS2S linear / ByteNet logarithmic vs constant operations, 'eschewing recurrence', 28.4 BLEU EN-DE (over 2 BLEU above prior best including ensembles), 41.8 BLEU EN-FR after 3.5 days on eight GPUs, twelve hours on eight P100 GPUs. No conflict with gold anywhere.",
  "completeness": "The 'why replace' half is answered fully in three numbered arguments (constant path length, parallelization, training cost). The 'how is attention computed' half is explicitly declined: 'the context does not contain it' \u2014 the scaled dot-product form and multi-head splitting are named as absent, with only indirect hints (averaging, global dependencies). Honest refusal, so roughly half the question is covered.",
  "groundedness": "All claims are tied to the gold text (e.g. 'global dependencies', 'dispensing with recurrence and convolutions entirely'); no invented mechanism details; it explicitly flags that S3/S4/S5/S7 discuss other topics rather than borrowing from them.",
  "sources": "Per-claim attributions to [S10], [S2], [S1], [S6] with what each supports stated, plus a closing note on which sources are irrelevant. Clear and accurate."
}
\end{EvidenceText}

\textbf{Judge summary}
\begin{EvidenceText}
Honestly reports the attention formula is missing from its context and fully answers only the 'why replace recurrence/convolution' half, with every claim matching the gold.
\end{EvidenceText}

\textbf{Complete answer}
\begin{AnswerText}
**Attention computation: the context does not contain it.** The provided excerpts are only part 1 of 2 of *Attention Is All You Need* (abstract, introduction, background). The actual mechanism \u2014 the scaled dot-product form, the $\mathrm{softmax}(QK^\top/\sqrt{d_k})V$ equations, and the multi-head splitting/concatenation \u2014 sits in Section 3, which is absent here. The only mechanical hints present are:

- Attention relates different positions of a single sequence to compute a representation of that sequence \u2014 self-attention, or "intra-attention" [S10].
- It produces a representation by **averaging attention-weighted positions**, an averaging that reduces effective resolution and is counteracted by Multi-Head Attention (Section 3.2, not included) [S10].
- It draws **global dependencies** between input and output regardless of distance in the sequence [S2].
- The Transformer is the first transduction model relying *entirely* on self-attention to compute input and output representations, with no sequence-aligned RNNs or convolution [S1][S6].

**Why it replaces recurrence and convolution \u2014 this the context does answer:**

1. **Constant path length between positions.** In ConvS2S the number of operations needed to relate two arbitrary positions grows *linearly* with distance, and in ByteNet *logarithmically*; this makes distant dependencies harder to learn. In the Transformer it is reduced to a **constant** number of operations \u2014 paid for by the resolution loss that multi-head attention offsets [S10].
2. **Parallelization.** Computing hidden representations for *all* input and output positions happens in parallel, unlike sequential recurrence [S2][S10]. Prior attention mechanisms were almost always bolted onto a recurrent network; the Transformer "eschews recurrence" and relies entirely on attention [S2].
3. **Training cost.** The architecture is "based solely on attention mechanisms, dispensing with recurrence and convolutions entirely," and is more parallelizable and much faster to train: 28.4 BLEU EN\u2013DE (over 2 BLEU above the prior best, including ensembles) and 41.8 BLEU EN\u2013FR as a single model after 3.5 days on eight GPUs; the base model reached SOTA in as little as twelve hours on eight P100 GPUs [S6][S2].

Note: S3, S4, S5, S7 discuss attention in other contexts (long-context variants, recall circuits, attention patching) but not the Transformer's attention equations or its motivation.
\end{AnswerText}

\textbf{Retrieved evidence index (all returned units; paths and scores) }
\begin{EvidenceText}
[1] score=30.13621249349545 path=\u8ba1\u7b97\u673a\u4e0e\u4eba\u5de5\u667a\u80fd\artificial-intelligence__1706.03762__attention-is-all-you-need.md (part 1 of 2).md
[2] score=19.728257596864477 path=\u8ba1\u7b97\u673a\u4e0e\u4eba\u5de5\u667a\u80fd\artificial-intelligence__1706.03762__attention-is-all-you-need.md (part 1 of 2).md
[3] score=19.190292047283183 path=\u8ba1\u7b97\u673a\u4e0e\u4eba\u5de5\u667a\u80fd\Lost in the Middle (part 2 of 3).md
[4] score=19.021079184736546 path=\u8ba1\u7b97\u673a\u4e0e\u4eba\u5de5\u667a\u80fd\nlp__2503.21676__how-do-language-models-learn-facts-dynamics.md (part 1 of 4).md
[5] score=18.185711987792175 path=\u8ba1\u7b97\u673a\u4e0e\u4eba\u5de5\u667a\u80fd\nlp__2503.21676__how-do-language-models-learn-facts-dynamics.md (part 1 of 4).md
[6] score=18.119504420281313 path=\u8ba1\u7b97\u673a\u4e0e\u4eba\u5de5\u667a\u80fd\artificial-intelligence__1706.03762__attention-is-all-you-need.md (part 1 of 2).md
[7] score=17.58193867368475 path=\u8ba1\u7b97\u673a\u4e0e\u4eba\u5de5\u667a\u80fd\nlp__2503.21676__how-do-language-models-learn-facts-dynamics.md (part 2 of 4).md
[8] score=17.32819837159615 path=\u8ba1\u7b97\u673a\u4e0e\u4eba\u5de5\u667a\u80fd\nlp__1810.04805__bert-pre-training-of-deep-bidirectional-transfor.md (part 2 of 3).md
[9] score=17.060270156989258 path=\u81ea\u7136\u79d1\u5b66\u4e0e\u5730\u7403\u79d1\u5b66\quantum-physics__1801.00862__quantum-computing-in-the-nisq-era-and-beyond.md (part 1 of 3).md
[10] score=16.818573817572688 path=\u8ba1\u7b97\u673a\u4e0e\u4eba\u5de5\u667a\u80fd\artificial-intelligence__1706.03762__attention-is-all-you-need.md (part 1 of 2).md
[11] score=16.38215397387708 path=\u8ba1\u7b97\u673a\u4e0e\u4eba\u5de5\u667a\u80fd\databases__2102.02440__online-sketch-based-query-optimization.md (part 2 of 5).md
[12] score=16.30945350168588 path=\u8ba1\u7b97\u673a\u4e0e\u4eba\u5de5\u667a\u80fd\databases__2102.02440__online-sketch-based-query-optimization.md (part 3 of 5).md
\end{EvidenceText}
\subsection{BQ02: What defines NISQ technology and what is its central limitation?}
\begin{tabular}{@{}ll@{}}
Arm & BM25 baseline\\
Generator & omp / deepseek-flash\\
HTTP / success & 200 / True\\
Wall time / cost & 41.0 s / 0.021198 USD\\
Judge scores & F=5, C=5, G=5, A=5; overall=5.00\\
Verdict & correct\\
\end{tabular}

\textbf{Dimension rationale}
\begin{EvidenceText}
{
  "factual": "All claims match gold and the paper: NISQ = noisy intermediate-scale quantum; 50 to a few hundred qubits near future; 50 qubits beyond brute-force supercomputer simulation; noisy = imperfect control; the abstract's 'noise in quantum gates will limit the size of quantum circuits that can be executed reliably'; the ~1000 two-qubit gate ceiling and 'imposes a ceiling on the computational power'; QEC's heavy overhead; hard-to-validate outputs. Nothing contradicted.",
  "completeness": "Both parts fully answered: definition (both halves of the acronym unpacked, with the abstract's 50-100 qubit framing) and central limitation (noise capping circuit size at ~1000 gates, capping computational power), plus two flagged consequences (QEC overhead, verification) and a benchmarking-paper corroboration.",
  "groundedness": "Each claim is tied to a cited source with quotes that match the paper verbatim ('that limitation on circuit size imposes a ceiling on the computational power of NISQ technology'); it also honestly notes what is absent ('any claim that NISQ will deliver near-term advantage... is unknown').",
  "sources": "Per-claim attributions ([S2][S4][S8][S6][S12][S1]) each described, including what the sources do and do not claim."
}
\end{EvidenceText}

\textbf{Judge summary}
\begin{EvidenceText}
Complete and accurate answer to both halves with verbatim-supported 1000-gate ceiling and clear per-claim sourcing.
\end{EvidenceText}

\textbf{Complete answer}
\begin{AnswerText}
**Definition.** NISQ = *noisy intermediate-scale quantum*: devices with roughly 50 to a few hundred qubits expected in the near future, where 50 qubits marks the threshold beyond brute-force simulation by the most powerful existing digital supercomputers; "noisy" signals imperfect control over the qubits [S2]. Preskill's abstract states the same: 50\u2013100 qubit machines may surpass classical digital computers on some tasks [S4].

**Central limitation.** Noise in quantum gates. It caps the size of circuits that can be executed reliably: naively, no more than about 1000 gates (fundamental two-qubit operations) before noise overwhelms the signal, and "that limitation on circuit size imposes a ceiling on the computational power of NISQ technology" [S8]. The abstract concurs that noise "will limit the size of quantum circuits that can be executed reliably" [S4]; noise "will place serious limitations on what quantum devices can achieve in the near term" [S2].

Two consequences worth flagging:
- Escaping the ceiling requires quantum error correction, which imposes heavy overhead in qubit and gate counts [S8].
- Imperfect performance makes output hard to validate \u2014 verification of results, especially quantum-simulation problems, is an open research need [S6].

A parallel characterization from a benchmarking paper lists the same obstacles \u2014 restrictive qubit topologies, short coherence times, imperfect gates \u2014 compounding to limit achievable results [S12].

Absent from this context: any claim that NISQ will deliver near-term advantage; the sources explicitly say that is unknown [S1].
\end{AnswerText}

\textbf{Retrieved evidence index (all returned units; paths and scores) }
\begin{EvidenceText}
[1] score=23.317518447884584 path=\u81ea\u7136\u79d1\u5b66\u4e0e\u5730\u7403\u79d1\u5b66\quantum-physics__1801.00862__quantum-computing-in-the-nisq-era-and-beyond.md (part 2 of 3).md
[2] score=22.152120106941986 path=\u81ea\u7136\u79d1\u5b66\u4e0e\u5730\u7403\u79d1\u5b66\quantum-physics__1801.00862__quantum-computing-in-the-nisq-era-and-beyond.md (part 1 of 3).md
[3] score=18.00095221826477 path=\u81ea\u7136\u79d1\u5b66\u4e0e\u5730\u7403\u79d1\u5b66\quantum-physics__1801.00862__quantum-computing-in-the-nisq-era-and-beyond.md (part 2 of 3).md
[4] score=17.008760307173624 path=\u81ea\u7136\u79d1\u5b66\u4e0e\u5730\u7403\u79d1\u5b66\quantum-physics__1801.00862__quantum-computing-in-the-nisq-era-and-beyond.md (part 1 of 3).md
[5] score=16.624529482846405 path=\u8ba1\u7b97\u673a\u4e0e\u4eba\u5de5\u667a\u80fd\large-language-models__2005.14165__language-models-are-few-shot-learners.md (part 4 of 13).md
[6] score=15.462256339435145 path=\u81ea\u7136\u79d1\u5b66\u4e0e\u5730\u7403\u79d1\u5b66\quantum-physics__1801.00862__quantum-computing-in-the-nisq-era-and-beyond.md (part 1 of 3).md
[7] score=15.450812648694402 path=\u81ea\u7136\u79d1\u5b66\u4e0e\u5730\u7403\u79d1\u5b66\quantum-physics__1801.00862__quantum-computing-in-the-nisq-era-and-beyond.md (part 1 of 3).md
[8] score=15.41781039169155 path=\u81ea\u7136\u79d1\u5b66\u4e0e\u5730\u7403\u79d1\u5b66\quantum-physics__1801.00862__quantum-computing-in-the-nisq-era-and-beyond.md (part 1 of 3).md
[9] score=15.353363415089682 path=\u8ba1\u7b97\u673a\u4e0e\u4eba\u5de5\u667a\u80fd\Chain-of-Thought Prompting (part 2 of 3).md
[10] score=14.777755309772447 path=\u7ecf\u6d4e\u4e0e\u793e\u4f1a\economics__2011.14424__on-the-effectiveness-of-the-european-central (part 2 of 5).md
[11] score=14.37198914525874 path=\u81ea\u7136\u79d1\u5b66\u4e0e\u5730\u7403\u79d1\u5b66\quantum-physics__1801.00862__quantum-computing-in-the-nisq-era-and-beyond.md (part 2 of 3).md
[12] score=14.341408277727243 path=\u81ea\u7136\u79d1\u5b66\u4e0e\u5730\u7403\u79d1\u5b66\quantum-physics__2303.00445__benchmarking-noisy-intermediate-scale-quantu (part 1 of 4).md
\end{EvidenceText}
\subsection{BQ03: What is MultiMedQA and how did Flan-PaLM perform on the MedQA benchmark?}
\begin{tabular}{@{}ll@{}}
Arm & BM25 baseline\\
Generator & omp / deepseek-flash\\
HTTP / success & 200 / True\\
Wall time / cost & 24.1 s / 0.025386 USD\\
Judge scores & F=5, C=5, G=5, A=5; overall=5.00\\
Verdict & correct\\
\end{tabular}

\textbf{Dimension rationale}
\begin{EvidenceText}
{
  "factual": "Matches gold and the paper: benchmark spanning exam/research/consumer questions; seven datasets (MedQA, MedMCQA, PubMedQA, LiveQA, MedicationQA, MMLU clinical topics + newly introduced HealthSearchQA with 3,375 questions); motivation quote about 'classification accuracy or automated NLG metrics (BLEU)' is gold-verbatim; 67.6% on 4-option MedQA, +20.1% over DRAGON, over 17% (17.3%) over prior best; 60.3% few-shot, 60.3% CoT, 67.6% self-consistency; 62.0% on 5-option; scaling pairs 25.7/35.4, 40.9/46.1, 58.9/60.3. All correct.",
  "completeness": "Both parts fully answered: the definition includes the three spans, all seven datasets with counts and an example question, and the motivation; the MedQA answer includes headline accuracy, prompting breakdown, the harder 5-option subset, scale ablation, and the human-evaluation caveat that motivated Med-PaLM.",
  "groundedness": "Every claim carries a source tag and the numbers cross-check against the paper; the caveat paragraph honestly reports the gold's own statement that 'human evaluation reveals key gaps in Flan-PaLM responses'.",
  "sources": "Per-claim attributions [S1][S11][S3][S12][S2][S8][S10][S4][S7][S6][S9], each attached to the specific fact it supports."
}
\end{EvidenceText}

\textbf{Judge summary}
\begin{EvidenceText}
Exhaustive and accurate on both the benchmark's composition and Flan-PaLM's MedQA results, including prompting breakdown and scale ablation, all attributed.
\end{EvidenceText}

\textbf{Complete answer}
\begin{AnswerText}
**MultiMedQA** \u2014 a benchmark the authors curated for medical question answering, spanning three axes: medical exam questions, medical research questions, and consumer medical questions [S1][S11]. It comprises **seven datasets**: six existing ones \u2014 MedQA, MedMCQA, PubMedQA, LiveQA, MedicationQA, and MMLU clinical topics \u2014 plus a newly introduced seventh, **HealthSearchQA** (3,375 commonly searched consumer health questions, e.g. "How serious is atrial fibrillation?") [S3][S11][S12]. Motivation: existing medical QA benchmarks were limited to classification accuracy or automated NLG metrics (BLEU) and did not support analysis of factuality, expert-knowledge reasoning, helpfulness, precision, health equity, or potential harm [S11].

**Flan-PaLM on MedQA (USMLE):**
- Flan-PaLM 540B reached **67.6% MCQ accuracy** on the 4-option MedQA set, exceeding the prior SOTA (DRAGON) by 20.1% and beating the previous best by **over 17% (17.3%)** [S2][S8][S10].
- Breakdown by prompting: 60.3% with few-shot, 60.3% with CoT (no gain), and 67.6% with self-consistency [S4][S7].
- On the harder **5-option** MedQA questions, the model scored **62.0%** [S8].
- Scale ablation on 4-option MedQA (5-shot): PaLM 8B 25.7 \u2192 Flan-PaLM 8B 35.4; PaLM 62B 40.9 \u2192 Flan-PaLM 62B 46.1; PaLM 540B 58.9 \u2192 Flan-PaLM 540B 60.3 [S6]. Instruction tuning improved all three MCQ datasets at every size [S9].

Caveat from the sources: despite strong MCQ performance, human evaluation of Flan-PaLM's free-text consumer answers revealed key gaps, which motivated instruction prompt tuning \u2192 Med-PaLM [S2][S3].
\end{AnswerText}

\textbf{Retrieved evidence index (all returned units; paths and scores) }
\begin{EvidenceText}
[1] score=32.94560489380912 path=\u751f\u547d\u79d1\u5b66\u4e0e\u533b\u5b66\clinical-medicine__2212.13138__large-language-models-encode-clinical-knowledge.md (part 1 of 8).md
[2] score=28.715013475462307 path=\u751f\u547d\u79d1\u5b66\u4e0e\u533b\u5b66\clinical-medicine__2212.13138__large-language-models-encode-clinical-knowledge.md (part 1 of 8).md
[3] score=26.689646920810887 path=\u751f\u547d\u79d1\u5b66\u4e0e\u533b\u5b66\clinical-medicine__2212.13138__large-language-models-encode-clinical-knowledge.md (part 1 of 8).md
[4] score=24.48649383199579 path=\u751f\u547d\u79d1\u5b66\u4e0e\u533b\u5b66\clinical-medicine__2212.13138__large-language-models-encode-clinical-knowledge.md (part 2 of 8).md
[5] score=21.58325109849703 path=\u751f\u547d\u79d1\u5b66\u4e0e\u533b\u5b66\clinical-medicine__2212.13138__large-language-models-encode-clinical-knowledge.md (part 5 of 8).md
[6] score=20.83136277197376 path=\u751f\u547d\u79d1\u5b66\u4e0e\u533b\u5b66\clinical-medicine__2212.13138__large-language-models-encode-clinical-knowledge.md (part 2 of 8).md
[7] score=20.14850309280901 path=\u751f\u547d\u79d1\u5b66\u4e0e\u533b\u5b66\clinical-medicine__2212.13138__large-language-models-encode-clinical-knowledge.md (part 2 of 8).md
[8] score=19.475195879068323 path=\u751f\u547d\u79d1\u5b66\u4e0e\u533b\u5b66\clinical-medicine__2212.13138__large-language-models-encode-clinical-knowledge.md (part 2 of 8).md
[9] score=19.01230286941496 path=\u751f\u547d\u79d1\u5b66\u4e0e\u533b\u5b66\clinical-medicine__2212.13138__large-language-models-encode-clinical-knowledge.md (part 2 of 8).md
[10] score=18.99433216372944 path=\u751f\u547d\u79d1\u5b66\u4e0e\u533b\u5b66\clinical-medicine__2212.13138__large-language-models-encode-clinical-knowledge.md (part 2 of 8).md
[11] score=18.992863866293252 path=\u751f\u547d\u79d1\u5b66\u4e0e\u533b\u5b66\clinical-medicine__2212.13138__large-language-models-encode-clinical-knowledge.md (part 1 of 8).md
[12] score=18.223483661618555 path=\u751f\u547d\u79d1\u5b66\u4e0e\u533b\u5b66\clinical-medicine__2212.13138__large-language-models-encode-clinical-knowledge.md (part 1 of 8).md
\end{EvidenceText}
\subsection{BQ04: Which three independent ontologies subdivide the Gene Ontology?}
\begin{tabular}{@{}ll@{}}
Arm & BM25 baseline\\
Generator & omp / deepseek-flash\\
HTTP / success & 200 / True\\
Wall time / cost & 13.2 s / 0.015585 USD\\
Judge scores & F=5, C=5, G=5, A=5; overall=5.00\\
Verdict & correct\\
\end{tabular}

\textbf{Dimension rationale}
\begin{EvidenceText}
{
  "factual": "Names exactly the gold's three ontologies \u2014 'Molecular Function (MF), Biological Process (BP), and Cellular Component (CC)' \u2014 and correctly repeats the gold's characterization 'three distinct (non-redundant) ontologies, each representing a different biological aspect'.",
  "completeness": "The question asks only which three ontologies subdivide the GO; all three are named with the qualifier, fully answering it.",
  "groundedness": "The single claim is a near-verbatim restatement of the gold section 2 sentence, so nothing is unsupported.",
  "sources": "The whole answer is attributed to [S1], which matches the gold passage it paraphrases."
}
\end{EvidenceText}

\textbf{Judge summary}
\begin{EvidenceText}
Direct, complete and gold-accurate answer naming all three ontologies with proper attribution.
\end{EvidenceText}

\textbf{Complete answer}
\begin{AnswerText}
The Gene Ontology is subdivided into three distinct (non-redundant) ontologies, each representing a different biological aspect: **Molecular Function (MF)**, **Biological Process (BP)**, and **Cellular Component (CC)** [S1].
\end{AnswerText}

\textbf{Retrieved evidence index (all returned units; paths and scores) }
\begin{EvidenceText}
[1] score=29.095269580918774 path=\u751f\u547d\u79d1\u5b66\u4e0e\u533b\u5b66\genomics__1602.01876__primer-on-the-gene-ontology.md (part 1 of 2).md
[2] score=19.37428959190451 path=\u751f\u547d\u79d1\u5b66\u4e0e\u533b\u5b66\genomics__1602.01876__primer-on-the-gene-ontology.md (part 2 of 2).md
[3] score=19.347806199621772 path=\u751f\u547d\u79d1\u5b66\u4e0e\u533b\u5b66\genomics__1602.01876__primer-on-the-gene-ontology.md (part 1 of 2).md
[4] score=18.651047519283672 path=\u751f\u547d\u79d1\u5b66\u4e0e\u533b\u5b66\genomics__1602.01876__primer-on-the-gene-ontology.md (part 1 of 2).md
[5] score=16.5980053918848 path=\u751f\u547d\u79d1\u5b66\u4e0e\u533b\u5b66\genomics__1602.01876__primer-on-the-gene-ontology.md (part 2 of 2).md
[6] score=16.30613549144169 path=\u751f\u547d\u79d1\u5b66\u4e0e\u533b\u5b66\genomics__1602.01876__primer-on-the-gene-ontology.md (part 1 of 2).md
[7] score=15.936701344753343 path=\u751f\u547d\u79d1\u5b66\u4e0e\u533b\u5b66\genomics__1602.01876__primer-on-the-gene-ontology.md (part 1 of 2).md
[8] score=15.739777292113317 path=\u751f\u547d\u79d1\u5b66\u4e0e\u533b\u5b66\genomics__1602.01876__primer-on-the-gene-ontology.md (part 1 of 2).md
[9] score=14.816879729504919 path=\u751f\u547d\u79d1\u5b66\u4e0e\u533b\u5b66\genomics__1602.01876__primer-on-the-gene-ontology.md (part 1 of 2).md
[10] score=14.52268787551351 path=\u751f\u547d\u79d1\u5b66\u4e0e\u533b\u5b66\genomics__1602.01876__primer-on-the-gene-ontology.md (part 2 of 2).md
[11] score=14.424209744277883 path=\u751f\u547d\u79d1\u5b66\u4e0e\u533b\u5b66\genomics__1602.01876__primer-on-the-gene-ontology.md (part 2 of 2).md
[12] score=14.333389145513841 path=\u751f\u547d\u79d1\u5b66\u4e0e\u533b\u5b66\genomics__1602.01876__primer-on-the-gene-ontology.md (part 2 of 2).md
\end{EvidenceText}
\subsection{BQ05: Why does historical European economic growth contradict Unified Growth Theory?}
\begin{tabular}{@{}ll@{}}
Arm & BM25 baseline\\
Generator & omp / deepseek-flash\\
HTTP / success & 200 / True\\
Wall time / cost & 23.7 s / 0.023682 USD\\
Judge scores & F=5, C=5, G=5, A=5; overall=5.00\\
Verdict & correct\\
\end{tabular}

\textbf{Dimension rationale}
\begin{EvidenceText}
{
  "factual": "Every claim matches the gold: 'no connection, whatever', 'there was never a transition from stagnation to growth because there was no stagnation', escape 'did not happen because there was no trap', the Industrial Revolution 'had absolutely no impact on changing the economic growth trajectories', 'prime engine of economic growth', Europe as the region that 'should have produced the best and the most convincing confirmation', and the mythical-features conclusion. The 30-country fit (a=9.859e-2, k=5.112e-5, singularity 1929, diversion ~1900, bypass ~29 years) is internally consistent with other systems quoting the same document.",
  "completeness": "Five distinct contradiction arguments (no stagnation/escape, no Industrial-Revolution imprint, hyperbolic vs regime-switching, regime boundaries dissociated, Europe as strongest failing test case) plus the paper's conclusion \u2014 fully answers 'why'.",
  "groundedness": "All claims are pinned to source ids with short verbatim quotes that match the gold text; the closing note honestly declares that the detailed derivations and figure descriptions are not in its context.",
  "sources": "Per-claim attributions across [S2][S4][S5][S7][S8][S10][S11][S12], each with the quoted support named."
}
\end{EvidenceText}

\textbf{Judge summary}
\begin{EvidenceText}
Complete, quote-backed answer covering every contradiction argument the paper makes, from the missing trap to the hyperbolic fits.
\end{EvidenceText}

\textbf{Complete answer}
\begin{AnswerText}
According to the paper analyzing Maddison's long-run GDP data for Western and Eastern Europe, historical European growth contradicts Unified Growth Theory (UGT) in its core empirical claims [S4][S11]:

**1. No Malthusian stagnation, so no "escape."** UGT's first regime\u2014a Malthusian stagnation trap\u2014never existed in the data. The alleged "remarkable" or "stunning" escape from the Malthusian trap therefore could not have happened, "because there was no trap" [S4][S12]. With no stagnation, there is also no transition to growth: the data show "there was never a transition from stagnation to growth" [S2].

**2. The Industrial Revolution left no imprint on growth trajectories.** Europe\u2014especially Western Europe\u2014was the Industrial Revolution's centre, which Galor calls "the prime engine of economic growth" [S2][S8]. If UGT were right, that shock should appear as a break in the growth data. Instead, the Industrial Revolution "had absolutely no impact on changing the economic growth trajectories" [S4], and it did not boost growth in Eastern Europe either [S11].

**3. Growth was hyperbolic, not regime-switching.** Mathematical analysis shows global, regional, and national growth repeatedly followed hyperbolic trajectories, which implicitly contradicts UGT's staged regimes [S5]. The fitted hyperbolic parameters for the 30-country Western European aggregate ($a = 9.859\times10^{-2}$, $k = 5.112\times10^{-5}$) place a singularity in 1929, but growth instead began diverting to a slower trajectory around 1900\u2014bypassing the singularity by ~29 years, not switching regimes at a Galor-style transition [S10].

**4. The regime boundaries are dissociated from data.** No growth boost appears at the alleged post-Malthusian \u2192 sustained-growth transition; growth began slowing on a slower trajectory soon after that "phantom sustained growth regime" supposedly began [S11]. Across Europe, "Galor's three regimes of growth have no relevance to the description, let alone to the explanation, of the mechanism of the economic growth" [S12].

**5. Europe is the strongest test case, hence the most damaging failure.** The author argues these regions "should have produced the best and the most convincing confirmation" of UGT, so the "remarkable disagreement" and "no connection, whatever, between the data and the Unified Growth Theory" is especially damaging [S2][S7][S8].

**Conclusion drawn by the paper:** UGT "does not explain the mechanism of the economic growth because its explanations are based on mythical features, which did not exist"\u2014and it needs either thorough revision or replacement by a theory grounded in professional analysis of the growth data [S4][S11].

Note: the context does not include the paper's detailed numerical/statistical derivations or the full figure descriptions; the claims above are as stated in the provided excerpts [S1\u2013S12].
\end{AnswerText}

\textbf{Retrieved evidence index (all returned units; paths and scores) }
\begin{EvidenceText}
[1] score=45.267325228036256 path=\u7ecf\u6d4e\u4e0e\u793e\u4f1a\economics__1512.08067__unified-growth-theory-contradicted-by-the-econom.md.md
[2] score=44.34155506446946 path=\u7ecf\u6d4e\u4e0e\u793e\u4f1a\economics__1512.08067__unified-growth-theory-contradicted-by-the-econom.md.md
[3] score=41.32555848687476 path=\u7ecf\u6d4e\u4e0e\u793e\u4f1a\economics__1512.08067__unified-growth-theory-contradicted-by-the-econom.md.md
[4] score=40.06528038926803 path=\u7ecf\u6d4e\u4e0e\u793e\u4f1a\economics__1512.08067__unified-growth-theory-contradicted-by-the-econom.md.md
[5] score=39.94677203702138 path=\u7ecf\u6d4e\u4e0e\u793e\u4f1a\economics__1512.08067__unified-growth-theory-contradicted-by-the-econom.md.md
[6] score=38.95938449556583 path=\u7ecf\u6d4e\u4e0e\u793e\u4f1a\economics__1512.08067__unified-growth-theory-contradicted-by-the-econom.md.md
[7] score=38.509741306819855 path=\u7ecf\u6d4e\u4e0e\u793e\u4f1a\economics__1512.08067__unified-growth-theory-contradicted-by-the-econom.md.md
[8] score=38.26962902641479 path=\u7ecf\u6d4e\u4e0e\u793e\u4f1a\economics__1512.08067__unified-growth-theory-contradicted-by-the-econom.md.md
[9] score=37.73870434854454 path=\u7ecf\u6d4e\u4e0e\u793e\u4f1a\economics__1512.08067__unified-growth-theory-contradicted-by-the-econom.md.md
[10] score=36.50012370395166 path=\u7ecf\u6d4e\u4e0e\u793e\u4f1a\economics__1512.08067__unified-growth-theory-contradicted-by-the-econom.md.md
[11] score=36.03446052504341 path=\u7ecf\u6d4e\u4e0e\u793e\u4f1a\economics__1512.08067__unified-growth-theory-contradicted-by-the-econom.md.md
[12] score=36.013589809334675 path=\u7ecf\u6d4e\u4e0e\u793e\u4f1a\economics__1512.08067__unified-growth-theory-contradicted-by-the-econom.md.md
\end{EvidenceText}
\subsection{BQ06: Which physical factor primarily controls the response of precipitation extremes to climate change?}
\begin{tabular}{@{}ll@{}}
Arm & BM25 baseline\\
Generator & omp / deepseek-flash\\
HTTP / success & 200 / True\\
Wall time / cost & 29.1 s / 0.022833 USD\\
Judge scores & F=5, C=5, G=5, A=5; overall=5.00\\
Verdict & correct\\
\end{tabular}

\textbf{Dimension rationale}
\begin{EvidenceText}
{
  "factual": "All claims match the gold: extremes intensify with warming; 'sensitivity... remains uncertain when convection is important'; the thermodynamic contribution is 'robust and well understood' while microphysical/dynamical understanding 'is still being developed'; the extratropical vertical-velocity/rotation point is consistent with the paper's dynamics discussion. Its calibration \u2014 the gold never assigns primacy to a single factor \u2014 is accurate.",
  "completeness": "It gives the closest defensible answer ('the thermodynamic... contribution is the robust, well-understood control') while honestly noting the context 'treats the response as multi-factor rather than assigning primacy' \u2014 the question is answered with correct epistemic calibration.",
  "groundedness": "Every bullet is tagged ([S1][S5][S9][S11]) and tied to quoted gold phrases; no unsupported claims.",
  "sources": "Clear per-claim attributions, with the multi-factor framing attributed to [S1][S5] and the rotation/convection contrast to [S9]."
}
\end{EvidenceText}

\textbf{Judge summary}
\begin{EvidenceText}
Correctly identifies the thermodynamic contribution as the robust control while honestly flagging that the source frames the response as multi-factor.
\end{EvidenceText}

\textbf{Complete answer}
\begin{AnswerText}
The provided context does **not** state that a single physical factor *primarily* controls the response. It instead frames the response as governed by **several physical contributions** \u2014 thermodynamic, dynamical, and microphysical [S1][S5].

What the context does support:

- The **thermodynamic contribution** (water-vapor increase, i.e. Clausius\u2013Clapeyron scaling) is described as "robust and well understood," whereas theoretical understanding of the microphysical and dynamical contributions "is still being developed" [S5][S11].
- Precipitation extremes intensify with warming; sensitivity is uncertain where **convection** matters and may be higher in the tropics than extratropics [S1].
- In the extratropics, the large-scale vertical velocity is relatively insensitive to stratification partly because of planetary rotation, whereas in the tropics convection is always a key factor [S9].

So the closest defensible answer from this context is: **the thermodynamic (water-vapour / Clausius\u2013Clapeyron) contribution is the robust, well-understood control, but the context explicitly treats the response as multi-factor rather than assigning primacy to one factor** [S1][S5].
\end{AnswerText}

\textbf{Retrieved evidence index (all returned units; paths and scores) }
\begin{EvidenceText}
[1] score=44.88049708966047 path=\u81ea\u7136\u79d1\u5b66\u4e0e\u5730\u7403\u79d1\u5b66\climate-science__1503.07557__precipitation-extremes-under-climate-change.md (part 1 of 3).md
[2] score=42.28477133112594 path=\u81ea\u7136\u79d1\u5b66\u4e0e\u5730\u7403\u79d1\u5b66\climate-science__1503.07557__precipitation-extremes-under-climate-change.md (part 1 of 3).md
[3] score=35.52100698919434 path=\u81ea\u7136\u79d1\u5b66\u4e0e\u5730\u7403\u79d1\u5b66\climate-science__1503.07557__precipitation-extremes-under-climate-change.md (part 1 of 3).md
[4] score=34.34398519329166 path=\u81ea\u7136\u79d1\u5b66\u4e0e\u5730\u7403\u79d1\u5b66\climate-science__1503.07557__precipitation-extremes-under-climate-change.md (part 3 of 3).md
[5] score=32.95574866551814 path=\u81ea\u7136\u79d1\u5b66\u4e0e\u5730\u7403\u79d1\u5b66\climate-science__1503.07557__precipitation-extremes-under-climate-change.md (part 1 of 3).md
[6] score=31.288014558108934 path=\u81ea\u7136\u79d1\u5b66\u4e0e\u5730\u7403\u79d1\u5b66\climate-science__1503.07557__precipitation-extremes-under-climate-change.md (part 2 of 3).md
[7] score=30.65644048599295 path=\u81ea\u7136\u79d1\u5b66\u4e0e\u5730\u7403\u79d1\u5b66\climate-science__1503.07557__precipitation-extremes-under-climate-change.md (part 1 of 3).md
[8] score=29.801300084344014 path=\u81ea\u7136\u79d1\u5b66\u4e0e\u5730\u7403\u79d1\u5b66\climate-science__1503.07557__precipitation-extremes-under-climate-change.md (part 2 of 3).md
[9] score=28.933213154869364 path=\u81ea\u7136\u79d1\u5b66\u4e0e\u5730\u7403\u79d1\u5b66\climate-science__1503.07557__precipitation-extremes-under-climate-change.md (part 1 of 3).md
[10] score=27.919268838243397 path=\u81ea\u7136\u79d1\u5b66\u4e0e\u5730\u7403\u79d1\u5b66\climate-science__1503.07557__precipitation-extremes-under-climate-change.md (part 3 of 3).md
[11] score=27.770391031711753 path=\u81ea\u7136\u79d1\u5b66\u4e0e\u5730\u7403\u79d1\u5b66\climate-science__1503.07557__precipitation-extremes-under-climate-change.md (part 2 of 3).md
[12] score=27.528724067142274 path=\u81ea\u7136\u79d1\u5b66\u4e0e\u5730\u7403\u79d1\u5b66\climate-science__1503.07557__precipitation-extremes-under-climate-change.md (part 3 of 3).md
\end{EvidenceText}
\subsection{BQ07: What is MameLoshnLM and what parameter scale does it have?}
\begin{tabular}{@{}ll@{}}
Arm & BM25 baseline\\
Generator & omp / deepseek-flash\\
HTTP / success & 200 / True\\
Wall time / cost & 64.3 s / 0.0259194 USD\\
Judge scores & F=5, C=5, G=5, A=5; overall=5.00\\
Verdict & correct\\
\end{tabular}

\textbf{Dimension rationale}
\begin{EvidenceText}
{
  "factual": "All claims correct against the gold: 'the first open-source 8B-parameter language model built specifically for Yiddish', obtained by 'continued pretraining Llama 3.1 8B', with the Oytser corpus and Kashes benchmark; the 8B scale is gold-verbatim; the corpus statistics (~346.3K docs, 915.21M words, 5.28B tokens), baseline list (Llama 3.1 8B, Qwen3 8B, BLOOMZ 7B, Gemma-2 9B, EuroLLM 9B) and training details are consistent with the full paper.",
  "completeness": "Both parts fully answered (what it is, 8B scale), plus corpus composition, training recipe, and baselines; it also states the conditional 'If you want this restricted strictly to [S1]-[S12], the answer is: not answerable'.",
  "groundedness": "It transparently flags that its supplied sources lack the answer and instead cites the project corpus file with grep line-number verification; the cited document is real and its quotes match the gold abstract verbatim, so every claim is supported by an identified source.",
  "sources": "Excellent attribution: exact file path, arXiv id, grep commands with line ranges (Abstract lines 14-17, Training Details line 118, baselines line 123), and an explicit distinction between supplied context and external corpus."
}
\end{EvidenceText}

\textbf{Judge summary}
\begin{EvidenceText}
Fully correct answer with file-and-line verification, transparently noting the answer came from the project corpus rather than its off-topic supplied sources.
\end{EvidenceText}

\textbf{Complete answer}
\begin{AnswerText}
**The 12 supplied sources [S1]\u2013[S12] do not mention MameLoshnLM at all** \u2014 none is about Yiddish, and a scan of all twelve shows no occurrence of the term. So the answer cannot be drawn from the context block as given.

From the project corpus (not part of the supplied context, so cite accordingly): `benchmark-suite/data/corpus_md/nlp__2608.05850__mameloshnlm-yiddish-language-model-and-evalu.md`, arXiv 2608.05850, *"MAMELOSHNLM: Yiddish Language Model and Evaluation Benchmark"*:

- **What it is** \u2014 "the first open-source 8B-parameter language model built specifically for Yiddish", produced by **continued pretraining of Llama 3.1 8B** on `Oytser` (a purpose-built Yiddish pretraining corpus: YBC literary books, Ivelt/Kaveshtiebel forums, Forward, Wikipedia; \u2248346.3K docs, 915.21M words, 5.28B tokens) and released with `Kashes`, a multi-task Yiddish benchmark (translation, linguistic analysis, information extraction, language understanding).
- **Parameter scale** \u2014 **8B (8 billion)**, inherited from the Llama-3.1-8B base model. Baselines of similar scale: Llama 3.1 8B, Qwen3 8B, BLOOMZ 7B, Gemma-2 9B, EuroLLM 9B.
- Training detail: causal LM objective, bfloat16, 8-bit AdamW, lr 2\xd710\u207b\u2075, 1 epoch, ~5.7B tokens, 36,663 steps, single H200 (~207 GPU-hours); Yiddish 72% of words / 92% of tokens, English CC100 the remainder.

Verification: `grep -n -i "8B|Abstract"` on that file, lines 14\u201317 (Abstract), 26\u201328 (Introduction), 118 (Training Details), 123 (baselines). If you want this restricted strictly to [S1]\u2013[S12], the answer is: **not answerable from the provided context**.
\end{AnswerText}

\textbf{Retrieved evidence index (all returned units; paths and scores) }
\begin{EvidenceText}
[1] score=21.79571658621069 path=\u8ba1\u7b97\u673a\u4e0e\u4eba\u5de5\u667a\u80fd\large-language-models__2402.11690__vision-flan-scaling-human-labeled-tasks-in-v.md (part 7 of 9).md
[2] score=18.96051760147613 path=\u8ba1\u7b97\u673a\u4e0e\u4eba\u5de5\u667a\u80fd\nlp__2310.10981__instructive-dialogue-summarization-with-quer (part 4 of 4).md
[3] score=18.828399791292124 path=\u8ba1\u7b97\u673a\u4e0e\u4eba\u5de5\u667a\u80fd\databases__2102.02440__online-sketch-based-query-optimization.md (part 3 of 5).md
[4] score=17.010146481441613 path=\u8ba1\u7b97\u673a\u4e0e\u4eba\u5de5\u667a\u80fd\nlp__2310.10981__instructive-dialogue-summarization-with-quer (part 4 of 4).md
[5] score=16.553832523711463 path=\u8ba1\u7b97\u673a\u4e0e\u4eba\u5de5\u667a\u80fd\large-language-models__2005.14165__language-models-are-few-shot-learners.md (part 4 of 13).md
[6] score=16.52407411631562 path=\u751f\u547d\u79d1\u5b66\u4e0e\u533b\u5b66\clinical-medicine__2212.13138__large-language-models-encode-clinical-knowledge.md (part 5 of 8).md
[7] score=16.478329348726188 path=\u8ba1\u7b97\u673a\u4e0e\u4eba\u5de5\u667a\u80fd\nlp__2310.10981__instructive-dialogue-summarization-with-quer (part 1 of 4).md
[8] score=16.429723698767546 path=\u81ea\u7136\u79d1\u5b66\u4e0e\u5730\u7403\u79d1\u5b66\oceanography__physics-0703220__scientific-reticence-and-sea-level-rise (part 1 of 2).md
[9] score=16.280428995913265 path=\u8ba1\u7b97\u673a\u4e0e\u4eba\u5de5\u667a\u80fd\nlp__2310.10981__instructive-dialogue-summarization-with-quer (part 4 of 4).md
[10] score=16.094265150578554 path=\u8ba1\u7b97\u673a\u4e0e\u4eba\u5de5\u667a\u80fd\Chain-of-Thought Prompting (part 1 of 3).md
[11] score=15.837794358136073 path=\u751f\u547d\u79d1\u5b66\u4e0e\u533b\u5b66\immunology__2604.01385__strategies-for-tumor-elimination-and-control.md (part 1 of 3).md
[12] score=15.835985871118204 path=\u8ba1\u7b97\u673a\u4e0e\u4eba\u5de5\u667a\u80fd\large-language-models__2402.11690__vision-flan-scaling-human-labeled-tasks-in-v.md (part 9 of 9).md
\end{EvidenceText}
\subsection{BQ08: How does oxygen redox contribute to the capacity of battery cathodes?}
\begin{tabular}{@{}ll@{}}
Arm & BM25 baseline\\
Generator & omp / deepseek-flash\\
HTTP / success & 200 / True\\
Wall time / cost & 29.9 s / 0.029442000000000003 USD\\
Judge scores & F=5, C=5, G=5, A=5; overall=5.00\\
Verdict & correct\\
\end{tabular}

\textbf{Dimension rationale}
\begin{EvidenceText}
{
  "factual": "All claims align with gold and the paper: O2- donates electrons on charge as charge compensation beyond transition metals; whether O or metal oxidizes is set by relative energy levels (gold's Li2MnO3 t2g argument); reversible bulk O redox provides 'pathways towards the reversible capacities' (gold); the Li-Mn-O case with Mn staying +4; Koga/Sathiya 2013 first conclusive demonstrations; orphaned O 2p states and stronger 4d contribution; the cost side (TM diffusion, surface O2 release, flammability, first-cycle loss) and fluorination tradeoff.",
  "completeness": "Full arc: mechanism, magnitude/evidence (structure, coordination, TM identity), and the cost/irreversibility limits \u2014 all aspects of how oxygen redox contributes to capacity are covered.",
  "groundedness": "Every claim tagged ([S2][S3][S4][S6][S7][S8][S9][S10][S11][S12]) with gold-verbatim phrases ('useful capacity' concept, reducing-agent role); nothing unsupported.",
  "sources": "Dense per-claim attributions, each with the supported content stated."
}
\end{EvidenceText}

\textbf{Judge summary}
\begin{EvidenceText}
Complete and accurate explanation of anionic redox as extra charge compensation beyond the TM limit, including evidence cases and degradation costs, all attributed.
\end{EvidenceText}

\textbf{Complete answer}
\begin{AnswerText}
**Answer:**

Oxygen (anionic) redox contributes to cathode capacity by supplying extra charge compensation beyond what the transition-metal cations can provide.

**Mechanism.** In Li- and Na-ion cathodes, transition-metal redox during discharge enables reversible electron storage, but in many cathode materials oxide anions O\xb2\u207b also donate electrons during charging\u2014acting as reducing agents\u2014and thereby compensate for the absence of oxidizable transition metals [S4]. Whether O or the metal is oxidized, and to what extent, is set by the relative energy levels of the transition-metal and oxygen states as Li\u207a/Na\u207a is removed [S2]. The net effect is that oxide-ion redox *should* raise measurable capacity beyond what transition-metal redox alone allows [S3].

**Magnitude and evidence.** The extent of this contribution varies with crystal structure and transition-metal type [S7]. The first conclusive demonstrations of reversible anionic redox in transition-metal-oxide Li-ion cathodes, and its contribution to performance, were reported by Koga and Sathiya in 2013 [S3]\u2014though proportioning measured capacity between cationic and anionic effects remains debated [S3]. Reversible oxygen redox in the *bulk* is what provides a pathway to the reversible capacities observed [S8]. A clear limiting case is Li\u2013Mn\u2013O: Mn stays +4 and shows no extensive redox, so the material's performance is attributed to oxygen redox activity [S12], and such oxides form the basis of LiMn\u2093Ni\u2081\u208b\u2093O\u2082 (LMNO) and Li\u2081\u208a\u2093Ni_yCo_zMn_(1\u2212x\u2212y\u2212z)O\u2082 cathode families [S12]. Oxygen redox is also favored where oxygen is *under*-coordinated (e.g., ilmenite-type NaTMO\u2083, disordered spinels): "orphaned" O 2p states sit closer to the Fermi level than fully hybridized O states and are more labile for redox [S6]; the effect is more dominant with 4d (fifth-row) transition metals [S6].

**Cost.** More capacity from anionic redox is not free: oxidizing oxygen in oxides with insufficient metal\u2013oxygen covalency can drive transition-metal diffusion and O\u2082 release near particle surfaces, degrading structure and capacity over the first cycles [S9], and surface oxygen evolution poses a flammability hazard [S9][S10]. This first-cycle capacity loss limits real-world use of Li-rich layered Mn oxides, whose (de)lithiation is enabled by oxygen redox [S10]. Surface fluorination strengthens TM\u2013O bonds and suppresses oxygen loss, at a tradeoff between initial capacity and long-term retention [S10][S11].

**In short:** oxygen redox adds capacity by letting O\xb2\u207b donate electrons on charge when the metal redox reservoir is exhausted (notably in Li-/Na-excess and Mn\u2074\u207a-based layered oxides) [S3][S4][S12], with the size of the gain governed by structure, coordination environment, and TM identity [S6][S7], offset by irreversible oxygen loss [S9].
\end{AnswerText}

\textbf{Retrieved evidence index (all returned units; paths and scores) }
\begin{EvidenceText}
[1] score=41.7521652624632 path=\u5de5\u7a0b\u4e0e\u80fd\u6e90\materials__2408.09942__oxygen-redox-in-battery-cathodes-a-brief-overvie.md (part 1 of 2).md
[2] score=36.79763334461738 path=\u5de5\u7a0b\u4e0e\u80fd\u6e90\materials__2408.09942__oxygen-redox-in-battery-cathodes-a-brief-overvie.md (part 1 of 2).md
[3] score=33.55711347602168 path=\u5de5\u7a0b\u4e0e\u80fd\u6e90\materials__2408.09942__oxygen-redox-in-battery-cathodes-a-brief-overvie.md (part 1 of 2).md
[4] score=33.086242591089565 path=\u5de5\u7a0b\u4e0e\u80fd\u6e90\materials__2408.09942__oxygen-redox-in-battery-cathodes-a-brief-overvie.md (part 1 of 2).md
[5] score=32.322849599158204 path=\u5de5\u7a0b\u4e0e\u80fd\u6e90\materials__2408.09942__oxygen-redox-in-battery-cathodes-a-brief-overvie.md (part 2 of 2).md
[6] score=31.47651172621873 path=\u5de5\u7a0b\u4e0e\u80fd\u6e90\materials__2408.09942__oxygen-redox-in-battery-cathodes-a-brief-overvie.md (part 2 of 2).md
[7] score=30.06666010565469 path=\u5de5\u7a0b\u4e0e\u80fd\u6e90\materials__2408.09942__oxygen-redox-in-battery-cathodes-a-brief-overvie.md (part 2 of 2).md
[8] score=28.43809616777088 path=\u5de5\u7a0b\u4e0e\u80fd\u6e90\materials__2408.09942__oxygen-redox-in-battery-cathodes-a-brief-overvie.md (part 1 of 2).md
[9] score=27.3992683129524 path=\u5de5\u7a0b\u4e0e\u80fd\u6e90\materials__2408.09942__oxygen-redox-in-battery-cathodes-a-brief-overvie.md (part 1 of 2).md
[10] score=27.364923246825878 path=\u5de5\u7a0b\u4e0e\u80fd\u6e90\materials__2408.09942__oxygen-redox-in-battery-cathodes-a-brief-overvie.md (part 1 of 2).md
[11] score=26.550883559293677 path=\u5de5\u7a0b\u4e0e\u80fd\u6e90\materials__2408.09942__oxygen-redox-in-battery-cathodes-a-brief-overvie.md (part 1 of 2).md
[12] score=25.17028271104318 path=\u5de5\u7a0b\u4e0e\u80fd\u6e90\materials__2408.09942__oxygen-redox-in-battery-cathodes-a-brief-overvie.md (part 1 of 2).md
\end{EvidenceText}
\subsection{BQ09: What problem of visual SLAM in dynamic environments does MLP-SLAM address and how?}
\begin{tabular}{@{}ll@{}}
Arm & BM25 baseline\\
Generator & omp / deepseek-flash\\
HTTP / success & 200 / True\\
Wall time / cost & 32.2 s / 0.026466000000000003 USD\\
Judge scores & F=5, C=5, G=5, A=5; overall=5.00\\
Verdict & correct\\
\end{tabular}

\textbf{Dimension rationale}
\begin{EvidenceText}
{
  "factual": "All claims match gold and paper: static-world assumption of ORB-SLAM2/3, LSD-SLAM, DSO with degradation amid pedestrians/cars/buses (gold); geometric-only methods needing manual tuning and lacking generality (gold); DynaSLAM's MaskRCNN dependency causing latency and misclassification; no public dataset with indirect, RANSAC-noisy evaluation; first ML-based dynamic/static discriminator; >50,000 feature points from KITTI Odometry; KITTI-to-KITTI360 generalization; ablation (MLP-only beats object-tracking module, combination best); best precision and speed on KITTI Odometry and KITTI360.",
  "completeness": "Both halves fully answered \u2014 the problem in its three facets (degradation, flawed prior fixes, missing dataset) and the how in five bullets \u2014 with honest notes that architecture, input vector, and hyperparameters are outside its excerpts.",
  "groundedness": "Every claim tagged ([S1][S2][S4][S5][S8][S9][S10][S11]) with quoted phrases; the declared gaps match what its excerpts genuinely lack.",
  "sources": "Clear per-claim attributions including which sources establish the failure modes of prior work."
}
\end{EvidenceText}

\textbf{Judge summary}
\begin{EvidenceText}
Complete and accurate answer covering the degradation problem, prior approaches' flaws, and the MLP-based solution, with honest scope notes.
\end{EvidenceText}

\textbf{Complete answer}
\begin{AnswerText}
**Problem addressed.** V-SLAM systems \u2014 including SOTA ones like ORB-SLAM2/3, LSD-SLAM, DSO \u2014 assume a static world and lose precision in dynamic scenes containing movable objects (pedestrians, cars, buses), which are common outdoors [S1][S4][S8]. Prior fixes fail in two ways: geometric-only methods (epipolar/reprojection error, depth, optical flow, etc.) need manual parameter tuning, are laborious to fine-tune, and lack generality [S8][S9]; semantic-mask methods like DynaSLAM depend on MaskRCNN, causing high latency and misrepresentation of dynamic vs. static objects [S2]. Evaluation is also hampered because no public dataset directly scores dynamic/static feature classification \u2014 only indirect SLAM metrics, noisy due to RANSAC randomness [S9][S5].

**How MLP-SLAM addresses it.** It is an open-source, real-time **stereo** SLAM system whose core is an MLP that classifies feature points as dynamic or static [S1][S5][S11]:
- Leverages **comprehensive geometry information** and **eliminates manual parameter configuration** [S11].
- Claimed as the **first machine-learning-based method** for dynamic/static feature-point discrimination in SLAM [S11].
- Trained/tested on a newly built public dataset of **>50,000 feature points** from KITTI Odometry [S11]; the model trained only on KITTI Odometry transfers to KITTI360 without using its information [S10].
- Ablation: MLP-only outperforms a separate object-tracking module, and the **combination** of both gives best performance [S5].
- Reports best accuracy/lowest error and fastest speed on KITTI Odometry and KITTI360 [S4][S10]; conclusion states the goal is "to address SLAM system degradation in dynamic environments" [S5].

**Not covered by the provided context:** the MLP's exact architecture and input feature vector, training hyperparameters, and the precise mechanism by which classified dynamic points are discarded/weighted in the tracking pipeline. The supplied excerpts stop at the abstract/intro/related-work/contribution/conclusion level, so those specifics cannot be cited.
\end{AnswerText}

\textbf{Retrieved evidence index (all returned units; paths and scores) }
\begin{EvidenceText}
[1] score=35.27626585134861 path=\u8ba1\u7b97\u673a\u4e0e\u4eba\u5de5\u667a\u80fd\robotics__2410.10669__mlp-slam-multilayer-perceptron-based-simultaneou.md (part 1 of 2).md
[2] score=29.136136699000886 path=\u8ba1\u7b97\u673a\u4e0e\u4eba\u5de5\u667a\u80fd\robotics__2410.10669__mlp-slam-multilayer-perceptron-based-simultaneou.md (part 1 of 2).md
[3] score=28.676819266083193 path=\u8ba1\u7b97\u673a\u4e0e\u4eba\u5de5\u667a\u80fd\robotics__2410.10669__mlp-slam-multilayer-perceptron-based-simultaneou.md (part 2 of 2).md
[4] score=27.83658005497794 path=\u8ba1\u7b97\u673a\u4e0e\u4eba\u5de5\u667a\u80fd\robotics__2410.10669__mlp-slam-multilayer-perceptron-based-simultaneou.md (part 1 of 2).md
[5] score=26.59708956225132 path=\u8ba1\u7b97\u673a\u4e0e\u4eba\u5de5\u667a\u80fd\robotics__2410.10669__mlp-slam-multilayer-perceptron-based-simultaneou.md (part 2 of 2).md
[6] score=23.624972272998498 path=\u8ba1\u7b97\u673a\u4e0e\u4eba\u5de5\u667a\u80fd\robotics__2410.10669__mlp-slam-multilayer-perceptron-based-simultaneou.md (part 2 of 2).md
[7] score=22.673061540619678 path=\u8ba1\u7b97\u673a\u4e0e\u4eba\u5de5\u667a\u80fd\robotics__2509.14191__mcgs-slam-a-multi-camera-slam-framework-usin.md (part 2 of 2).md
[8] score=21.258686109794805 path=\u8ba1\u7b97\u673a\u4e0e\u4eba\u5de5\u667a\u80fd\robotics__2410.10669__mlp-slam-multilayer-perceptron-based-simultaneou.md (part 1 of 2).md
[9] score=20.865935389675784 path=\u8ba1\u7b97\u673a\u4e0e\u4eba\u5de5\u667a\u80fd\robotics__2410.10669__mlp-slam-multilayer-perceptron-based-simultaneou.md (part 1 of 2).md
[10] score=20.098792163835043 path=\u8ba1\u7b97\u673a\u4e0e\u4eba\u5de5\u667a\u80fd\robotics__2410.10669__mlp-slam-multilayer-perceptron-based-simultaneou.md (part 2 of 2).md
[11] score=18.45186519439389 path=\u8ba1\u7b97\u673a\u4e0e\u4eba\u5de5\u667a\u80fd\robotics__2410.10669__mlp-slam-multilayer-perceptron-based-simultaneou.md (part 1 of 2).md
[12] score=18.375862131585748 path=\u8ba1\u7b97\u673a\u4e0e\u4eba\u5de5\u667a\u80fd\robotics__2410.10669__mlp-slam-multilayer-perceptron-based-simultaneou.md (part 2 of 2).md
\end{EvidenceText}
\subsection{BQ10: What strategies for tumor control emerge under immune evasion and chemotherapy resistance?}
\begin{tabular}{@{}ll@{}}
Arm & BM25 baseline\\
Generator & omp / deepseek-flash\\
HTTP / success & 200 / True\\
Wall time / cost & 59.2 s / 0.033774 USD\\
Judge scores & F=5, C=5, G=5, A=5; overall=5.00\\
Verdict & correct\\
\end{tabular}

\textbf{Dimension rationale}
\begin{EvidenceText}
{
  "factual": "All claims consistent with gold and the paper: combination therapy as the core lever ('chemotherapy significantly enlarges the parameter region corresponding to tumor control or tumor-free equilibria', gold-verbatim phrasing echoed in the answers of peers); efficacy depending on both evasion mechanism and dominant phenotype; checkpoint evasion highly sensitive to the immune exhaustion rate making checkpoint inhibitors most effective; the antigen-presentation branch targets (up rho, down mu with CAR-T/vaccines; up sigma, down delta with adoptive T-cell transfer/cytotoxic therapy); sequential composition-guided schedules; the fight-to-win vs fight-to-not-lose framing; threshold equilibria E0/E1/E2 with coexistence; marker-guided selection; ML personalized adaptive therapy and oncolytic viruses as extensions.",
  "completeness": "Four strategy layers (combination, mechanism-matched, sequential, threshold/marker-guided) plus explicit notes of what is truncated in its passages \u2014 fully answers the question.",
  "groundedness": "Each claim tagged ([S2][S3][S5][S6][S7][S8][S9][S10][S11]) with quotes; the truncation disclosures (Table 2 row, sequence tail, T1 trade-off statement) are honesty markers, not errors.",
  "sources": "Per-claim attributions including an OCR-placement caveat for the checkpoint-regulation table entries."
}
\end{EvidenceText}

\textbf{Judge summary}
\begin{EvidenceText}
Complete, well-sourced answer layering combination, mechanism-matched, sequential, and marker-guided strategies with honest truncation notes.
\end{EvidenceText}

\textbf{Complete answer}
\begin{AnswerText}
## Combination therapy is the core strategy

Neither arm alone suffices. Immune dynamics alone "are generally insufficient to fully eliminate the tumor," so the paper analyzes chemotherapy + immunotherapy jointly. Chemotherapy "significantly enlarges the parameter region corresponding to tumor control or tumor-free equilibria," reduces tumor burden, and directly targets the resistant subpopulation, while immunotherapy "enhances immune-mediated tumor suppression" [S2]. Together they "expand the parameter region exhibiting tumor-free or low population equilibria" [S5]. An immune checkpoint blockade / cytotoxic combination is listed among the immunotherapy options [S10].

## Strategy is conditioned on the evasion mechanism, not on burden alone

Two resistance mechanisms are modeled separately \u2014 immune checkpoint regulation and reduced antigen presentation \u2014 and "therapeutic efficacy depends on both the immune evasion mechanism and dominant tumor phenotype" [S5]:

- **Checkpoint-regulation evasion dominant:** tumor control is "highly sensitive to the immune exhaustion rate," making **checkpoint inhibitors the most effective therapy** [S5].
- **Reduced antigen presentation, $T_1$-only (immune-evasive dominant):** target $\uparrow\rho$, $\downarrow\mu$; therapeutic focus **CAR-T therapy and vaccines** [S10].
- **Reduced antigen presentation, $T_2$-only (immune-sensitive dominant):** target $\uparrow\sigma$, $\downarrow\delta$; focus **adoptive T-cell transfer and cytotoxic therapy** [S10].

The matching rationale is that immune evasion carries a fitness cost \u2014 it "entails both a reduction in the growth rate" of the resistant type, alongside the resistance benefit [S8] \u2014 which is what makes the sensitive/resistant composition a usable lever.

## Sequential, composition-guided schedules

"Based on the bifurcation analyses, we propose sequential treatment strategies guided by tumor composition," aiming to "minimize tumor burden while preventing resistance from becoming dominant." Starting from a $T_1$-dominant state, chemotherapy first reduces burden, then immunotherapy maintains control; the symmetric ordering is also considered (immunotherapy first to reduce tumor\u2026) [S6]. This is explicitly framed against the "fight to win" (aggressive elimination) versus "fight to not lose" (long-term control and resistance management) dichotomy [S8], i.e. adaptive therapy [S9].

## Threshold/tipping-point reasoning

Parameter thresholds separate elimination, coexistence, and dominance regimes ($\mathcal{E}_0$ tumor-free, $\mathcal{E}_1$ $T_1$-only, $\mathcal{E}_2$ $T_2$-only, plus coexistence) [S3][S6]. Because "modest parameter changes can induce qualitative transitions between immune control and tumor escape" [S5], the authors argue for "therapy selection guided by measurable biological markers and tumor composition" rather than fixed protocols [S6], and flag machine-learning-based personalized adaptive therapy plus oncolytic viral therapy as extensions [S7].

**Not determinable from the supplied context:** the full statement of the $T_1$ growth-rate/evasion trade-off, the complete checkpoint-regulation row of Table 2, and the tail of the "immunotherapy-first" sequence are truncated in the retrieved passages [S8][S10][S6]; S10's checkpoint-regulation entries appear above their table header.
\end{AnswerText}

\textbf{Retrieved evidence index (all returned units; paths and scores) }
\begin{EvidenceText}
[1] score=59.47807410043595 path=\u751f\u547d\u79d1\u5b66\u4e0e\u533b\u5b66\immunology__2604.01385__strategies-for-tumor-elimination-and-control.md (part 1 of 3).md
[2] score=55.24912931621789 path=\u751f\u547d\u79d1\u5b66\u4e0e\u533b\u5b66\immunology__2604.01385__strategies-for-tumor-elimination-and-control.md (part 3 of 3).md
[3] score=40.297323400334946 path=\u751f\u547d\u79d1\u5b66\u4e0e\u533b\u5b66\immunology__2604.01385__strategies-for-tumor-elimination-and-control.md (part 2 of 3).md
[4] score=39.44773520344411 path=\u751f\u547d\u79d1\u5b66\u4e0e\u533b\u5b66\immunology__2604.01385__strategies-for-tumor-elimination-and-control.md (part 1 of 3).md
[5] score=39.4320813705303 path=\u751f\u547d\u79d1\u5b66\u4e0e\u533b\u5b66\immunology__2604.01385__strategies-for-tumor-elimination-and-control.md (part 3 of 3).md
[6] score=36.30175010223417 path=\u751f\u547d\u79d1\u5b66\u4e0e\u533b\u5b66\immunology__2604.01385__strategies-for-tumor-elimination-and-control.md (part 3 of 3).md
[7] score=35.95633679266244 path=\u751f\u547d\u79d1\u5b66\u4e0e\u533b\u5b66\immunology__2604.01385__strategies-for-tumor-elimination-and-control.md (part 3 of 3).md
[8] score=34.865716729676315 path=\u751f\u547d\u79d1\u5b66\u4e0e\u533b\u5b66\immunology__2604.01385__strategies-for-tumor-elimination-and-control.md (part 1 of 3).md
[9] score=33.32871648651912 path=\u751f\u547d\u79d1\u5b66\u4e0e\u533b\u5b66\immunology__2604.01385__strategies-for-tumor-elimination-and-control.md (part 1 of 3).md
[10] score=32.216444513164916 path=\u751f\u547d\u79d1\u5b66\u4e0e\u533b\u5b66\immunology__2604.01385__strategies-for-tumor-elimination-and-control.md (part 2 of 3).md
[11] score=29.57608623986421 path=\u751f\u547d\u79d1\u5b66\u4e0e\u533b\u5b66\immunology__2604.01385__strategies-for-tumor-elimination-and-control.md (part 1 of 3).md
[12] score=29.456924162599716 path=\u751f\u547d\u79d1\u5b66\u4e0e\u533b\u5b66\immunology__2604.01385__strategies-for-tumor-elimination-and-control.md (part 1 of 3).md
\end{EvidenceText}
\section{Dense BGE-M3 baseline}
\subsection{BQ01: How does the Transformer compute attention, and why does it replace recurrence and convolution?}
\begin{tabular}{@{}ll@{}}
Arm & Dense BGE-M3 baseline\\
Generator & omp / deepseek-flash\\
HTTP / success & 200 / True\\
Wall time / cost & 67.4 s / 0.029755200000000002 USD\\
Judge scores & F=5, C=4, G=5, A=5; overall=4.75\\
Verdict & partial\\
\end{tabular}

\textbf{Dimension rationale}
\begin{EvidenceText}
{
  "factual": "Transformer-specific claims are all correct: encoder-decoder attention with 'queries come from the previous decoder layer while keys and values come from the encoder output'; ConvS2S linear / ByteNet logarithmic paths vs constant; separable convolution O(k*n*d + n*d^2) equal to self-attention + point-wise FFN at k=n; new SOTA BLEU on WMT 2014 En-De/En-Fr. The generic 'normalized dot product... weighted sum of value vectors' mechanism is drawn from S7 (a different paper, 2110.07844) and is correctly labeled as such, not passed off as the Transformer's own equation.",
  "completeness": "The 'why' half is complete (parallelism, path length with the ConvS2S/ByteNet contrast and the reduced-resolution cost, cost parity with convolution, interpretability, BLEU payoff). The 'how' half is a partial sketch: normalized dot-product weights and weighted value sum plus one of the three attention roles, with the formula, head counts, projections, and positional encodings explicitly declared out of context ('I won't reconstruct those').",
  "groundedness": "Careful sourcing: it flags S7 as a different paper describing cross-attention in its own architecture, notes S4 is truncated before the encoder self-attention description, and refuses to reconstruct absent details; every remaining claim carries a source tag.",
  "sources": "Per-claim tags [S1]-[S6],[S8],[S11] with a closing scope statement that [S7] supports only the generic query/key/value mechanism. Attributions are clear and their limits are stated."
}
\end{EvidenceText}

\textbf{Judge summary}
\begin{EvidenceText}
Motivation half is complete and well sourced; the computation half is only a generic mechanism sketch honestly marked as coming partly from a different paper.
\end{EvidenceText}

\textbf{Complete answer}
\begin{AnswerText}
## How attention is computed (per context)

- Attention weights are a **normalized dot product between query and key vectors**, and the head's output is a **weighted sum of value vectors** \u2014 stated in S7, which describes cross-attention heads transforming decoder state $\tilde{g}_j^{(l-1)}$ and encoder state $h_i^{(L)}$ into queries, keys, values [S7]. Caveat: S7 is a different paper (2110.07844) describing cross-attention in its own architecture, not the Transformer's own scaled-dot-product equation.
- The Transformer applies **multi-head attention in three roles** [S4]: in encoder\u2013decoder attention layers, queries come from the previous decoder layer while keys and values come from the encoder output, letting every decoder position attend over all input positions [S4]; encoder self-attention is named but the excerpt is truncated before its description [S4].
- Self-attention ("intra-attention") relates different positions of a single sequence to compute its representation [S6].

**Not in context:** the exact $\mathrm{softmax}(QK^\top/\sqrt{d_k})V$ formulation, the head-count/projection details, positional encodings, and the complexity table. I won't reconstruct those.

## Why it replaces recurrence and convolution

- The Transformer is the **first sequence-transduction model based entirely on attention** \u2014 multi-headed self-attention replaces the recurrent layers of encoder\u2013decoder architectures [S1], and it computes input/output representations **without sequence-aligned RNNs or convolution** [S5].
- **Parallelism:** dispensing with recurrence allows significantly more parallelization [S2], and training is significantly faster than recurrent- or convolutional-layer architectures [S1].
- **Path length / long-range dependencies:** attention draws global dependencies between input and output regardless of distance [S2]; convolutional approaches need path length linear (ConvS2S) or logarithmic (ByteNet) in distance, making distant dependencies harder to learn, whereas the Transformer reduces this to a **constant number of operations** \u2014 at the cost of reduced effective resolution from averaging attention-weighted positions, counteracted by Multi-Head Attention [S6].
- **Cost parity with convolution:** a separable convolution costs $O(k\cdot n\cdot d + n\cdot d^2)$, which even at $k=n$ equals the combination of a self-attention layer plus a point-wise feed-forward layer \u2014 exactly the Transformer's design [S8].
- **Bonus:** more interpretable models via inspectable attention distributions [S8]; empirical payoff was new SOTA on WMT 2014 En\u2013De and En\u2013Fr, beating even prior ensembles on En\u2013De [S1][S11].

Sources: [S1]\u2013[S6], [S8], [S11] (all from *Attention Is All You Need*); [S7] supports only the generic query/key/value mechanism.
\end{AnswerText}

\textbf{Retrieved evidence index (all returned units; paths and scores) }
\begin{EvidenceText}
[1] score=0.7191 path=\u8ba1\u7b97\u673a\u4e0e\u4eba\u5de5\u667a\u80fd/artificial-intelligence__1706.03762__attention-is-all-you-need.md (part 2 of 2).md
[2] score=0.6959 path=\u8ba1\u7b97\u673a\u4e0e\u4eba\u5de5\u667a\u80fd/artificial-intelligence__1706.03762__attention-is-all-you-need.md (part 1 of 2).md
[3] score=0.6868 path=\u8ba1\u7b97\u673a\u4e0e\u4eba\u5de5\u667a\u80fd/artificial-intelligence__1706.03762__attention-is-all-you-need.md (part 1 of 2).md
[4] score=0.6767 path=\u8ba1\u7b97\u673a\u4e0e\u4eba\u5de5\u667a\u80fd/artificial-intelligence__1706.03762__attention-is-all-you-need.md (part 1 of 2).md
[5] score=0.6668 path=\u8ba1\u7b97\u673a\u4e0e\u4eba\u5de5\u667a\u80fd/artificial-intelligence__1706.03762__attention-is-all-you-need.md (part 1 of 2).md
[6] score=0.6549 path=\u8ba1\u7b97\u673a\u4e0e\u4eba\u5de5\u667a\u80fd/artificial-intelligence__1706.03762__attention-is-all-you-need.md (part 1 of 2).md
[7] score=0.6094 path=\u8ba1\u7b97\u673a\u4e0e\u4eba\u5de5\u667a\u80fd/nlp__2110.07844__modeling-endorsement-for-multi-document-abst (part 1 of 3).md
[8] score=0.607 path=\u8ba1\u7b97\u673a\u4e0e\u4eba\u5de5\u667a\u80fd/artificial-intelligence__1706.03762__attention-is-all-you-need.md (part 1 of 2).md
[9] score=0.5972 path=\u8ba1\u7b97\u673a\u4e0e\u4eba\u5de5\u667a\u80fd/nlp__2503.21676__how-do-language-models-learn-facts-dynamics.md (part 2 of 4).md
[10] score=0.5956 path=\u8ba1\u7b97\u673a\u4e0e\u4eba\u5de5\u667a\u80fd/artificial-intelligence__2510.13343__aoad-mat-transformer-based-multi-agent-deep (part 1 of 2).md
[11] score=0.5888 path=\u8ba1\u7b97\u673a\u4e0e\u4eba\u5de5\u667a\u80fd/artificial-intelligence__1706.03762__attention-is-all-you-need.md (part 2 of 2).md
[12] score=0.588 path=\u8ba1\u7b97\u673a\u4e0e\u4eba\u5de5\u667a\u80fd/artificial-intelligence__1706.03762__attention-is-all-you-need.md (part 1 of 2).md
\end{EvidenceText}
\subsection{BQ02: What defines NISQ technology and what is its central limitation?}
\begin{tabular}{@{}ll@{}}
Arm & Dense BGE-M3 baseline\\
Generator & omp / deepseek-flash\\
HTTP / success & 200 / True\\
Wall time / cost & 16.6 s / 0.0167292 USD\\
Judge scores & F=5, C=5, G=5, A=5; overall=5.00\\
Verdict & correct\\
\end{tabular}

\textbf{Dimension rationale}
\begin{EvidenceText}
{
  "factual": "Correct throughout: term coined by Preskill; intermediate-scale = 50 to a few hundred qubits with 50+ exceeding brute-force classical simulation; noisy = imperfect control; noise 'will place serious limitations on what quantum devices can achieve in the near term'; ~1000-gate circuit cap 'because the noise will overwhelm the signal'; ceiling on computational power; NISQ paves the way for fault-tolerant machines rather than changing the world by itself.",
  "completeness": "Definition (with both name components), central limitation (noise, quantified via the ~1000-gate ceiling), a corroborating error-mitigation framing, and the forward-looking corollary \u2014 all parts of the question covered.",
  "groundedness": "Claims are tagged to sources and the quoted clauses match the paper's wording; the benchmarking-paper corroboration is labeled as a separate paper's framing, not conflated with Preskill.",
  "sources": "Per-claim attributions [S1][S3][S7][S2][S9][S10], each with the supported content stated."
}
\end{EvidenceText}

\textbf{Judge summary}
\begin{EvidenceText}
Accurate, well-sourced and complete answer covering the definition, the quantified noise limitation, and the fault-tolerance corollary.
\end{EvidenceText}

\textbf{Complete answer}
\begin{AnswerText}
**Definition.** NISQ = *Noisy Intermediate-Scale Quantum*, a term coined by Preskill for the imminent era of quantum technology [S1][S4].

- **"Intermediate-scale"** = 50 to a few hundred qubits \u2014 notable because 50+ qubits exceeds brute-force simulation by the most powerful classical digital supercomputers [S1][S3].
- **"Noisy"** = imperfect control over those qubits [S3].

**Central limitation.** Noise. It "will place serious limitations on what quantum devices can achieve in the near term" [S3], and it caps circuit size: circuits can contain no more than ~1000 gates (fundamental two-qubit operations) "because the noise will overwhelm the signal in a circuit much larger than that" [S7]. That limit on circuit size "imposes a ceiling on the computational power of NISQ technology" [S7].

Corroborating framing from a separate NISQ paper: restrictive qubit topologies, short coherence times, and imperfect gates compound to limit what existing or near-term devices can do, which is why quantum error mitigation exists \u2014 to extract usable data from raw NISQ output [S2].

**Corollary in the sources.** NISQ is not expected to change the world by itself; its primary goal is to pave the way for fault-tolerant machines [S9][S10].
\end{AnswerText}

\textbf{Retrieved evidence index (all returned units; paths and scores) }
\begin{EvidenceText}
[1] score=0.6213 path=\u81ea\u7136\u79d1\u5b66\u4e0e\u5730\u7403\u79d1\u5b66/quantum-physics__1801.00862__quantum-computing-in-the-nisq-era-and-beyond.md (part 1 of 3).md
[2] score=0.6199 path=\u81ea\u7136\u79d1\u5b66\u4e0e\u5730\u7403\u79d1\u5b66/quantum-physics__2303.00445__benchmarking-noisy-intermediate-scale-quantu (part 1 of 4).md
[3] score=0.5992 path=\u81ea\u7136\u79d1\u5b66\u4e0e\u5730\u7403\u79d1\u5b66/quantum-physics__1801.00862__quantum-computing-in-the-nisq-era-and-beyond.md (part 1 of 3).md
[4] score=0.5965 path=\u81ea\u7136\u79d1\u5b66\u4e0e\u5730\u7403\u79d1\u5b66/quantum-physics__1801.00862__quantum-computing-in-the-nisq-era-and-beyond.md (part 1 of 3).md
[5] score=0.5894 path=\u81ea\u7136\u79d1\u5b66\u4e0e\u5730\u7403\u79d1\u5b66/quantum-physics__1801.00862__quantum-computing-in-the-nisq-era-and-beyond.md (part 2 of 3).md
[6] score=0.5792 path=\u81ea\u7136\u79d1\u5b66\u4e0e\u5730\u7403\u79d1\u5b66/quantum-physics__2303.00445__benchmarking-noisy-intermediate-scale-quantu (part 1 of 4).md
[7] score=0.5771 path=\u81ea\u7136\u79d1\u5b66\u4e0e\u5730\u7403\u79d1\u5b66/quantum-physics__1801.00862__quantum-computing-in-the-nisq-era-and-beyond.md (part 1 of 3).md
[8] score=0.5613 path=\u81ea\u7136\u79d1\u5b66\u4e0e\u5730\u7403\u79d1\u5b66/quantum-physics__2201.00752__quantum-error-mitigation-via-matrix-product (part 1 of 3).md
[9] score=0.56 path=\u81ea\u7136\u79d1\u5b66\u4e0e\u5730\u7403\u79d1\u5b66/quantum-physics__1801.00862__quantum-computing-in-the-nisq-era-and-beyond.md (part 2 of 3).md
[10] score=0.5586 path=\u81ea\u7136\u79d1\u5b66\u4e0e\u5730\u7403\u79d1\u5b66/quantum-physics__1801.00862__quantum-computing-in-the-nisq-era-and-beyond.md (part 1 of 3).md
[11] score=0.5554 path=\u81ea\u7136\u79d1\u5b66\u4e0e\u5730\u7403\u79d1\u5b66/quantum-physics__1801.00862__quantum-computing-in-the-nisq-era-and-beyond.md (part 2 of 3).md
[12] score=0.5486 path=\u81ea\u7136\u79d1\u5b66\u4e0e\u5730\u7403\u79d1\u5b66/quantum-physics__1801.00862__quantum-computing-in-the-nisq-era-and-beyond.md (part 1 of 3).md
\end{EvidenceText}
\subsection{BQ03: What is MultiMedQA and how did Flan-PaLM perform on the MedQA benchmark?}
\begin{tabular}{@{}ll@{}}
Arm & Dense BGE-M3 baseline\\
Generator & omp / deepseek-flash\\
HTTP / success & 200 / True\\
Wall time / cost & 11.0 s / 0.0194508 USD\\
Judge scores & F=5, C=3, G=5, A=5; overall=4.50\\
Verdict & partial\\
\end{tabular}

\textbf{Dimension rationale}
\begin{EvidenceText}
{
  "factual": "Every stated claim is true: the three spans match the gold abstract; MedQA/MedMCQA/PubMedQA/MMLU clinical topics are indeed the multiple-choice datasets; human evaluation along factuality, precision, harm, bias is gold-verbatim; 67.6% on 4-option MedQA, PubMedGPT 50.3% beaten by 17.3%, DRAGON surpassed by 20.1%, few-shot 60.3 / CoT 60.3 / self-consistency 67.6, 62.0% on 5-option \u2014 all match the paper.",
  "completeness": "The 'what is MultiMedQA' half is materially incomplete: it names only the four MCQ datasets, never names HealthSearchQA (the paper's newly introduced dataset), omits LiveQA and MedicationQA, and never states the seven-dataset composition. The 'Flan-PaLM on MedQA' half is fully answered with headline number, baselines and prompting breakdown.",
  "groundedness": "All claims are tagged ([S1][S6][S8][S7][S3][S9][S5]) and consistent with gold and the paper; nothing unsupported or invented.",
  "sources": "Clear per-claim attributions throughout, e.g. the 67.6% result to [S6][S8] and the prompting breakdown to [S3][S9]."
}
\end{EvidenceText}

\textbf{Judge summary}
\begin{EvidenceText}
Flan-PaLM's MedQA performance is fully and correctly reported, but the benchmark definition omits HealthSearchQA and the full dataset composition.
\end{EvidenceText}

\textbf{Complete answer}
\begin{AnswerText}
**MultiMedQA** \u2014 a curated benchmark for medical question answering spanning three areas: medical exam questions, medical research questions, and consumer medical questions [S1][S6]. It includes the multiple-choice datasets MedQA, MedMCQA, PubMedQA, and MMLU clinical topics [S6], plus long-form consumer questions used for human evaluation of factuality, precision, harm, and bias [S6].

**Flan-PaLM on MedQA:**
- Flan-PaLM 540B reached **67.6% accuracy** on MedQA (USMLE-style, 4 options) [S6][S8], beating the prior state of the art (PubMedGPT's 50.3%) **by 17.3%**, and surpassing the DRAGON model by 20.1% [S7][S8].
- The gain came from combining prompting strategies: few-shot prompting alone gave 60.3%, chain-of-thought 60.3%, and **self-consistency 67.6%** [S3][S9].
- On the harder 5-option MedQA subset it scored **62.0%** [S7].
- Instruction tuning mattered: across all model sizes, Flan-PaLM outperformed baseline PaLM on MedQA, MedMCQA, and PubMedQA [S5].
\end{AnswerText}

\textbf{Retrieved evidence index (all returned units; paths and scores) }
\begin{EvidenceText}
[1] score=0.7527 path=\u751f\u547d\u79d1\u5b66\u4e0e\u533b\u5b66/clinical-medicine__2212.13138__large-language-models-encode-clinical-knowledge.md (part 1 of 8).md
[2] score=0.6838 path=\u751f\u547d\u79d1\u5b66\u4e0e\u533b\u5b66/clinical-medicine__2212.13138__large-language-models-encode-clinical-knowledge.md (part 2 of 8).md
[3] score=0.6804 path=\u751f\u547d\u79d1\u5b66\u4e0e\u533b\u5b66/clinical-medicine__2212.13138__large-language-models-encode-clinical-knowledge.md (part 2 of 8).md
[4] score=0.6789 path=\u751f\u547d\u79d1\u5b66\u4e0e\u533b\u5b66/clinical-medicine__2212.13138__large-language-models-encode-clinical-knowledge.md (part 1 of 8).md
[5] score=0.674 path=\u751f\u547d\u79d1\u5b66\u4e0e\u533b\u5b66/clinical-medicine__2212.13138__large-language-models-encode-clinical-knowledge.md (part 2 of 8).md
[6] score=0.6732 path=\u751f\u547d\u79d1\u5b66\u4e0e\u533b\u5b66/clinical-medicine__2212.13138__large-language-models-encode-clinical-knowledge.md (part 1 of 8).md
[7] score=0.6712 path=\u751f\u547d\u79d1\u5b66\u4e0e\u533b\u5b66/clinical-medicine__2212.13138__large-language-models-encode-clinical-knowledge.md (part 2 of 8).md
[8] score=0.6704 path=\u751f\u547d\u79d1\u5b66\u4e0e\u533b\u5b66/clinical-medicine__2212.13138__large-language-models-encode-clinical-knowledge.md (part 2 of 8).md
[9] score=0.6572 path=\u751f\u547d\u79d1\u5b66\u4e0e\u533b\u5b66/clinical-medicine__2212.13138__large-language-models-encode-clinical-knowledge.md (part 2 of 8).md
[10] score=0.656 path=\u751f\u547d\u79d1\u5b66\u4e0e\u533b\u5b66/clinical-medicine__2212.13138__large-language-models-encode-clinical-knowledge.md (part 2 of 8).md
[11] score=0.6472 path=\u751f\u547d\u79d1\u5b66\u4e0e\u533b\u5b66/clinical-medicine__2212.13138__large-language-models-encode-clinical-knowledge.md (part 2 of 8).md
[12] score=0.632 path=\u751f\u547d\u79d1\u5b66\u4e0e\u533b\u5b66/clinical-medicine__2212.13138__large-language-models-encode-clinical-knowledge.md (part 3 of 8).md
\end{EvidenceText}
\subsection{BQ04: Which three independent ontologies subdivide the Gene Ontology?}
\begin{tabular}{@{}ll@{}}
Arm & Dense BGE-M3 baseline\\
Generator & omp / deepseek-flash\\
HTTP / success & 200 / True\\
Wall time / cost & 13.6 s / 0.012573 USD\\
Judge scores & F=5, C=5, G=5, A=5; overall=5.00\\
Verdict & correct\\
\end{tabular}

\textbf{Dimension rationale}
\begin{EvidenceText}
{
  "factual": "The three names (MF, BP, CC) are exactly the gold's answer, and the quoted 'non-redundant and share a common space of identifiers and a well-specified syntax' plus 'terms linked by relations into a hierarchical vocabulary' are verbatim gold phrases.",
  "completeness": "All three ontologies named plus the gold's own qualifiers \u2014 fully answers the question.",
  "groundedness": "Everything is either the requested list or a direct quote from the gold passage; no unsupported content.",
  "sources": "Attributed to [S1] with quoted phrases, an accurate pointer to the gold passage."
}
\end{EvidenceText}

\textbf{Judge summary}
\begin{EvidenceText}
Concise, fully correct answer quoting the gold's exact wording with a valid source tag.
\end{EvidenceText}

\textbf{Complete answer}
\begin{AnswerText}
The three ontologies are **Molecular Function (MF)**, **Biological Process (BP)**, and **Cellular Component (CC)** \u2014 "non-redundant and share a common space of identifiers and a well-specified syntax," with terms linked by relations into a hierarchical vocabulary [S1].
\end{AnswerText}

\textbf{Retrieved evidence index (all returned units; paths and scores) }
\begin{EvidenceText}
[1] score=0.6548 path=\u751f\u547d\u79d1\u5b66\u4e0e\u533b\u5b66/genomics__1602.01876__primer-on-the-gene-ontology.md (part 1 of 2).md
[2] score=0.5409 path=\u751f\u547d\u79d1\u5b66\u4e0e\u533b\u5b66/genomics__1602.01876__primer-on-the-gene-ontology.md (part 2 of 2).md
[3] score=0.5307 path=\u751f\u547d\u79d1\u5b66\u4e0e\u533b\u5b66/genomics__1602.01876__primer-on-the-gene-ontology.md (part 2 of 2).md
[4] score=0.5298 path=\u751f\u547d\u79d1\u5b66\u4e0e\u533b\u5b66/genomics__1602.01876__primer-on-the-gene-ontology.md (part 1 of 2).md
[5] score=0.5136 path=\u751f\u547d\u79d1\u5b66\u4e0e\u533b\u5b66/genomics__1602.01876__primer-on-the-gene-ontology.md (part 1 of 2).md
[6] score=0.5107 path=\u751f\u547d\u79d1\u5b66\u4e0e\u533b\u5b66/genomics__1602.01876__primer-on-the-gene-ontology.md (part 1 of 2).md
[7] score=0.5106 path=\u751f\u547d\u79d1\u5b66\u4e0e\u533b\u5b66/genomics__2301.08436__spacetx-a-roadmap-for-benchmarking-spatial-t.md (part 2 of 2).md
[8] score=0.4974 path=\u751f\u547d\u79d1\u5b66\u4e0e\u533b\u5b66/genomics__1602.01876__primer-on-the-gene-ontology.md (part 1 of 2).md
[9] score=0.4926 path=\u751f\u547d\u79d1\u5b66\u4e0e\u533b\u5b66/genomics__1602.01876__primer-on-the-gene-ontology.md (part 1 of 2).md
[10] score=0.4924 path=\u751f\u547d\u79d1\u5b66\u4e0e\u533b\u5b66/genomics__1602.01876__primer-on-the-gene-ontology.md (part 1 of 2).md
[11] score=0.4776 path=\u751f\u547d\u79d1\u5b66\u4e0e\u533b\u5b66/genomics__1602.01876__primer-on-the-gene-ontology.md (part 1 of 2).md
[12] score=0.472 path=\u751f\u547d\u79d1\u5b66\u4e0e\u533b\u5b66/genomics__1602.01876__primer-on-the-gene-ontology.md (part 1 of 2).md
\end{EvidenceText}
\subsection{BQ05: Why does historical European economic growth contradict Unified Growth Theory?}
\begin{tabular}{@{}ll@{}}
Arm & Dense BGE-M3 baseline\\
Generator & omp / deepseek-flash\\
HTTP / success & 200 / True\\
Wall time / cost & 63.7 s / 0.024572999999999998 USD\\
Judge scores & F=5, C=5, G=5, A=5; overall=5.00\\
Verdict & correct\\
\end{tabular}

\textbf{Dimension rationale}
\begin{EvidenceText}
{
  "factual": "All claims match the gold: Europe as the region where UGT 'should be prominently displayed'; the Industrial Revolution as 'prime engine of economic growth'; Maddison (2010) data; hyperbolic growth from 'mathematical analysis'; rejection/revision principle; and the paradox that UGT is 'contradicted by the same data which were used during its development'. The reciprocal-1/GDP signature argument and the 12/30-country/Eastern Europe coverage are consistent with the document.",
  "completeness": "Four numbered contradiction arguments plus the scope statement (12 countries, 30 countries, Eastern Europe, both Maddison 2001 and 2010) and the author's broader 'diametrically opposite' story framing \u2014 fully answers the question.",
  "groundedness": "Each numbered point carries source tags with quoted phrases ('Galor's three regimes of growth have no relevance to the description, let alone to the explanation...'), consistent with the gold and the paper; nothing unsupported.",
  "sources": "Dense per-claim attributions [S1][S3][S4][S5][S6][S7][S8][S9][S11], including which dataset each claim applies to."
}
\end{EvidenceText}

\textbf{Judge summary}
\begin{EvidenceText}
Accurate, well-structured answer covering the missing transition signature, absent regimes, no-trap argument, and hyperbolic growth with per-dataset sourcing.
\end{EvidenceText}

\textbf{Complete answer}
\begin{AnswerText}
Answer: Europe is the hardest case for Unified Growth Theory (UGT), not the easiest \u2014 it should confirm the theory most strongly because Western Europe was the centre of the Industrial Revolution, the alleged "prime engine of economic growth," so any real regime mechanism would be prominently visible there [S1][S6]. Using Maddison's data (2010), Nielsen finds that instead of confirming UGT, European growth contradicts it on four counts:

1. **The predicted transition signature is absent.** UGT requires a visible boost in growth at the Industrial Revolution \u2014 a clear, prominent downward change in the trajectory of reciprocal GDP values (1/GDP). That signature is missing in the European data [S8], as shown for the 12 leading Western European countries [S5] and for all 30 Western European countries [S7].

2. **The three regimes are not observable.** Galor's three regimes of growth "have no relevance to the description, let alone to the explanation, of the mechanism of the economic growth"; there was no transition between regimes [S5][S7].

3. **The Malthusian regime did not exist, so there was no trap to escape.** The "remarkable"/"stunning" escape from the Malthusian trap did not happen because there was no trap, and the Industrial Revolution did not boost growth in Eastern Europe [S11][S12].

4. **Growth was hyperbolic, not regime-switching.** Mathematical analysis of historical growth shows global, regional, and national growths were hyperbolic [S9], so any theory contradicted by even one good dataset should be rejected or revised \u2014 and UGT is contradicted repeatedly, including by the same data used to develop it [S9].

The contradiction holds across the 12 leading countries, the total of 30 Western European countries [S4], and Eastern Europe [S11], and applies to both Maddison (2001) and Maddison (2010) [S3]. Nielsen's broader point is that the data "tell one story while the Unified Growth Theory tells another and diametrically opposite" one [S8], and that Galor's explanations are "dissociated from reality" \u2014 they may describe something else, but not the mechanism of economic growth [S4][S6].
\end{AnswerText}

\textbf{Retrieved evidence index (all returned units; paths and scores) }
\begin{EvidenceText}
[1] score=0.7615 path=\u7ecf\u6d4e\u4e0e\u793e\u4f1a/economics__1512.08067__unified-growth-theory-contradicted-by-the-econom.md.md
[2] score=0.7605 path=\u7ecf\u6d4e\u4e0e\u793e\u4f1a/economics__1512.08067__unified-growth-theory-contradicted-by-the-econom.md.md
[3] score=0.7447 path=\u7ecf\u6d4e\u4e0e\u793e\u4f1a/economics__1512.08067__unified-growth-theory-contradicted-by-the-econom.md.md
[4] score=0.7178 path=\u7ecf\u6d4e\u4e0e\u793e\u4f1a/economics__1512.08067__unified-growth-theory-contradicted-by-the-econom.md.md
[5] score=0.7084 path=\u7ecf\u6d4e\u4e0e\u793e\u4f1a/economics__1512.08067__unified-growth-theory-contradicted-by-the-econom.md.md
[6] score=0.701 path=\u7ecf\u6d4e\u4e0e\u793e\u4f1a/economics__1512.08067__unified-growth-theory-contradicted-by-the-econom.md.md
[7] score=0.7 path=\u7ecf\u6d4e\u4e0e\u793e\u4f1a/economics__1512.08067__unified-growth-theory-contradicted-by-the-econom.md.md
[8] score=0.6883 path=\u7ecf\u6d4e\u4e0e\u793e\u4f1a/economics__1512.08067__unified-growth-theory-contradicted-by-the-econom.md.md
[9] score=0.6794 path=\u7ecf\u6d4e\u4e0e\u793e\u4f1a/economics__1512.08067__unified-growth-theory-contradicted-by-the-econom.md.md
[10] score=0.679 path=\u7ecf\u6d4e\u4e0e\u793e\u4f1a/economics__1512.08067__unified-growth-theory-contradicted-by-the-econom.md.md
[11] score=0.6696 path=\u7ecf\u6d4e\u4e0e\u793e\u4f1a/economics__1512.08067__unified-growth-theory-contradicted-by-the-econom.md.md
[12] score=0.663 path=\u7ecf\u6d4e\u4e0e\u793e\u4f1a/economics__1512.08067__unified-growth-theory-contradicted-by-the-econom.md.md
\end{EvidenceText}
\subsection{BQ06: Which physical factor primarily controls the response of precipitation extremes to climate change?}
\begin{tabular}{@{}ll@{}}
Arm & Dense BGE-M3 baseline\\
Generator & omp / deepseek-flash\\
HTTP / success & 200 / True\\
Wall time / cost & 23.2 s / 0.024976200000000004 USD\\
Judge scores & F=5, C=5, G=5, A=5; overall=5.00\\
Verdict & correct\\
\end{tabular}

\textbf{Dimension rationale}
\begin{EvidenceText}
{
  "factual": "Correct throughout: the review's focus quote; the abstract's split of 'robust and well understood' thermodynamic vs developing microphysical/dynamical understanding; uncertainty 'when convection is important'; and the governing approximation P = -eps*omega*S with S(T,p) as the saturation-specific-humidity derivative \u2014 all match gold and the paper.",
  "completeness": "Names the thermodynamic (moisture) contribution as the robust factor, then explains the three-term decomposition and honestly notes that no excerpt 'crowns' moisture as exclusively primary, offering the defensible reading that moisture is the well-constrained backbone \u2014 the question is answered with appropriate calibration.",
  "groundedness": "Claims are attributed to [S1][S3][S8] with the equation reproduced; the hedge itself is grounded in the verified absence of a primacy sentence in the excerpts.",
  "sources": "Specific attributions including the equation-bearing source [S8] and abstract sources [S1][S3]."
}
\end{EvidenceText}

\textbf{Judge summary}
\begin{EvidenceText}
Answers with the thermodynamic/moisture factor, correctly reproducing the paper's decomposition and honestly calibrating the 'primary' claim to its sources.
\end{EvidenceText}

\textbf{Complete answer}
\begin{AnswerText}
Short answer: the **thermodynamic (moisture) contribution \u2014 the increase in atmospheric water vapor with warming** \u2014 is the robust, well-understood factor, but the excerpts do not explicitly crown it "the primary" control, and the paper's theoretical framing makes the response a product of several factors.

Evidence:

- The review's stated focus is "the physical factors that control the response," and the abstract/executive summary separates three: a **thermodynamic contribution** that is "robust and well understood," versus **microphysical and dynamical contributions** whose "theoretical understanding \u2026 is still being developed"; sensitivity "remains uncertain when convection is important" [S1, S3].
- The governing approximation makes the dependence explicit:  
  $$P \simeq -\epsilon\,\{\omega(p)\,S(T,p)\}$$  
  where $S(T,p)=dq_s/dp|_{\theta_e^*}$ is the derivative of saturation specific humidity (the thermodynamic/moisture factor), $\omega$ the vertical velocity in pressure coordinates (dynamics), and $\epsilon$ the precipitation efficiency (microphysics) [S8].
- So the intensity response is controlled by changes in moisture $S$, vertical velocity $\omega$, and efficiency $\epsilon$ \u2014 with only the moisture term described as robustly understood [S3, S8].

Caveat: the excerpts contain no sentence identifying one factor as exclusively "primarily" responsible; claiming moisture alone would overstate the source. The defensible reading is that moisture/thermodynamics is the well-constrained backbone of the response, while dynamical and microphysical changes explain the residual spread (notably where convection matters) [S1, S3].
\end{AnswerText}

\textbf{Retrieved evidence index (all returned units; paths and scores) }
\begin{EvidenceText}
[1] score=0.7273 path=\u81ea\u7136\u79d1\u5b66\u4e0e\u5730\u7403\u79d1\u5b66/climate-science__1503.07557__precipitation-extremes-under-climate-change.md (part 1 of 3).md
[2] score=0.6831 path=\u81ea\u7136\u79d1\u5b66\u4e0e\u5730\u7403\u79d1\u5b66/climate-science__1503.07557__precipitation-extremes-under-climate-change.md (part 1 of 3).md
[3] score=0.6695 path=\u81ea\u7136\u79d1\u5b66\u4e0e\u5730\u7403\u79d1\u5b66/climate-science__1503.07557__precipitation-extremes-under-climate-change.md (part 1 of 3).md
[4] score=0.6553 path=\u81ea\u7136\u79d1\u5b66\u4e0e\u5730\u7403\u79d1\u5b66/climate-science__1503.07557__precipitation-extremes-under-climate-change.md (part 3 of 3).md
[5] score=0.6518 path=\u81ea\u7136\u79d1\u5b66\u4e0e\u5730\u7403\u79d1\u5b66/climate-science__1503.07557__precipitation-extremes-under-climate-change.md (part 1 of 3).md
[6] score=0.6516 path=\u81ea\u7136\u79d1\u5b66\u4e0e\u5730\u7403\u79d1\u5b66/climate-science__1503.07557__precipitation-extremes-under-climate-change.md (part 1 of 3).md
[7] score=0.6451 path=\u81ea\u7136\u79d1\u5b66\u4e0e\u5730\u7403\u79d1\u5b66/climate-science__1503.07557__precipitation-extremes-under-climate-change.md (part 3 of 3).md
[8] score=0.6398 path=\u81ea\u7136\u79d1\u5b66\u4e0e\u5730\u7403\u79d1\u5b66/climate-science__1503.07557__precipitation-extremes-under-climate-change.md (part 1 of 3).md
[9] score=0.6391 path=\u81ea\u7136\u79d1\u5b66\u4e0e\u5730\u7403\u79d1\u5b66/climate-science__1503.07557__precipitation-extremes-under-climate-change.md (part 2 of 3).md
[10] score=0.6323 path=\u81ea\u7136\u79d1\u5b66\u4e0e\u5730\u7403\u79d1\u5b66/climate-science__1503.07557__precipitation-extremes-under-climate-change.md (part 1 of 3).md
[11] score=0.6307 path=\u81ea\u7136\u79d1\u5b66\u4e0e\u5730\u7403\u79d1\u5b66/climate-science__1503.07557__precipitation-extremes-under-climate-change.md (part 1 of 3).md
[12] score=0.6306 path=\u81ea\u7136\u79d1\u5b66\u4e0e\u5730\u7403\u79d1\u5b66/climate-science__2007.10574__impacts-of-the-madden-julian-oscillation-on (part 1 of 2).md
\end{EvidenceText}
\subsection{BQ07: What is MameLoshnLM and what parameter scale does it have?}
\begin{tabular}{@{}ll@{}}
Arm & Dense BGE-M3 baseline\\
Generator & omp / deepseek-flash\\
HTTP / success & 200 / True\\
Wall time / cost & 13.3 s / 0.0166752 USD\\
Judge scores & F=5, C=5, G=5, A=5; overall=5.00\\
Verdict & correct\\
\end{tabular}

\textbf{Dimension rationale}
\begin{EvidenceText}
{
  "factual": "All correct: 'first open-source large language model built specifically for Yiddish'; mame-loshn as 'mother tongue'; continued pretraining of Llama-3.1-8B on Oytser with a causal LM objective; 8B scale ('open-source 8B-parameter language model') quoted from abstract and conclusion; the five same-scale baselines and the Kashes averages (62.6 vs Gemma-2 9B 57.0, Llama 3.1 8B 56.8, Qwen3 8B 54.7) are consistent across independent answers; training hyperparameters and the non-commercial license are plausible paper details.",
  "completeness": "Both parts fully answered with definition, name origin, training recipe, both companion resources, parameter scale, baselines, benchmark results, and licensing.",
  "groundedness": "Every claim carries a source tag ([S1][S2][S3][S5][S7][S8][S9][S12]) and the quotes match the paper; no unsupported content.",
  "sources": "Dense per-claim attributions, e.g. the 8B scale to [S5][S3], the baseline list to [S9], and the benchmark scores to [S7]."
}
\end{EvidenceText}

\textbf{Judge summary}
\begin{EvidenceText}
Complete, accurate and thoroughly sourced answer covering definition, 8B scale, training setup, and benchmark standing.
\end{EvidenceText}

\textbf{Complete answer}
\begin{AnswerText}
**MameLoshnLM** is the first open-source large language model built specifically for Yiddish, named after the traditional Yiddish term *mame-loshn* ("mother tongue") [S1][S5]. It is produced by continued pretraining of Llama-3.1-8B on the Yiddish corpus *Oytser*, using a causal language modeling objective [S2]. Alongside the model, the authors release two supporting resources: the Oytser pretraining corpus and the Kashes evaluation benchmark [S3][S8].

**Parameter scale: 8B.** The abstract and conclusion both describe it as an "open-source 8B-parameter language model" [S5][S3], and the experimental setup compares it against five open-weight baselines "of similar scale" \u2014 Llama 3.1 8B, Qwen3 8B, BLOOMZ 7B, Gemma-2 9B, EuroLLM 9B [S9]. On the Kashes benchmark it achieves the best average score (62.6), ahead of Gemma-2 9B (57.0), Llama 3.1 8B (56.8), and Qwen3 8B (54.7) [S7].

Training specifics: bfloat16 precision, 8-bit AdamW, learning rate $2\times10^{-5}$, cosine schedule with 2% warmup, weight decay 0.01, one epoch, max sequence length 1,024 [S2]. Released under a non-commercial license [S12].
\end{AnswerText}

\textbf{Retrieved evidence index (all returned units; paths and scores) }
\begin{EvidenceText}
[1] score=0.5682 path=\u8ba1\u7b97\u673a\u4e0e\u4eba\u5de5\u667a\u80fd/nlp__2608.05850__mameloshnlm-yiddish-language-model-and-evalu.md (part 1 of 4).md
[2] score=0.5574 path=\u8ba1\u7b97\u673a\u4e0e\u4eba\u5de5\u667a\u80fd/nlp__2608.05850__mameloshnlm-yiddish-language-model-and-evalu.md (part 1 of 4).md
[3] score=0.5541 path=\u8ba1\u7b97\u673a\u4e0e\u4eba\u5de5\u667a\u80fd/nlp__2608.05850__mameloshnlm-yiddish-language-model-and-evalu.md (part 1 of 4).md
[4] score=0.5528 path=\u8ba1\u7b97\u673a\u4e0e\u4eba\u5de5\u667a\u80fd/nlp__2608.05850__mameloshnlm-yiddish-language-model-and-evalu.md (part 1 of 4).md
[5] score=0.5362 path=\u8ba1\u7b97\u673a\u4e0e\u4eba\u5de5\u667a\u80fd/nlp__2608.05850__mameloshnlm-yiddish-language-model-and-evalu.md (part 1 of 4).md
[6] score=0.5315 path=\u8ba1\u7b97\u673a\u4e0e\u4eba\u5de5\u667a\u80fd/nlp__2608.05850__mameloshnlm-yiddish-language-model-and-evalu.md (part 4 of 4).md
[7] score=0.5256 path=\u8ba1\u7b97\u673a\u4e0e\u4eba\u5de5\u667a\u80fd/nlp__2608.05850__mameloshnlm-yiddish-language-model-and-evalu.md (part 1 of 4).md
[8] score=0.5251 path=\u8ba1\u7b97\u673a\u4e0e\u4eba\u5de5\u667a\u80fd/nlp__2608.05850__mameloshnlm-yiddish-language-model-and-evalu.md (part 1 of 4).md
[9] score=0.5199 path=\u8ba1\u7b97\u673a\u4e0e\u4eba\u5de5\u667a\u80fd/nlp__2608.05850__mameloshnlm-yiddish-language-model-and-evalu.md (part 1 of 4).md
[10] score=0.5154 path=\u8ba1\u7b97\u673a\u4e0e\u4eba\u5de5\u667a\u80fd/nlp__2608.05850__mameloshnlm-yiddish-language-model-and-evalu.md (part 1 of 4).md
[11] score=0.5135 path=\u8ba1\u7b97\u673a\u4e0e\u4eba\u5de5\u667a\u80fd/nlp__2608.05850__mameloshnlm-yiddish-language-model-and-evalu.md (part 2 of 4).md
[12] score=0.5131 path=\u8ba1\u7b97\u673a\u4e0e\u4eba\u5de5\u667a\u80fd/nlp__2608.05850__mameloshnlm-yiddish-language-model-and-evalu.md (part 2 of 4).md
\end{EvidenceText}
\subsection{BQ08: How does oxygen redox contribute to the capacity of battery cathodes?}
\begin{tabular}{@{}ll@{}}
Arm & Dense BGE-M3 baseline\\
Generator & omp / deepseek-flash\\
HTTP / success & 200 / True\\
Wall time / cost & 21.2 s / 0.02151 USD\\
Judge scores & F=5, C=5, G=5, A=5; overall=5.00\\
Verdict & correct\\
\end{tabular}

\textbf{Dimension rationale}
\begin{EvidenceText}
{
  "factual": "All correct: O2- donating electrons as reducing agents; TM cation redox setting the 'ultimate limit'; the gold's overcapacity statement ('achieved and even exceeded... charge density'); oxygen participation 'governing the useful capacity' (gold-verbatim); Li-excess disordered spinels as the first significant contribution; Li2MnO3 with Mn staying +4; reversibility as the crux with capacity fade after initial cycling (gold); O2 release and solvent decomposition failure modes; DFT as the only probing route.",
  "completeness": "Mechanism, clearest material cases (spinel, Li2MnO3), and caveats (reversibility, O2 release, flammability) \u2014 fully answers how oxygen redox contributes to capacity.",
  "groundedness": "Every claim tagged ([S1][S2][S3][S5][S6][S7][S9][S10][S12]) with quotes; the closing note honestly states its context has no quantitative capacity figures.",
  "sources": "Per-claim attributions with an explicit statement of what the context does not quantify."
}
\end{EvidenceText}

\textbf{Judge summary}
\begin{EvidenceText}
Accurate, well-sourced answer covering the extra-capacity mechanism, its clearest material cases, and reversibility caveats.
\end{EvidenceText}

\textbf{Complete answer}
\begin{AnswerText}
**Answer**

Oxygen (anionic) redox adds capacity beyond what transition-metal cations alone can supply, by having oxide anions donate electrons during charge.

**Mechanism and why it adds capacity**
- In many cathodes, O\xb2\u207b anions donate electrons during charging \u2014 effectively acting as reducing agents \u2014 thereby compensating for the absence of oxidizable transition metals [S2].
- Transition-metal cation redox normally "sets the ultimate limit" on cathode performance, accommodating reversible Li/Na insertion and removal [S5]. Oxygen redox is what lets the electrode exceed that limit: oxygen or other anions participating in redox permits greater energy storage than would otherwise be theoretically predicted [S2].
- Consistent with this, real cathodes have "achieved and even exceeded" the charge density of the transition-metal oxidation limit, delivering overcapacity; the atomic-scale mechanism of this overcapacity still requires elucidation [S10].
- Hence oxygen participation is described as governing the "useful capacity" of Li- and Na-ion cathodes [S3].

**Where the contribution is clearest**
- Oxygen redox was first found to contribute significantly to overall capacity in Li-excess disordered spinel-type cathodes, where the O redox contribution is high [S1].
- The canonical case is Li\u2082MnO\u2083: Mn stays +4 and shows no extensive redox, so the material's performance is attributed to oxygen redox activity [S7][S12]. Charge compensation in Li\u2082MnO\u2083 during delithiation drove the concept into the field's foreground [S12].

**Caveats**
- Reversibility is the crux: reversible oxygen redox in the cathode bulk offers a pathway to reversible capacity, whereas irreversible or poorly reversible processes cause the capacity fade typically seen after initial cycling [S6].
- Failure modes include oxygen removal/oxygen gas evolution and electrolyte-solvent decomposition; O\u2082 release (mainly at surfaces) is an ignition/flammability hazard, and along with large first-cycle capacity loss it limits real-world use [S6][S9].
- Direct probing of anionic redox is not possible; DFT-based computational analysis is the main approach to understanding and harnessing it [S3].

Note: the context does not quantify the capacity contribution (e.g., mAh/g figures or % of capacity from oxygen) \u2014 it only states qualitative contributions and the Li-excess spinel and Li\u2082MnO\u2083 cases [S1][S7][S12].
\end{AnswerText}

\textbf{Retrieved evidence index (all returned units; paths and scores) }
\begin{EvidenceText}
[1] score=0.7297 path=\u5de5\u7a0b\u4e0e\u80fd\u6e90/materials__2408.09942__oxygen-redox-in-battery-cathodes-a-brief-overvie.md (part 2 of 2).md
[2] score=0.7292 path=\u5de5\u7a0b\u4e0e\u80fd\u6e90/materials__2408.09942__oxygen-redox-in-battery-cathodes-a-brief-overvie.md (part 1 of 2).md
[3] score=0.7166 path=\u5de5\u7a0b\u4e0e\u80fd\u6e90/materials__2408.09942__oxygen-redox-in-battery-cathodes-a-brief-overvie.md (part 1 of 2).md
[4] score=0.6976 path=\u5de5\u7a0b\u4e0e\u80fd\u6e90/materials__2408.09942__oxygen-redox-in-battery-cathodes-a-brief-overvie.md (part 1 of 2).md
[5] score=0.6877 path=\u5de5\u7a0b\u4e0e\u80fd\u6e90/materials__2408.09942__oxygen-redox-in-battery-cathodes-a-brief-overvie.md (part 1 of 2).md
[6] score=0.6784 path=\u5de5\u7a0b\u4e0e\u80fd\u6e90/materials__2408.09942__oxygen-redox-in-battery-cathodes-a-brief-overvie.md (part 1 of 2).md
[7] score=0.6761 path=\u5de5\u7a0b\u4e0e\u80fd\u6e90/materials__2408.09942__oxygen-redox-in-battery-cathodes-a-brief-overvie.md (part 1 of 2).md
[8] score=0.6757 path=\u5de5\u7a0b\u4e0e\u80fd\u6e90/materials__2408.09942__oxygen-redox-in-battery-cathodes-a-brief-overvie.md (part 1 of 2).md
[9] score=0.6752 path=\u5de5\u7a0b\u4e0e\u80fd\u6e90/materials__2408.09942__oxygen-redox-in-battery-cathodes-a-brief-overvie.md (part 1 of 2).md
[10] score=0.6618 path=\u5de5\u7a0b\u4e0e\u80fd\u6e90/materials__2408.09942__oxygen-redox-in-battery-cathodes-a-brief-overvie.md (part 1 of 2).md
[11] score=0.6613 path=\u5de5\u7a0b\u4e0e\u80fd\u6e90/materials__2408.09942__oxygen-redox-in-battery-cathodes-a-brief-overvie.md (part 1 of 2).md
[12] score=0.6609 path=\u5de5\u7a0b\u4e0e\u80fd\u6e90/materials__2408.09942__oxygen-redox-in-battery-cathodes-a-brief-overvie.md (part 1 of 2).md
\end{EvidenceText}
\subsection{BQ09: What problem of visual SLAM in dynamic environments does MLP-SLAM address and how?}
\begin{tabular}{@{}ll@{}}
Arm & Dense BGE-M3 baseline\\
Generator & omp / deepseek-flash\\
HTTP / success & 200 / True\\
Wall time / cost & 21.2 s / 0.021453 USD\\
Judge scores & F=5, C=5, G=5, A=5; overall=5.00\\
Verdict & correct\\
\end{tabular}

\textbf{Dimension rationale}
\begin{EvidenceText}
{
  "factual": "All correct: the static-world assumption of SOTA systems and degradation with movable objects (gold); DynaSLAM discarding in-mask features with misclassification and MaskRCNN latency; the learned binary discriminator replacing semantic-mask gating inside a real-time stereo ORB-SLAM2-based system; the first-MLP-discriminator claim (gold-verbatim); KITTI-only training generalizing to KITTI360; the released 50,000+ point dataset; the synthetic moving-car and KITTI360 checks indicating genuine dynamic-object recognition; stereo-only scope with claimed extendability.",
  "completeness": "Problem, mechanism, two supporting gaps (dataset; genuine-dynamic-object evidence), reported outcomes, and scope limits \u2014 fully answered, with honest notes on what its excerpts lack (input vector, hyperparameters, truncated table figures).",
  "groundedness": "Every claim tagged ([S1][S2][S3][S4][S5][S7][S8][S9][S10][S11][S12]); the declared gaps match its excerpt limits.",
  "sources": "Per-claim attributions including which sources support the prior-work criticisms."
}
\end{EvidenceText}

\textbf{Judge summary}
\begin{EvidenceText}
Complete, well-sourced answer linking the static-world problem to the learned discriminator solution, plus dataset and scope caveats.
\end{EvidenceText}

\textbf{Complete answer}
\begin{AnswerText}
**Problem.** MLP-SLAM targets the static-world assumption underlying V-SLAM: most systems \u2014 including SOTA ones such as ORB-SLAM2, ORB-SLAM3, LSD-SLAM \u2014 assume a static world, so accuracy holds up only in scenes with few dynamic objects and degrades where movable objects are numerous (pedestrians, cars, buses, typical of outdoor driving scenes) [S7][S1]. Prior dynamic-SLAM work (e.g., DynaSLAM) discards feature points inside semantic masks of potentially moving objects, which can misclassify dynamic vs. static objects and incurs high latency from the MaskRCNN dependency [S9].

**How.** The core idea is to replace semantic-mask gating with a learned binary discriminator: a multilayer perceptron classifies each feature point as dynamic or static, and dynamic points are excluded from tracking, inside a real-time stereo SLAM system built on ORB-SLAM2 [S8][S3][S2]. Per the authors, this is the first use of an MLP as a dynamic/static feature-point discriminator in SLAM [S8]. The MLP was trained only on KITTI Odometry, without any KITTI360 information, yet still generalized on the KITTI360 comparison [S3].

**Two supporting gaps it also fills.** (i) No public dataset existed for directly evaluating dynamic/static feature classification; prior evaluation was indirect via SLAM metrics and unreliable because of RANSAC randomness in pose estimation \u2014 so they released a dataset of 50,000+ feature points [S11][S4][S2]. (ii) Both a synthetic "moving car" sequence and KITTI360 checks indicate the model identifies genuinely dynamic objects rather than memorized semantics [S5].

**Reported outcome.** Highest average precision on KITTI Odometry and the fastest speed among the compared dynamic SLAM systems, with the discriminator outperforming other methods on their dataset [S10][S12][S4]. Scope limits stated in the paper: experiments are stereo-only, though the method is claimed extendable to monocular/RGB-D [S5].

**Not in context:** the MLP's input feature vector, layer sizes/training hyperparameters, the exact dynamic-point rejection policy inside the tracker, and the quantitative margins in Tables II/III (figures truncated).
\end{AnswerText}

\textbf{Retrieved evidence index (all returned units; paths and scores) }
\begin{EvidenceText}
[1] score=0.6879 path=\u8ba1\u7b97\u673a\u4e0e\u4eba\u5de5\u667a\u80fd/robotics__2410.10669__mlp-slam-multilayer-perceptron-based-simultaneou.md (part 1 of 2).md
[2] score=0.6821 path=\u8ba1\u7b97\u673a\u4e0e\u4eba\u5de5\u667a\u80fd/robotics__2410.10669__mlp-slam-multilayer-perceptron-based-simultaneou.md (part 2 of 2).md
[3] score=0.6606 path=\u8ba1\u7b97\u673a\u4e0e\u4eba\u5de5\u667a\u80fd/robotics__2410.10669__mlp-slam-multilayer-perceptron-based-simultaneou.md (part 2 of 2).md
[4] score=0.6549 path=\u8ba1\u7b97\u673a\u4e0e\u4eba\u5de5\u667a\u80fd/robotics__2410.10669__mlp-slam-multilayer-perceptron-based-simultaneou.md (part 1 of 2).md
[5] score=0.6549 path=\u8ba1\u7b97\u673a\u4e0e\u4eba\u5de5\u667a\u80fd/robotics__2410.10669__mlp-slam-multilayer-perceptron-based-simultaneou.md (part 2 of 2).md
[6] score=0.6528 path=\u8ba1\u7b97\u673a\u4e0e\u4eba\u5de5\u667a\u80fd/robotics__2509.14191__mcgs-slam-a-multi-camera-slam-framework-usin.md (part 1 of 2).md
[7] score=0.6505 path=\u8ba1\u7b97\u673a\u4e0e\u4eba\u5de5\u667a\u80fd/robotics__2410.10669__mlp-slam-multilayer-perceptron-based-simultaneou.md (part 1 of 2).md
[8] score=0.6453 path=\u8ba1\u7b97\u673a\u4e0e\u4eba\u5de5\u667a\u80fd/robotics__2410.10669__mlp-slam-multilayer-perceptron-based-simultaneou.md (part 1 of 2).md
[9] score=0.6368 path=\u8ba1\u7b97\u673a\u4e0e\u4eba\u5de5\u667a\u80fd/robotics__2410.10669__mlp-slam-multilayer-perceptron-based-simultaneou.md (part 1 of 2).md
[10] score=0.6356 path=\u8ba1\u7b97\u673a\u4e0e\u4eba\u5de5\u667a\u80fd/robotics__2410.10669__mlp-slam-multilayer-perceptron-based-simultaneou.md (part 2 of 2).md
[11] score=0.6335 path=\u8ba1\u7b97\u673a\u4e0e\u4eba\u5de5\u667a\u80fd/robotics__2410.10669__mlp-slam-multilayer-perceptron-based-simultaneou.md (part 1 of 2).md
[12] score=0.6284 path=\u8ba1\u7b97\u673a\u4e0e\u4eba\u5de5\u667a\u80fd/robotics__2410.10669__mlp-slam-multilayer-perceptron-based-simultaneou.md (part 1 of 2).md
\end{EvidenceText}
\subsection{BQ10: What strategies for tumor control emerge under immune evasion and chemotherapy resistance?}
\begin{tabular}{@{}ll@{}}
Arm & Dense BGE-M3 baseline\\
Generator & omp / deepseek-flash\\
HTTP / success & 200 / True\\
Wall time / cost & 33.3 s / 0.0219732 USD\\
Judge scores & F=5, C=4, G=5, A=5; overall=4.75\\
Verdict & correct\\
\end{tabular}

\textbf{Dimension rationale}
\begin{EvidenceText}
{
  "factual": "All claims correct: combination therapy enlarging the tumor-control parameter region (gold-verbatim); immunotherapy enhancing immune-mediated suppression; nonlinear tipping between immune control and escape; efficacy depending on evasion mechanism and dominant phenotype (gold-aligned); checkpoint-regulation dominance making control highly sensitive to the immune exhaustion rate, pointing to checkpoint inhibition (truncation honestly noted); the sequencing rules for T2-dominant and coexisting-resistance tumors; the T1-mutant resurgence caveat; threshold conditions and competition/cooperation modeling (abstract); future-work deferral of schedule optimization.",
  "completeness": "Covers combination, checkpoint-branch mechanism matching, sequencing, and thresholds, but omits the reduced-antigen-presentation branch targets (up rho, down mu with CAR-T/vaccines; adoptive T-cell transfer focus) and does not mention adaptive therapy explicitly.",
  "groundedness": "Every claim tagged ([S2][S5][S6][S10][S11][S12]) with quotes and truncation disclosure; nothing unsupported.",
  "sources": "Clear per-claim attributions including the future-work passage."
}
\end{EvidenceText}

\textbf{Judge summary}
\begin{EvidenceText}
Accurate, well-sourced answer covering combination, checkpoint-matched, and sequential strategies, though thinner on the antigen-presentation branch.
\end{EvidenceText}

\textbf{Complete answer}
\begin{AnswerText}
**Combination therapy is the core lever.** Immune dynamics alone are generally insufficient to fully eliminate a tumor; adding chemotherapy alongside immunotherapy significantly enlarges the parameter region corresponding to tumor control or tumor-free equilibria, because chemo reduces burden and directly targets resistant subpopulations [S5], while immunotherapy enhances immune-mediated suppression [S11]. Together they expand the region of tumor-free/low-population equilibria, and modest parameter shifts can tip the system qualitatively between immune control and escape \u2014 the dynamics are explicitly nonlinear [S11].

**Efficacy is phenotype- and mechanism-dependent, not universal.** Bifurcation analyses show therapeutic efficacy depends on both the immune-evasion mechanism and the dominant tumor phenotype [S2]. Two resistance strategies were modeled \u2014 immune checkpoint regulation and reduced antigen presentation \u2014 and when the checkpoint-regulating immune-resistant subpopulation dominates, tumor control becomes highly sensitive to the immune exhaustion rate, which points toward checkpoint inhibition [S2] (the sentence is truncated in the provided excerpt, so the full claim cannot be verified).

**Sequencing rules that emerge** [S10]:
- **T\u2082-dominant tumors:** chemotherapy may reduce burden but can also drive resistant escape; immunotherapy is therefore typically applied *afterward* to stabilize the system.
- **Predicted coexistence of both resistance types:** apply chemotherapy first to eliminate T\u2081 cells, leaving a low T\u2082-only equilibrium, then apply immunotherapy to potentially eliminate the tumor.
- **Caveat:** if the tumor is only *controlled* rather than eliminated, emergence of T\u2081 mutants can cause resurgence \u2014 control is not durable in itself.

**Threshold conditions** govern persistence, elimination, and phenotype dominance as therapeutic intensity varies, and the models include competition and cooperation between resistant and sensitive subpopulations \u2014 the stated purpose being a theoretical framework for targeted and combination therapy design [S6].

**Not in the provided context:** specific numerical thresholds, doses, or schedules. The paper defers treatment-schedule optimization (minimizing burden while reducing toxicity) to future work via optimal control methods, along with alternative immune-evasion strategies [S12].
\end{AnswerText}

\textbf{Retrieved evidence index (all returned units; paths and scores) }
\begin{EvidenceText}
[1] score=0.7545 path=\u751f\u547d\u79d1\u5b66\u4e0e\u533b\u5b66/immunology__2604.01385__strategies-for-tumor-elimination-and-control.md (part 1 of 3).md
[2] score=0.7277 path=\u751f\u547d\u79d1\u5b66\u4e0e\u533b\u5b66/immunology__2604.01385__strategies-for-tumor-elimination-and-control.md (part 3 of 3).md
[3] score=0.7167 path=\u751f\u547d\u79d1\u5b66\u4e0e\u533b\u5b66/immunology__2604.01385__strategies-for-tumor-elimination-and-control.md (part 3 of 3).md
[4] score=0.716 path=\u751f\u547d\u79d1\u5b66\u4e0e\u533b\u5b66/immunology__2604.01385__strategies-for-tumor-elimination-and-control.md (part 2 of 3).md
[5] score=0.6776 path=\u751f\u547d\u79d1\u5b66\u4e0e\u533b\u5b66/immunology__2604.01385__strategies-for-tumor-elimination-and-control.md (part 3 of 3).md
[6] score=0.6746 path=\u751f\u547d\u79d1\u5b66\u4e0e\u533b\u5b66/immunology__2604.01385__strategies-for-tumor-elimination-and-control.md (part 1 of 3).md
[7] score=0.6737 path=\u751f\u547d\u79d1\u5b66\u4e0e\u533b\u5b66/immunology__2604.01385__strategies-for-tumor-elimination-and-control.md (part 1 of 3).md
[8] score=0.6688 path=\u751f\u547d\u79d1\u5b66\u4e0e\u533b\u5b66/immunology__2604.01385__strategies-for-tumor-elimination-and-control.md (part 1 of 3).md
[9] score=0.6617 path=\u751f\u547d\u79d1\u5b66\u4e0e\u533b\u5b66/immunology__2604.01385__strategies-for-tumor-elimination-and-control.md (part 1 of 3).md
[10] score=0.6608 path=\u751f\u547d\u79d1\u5b66\u4e0e\u533b\u5b66/immunology__2604.01385__strategies-for-tumor-elimination-and-control.md (part 3 of 3).md
[11] score=0.6491 path=\u751f\u547d\u79d1\u5b66\u4e0e\u533b\u5b66/immunology__2604.01385__strategies-for-tumor-elimination-and-control.md (part 3 of 3).md
[12] score=0.6463 path=\u751f\u547d\u79d1\u5b66\u4e0e\u533b\u5b66/immunology__2604.01385__strategies-for-tumor-elimination-and-control.md (part 3 of 3).md
\end{EvidenceText}
\section{RRF baseline}
\subsection{BQ01: How does the Transformer compute attention, and why does it replace recurrence and convolution?}
\begin{tabular}{@{}ll@{}}
Arm & RRF baseline\\
Generator & omp / deepseek-flash\\
HTTP / success & 200 / True\\
Wall time / cost & 27.7 s / 0.025184400000000003 USD\\
Judge scores & F=5, C=3, G=5, A=5; overall=4.50\\
Verdict & partial\\
\end{tabular}

\textbf{Dimension rationale}
\begin{EvidenceText}
{
  "factual": "Every claim matches the GOLD excerpt: self-attention definition; 'the first transduction model relying entirely on self-attention'; 'global dependencies between input and output' regardless of distance; the encoder-decoder connection itself being an attention mechanism; ConvS2S linear / ByteNet logarithmic vs constant operations; 28.4 BLEU (over 2 above prior best including ensembles) and 41.8 BLEU after 3.5 days on eight GPUs, twelve hours on eight P100 GPUs. No errors.",
  "completeness": "The 'why' half is fully answered in three points mirroring gold (constant path length, parallelization, empirical payoff). The 'how' half is an honest refusal: it states the excerpts cover only abstract/introduction/Background and that the formula, multi-head concatenation and positional encodings live in absent section 3.1-3.5, offering only definitional hints \u2014 similar coverage to System U.",
  "groundedness": "All claims are pinned to [S1]-[S4], and it explicitly states '[S5]-[S12] are unrelated documents... and were not used', so no content from irrelevant sources leaked in; nothing unsupported.",
  "sources": "Per-claim attributions plus an explicit disclosure of which sources were discarded and why (their topics are named). Clear and accurate."
}
\end{EvidenceText}

\textbf{Judge summary}
\begin{EvidenceText}
With the formula absent from its excerpts it answers honestly: the motivation half is complete and gold-accurate, the computation half is a declared no-answer with all available hints given.
\end{EvidenceText}

\textbf{Complete answer}
\begin{AnswerText}
**The exact attention computation is not in the provided context.** The excerpts come from the paper's abstract, introduction, and \xa72 Background \u2014 the scaled dot-product formula ($\mathrm{softmax}(QK^\top/\sqrt{d_k})V$), multi-head concatenation, and positional encodings live in \xa73.1\u20133.5, which are absent here. Below is only what the supplied sources support.

## How it computes attention (as far as the context states)

- Self-attention ("intra-attention") is "an attention mechanism relating different positions of a single sequence in order to compute a representation of the sequence" [S4].
- The Transformer is "the first transduction model relying entirely on self-attention to compute representations of its input and output without using sequence-aligned RNNs or convolution" [S1].
- Attention draws "global dependencies between input and output" directly, without regard to distance in the sequences [S2].
- The encoder\u2013decoder connection is itself an attention mechanism [S3].
- The averaging over attention-weighted positions reduces effective resolution; the paper counteracts this "with Multi-Head Attention as described in section 3.2" [S4] \u2014 the mechanism's details are outside the provided context.

## Why it replaces recurrence and convolution

1. **Constant path length between any two positions.** In prior convolutional seq2seq models (ConvS2S, ByteNet, Extended Neural GPU) the number of operations to relate two arbitrary positions grows linearly (ConvS2S) or logarithmically (ByteNet) with distance, "which makes it more difficult to learn dependencies between distant positions." In the Transformer this is reduced to a **constant** number of operations [S4].
2. **Parallelization / less sequential computation.** Existing attention mechanisms were almost always "used in conjunction with a recurrent network" [S2]; eschewing recurrence removes the sequential bottleneck, so the model "allows for significantly more parallelization" [S2]. Reducing sequential computation is the same motivation behind the convolutional alternatives [S2].
3. **Empirical payoff.** Superior translation quality while being more parallelizable and requiring significantly less training time: 28.4 BLEU EN\u2192DE (over 2 BLEU above the prior best, including ensembles), and 41.8 BLEU EN\u2192FR single-model SOTA after 3.5 days on eight GPUs; base models reached SOTA "after being trained for as little as twelve hours on eight P100 GPUs" [S2, S3].

Note: [S5]\u2013[S12] are unrelated documents (long-context LMs, fact-learning dynamics, BERT, quantum computing, sketch-based query optimization) and were not used.
\end{AnswerText}

\textbf{Retrieved evidence index (all returned units; paths and scores) }
\begin{EvidenceText}
[1] score=30.13621249349545 path=\u8ba1\u7b97\u673a\u4e0e\u4eba\u5de5\u667a\u80fd\artificial-intelligence__1706.03762__attention-is-all-you-need.md (part 1 of 2).md
[2] score=19.728257596864477 path=\u8ba1\u7b97\u673a\u4e0e\u4eba\u5de5\u667a\u80fd\artificial-intelligence__1706.03762__attention-is-all-you-need.md (part 1 of 2).md
[3] score=18.119504420281313 path=\u8ba1\u7b97\u673a\u4e0e\u4eba\u5de5\u667a\u80fd\artificial-intelligence__1706.03762__attention-is-all-you-need.md (part 1 of 2).md
[4] score=16.818573817572688 path=\u8ba1\u7b97\u673a\u4e0e\u4eba\u5de5\u667a\u80fd\artificial-intelligence__1706.03762__attention-is-all-you-need.md (part 1 of 2).md
[5] score=19.190292047283183 path=\u8ba1\u7b97\u673a\u4e0e\u4eba\u5de5\u667a\u80fd\Lost in the Middle (part 2 of 3).md
[6] score=19.021079184736546 path=\u8ba1\u7b97\u673a\u4e0e\u4eba\u5de5\u667a\u80fd\nlp__2503.21676__how-do-language-models-learn-facts-dynamics.md (part 1 of 4).md
[7] score=18.185711987792175 path=\u8ba1\u7b97\u673a\u4e0e\u4eba\u5de5\u667a\u80fd\nlp__2503.21676__how-do-language-models-learn-facts-dynamics.md (part 1 of 4).md
[8] score=17.58193867368475 path=\u8ba1\u7b97\u673a\u4e0e\u4eba\u5de5\u667a\u80fd\nlp__2503.21676__how-do-language-models-learn-facts-dynamics.md (part 2 of 4).md
[9] score=17.32819837159615 path=\u8ba1\u7b97\u673a\u4e0e\u4eba\u5de5\u667a\u80fd\nlp__1810.04805__bert-pre-training-of-deep-bidirectional-transfor.md (part 2 of 3).md
[10] score=17.060270156989258 path=\u81ea\u7136\u79d1\u5b66\u4e0e\u5730\u7403\u79d1\u5b66\quantum-physics__1801.00862__quantum-computing-in-the-nisq-era-and-beyond.md (part 1 of 3).md
[11] score=16.38215397387708 path=\u8ba1\u7b97\u673a\u4e0e\u4eba\u5de5\u667a\u80fd\databases__2102.02440__online-sketch-based-query-optimization.md (part 2 of 5).md
[12] score=16.30945350168588 path=\u8ba1\u7b97\u673a\u4e0e\u4eba\u5de5\u667a\u80fd\databases__2102.02440__online-sketch-based-query-optimization.md (part 3 of 5).md
\end{EvidenceText}
\subsection{BQ02: What defines NISQ technology and what is its central limitation?}
\begin{tabular}{@{}ll@{}}
Arm & RRF baseline\\
Generator & omp / deepseek-flash\\
HTTP / success & 200 / True\\
Wall time / cost & 18.8 s / 0.0217368 USD\\
Judge scores & F=5, C=5, G=5, A=5; overall=5.00\\
Verdict & correct\\
\end{tabular}

\textbf{Dimension rationale}
\begin{EvidenceText}
{
  "factual": "All claims correct: intermediate-scale = 50 to a few hundred qubits exceeding brute-force supercomputer simulation; noisy = imperfect control; the abstract's programmatic statement and the many-body-physics use case; the verbatim 1000-gate passage with 'ceiling on the computational power'; QEC heavy overhead; noise's serious near-term limitations; hard-to-validate outputs especially for quantum simulation.",
  "completeness": "Both parts answered completely: the definition via the acronym's two components plus the abstract framing and a corroborating benchmarking source; the central limitation quantified (1000-gate ceiling) with the QEC-overhead escape route and two corollary limits.",
  "groundedness": "Every claim is sourced ([S4][S5][S1][S12][S8][S6]) with quotes matching the paper; the 'Not in context' section honestly declines to supply a formal error-rate threshold or any quantitative definition beyond the 50-few-hundred range.",
  "sources": "Clear per-claim attributions with the closing note specifying exactly what its sources do not contain."
}
\end{EvidenceText}

\textbf{Judge summary}
\begin{EvidenceText}
Complete and accurate on definition and limitation, with the verbatim 1000-gate quote, QEC-overhead escape route, and honest disclosure of what its context lacks.
\end{EvidenceText}

\textbf{Complete answer}
\begin{AnswerText}
**Definition (Noisy Intermediate-Scale Quantum).** NISQ names devices that are *intermediate-scale* \u2014 roughly 50 to a few hundred qubits, a range that already exceeds what the most powerful classical supercomputers can simulate by brute force \u2014 and *noisy*, meaning we will have imperfect control over those qubits [S4]. The paper's abstract states the same programmatically: quantum computers with 50\u2013100 qubits may perform tasks surpassing classical digital computers, and NISQ devices will be useful for exploring many-body quantum physics, possibly with other applications [S5]. The acronym is defined simply as "NISQ = noisy intermediate-scale quantum," an era the field is just entering [S1]. A benchmarking paper confirms the characterization as "various obstacles including restrictive qubit topologies, short coherence times and imperfect quantum gates" [S12].

**Central limitation.** Noise in quantum gates limits the size of circuits that can be executed reliably [S5]. Concretely: "with these noisy devices we don't expect to be able to execute a circuit that contains many more than about 1000 gates \u2014 that is, 1000 fundamental two-qubit operations \u2014 because the noise will overwhelm the signal in a circuit much larger than that. That limitation on circuit size imposes a ceiling on the computational power of NISQ technology" [S8]. Scaling past that ceiling requires quantum error correction, which carries a heavy overhead in qubit and gate counts [S8].

Two corollary limits in the context: noise "will place serious limitations on what quantum devices can achieve in the near term" [S4], and because of imperfect NISQ performance it "may be hard to validate that a quantum computer is really giving the right answer," especially for quantum-simulation problems [S6].

**Not in context:** a formal error-rate threshold for NISQ operation, or a quantitative definition of "intermediate scale" beyond the 50\u2013few-hundred qubit range given [S4][S5].
\end{AnswerText}

\textbf{Retrieved evidence index (all returned units; paths and scores) }
\begin{EvidenceText}
[1] score=23.317518447884584 path=\u81ea\u7136\u79d1\u5b66\u4e0e\u5730\u7403\u79d1\u5b66\quantum-physics__1801.00862__quantum-computing-in-the-nisq-era-and-beyond.md (part 2 of 3).md
[2] score=18.00095221826477 path=\u81ea\u7136\u79d1\u5b66\u4e0e\u5730\u7403\u79d1\u5b66\quantum-physics__1801.00862__quantum-computing-in-the-nisq-era-and-beyond.md (part 2 of 3).md
[3] score=14.37198914525874 path=\u81ea\u7136\u79d1\u5b66\u4e0e\u5730\u7403\u79d1\u5b66\quantum-physics__1801.00862__quantum-computing-in-the-nisq-era-and-beyond.md (part 2 of 3).md
[4] score=22.152120106941986 path=\u81ea\u7136\u79d1\u5b66\u4e0e\u5730\u7403\u79d1\u5b66\quantum-physics__1801.00862__quantum-computing-in-the-nisq-era-and-beyond.md (part 1 of 3).md
[5] score=17.008760307173624 path=\u81ea\u7136\u79d1\u5b66\u4e0e\u5730\u7403\u79d1\u5b66\quantum-physics__1801.00862__quantum-computing-in-the-nisq-era-and-beyond.md (part 1 of 3).md
[6] score=15.462256339435145 path=\u81ea\u7136\u79d1\u5b66\u4e0e\u5730\u7403\u79d1\u5b66\quantum-physics__1801.00862__quantum-computing-in-the-nisq-era-and-beyond.md (part 1 of 3).md
[7] score=15.450812648694402 path=\u81ea\u7136\u79d1\u5b66\u4e0e\u5730\u7403\u79d1\u5b66\quantum-physics__1801.00862__quantum-computing-in-the-nisq-era-and-beyond.md (part 1 of 3).md
[8] score=15.41781039169155 path=\u81ea\u7136\u79d1\u5b66\u4e0e\u5730\u7403\u79d1\u5b66\quantum-physics__1801.00862__quantum-computing-in-the-nisq-era-and-beyond.md (part 1 of 3).md
[9] score=16.624529482846405 path=\u8ba1\u7b97\u673a\u4e0e\u4eba\u5de5\u667a\u80fd\large-language-models__2005.14165__language-models-are-few-shot-learners.md (part 4 of 13).md
[10] score=15.353363415089682 path=\u8ba1\u7b97\u673a\u4e0e\u4eba\u5de5\u667a\u80fd\Chain-of-Thought Prompting (part 2 of 3).md
[11] score=14.777755309772447 path=\u7ecf\u6d4e\u4e0e\u793e\u4f1a\economics__2011.14424__on-the-effectiveness-of-the-european-central (part 2 of 5).md
[12] score=14.341408277727243 path=\u81ea\u7136\u79d1\u5b66\u4e0e\u5730\u7403\u79d1\u5b66\quantum-physics__2303.00445__benchmarking-noisy-intermediate-scale-quantu (part 1 of 4).md
\end{EvidenceText}
\subsection{BQ03: What is MultiMedQA and how did Flan-PaLM perform on the MedQA benchmark?}
\begin{tabular}{@{}ll@{}}
Arm & RRF baseline\\
Generator & omp / deepseek-flash\\
HTTP / success & 200 / True\\
Wall time / cost & 15.3 s / 0.0230232 USD\\
Judge scores & F=5, C=5, G=5, A=5; overall=5.00\\
Verdict & correct\\
\end{tabular}

\textbf{Dimension rationale}
\begin{EvidenceText}
{
  "factual": "All claims verified: seven-dataset composition with HealthSearchQA's 3,375 free-text open-domain questions; motivation beyond classification accuracy/BLEU probing factuality, expert reasoning, helpfulness, precision, health equity, harm (gold-aligned); 67.6% with self-consistency, 60.3% few-shot and CoT with no CoT gain; DRAGON +20.1% and 17.3% over the then-best; 62.0% on 5-option; full scaling ladder 25.7/35.4/40.9/46.1/58.9/60.3; the 'not exhaustive, plan to expand' admission is a real paper statement.",
  "completeness": "Both parts answered exhaustively: definition (composition, spans, purpose, stated limitations) and performance (headline, prompting breakdown, both baselines, 5-option set, full scaling ladder, SOTA claim, and the human-eval gap leading to Med-PaLM).",
  "groundedness": "Every figure and claim is attached to a source id ([S2][S3][S4][S5][S6][S7][S8][S9][S11][S1]) and matches gold or the paper; no invention detected.",
  "sources": "Dense per-claim attributions, e.g. scaling numbers to [S7], the CoT-no-gain observation to [S6][S8][S7], and the two SOTA comparisons to [S11][S9][S1]."
}
\end{EvidenceText}

\textbf{Judge summary}
\begin{EvidenceText}
Fully covers both halves with every number correct and individually sourced, including the benchmark's self-acknowledged limits and the scale ablation.
\end{EvidenceText}

\textbf{Complete answer}
\begin{AnswerText}
**MultiMedQA** is a medical question-answering benchmark curated by the authors, spanning medical exam questions, medical research questions, and consumer medical questions. It comprises seven datasets: six existing ones \u2014 MedQA, MedMCQA, PubMedQA, LiveQA, MedicationQA, and MMLU clinical topics \u2014 plus HealthSearchQA (newly introduced; 3,375 commonly searched health questions in free-text, open-domain format) [S4][S3][S5]. Its purpose was to broaden assessment beyond classification accuracy or BLEU-style metrics, probing factuality, expert-knowledge use in reasoning, helpfulness, precision, health equity, and potential harm [S4]. The authors acknowledge it is not exhaustive and plan to expand it [S5].

**Flan-PaLM on MedQA (USMLE-style, 4 options):**
- 67.6% MCQ accuracy with self-consistency prompting [S2][S6][S11]
- 60.3% with plain few-shot prompting (and also 60.3% with CoT \u2014 CoT gave no improvement over few-shot on MedQA/MedMCQA/PubMedQA) [S6][S8][S7]
- Surpassed prior SOTA (DRAGON) by 20.1%, and beat the then-best prior result by 17.3% overall \u2014 hence the "over 17% improvement over previous SOTA" claim [S11][S9][S1]
- On the harder 5-option MedQA set, the model scored 62.0% [S9]
- Scaling: MedQA 4-option few-shot accuracies were 25.7 (PaLM 8B), 35.4 (Flan-PaLM 8B), 40.9 (PaLM 62B), 46.1 (Flan-PaLM 62B), 58.9 (PaLM 540B), 60.3 (Flan-PaLM 540B) [S7]

This made Flan-PaLM reach SOTA on every MultiMedQA multiple-choice dataset (MedQA, MedMCQA, PubMedQA, MMLU clinical topics) [S2][S3], though human evaluation still revealed gaps in its free-text consumer-question answers \u2014 motivating instruction prompt tuning and Med-PaLM [S2][S1].
\end{AnswerText}

\textbf{Retrieved evidence index (all returned units; paths and scores) }
\begin{EvidenceText}
[1] score=32.94560489380912 path=\u751f\u547d\u79d1\u5b66\u4e0e\u533b\u5b66\clinical-medicine__2212.13138__large-language-models-encode-clinical-knowledge.md (part 1 of 8).md
[2] score=28.715013475462307 path=\u751f\u547d\u79d1\u5b66\u4e0e\u533b\u5b66\clinical-medicine__2212.13138__large-language-models-encode-clinical-knowledge.md (part 1 of 8).md
[3] score=26.689646920810887 path=\u751f\u547d\u79d1\u5b66\u4e0e\u533b\u5b66\clinical-medicine__2212.13138__large-language-models-encode-clinical-knowledge.md (part 1 of 8).md
[4] score=18.992863866293252 path=\u751f\u547d\u79d1\u5b66\u4e0e\u533b\u5b66\clinical-medicine__2212.13138__large-language-models-encode-clinical-knowledge.md (part 1 of 8).md
[5] score=18.223483661618555 path=\u751f\u547d\u79d1\u5b66\u4e0e\u533b\u5b66\clinical-medicine__2212.13138__large-language-models-encode-clinical-knowledge.md (part 1 of 8).md
[6] score=24.48649383199579 path=\u751f\u547d\u79d1\u5b66\u4e0e\u533b\u5b66\clinical-medicine__2212.13138__large-language-models-encode-clinical-knowledge.md (part 2 of 8).md
[7] score=20.83136277197376 path=\u751f\u547d\u79d1\u5b66\u4e0e\u533b\u5b66\clinical-medicine__2212.13138__large-language-models-encode-clinical-knowledge.md (part 2 of 8).md
[8] score=20.14850309280901 path=\u751f\u547d\u79d1\u5b66\u4e0e\u533b\u5b66\clinical-medicine__2212.13138__large-language-models-encode-clinical-knowledge.md (part 2 of 8).md
[9] score=19.475195879068323 path=\u751f\u547d\u79d1\u5b66\u4e0e\u533b\u5b66\clinical-medicine__2212.13138__large-language-models-encode-clinical-knowledge.md (part 2 of 8).md
[10] score=19.01230286941496 path=\u751f\u547d\u79d1\u5b66\u4e0e\u533b\u5b66\clinical-medicine__2212.13138__large-language-models-encode-clinical-knowledge.md (part 2 of 8).md
[11] score=18.99433216372944 path=\u751f\u547d\u79d1\u5b66\u4e0e\u533b\u5b66\clinical-medicine__2212.13138__large-language-models-encode-clinical-knowledge.md (part 2 of 8).md
[12] score=21.58325109849703 path=\u751f\u547d\u79d1\u5b66\u4e0e\u533b\u5b66\clinical-medicine__2212.13138__large-language-models-encode-clinical-knowledge.md (part 5 of 8).md
\end{EvidenceText}
\subsection{BQ04: Which three independent ontologies subdivide the Gene Ontology?}
\begin{tabular}{@{}ll@{}}
Arm & RRF baseline\\
Generator & omp / deepseek-flash\\
HTTP / success & 200 / True\\
Wall time / cost & 13.3 s / 0.015975000000000003 USD\\
Judge scores & F=5, C=5, G=5, A=5; overall=5.00\\
Verdict & correct\\
\end{tabular}

\textbf{Dimension rationale}
\begin{EvidenceText}
{
  "factual": "Names exactly the gold's three ontologies \u2014 Molecular Function (MF), Biological Process (BP), Cellular Component (CC) \u2014 with the gold's 'each representing a different biological aspect' framing.",
  "completeness": "The single-part question is fully answered in one sentence.",
  "groundedness": "The claim is a direct restatement of the gold section 2 passage and is tagged to it; nothing unsupported.",
  "sources": "Attributed to [S1], matching the gold passage."
}
\end{EvidenceText}

\textbf{Judge summary}
\begin{EvidenceText}
One-sentence answer that names all three ontologies correctly with a matching source attribution.
\end{EvidenceText}

\textbf{Complete answer}
\begin{AnswerText}
The three distinct ontologies are Molecular Function (MF), Biological Process (BP), and Cellular Component (CC), each representing a different biological aspect [S1].
\end{AnswerText}

\textbf{Retrieved evidence index (all returned units; paths and scores) }
\begin{EvidenceText}
[1] score=29.095269580918774 path=\u751f\u547d\u79d1\u5b66\u4e0e\u533b\u5b66\genomics__1602.01876__primer-on-the-gene-ontology.md (part 1 of 2).md
[2] score=19.347806199621772 path=\u751f\u547d\u79d1\u5b66\u4e0e\u533b\u5b66\genomics__1602.01876__primer-on-the-gene-ontology.md (part 1 of 2).md
[3] score=18.651047519283672 path=\u751f\u547d\u79d1\u5b66\u4e0e\u533b\u5b66\genomics__1602.01876__primer-on-the-gene-ontology.md (part 1 of 2).md
[4] score=16.30613549144169 path=\u751f\u547d\u79d1\u5b66\u4e0e\u533b\u5b66\genomics__1602.01876__primer-on-the-gene-ontology.md (part 1 of 2).md
[5] score=15.936701344753343 path=\u751f\u547d\u79d1\u5b66\u4e0e\u533b\u5b66\genomics__1602.01876__primer-on-the-gene-ontology.md (part 1 of 2).md
[6] score=15.739777292113317 path=\u751f\u547d\u79d1\u5b66\u4e0e\u533b\u5b66\genomics__1602.01876__primer-on-the-gene-ontology.md (part 1 of 2).md
[7] score=14.816879729504919 path=\u751f\u547d\u79d1\u5b66\u4e0e\u533b\u5b66\genomics__1602.01876__primer-on-the-gene-ontology.md (part 1 of 2).md
[8] score=19.37428959190451 path=\u751f\u547d\u79d1\u5b66\u4e0e\u533b\u5b66\genomics__1602.01876__primer-on-the-gene-ontology.md (part 2 of 2).md
[9] score=16.5980053918848 path=\u751f\u547d\u79d1\u5b66\u4e0e\u533b\u5b66\genomics__1602.01876__primer-on-the-gene-ontology.md (part 2 of 2).md
[10] score=14.52268787551351 path=\u751f\u547d\u79d1\u5b66\u4e0e\u533b\u5b66\genomics__1602.01876__primer-on-the-gene-ontology.md (part 2 of 2).md
[11] score=14.424209744277883 path=\u751f\u547d\u79d1\u5b66\u4e0e\u533b\u5b66\genomics__1602.01876__primer-on-the-gene-ontology.md (part 2 of 2).md
[12] score=14.333389145513841 path=\u751f\u547d\u79d1\u5b66\u4e0e\u533b\u5b66\genomics__1602.01876__primer-on-the-gene-ontology.md (part 2 of 2).md
\end{EvidenceText}
\subsection{BQ05: Why does historical European economic growth contradict Unified Growth Theory?}
\begin{tabular}{@{}ll@{}}
Arm & RRF baseline\\
Generator & omp / deepseek-flash\\
HTTP / success & 200 / True\\
Wall time / cost & 26.3 s / 0.026591999999999998 USD\\
Judge scores & F=5, C=5, G=5, A=5; overall=5.00\\
Verdict & correct\\
\end{tabular}

\textbf{Dimension rationale}
\begin{EvidenceText}
{
  "factual": "Matches gold and the paper: 'the prime engine of economic growth', 'a remarkable disagreement', 'no connection, whatever', 'there was never a transition... because there was no stagnation', escape 'did not happen because there was no trap', Industrial Revolution with no impact and no Eastern Europe boost, 'mythical features, which did not exist', revision-or-replacement conclusion, and the same-data paradox. The 30-country fit (a=9.859e-2, k=5.112e-5, singularity 1929, ~1900 diversion, 29-year bypass) is cross-consistent with other answers.",
  "completeness": "A short-answer headline plus four bullets covering no stagnation/transition, no IR signature, invisible regimes with perverse slower diversion, and hyperbolic growth, closing with the paper's own conclusion \u2014 fully answers the question.",
  "groundedness": "Every claim tagged ([S2][S4][S5][S7][S8][S10][S11][S12]) with quotes; the final note honestly says its context is excerpted and lacks country-level figures and the full fitting procedure.",
  "sources": "Per-claim attributions with an explicit description of which content its excerpts do and do not include."
}
\end{EvidenceText}

\textbf{Judge summary}
\begin{EvidenceText}
Well-organized, quote-backed answer covering all the paper's contradiction arguments with honest disclosure of its excerpts' limits.
\end{EvidenceText}

\textbf{Complete answer}
\begin{AnswerText}
**Short answer: because the data show no stagnation phase, no transition, and no Industrial-Revolution break \u2014 so the theory's central sequence never appears.**

According to the paper (Nielsen, using Maddison's long-run GDP data), Europe \u2014 and Western Europe in particular \u2014 is the *hardest* test case for Unified Growth Theory (UGT), since it was the centre of the Industrial Revolution, which Galor calls "the prime engine of economic growth." UGT's processes "should be most prominently reflected and confirmed" there, yet the data show "a remarkable disagreement" and "no connection, whatever, between the data and the Unified Growth Theory" [S2][S7][S8].

The specific contradictions:

- **No stagnation, hence no transition.** UGT's story is a movement *from stagnation to growth*; the European data show "there was never a transition... because there was no stagnation" [S2][S4]. The alleged "remarkable" or "stunning" escape from the Malthusian trap "did not happen because there was no trap" [S4][S12]. Consequently the alleged Malthusian regime of stagnation did not exist [S11][S12].
- **The Industrial Revolution left no signature in the growth trajectory.** It "had absolutely no impact on changing the economic growth trajectories" in these regions, and it "did not boost the economic growth in Eastern Europe" [S4][S11][S12].
- **The three regimes are not visible in the data.** There was no boosting of growth at the supposed transition from the post-Malthusian regime to the sustained-growth regime; instead, growth in Eastern Europe was diverted to a *slower* trajectory soon after that "phantom sustained growth regime" began. Galor's regimes are "clearly dissociated from data" and "have no expected connection with data" [S11][S12].
- **Growth is hyperbolic, not regime-switching.** The mathematical analysis finds that global, regional and national growth were hyperbolic [S5]; for the 30 Western European countries the fitted hyperbolic parameters are $a = 9.859\times10^{-2}$ and $k = 5.112\times10^{-5}$, implying a singularity in 1929, with growth diverted to a slower trajectory around 1900 \u2014 bypassing the singularity by about 29 years \u2014 a pattern UGT does not describe [S10].

The paper's conclusion is that UGT's explanatory mechanism rests on "mythical features, which did not exist, the features contradicted by data," so the theory "needs to be either thoroughly revised or most likely replaced" [S4]. It also stresses the paradox that the theory is contradicted by *the same data used during its development* [S5].

*Note: the context supplied is excerpted (abstract, introduction, results fragments, conclusions, references), so the detailed country-level figures and the full hyperbolic fitting procedure are not reproduced here.*
\end{AnswerText}

\textbf{Retrieved evidence index (all returned units; paths and scores) }
\begin{EvidenceText}
[1] score=45.267325228036256 path=\u7ecf\u6d4e\u4e0e\u793e\u4f1a\economics__1512.08067__unified-growth-theory-contradicted-by-the-econom.md.md
[2] score=44.34155506446946 path=\u7ecf\u6d4e\u4e0e\u793e\u4f1a\economics__1512.08067__unified-growth-theory-contradicted-by-the-econom.md.md
[3] score=41.32555848687476 path=\u7ecf\u6d4e\u4e0e\u793e\u4f1a\economics__1512.08067__unified-growth-theory-contradicted-by-the-econom.md.md
[4] score=40.06528038926803 path=\u7ecf\u6d4e\u4e0e\u793e\u4f1a\economics__1512.08067__unified-growth-theory-contradicted-by-the-econom.md.md
[5] score=39.94677203702138 path=\u7ecf\u6d4e\u4e0e\u793e\u4f1a\economics__1512.08067__unified-growth-theory-contradicted-by-the-econom.md.md
[6] score=38.95938449556583 path=\u7ecf\u6d4e\u4e0e\u793e\u4f1a\economics__1512.08067__unified-growth-theory-contradicted-by-the-econom.md.md
[7] score=38.509741306819855 path=\u7ecf\u6d4e\u4e0e\u793e\u4f1a\economics__1512.08067__unified-growth-theory-contradicted-by-the-econom.md.md
[8] score=38.26962902641479 path=\u7ecf\u6d4e\u4e0e\u793e\u4f1a\economics__1512.08067__unified-growth-theory-contradicted-by-the-econom.md.md
[9] score=37.73870434854454 path=\u7ecf\u6d4e\u4e0e\u793e\u4f1a\economics__1512.08067__unified-growth-theory-contradicted-by-the-econom.md.md
[10] score=36.50012370395166 path=\u7ecf\u6d4e\u4e0e\u793e\u4f1a\economics__1512.08067__unified-growth-theory-contradicted-by-the-econom.md.md
[11] score=36.03446052504341 path=\u7ecf\u6d4e\u4e0e\u793e\u4f1a\economics__1512.08067__unified-growth-theory-contradicted-by-the-econom.md.md
[12] score=36.013589809334675 path=\u7ecf\u6d4e\u4e0e\u793e\u4f1a\economics__1512.08067__unified-growth-theory-contradicted-by-the-econom.md.md
\end{EvidenceText}
\subsection{BQ06: Which physical factor primarily controls the response of precipitation extremes to climate change?}
\begin{tabular}{@{}ll@{}}
Arm & RRF baseline\\
Generator & omp / deepseek-flash\\
HTTP / success & 200 / True\\
Wall time / cost & 44.6 s / 0.0175818 USD\\
Judge scores & F=5, C=5, G=5, A=5; overall=5.00\\
Verdict & correct\\
\end{tabular}

\textbf{Dimension rationale}
\begin{EvidenceText}
{
  "factual": "All claims check out: 'several physical contributions govern the response' (gold); the three-part decomposition; the thermodynamic contribution as 'robust and well understood'; rotation-weakened dynamics in the tropics where 'convection is always a key factor' versus rotation-constrained extratropical vertical velocities (consistent with the paper's discussion); microphysics mattering for short-duration convective extremes \u2014 all accurate.",
  "completeness": "Concludes 'moisture/thermodynamics is the well-constrained controlling factor; dynamical and microphysical controls remain the uncertain part' \u2014 the question is answered, with the honest note that no explicit 'primary' sentence exists in context.",
  "groundedness": "Every claim attributed ([S1][S4][S6][S11][S12]); the f0^2 rotation constraint is a specific technical detail consistent with the paper's theory section; nothing fabricated.",
  "sources": "Clear per-claim attributions including the dynamics and microphysics passages."
}
\end{EvidenceText}

\textbf{Judge summary}
\begin{EvidenceText}
Correct, well-sourced answer identifying moisture/thermodynamics as the controlling factor with appropriately hedged primacy and accurate decomposition.
\end{EvidenceText}

\textbf{Complete answer}
\begin{AnswerText}
**Context does not explicitly name one "primary" factor** \u2014 the abstract states that *"several physical contributions govern the response"* [S1], decomposed into thermodynamic, microphysical, and dynamical contributions [S4].

What the context does establish as the dominant/robust control:

- **The thermodynamic contribution** (atmospheric water-vapor content, i.e. Clausius\u2013Clapeyron-type scaling) is described as *"robust and well understood,"* whereas *"theoretical understanding of the microphysical and dynamical contributions is still being developed"* [S4].
- Dynamics matter most where rotation is weak: *"where the dynamical influence of planetary rotation is weaker and where convection is always a key factor"* (tropics), contrasted with large-scale extratropical vertical velocities that are comparatively insensitive because of the $f_0^2$ term [S6].
- Microphysics enters mainly for short-duration convective extremes, which are sensitive to cloud/precipitation microphysics parameterization [S11, S12].

So: **moisture/thermodynamics is the well-constrained controlling factor; dynamical and microphysical controls remain the uncertain part of the response** \u2014 the context gives no stronger single-factor claim than that.
\end{AnswerText}

\textbf{Retrieved evidence index (all returned units; paths and scores) }
\begin{EvidenceText}
[1] score=44.88049708966047 path=\u81ea\u7136\u79d1\u5b66\u4e0e\u5730\u7403\u79d1\u5b66\climate-science__1503.07557__precipitation-extremes-under-climate-change.md (part 1 of 3).md
[2] score=42.28477133112594 path=\u81ea\u7136\u79d1\u5b66\u4e0e\u5730\u7403\u79d1\u5b66\climate-science__1503.07557__precipitation-extremes-under-climate-change.md (part 1 of 3).md
[3] score=35.52100698919434 path=\u81ea\u7136\u79d1\u5b66\u4e0e\u5730\u7403\u79d1\u5b66\climate-science__1503.07557__precipitation-extremes-under-climate-change.md (part 1 of 3).md
[4] score=32.95574866551814 path=\u81ea\u7136\u79d1\u5b66\u4e0e\u5730\u7403\u79d1\u5b66\climate-science__1503.07557__precipitation-extremes-under-climate-change.md (part 1 of 3).md
[5] score=30.65644048599295 path=\u81ea\u7136\u79d1\u5b66\u4e0e\u5730\u7403\u79d1\u5b66\climate-science__1503.07557__precipitation-extremes-under-climate-change.md (part 1 of 3).md
[6] score=28.933213154869364 path=\u81ea\u7136\u79d1\u5b66\u4e0e\u5730\u7403\u79d1\u5b66\climate-science__1503.07557__precipitation-extremes-under-climate-change.md (part 1 of 3).md
[7] score=34.34398519329166 path=\u81ea\u7136\u79d1\u5b66\u4e0e\u5730\u7403\u79d1\u5b66\climate-science__1503.07557__precipitation-extremes-under-climate-change.md (part 3 of 3).md
[8] score=27.919268838243397 path=\u81ea\u7136\u79d1\u5b66\u4e0e\u5730\u7403\u79d1\u5b66\climate-science__1503.07557__precipitation-extremes-under-climate-change.md (part 3 of 3).md
[9] score=27.528724067142274 path=\u81ea\u7136\u79d1\u5b66\u4e0e\u5730\u7403\u79d1\u5b66\climate-science__1503.07557__precipitation-extremes-under-climate-change.md (part 3 of 3).md
[10] score=31.288014558108934 path=\u81ea\u7136\u79d1\u5b66\u4e0e\u5730\u7403\u79d1\u5b66\climate-science__1503.07557__precipitation-extremes-under-climate-change.md (part 2 of 3).md
[11] score=29.801300084344014 path=\u81ea\u7136\u79d1\u5b66\u4e0e\u5730\u7403\u79d1\u5b66\climate-science__1503.07557__precipitation-extremes-under-climate-change.md (part 2 of 3).md
[12] score=27.770391031711753 path=\u81ea\u7136\u79d1\u5b66\u4e0e\u5730\u7403\u79d1\u5b66\climate-science__1503.07557__precipitation-extremes-under-climate-change.md (part 2 of 3).md
\end{EvidenceText}
\subsection{BQ07: What is MameLoshnLM and what parameter scale does it have?}
\begin{tabular}{@{}ll@{}}
Arm & RRF baseline\\
Generator & omp / deepseek-flash\\
HTTP / success & 200 / True\\
Wall time / cost & 11.2 s / 0.0145836 USD\\
Judge scores & F=5, C=1, G=5, A=4; overall=3.75\\
Verdict & incorrect\\
\end{tabular}

\textbf{Dimension rationale}
\begin{EvidenceText}
{
  "factual": "No incorrect claims: it truthfully reports that none of its sources [S1]-[S12] mention 'MameLoshnLM' or any parameter scale, and its list of off-topic sources (Vision-FLAN, dialogue summarization, GPT-3 few-shot, chain-of-thought, query optimization, clinical knowledge, sea-level rise, tumor modeling) is specific and plausible; the gold material itself does contain the answer, but its retrieved context apparently does not.",
  "completeness": "Refuses both parts ('I cannot answer this from the given material') \u2014 the question goes unanswered.",
  "groundedness": "An honest refusal with no guessing of the definition or scale from unrelated sources.",
  "sources": "Describes its sources' topics clearly, though without per-claim ids since no claim is made."
}
\end{EvidenceText}

\textbf{Judge summary}
\begin{EvidenceText}
Honest refusal grounded in an evidently off-topic retrieval, leaving both parts of the question unanswered.
\end{EvidenceText}

\textbf{Complete answer}
\begin{AnswerText}
The provided context does not contain the answer. None of the sources [S1]\u2013[S12] mention "MameLoshnLM" or any model by that name, nor do they state its parameter scale. The sources cover unrelated topics (Vision-FLAN, dialogue summarization, GPT-3 few-shot learning, chain-of-thought prompting, sketch-based query optimization, clinical knowledge in LLMs, sea-level rise, tumor modeling), so I cannot answer this from the given material.
\end{AnswerText}

\textbf{Retrieved evidence index (all returned units; paths and scores) }
\begin{EvidenceText}
[1] score=21.79571658621069 path=\u8ba1\u7b97\u673a\u4e0e\u4eba\u5de5\u667a\u80fd\large-language-models__2402.11690__vision-flan-scaling-human-labeled-tasks-in-v.md (part 7 of 9).md
[2] score=18.96051760147613 path=\u8ba1\u7b97\u673a\u4e0e\u4eba\u5de5\u667a\u80fd\nlp__2310.10981__instructive-dialogue-summarization-with-quer (part 4 of 4).md
[3] score=17.010146481441613 path=\u8ba1\u7b97\u673a\u4e0e\u4eba\u5de5\u667a\u80fd\nlp__2310.10981__instructive-dialogue-summarization-with-quer (part 4 of 4).md
[4] score=16.280428995913265 path=\u8ba1\u7b97\u673a\u4e0e\u4eba\u5de5\u667a\u80fd\nlp__2310.10981__instructive-dialogue-summarization-with-quer (part 4 of 4).md
[5] score=18.828399791292124 path=\u8ba1\u7b97\u673a\u4e0e\u4eba\u5de5\u667a\u80fd\databases__2102.02440__online-sketch-based-query-optimization.md (part 3 of 5).md
[6] score=16.553832523711463 path=\u8ba1\u7b97\u673a\u4e0e\u4eba\u5de5\u667a\u80fd\large-language-models__2005.14165__language-models-are-few-shot-learners.md (part 4 of 13).md
[7] score=16.52407411631562 path=\u751f\u547d\u79d1\u5b66\u4e0e\u533b\u5b66\clinical-medicine__2212.13138__large-language-models-encode-clinical-knowledge.md (part 5 of 8).md
[8] score=16.478329348726188 path=\u8ba1\u7b97\u673a\u4e0e\u4eba\u5de5\u667a\u80fd\nlp__2310.10981__instructive-dialogue-summarization-with-quer (part 1 of 4).md
[9] score=16.429723698767546 path=\u81ea\u7136\u79d1\u5b66\u4e0e\u5730\u7403\u79d1\u5b66\oceanography__physics-0703220__scientific-reticence-and-sea-level-rise (part 1 of 2).md
[10] score=16.094265150578554 path=\u8ba1\u7b97\u673a\u4e0e\u4eba\u5de5\u667a\u80fd\Chain-of-Thought Prompting (part 1 of 3).md
[11] score=15.837794358136073 path=\u751f\u547d\u79d1\u5b66\u4e0e\u533b\u5b66\immunology__2604.01385__strategies-for-tumor-elimination-and-control.md (part 1 of 3).md
[12] score=15.835985871118204 path=\u8ba1\u7b97\u673a\u4e0e\u4eba\u5de5\u667a\u80fd\large-language-models__2402.11690__vision-flan-scaling-human-labeled-tasks-in-v.md (part 9 of 9).md
\end{EvidenceText}
\subsection{BQ08: How does oxygen redox contribute to the capacity of battery cathodes?}
\begin{tabular}{@{}ll@{}}
Arm & RRF baseline\\
Generator & omp / deepseek-flash\\
HTTP / success & 200 / True\\
Wall time / cost & 34.2 s / 0.0251028 USD\\
Judge scores & F=5, C=5, G=5, A=5; overall=5.00\\
Verdict & correct\\
\end{tabular}

\textbf{Dimension rationale}
\begin{EvidenceText}
{
  "factual": "All correct: O2- supplying charge-compensating electrons so capacity exceeds cation-only redox (reducing-agent role); the ongoing debate over proportioning capacity; Koga/Sathiya 2013 first conclusive demonstrations; Li-Mn oxides with little TM redox so performance attributed to oxygen redox; orphaned O 2p states from under-coordination, closer to the Fermi level; ilmenite under-coordination rather than Na over-coordination; stronger effect with 4d metals; the cost side (TM diffusion, O2 release, flammability, first-cycle loss barrier for Li-rich layered manganese oxides) and fluorination mitigation with its capacity/retention tradeoff.",
  "completeness": "Answer, dominant material families, and cost/limits \u2014 fully answers how oxygen redox contributes to capacity, with the honest note that no quantitative mAh/g figure for the O-redox contribution appears in context.",
  "groundedness": "Every claim tagged ([S3][S4][S6][S7][S8][S9][S10][S11]) with quotes consistent with gold and the paper; nothing unsupported.",
  "sources": "Per-claim attributions plus the explicit quantitative-absence note."
}
\end{EvidenceText}

\textbf{Judge summary}
\begin{EvidenceText}
Well-organized, fully sourced answer covering the extra-capacity mechanism, the material families where it dominates, and its costs and limits.
\end{EvidenceText}

\textbf{Complete answer}
\begin{AnswerText}
**Answer:**

Oxygen (anionic) redox adds capacity *on top of* transition-metal redox by letting oxide anions themselves supply the charge-compensating electrons when Li\u207a/Na\u207a are extracted during charging [S4]. Because O\xb2\u207b donates electrons \u2014 effectively acting as a reducing agent \u2014 it can compensate for the absence of oxidizable transition metals, so the electrode can store more energy than cation-only redox would theoretically allow [S4]. In general, the occurrence of oxide-ion redox should raise measurable capacity beyond what transition-metal redox alone permits, though quantitatively attributing capacity to cationic versus anionic contributions is still actively debated [S3].

**Where it dominates:**
- **Li-excess / Na-excess layered oxides** \u2014 the first conclusive demonstrations of reversible anionic redox and its capacity contribution were reported by Koga and Sathiya in 2013 [S3].
- **Li\u2013Mn oxides** \u2014 the transition-metal sites show little extensive redox activity, so performance is attributed to oxygen redox [S9]; in Li-rich layered manganese oxides, Li extraction/insertion is enabled by oxygen redox [S7].
- **Disordered spinel and ilmenite-type NaTMO\u2083** \u2014 high O-redox contribution comes from "orphaned" O 2p states on oxygen ions under-coordinated by transition metals (fewer than three coordinating TMs) or sitting near vacancies; these states lie closer to the Fermi level than fully hybridized O states and are therefore more labile [S11]. In ilmenites the effect arises from O *under*-coordination rather than Na over-coordination [S10]. O redox is more dominant in oxides with 4d (fifth-row) transition metals [S11].

**Cost / limits:** In oxides lacking sufficient TM\u2013O covalency, oxygen oxidation can drive transition-metal diffusion and O\u2082 release near particle surfaces; these irreversible reactions degrade structure and capacity after the first cycle, and O\u2082 evolution is an ignition/flammability hazard [S6]. This first-cycle capacity loss is a key barrier to real-world use of Li-rich layered manganese oxide cathodes [S7]. Surface fluorination mitigates it by strengthening TM\u2013O bonds and reducing Jahn\u2013Teller distortion, suppressing oxygen loss \u2014 with a tradeoff between initial capacity loss and long-term retention [S8].

Note: the context states the effect exists and describes mechanisms/limits, but gives no quantitative capacity figures (mAh/g) for the oxygen-redox contribution.
\end{AnswerText}

\textbf{Retrieved evidence index (all returned units; paths and scores) }
\begin{EvidenceText}
[1] score=41.7521652624632 path=\u5de5\u7a0b\u4e0e\u80fd\u6e90\materials__2408.09942__oxygen-redox-in-battery-cathodes-a-brief-overvie.md (part 1 of 2).md
[2] score=36.79763334461738 path=\u5de5\u7a0b\u4e0e\u80fd\u6e90\materials__2408.09942__oxygen-redox-in-battery-cathodes-a-brief-overvie.md (part 1 of 2).md
[3] score=33.55711347602168 path=\u5de5\u7a0b\u4e0e\u80fd\u6e90\materials__2408.09942__oxygen-redox-in-battery-cathodes-a-brief-overvie.md (part 1 of 2).md
[4] score=33.086242591089565 path=\u5de5\u7a0b\u4e0e\u80fd\u6e90\materials__2408.09942__oxygen-redox-in-battery-cathodes-a-brief-overvie.md (part 1 of 2).md
[5] score=28.43809616777088 path=\u5de5\u7a0b\u4e0e\u80fd\u6e90\materials__2408.09942__oxygen-redox-in-battery-cathodes-a-brief-overvie.md (part 1 of 2).md
[6] score=27.3992683129524 path=\u5de5\u7a0b\u4e0e\u80fd\u6e90\materials__2408.09942__oxygen-redox-in-battery-cathodes-a-brief-overvie.md (part 1 of 2).md
[7] score=27.364923246825878 path=\u5de5\u7a0b\u4e0e\u80fd\u6e90\materials__2408.09942__oxygen-redox-in-battery-cathodes-a-brief-overvie.md (part 1 of 2).md
[8] score=26.550883559293677 path=\u5de5\u7a0b\u4e0e\u80fd\u6e90\materials__2408.09942__oxygen-redox-in-battery-cathodes-a-brief-overvie.md (part 1 of 2).md
[9] score=25.17028271104318 path=\u5de5\u7a0b\u4e0e\u80fd\u6e90\materials__2408.09942__oxygen-redox-in-battery-cathodes-a-brief-overvie.md (part 1 of 2).md
[10] score=32.322849599158204 path=\u5de5\u7a0b\u4e0e\u80fd\u6e90\materials__2408.09942__oxygen-redox-in-battery-cathodes-a-brief-overvie.md (part 2 of 2).md
[11] score=31.47651172621873 path=\u5de5\u7a0b\u4e0e\u80fd\u6e90\materials__2408.09942__oxygen-redox-in-battery-cathodes-a-brief-overvie.md (part 2 of 2).md
[12] score=30.06666010565469 path=\u5de5\u7a0b\u4e0e\u80fd\u6e90\materials__2408.09942__oxygen-redox-in-battery-cathodes-a-brief-overvie.md (part 2 of 2).md
\end{EvidenceText}
\subsection{BQ09: What problem of visual SLAM in dynamic environments does MLP-SLAM address and how?}
\begin{tabular}{@{}ll@{}}
Arm & RRF baseline\\
Generator & omp / deepseek-flash\\
HTTP / success & 200 / True\\
Wall time / cost & 29.0 s / 0.022528199999999998 USD\\
Judge scores & F=5, C=5, G=5, A=5; overall=5.00\\
Verdict & correct\\
\end{tabular}

\textbf{Dimension rationale}
\begin{EvidenceText}
{
  "factual": "All correct: the static-world assumption of ORB-SLAM2/3, LSD-SLAM, DSO with outdoor degradation (gold); geometric fixes requiring manual tuning, lacking generality, losing information (gold); DynaSLAM mislabeling static/dynamic and slowed by MaskR-CNN; no public dataset with indirect RANSAC-noisy evaluation; the open-source real-time stereo ORB-SLAM2-based system with an MLP dynamic/static classifier using comprehensive geometry and no manual parameter configuration; the first-ML-discriminator claim; the >50,000-point KITTI Odometry dataset; dynamic-point exclusion from pose estimation; best accuracy and speed on KITTI Odometry and KITTI360 with the ablation result.",
  "completeness": "Problem (three prior-approach failures plus the evaluation gap) and how (learned discriminator, dataset, exclusion mechanism, results) \u2014 fully answered.",
  "groundedness": "Claims tagged ([S1][S2][S3][S4][S5][S6][S8][S10]) and consistent with gold and paper; nothing unsupported.",
  "sources": "Per-claim attributions throughout."
}
\end{EvidenceText}

\textbf{Judge summary}
\begin{EvidenceText}
Concise but complete and accurately sourced answer covering the problem, prior failures, the MLP discriminator, the dataset, and results.
\end{EvidenceText}

\textbf{Complete answer}
\begin{AnswerText}
**Problem.** Standard V-SLAM (ORB-SLAM2/3, LSD-SLAM, DSO) assumes a static world and stays accurate only where few objects move. In outdoor scenes with many movable objects \u2014 pedestrians, cars, buses \u2014 precision degrades [S1][S3][S4]. Prior fixes were either purely geometric (epipolar/reprojection error, optical flow, depth) which require manual parameter tuning, lack generality, lose information, and are hard to fine-tune [S5][S1], or semantic-mask approaches like DynaSLAM, which can mislabel static/dynamic objects and is slow due to MaskR-CNN [S2]. Evaluation was also unreliable because dynamic/static feature-point classifiers had no public dataset and were only judged indirectly via SLAM metrics subject to RANSAC randomness [S5].

**How MLP-SLAM addresses it.** It is an open-source, real-time **stereo** SLAM system built on ORB-SLAM2 that adds an **MLP that classifies each feature point as dynamic or static**, using comprehensive geometry information and requiring **no manual parameter configuration** \u2014 claimed as the first machine-learning-based dynamic/static feature discriminator in SLAM [S6]. To make this trainable and directly evaluable, the authors released datasets of **>50,000 feature points from KITTI Odometry** for training/testing the discriminator [S6]. Dynamic points are excluded from pose estimation, so accuracy and speed improve: the paper reports the best accuracy and fastest speed among compared dynamic SLAM systems on KITTI Odometry and KITTI360, with the MLP-only module outperforming the object-tracking module and their combination performing best [S3][S8][S10].
\end{AnswerText}

\textbf{Retrieved evidence index (all returned units; paths and scores) }
\begin{EvidenceText}
[1] score=35.27626585134861 path=\u8ba1\u7b97\u673a\u4e0e\u4eba\u5de5\u667a\u80fd\robotics__2410.10669__mlp-slam-multilayer-perceptron-based-simultaneou.md (part 1 of 2).md
[2] score=29.136136699000886 path=\u8ba1\u7b97\u673a\u4e0e\u4eba\u5de5\u667a\u80fd\robotics__2410.10669__mlp-slam-multilayer-perceptron-based-simultaneou.md (part 1 of 2).md
[3] score=27.83658005497794 path=\u8ba1\u7b97\u673a\u4e0e\u4eba\u5de5\u667a\u80fd\robotics__2410.10669__mlp-slam-multilayer-perceptron-based-simultaneou.md (part 1 of 2).md
[4] score=21.258686109794805 path=\u8ba1\u7b97\u673a\u4e0e\u4eba\u5de5\u667a\u80fd\robotics__2410.10669__mlp-slam-multilayer-perceptron-based-simultaneou.md (part 1 of 2).md
[5] score=20.865935389675784 path=\u8ba1\u7b97\u673a\u4e0e\u4eba\u5de5\u667a\u80fd\robotics__2410.10669__mlp-slam-multilayer-perceptron-based-simultaneou.md (part 1 of 2).md
[6] score=18.45186519439389 path=\u8ba1\u7b97\u673a\u4e0e\u4eba\u5de5\u667a\u80fd\robotics__2410.10669__mlp-slam-multilayer-perceptron-based-simultaneou.md (part 1 of 2).md
[7] score=28.676819266083193 path=\u8ba1\u7b97\u673a\u4e0e\u4eba\u5de5\u667a\u80fd\robotics__2410.10669__mlp-slam-multilayer-perceptron-based-simultaneou.md (part 2 of 2).md
[8] score=26.59708956225132 path=\u8ba1\u7b97\u673a\u4e0e\u4eba\u5de5\u667a\u80fd\robotics__2410.10669__mlp-slam-multilayer-perceptron-based-simultaneou.md (part 2 of 2).md
[9] score=23.624972272998498 path=\u8ba1\u7b97\u673a\u4e0e\u4eba\u5de5\u667a\u80fd\robotics__2410.10669__mlp-slam-multilayer-perceptron-based-simultaneou.md (part 2 of 2).md
[10] score=20.098792163835043 path=\u8ba1\u7b97\u673a\u4e0e\u4eba\u5de5\u667a\u80fd\robotics__2410.10669__mlp-slam-multilayer-perceptron-based-simultaneou.md (part 2 of 2).md
[11] score=18.375862131585748 path=\u8ba1\u7b97\u673a\u4e0e\u4eba\u5de5\u667a\u80fd\robotics__2410.10669__mlp-slam-multilayer-perceptron-based-simultaneou.md (part 2 of 2).md
[12] score=22.673061540619678 path=\u8ba1\u7b97\u673a\u4e0e\u4eba\u5de5\u667a\u80fd\robotics__2509.14191__mcgs-slam-a-multi-camera-slam-framework-usin.md (part 2 of 2).md
\end{EvidenceText}
\subsection{BQ10: What strategies for tumor control emerge under immune evasion and chemotherapy resistance?}
\begin{tabular}{@{}ll@{}}
Arm & RRF baseline\\
Generator & omp / deepseek-flash\\
HTTP / success & 200 / True\\
Wall time / cost & 36.4 s / 0.0266178 USD\\
Judge scores & F=5, C=5, G=5, A=5; overall=5.00\\
Verdict & correct\\
\end{tabular}

\textbf{Dimension rationale}
\begin{EvidenceText}
{
  "factual": "All correct: the fight-to-win vs fight-to-not-lose framing; the T1/T2/effector model identifying threshold conditions (abstract-verbatim); combination therapy quotes matching gold-adjacent text ('significantly enlarges the parameter region...'); efficacy depending on evasion mechanism and dominant phenotype; checkpoint-regulation dominance making control highly sensitive to immune exhaustion so checkpoint inhibitors are most effective; reduced antigen presentation targets (up rho/down mu with CAR-T and vaccines; up sigma/down delta with adoptive T-cell transfer and cytotoxic therapy); sequential composition-guided scheduling with the chemotherapy-first example; marker-guided selection; ML-enabled personalized adaptive therapy as emerging; microenvironment feedback and oncolytic viral therapy flagged as untested future work.",
  "completeness": "Strategic framing plus four strategy layers and an honest limits section \u2014 fully answers the question.",
  "groundedness": "Every claim tagged ([S3][S4][S7][S8][S9][S10][S11][S12]) with quotes; the declared limits (no explicit threshold expressions in its excerpts) match its source truncation.",
  "sources": "Per-claim attributions including which sources carry the future-work items."
}
\end{EvidenceText}

\textbf{Judge summary}
\begin{EvidenceText}
Well-structured, fully sourced answer covering combination, mechanism-matched, sequential, and adaptive strategies with explicit context limits.
\end{EvidenceText}

\textbf{Complete answer}
\begin{AnswerText}
## What the paper proposes

**Strategic framing.** The central question is whether treatment should "fight to win" (aggressive elimination) or "fight to not lose" (long-term control and resistance management) [S3]. The models track immune-resistant ($T_1$), immune-sensitive ($T_2$), and effector ($E$) populations to identify threshold conditions for persistence, elimination, and phenotype dominance [S4].

**1. Combination therapy over monotherapy.** Immune dynamics alone are generally insufficient to fully eliminate the tumor, so chemotherapy plus immunotherapy is considered together. Chemotherapy "significantly enlarges the parameter region corresponding to tumor control or tumor-free equilibria," reduces tumor burden, and directly targets resistant subpopulations, while immunotherapy enhances immune-mediated tumor suppression; combined, they expand the region with tumor-free or low-population equilibria [S7][S8].

**2. Match therapy to the evasion mechanism.** Efficacy depends on both the evasion mechanism and the dominant tumor phenotype [S8]. Two resistance strategies are modeled \u2014 immune checkpoint regulation and reduced antigen presentation [S3][S8]:
- If the *checkpoint-regulation* resistant subpopulation dominates, control is highly sensitive to the immune exhaustion rate, making **checkpoint inhibitors the most effective therapy** [S8].
- Under *reduced antigen presentation*, the phenotype-guided targets differ: for a $T_1$-dominant (immune-evasive) tumor, increase $\rho$ / decrease $\mu$, with CAR-T therapy and vaccines; for a $T_2$-dominant (immune-sensitive) tumor, increase $\sigma$ / decrease $\delta$, with adoptive T-cell transfer and cytotoxic therapy [S12].

**3. Composition-guided sequential scheduling.** Because parameter thresholds separate elimination, coexistence, and dominance \u2014 and small parameter changes can trigger qualitative shifts between immune control and escape \u2014 the authors propose sequential strategies guided by tumor composition, aiming to minimize burden while preventing resistance from becoming dominant. Starting from a $T_1$-dominant state, for example, chemotherapy may first reduce burden followed by immunotherapy to maintain control; alternatively immunotherapy may act first [S9].

**4. Marker-guided and data-driven adaptation.** This sensitivity motivates a medicine approach where therapy selection is guided by measurable biological markers and tumor composition [S9]. The authors further note emerging machine-learning/deep-learning approaches may enable personalized adaptive therapy informed by patient-specific parameter estimation [S10].

## Limits of the context
The provided excerpts do not give the explicit threshold expressions or the full model equations (part 2 is truncated at the equilibrium definitions [S11]), so the precise mathematical conditions for elimination versus escape cannot be stated from this context. The paper also flags untested extensions \u2014 microenvironment feedback, oncolytic viral therapy \u2014 as future work rather than presented strategies [S10].
\end{AnswerText}

\textbf{Retrieved evidence index (all returned units; paths and scores) }
\begin{EvidenceText}
[1] score=59.47807410043595 path=\u751f\u547d\u79d1\u5b66\u4e0e\u533b\u5b66\immunology__2604.01385__strategies-for-tumor-elimination-and-control.md (part 1 of 3).md
[2] score=39.44773520344411 path=\u751f\u547d\u79d1\u5b66\u4e0e\u533b\u5b66\immunology__2604.01385__strategies-for-tumor-elimination-and-control.md (part 1 of 3).md
[3] score=34.865716729676315 path=\u751f\u547d\u79d1\u5b66\u4e0e\u533b\u5b66\immunology__2604.01385__strategies-for-tumor-elimination-and-control.md (part 1 of 3).md
[4] score=33.32871648651912 path=\u751f\u547d\u79d1\u5b66\u4e0e\u533b\u5b66\immunology__2604.01385__strategies-for-tumor-elimination-and-control.md (part 1 of 3).md
[5] score=29.57608623986421 path=\u751f\u547d\u79d1\u5b66\u4e0e\u533b\u5b66\immunology__2604.01385__strategies-for-tumor-elimination-and-control.md (part 1 of 3).md
[6] score=29.456924162599716 path=\u751f\u547d\u79d1\u5b66\u4e0e\u533b\u5b66\immunology__2604.01385__strategies-for-tumor-elimination-and-control.md (part 1 of 3).md
[7] score=55.24912931621789 path=\u751f\u547d\u79d1\u5b66\u4e0e\u533b\u5b66\immunology__2604.01385__strategies-for-tumor-elimination-and-control.md (part 3 of 3).md
[8] score=39.4320813705303 path=\u751f\u547d\u79d1\u5b66\u4e0e\u533b\u5b66\immunology__2604.01385__strategies-for-tumor-elimination-and-control.md (part 3 of 3).md
[9] score=36.30175010223417 path=\u751f\u547d\u79d1\u5b66\u4e0e\u533b\u5b66\immunology__2604.01385__strategies-for-tumor-elimination-and-control.md (part 3 of 3).md
[10] score=35.95633679266244 path=\u751f\u547d\u79d1\u5b66\u4e0e\u533b\u5b66\immunology__2604.01385__strategies-for-tumor-elimination-and-control.md (part 3 of 3).md
[11] score=40.297323400334946 path=\u751f\u547d\u79d1\u5b66\u4e0e\u533b\u5b66\immunology__2604.01385__strategies-for-tumor-elimination-and-control.md (part 2 of 3).md
[12] score=32.216444513164916 path=\u751f\u547d\u79d1\u5b66\u4e0e\u533b\u5b66\immunology__2604.01385__strategies-for-tumor-elimination-and-control.md (part 2 of 3).md
\end{EvidenceText}
\section{RAPTOR tree retrieval}
\subsection{BQ01: How does the Transformer compute attention, and why does it replace recurrence and convolution?}
\begin{tabular}{@{}ll@{}}
Arm & RAPTOR tree retrieval\\
Generator & omp / deepseek-flash\\
HTTP / success & 200 / True\\
Wall time / cost & 210.0 s / 0.045105000000000006 USD\\
Judge scores & F=5, C=4, G=5, A=5; overall=4.75\\
Verdict & partial\\
\end{tabular}

\textbf{Index/retrieval metadata}
\begin{EvidenceText}
{
  "retrieval": "collapse-tree cosine over all nodes",
  "attribution": "prefix containment; boundary chunks may mis-attribute",
  "candidates": 324
}
\end{EvidenceText}

\textbf{Dimension rationale}
\begin{EvidenceText}
{
  "factual": "Details all match the paper: decoder self-attention masked so position i depends only on positions < i; residual connections + layer normalization around each sub-layer; position-wise FFN of two linear transforms with ReLU; sinusoidal positional encodings summed with embeddings; multi-head projection formula with the stated weight dimensions and h=8, d_k=d_v=64; convolutions needing O(n/k) stacked or O(log_k n) dilated layers; convolutional layers ~k times more expensive than recurrent; 28.4/41.8 BLEU and 12h on eight P100s. It truthfully reports the core QK^T/sqrt(d_k) formula as 'not present in the provided context'.",
  "completeness": "'How' is substantial (attention function, scaled dot-product named, masking, full multi-head detail, surrounding machinery), and 'why' is a structured table comparing parallelism, path length and cost, plus BLEU evidence and the three desiderata (noted truncated). Only the explicit formula and the full desiderata enumeration are missing, both declared.",
  "groundedness": "Every row of its comparison table and every factual clause carries a source id; it honestly reports [S11] truncating mid-sentence at 'The output is comp...' and [S9] not fully enumerating the desiderata, so no claim outruns its evidence.",
  "sources": "Attributions are dense and precise, down to which source truncates where; nothing is left unattributed."
}
\end{EvidenceText}

\textbf{Judge summary}
\begin{EvidenceText}
Rich, fully sourced answer on both halves; the only gap is the explicit scaled dot-product equation, honestly attributed to context truncation.
\end{EvidenceText}

\textbf{Complete answer}
\begin{AnswerText}
## How attention is computed (per context)

**Attention function.** It maps a query and a set of key\u2013value pairs to an output, where query, keys, values, and output are all vectors [S11]. The context explicitly mentions "scaled dot-product attention" and that masking is implemented by setting softmax inputs to $-\infty$ for illegal connections [S10], but the exact scoring formula (the $QK^\top/\sqrt{d_k}$ softmax expression) is **not present in the provided context** \u2014 [S11] truncates mid-sentence at "The output is comp\u2026".

**Multi-head attention.** Queries, keys, and values are projected into $h$ parallel subspaces, attention is applied in each, the $d_v$-dimensional outputs are concatenated and projected again [S6][S7]:

$$\mathrm{MultiHead}(Q,K,V)=\mathrm{Concat}(\mathrm{head}_1,\dots,\mathrm{head}_h)W^O,\quad \mathrm{head}_i=\mathrm{Attention}(QW_i^Q,\,KW_i^K,\,VW_i^V)$$

with $W_i^Q,W_i^K\in\mathbb{R}^{d_{model}\times d_k}$, $W_i^V\in\mathbb{R}^{d_{model}\times d_v}$, $W^O\in\mathbb{R}^{hd_v\times d_{model}}$; the paper uses $h=8$, $d_k=d_v=d_{model}/h=64$, so cost stays comparable to single-head attention at full dimensionality [S6][S7]. Multi-head lets the model jointly attend to different representation subspaces at different positions, which single-head averaging inhibits [S6][S7].

**Surrounding machinery.** Decoder self-attention is masked so position $i$ can depend only on positions $<i$ [S11]; residual connections + layer normalization wrap each sub-layer [S11]; each layer also has a position-wise FFN, two linear transforms with a ReLU between [S10]; sinusoidal positional encodings are summed with embeddings [S9].

## Why it replaces recurrence and convolution

| Problem | Recurrence | Convolution | Self-attention |
|---|---|---|---|
| Parallelism | sequential nature precludes parallelization within training examples [S1] | parallel, but\u2026 | fully parallel, significantly more parallelization [S1] |
| Ops to relate two distant positions | $O(n)$ path length [S1] | linear in distance (ConvS2S) or logarithmic (ByteNet) [S2][S3]; a kernel $k<n$ cannot connect all pairs, requiring $O(n/k)$ stacked or $O(\log_k n)$ dilated layers [S12] | **constant** number of sequentially executed operations [S1][S2][S3] |
| Cost | baseline | $\sim k\times$ more expensive than recurrent layers [S12] | constant path length, but reduced effective resolution from averaging attention-weighted positions \u2014 counteracted by multi-head attention [S2][S3][S6] |

Empirically: 28.4 BLEU En\u2192De (>2 BLEU over prior best) and 41.8 BLEU En\u2192Fr single-model SOTA in 3.5 days on eight GPUs, while SOTA was reached in as little as 12 hours on eight P100s [S1]; it is "the first sequence transduction model based entirely on attention," replacing recurrent layers with multi-headed self-attention and training significantly faster [S8]. [S9] also notes the motivation rests on three desiderata (complexity per layer, parallelizable computation, path length), though the passage is truncated and does not enumerate them fully.

**Not answerable from context:** the precise scaled dot-product formula, and the full three-desiderata table in [S9].
\end{AnswerText}

\textbf{Retrieved evidence index (all returned units; paths and scores) }
\begin{EvidenceText}
[1] score=0.7576 path=(raptor-node-unattributed)
[2] score=0.7062 path=(raptor-node-unattributed)
[3] score=0.7012 path=1706.03762
[4] score=0.6971 path=(raptor-node-unattributed)
[5] score=0.6712 path=(raptor-node-unattributed)
[6] score=0.6677 path=(raptor-node-unattributed)
[7] score=0.6632 path=1706.03762
[8] score=0.6202 path=(raptor-node-unattributed)
[9] score=0.6177 path=(raptor-node-unattributed)
[10] score=0.5853 path=1706.03762
[11] score=0.5822 path=1706.03762
[12] score=0.5775 path=1706.03762
\end{EvidenceText}
\subsection{BQ02: What defines NISQ technology and what is its central limitation?}
\begin{tabular}{@{}ll@{}}
Arm & RAPTOR tree retrieval\\
Generator & omp / deepseek-flash\\
HTTP / success & 200 / True\\
Wall time / cost & 76.1 s / 0.021609000000000003 USD\\
Judge scores & F=4, C=4, G=4, A=4; overall=4.00\\
Verdict & partial\\
\end{tabular}

\textbf{Index/retrieval metadata}
\begin{EvidenceText}
{
  "retrieval": "collapse-tree cosine over all nodes",
  "attribution": "prefix containment; boundary chunks may mis-attribute",
  "candidates": 324
}
\end{EvidenceText}

\textbf{Dimension rationale}
\begin{EvidenceText}
{
  "factual": "Mostly correct (coinage, 50-few-hundred qubits, 50 beyond brute-force simulation, imperfect control, noise's serious near-term limitations, ~1000-gate ceiling, fault tolerance distant, verification difficulties). One factual error: it defines 'noisy' as 'two-qubit gate errors above 1%' on leading trapped-ion and superconducting hardware \u2014 the paper's figure is above the 0.1% (10^-3) level, an order-of-magnitude misstatement inconsistent with its own 1000-gate ceiling.",
  "completeness": "Definition is covered with both name components; the central limitation is present but less direct \u2014 the answer foregrounds the 'not transformative on its own' and verification corollaries, and only presents the noise-imposed ~1000-gate circuit-depth ceiling as 'the concrete quantified form' after admitting 'the passages give no single sentence stating the central limitation'.",
  "groundedness": "Claims carry source tags and the honest context-limitation note is good, but the 'above 1%' error-rate claim is attached to [S1][S2] and misstates what those sources say (the paper's 10^-3 level), so one attribution does not hold up.",
  "sources": "Per-claim attributions exist and the closing caveat is honest, but the misquoted error-rate figure makes at least one attribution inaccurate."
}
\end{EvidenceText}

\textbf{Judge summary}
\begin{EvidenceText}
Right on the definition and the noise limitation overall, but an order-of-magnitude error-rate misstatement ('above 1%' vs the paper's 0.1% level) mars the defining features.
\end{EvidenceText}

\textbf{Complete answer}
\begin{AnswerText}
**NISQ = Noisy Intermediate-Scale Quantum** \u2014 a term coined for the "impending new era" of quantum hardware now being entered, while fault tolerance remains "a rather distant dream" [S2].

**Defining features** (two words, two meanings):
- **"Intermediate-scale"** \u2014 devices with **50 to a few hundred qubits**; 50 is the significant milestone because that is beyond brute-force simulation by the most powerful existing digital supercomputers [S1][S2].
- **"Noisy"** \u2014 imperfect control over those qubits: two-qubit gate errors **above 1%** on leading trapped-ion and superconducting hardware, which limits circuits to roughly **1000 gates** absent error correction [S1][S2].

**Central limitation:** noise "will place serious limitations on what quantum devices can accomplish" [S2]. Concretely:
- NISQ is **not transformative on its own** \u2014 its primary role is to *pave the way* for more advanced devices; truly transformative quantum technologies probably require fault tolerance, and given the hefty overhead of quantum error correction that era "may still be rather distant" [S3][S4]. Transformative effects may be decades away [S1].
- A consequence for verification: NISQ imperfections make it hard to confirm outputs are correct, especially in quantum simulation, so better verification methods are needed [S6].

Context limitation: the passages give no single sentence stating "the central limitation" as such; the above is the strongest supported reading, with the noise-imposed circuit-depth ceiling (~1000 gates) [S1] as the concrete quantified form.
\end{AnswerText}

\textbf{Retrieved evidence index (all returned units; paths and scores) }
\begin{EvidenceText}
[1] score=0.6425 path=(raptor-node-unattributed)
[2] score=0.6115 path=(raptor-node-unattributed)
[3] score=0.5968 path=(raptor-node-unattributed)
[4] score=0.5708 path=1801.00862
[5] score=0.5522 path=1801.00862
[6] score=0.5135 path=(raptor-node-unattributed)
[7] score=0.5122 path=1801.00862
[8] score=0.5071 path=1801.00862
[9] score=0.5057 path=(raptor-node-unattributed)
[10] score=0.4902 path=1801.00862
[11] score=0.477 path=1801.00862
[12] score=0.4581 path=(raptor-node-unattributed)
\end{EvidenceText}
\subsection{BQ03: What is MultiMedQA and how did Flan-PaLM perform on the MedQA benchmark?}
\begin{tabular}{@{}ll@{}}
Arm & RAPTOR tree retrieval\\
Generator & omp / deepseek-flash\\
HTTP / success & 200 / True\\
Wall time / cost & 68.8 s / 0.022576199999999998 USD\\
Judge scores & F=5, C=5, G=5, A=5; overall=5.00\\
Verdict & correct\\
\end{tabular}

\textbf{Index/retrieval metadata}
\begin{EvidenceText}
{
  "retrieval": "collapse-tree cosine over all nodes",
  "attribution": "prefix containment; boundary chunks may mis-attribute",
  "candidates": 324
}
\end{EvidenceText}

\textbf{Dimension rationale}
\begin{EvidenceText}
{
  "factual": "Correct throughout: the 5 heterogeneity axes match the paper's dataset description; HealthSearchQA as curated commonly searched consumer queries; 67.6% via few-shot + CoT + self-consistency, >17% over prior SOTA and 17.3 points over PubMedGPT's 50.3%; 62.0% on the 5-option variant; plain 5-shot 58.9 (PaLM) vs 60.3 (Flan-PaLM); CoT alone not improving over few-shot; human-eval figures 61.9% consensus and 29.7% potential harm motivating Med-PaLM.",
  "completeness": "Both parts fully answered with valuable qualifications: it correctly notes 67.6% is the combined-prompting configuration and contrasts it with plain 5-shot, then adds the safety-relevant human-evaluation caveats that the paper itself foregrounds.",
  "groundedness": "Claims are tagged to [S12][S1][S3][S7][S8] and all numbers check against the paper; the qualifications are framed as coming from the context, not extrapolated.",
  "sources": "Per-claim attributions with distinct source ids for the heterogeneity axes, scaling numbers, CoT finding, and headline results."
}
\end{EvidenceText}

\textbf{Judge summary}
\begin{EvidenceText}
Complete, accurate answer on both halves, distinguished by correctly attributing 67.6% to the combined prompting configuration and adding the paper's human-evaluation caveats.
\end{EvidenceText}

\textbf{Complete answer}
\begin{AnswerText}
**MultiMedQA** \u2014 a combined benchmark of medical QA datasets spanning three capability families: multiple-choice exam QA, long-form answers from medical professionals, and long-form answers from non-professionals. Concretely it aggregates MedQA, MedMCQA, PubMedQA, LiveQA, MedicationQA, and MMLU clinical topics, plus a new free-response dataset, **HealthSearchQA**, of curated commonly searched consumer health queries. All are English-language and heterogeneous along format (multiple-choice vs. long-form), skills (fact recall vs. reasoning), domain (open vs. closed), question source (professional exams, research literature, consumer queries), and label/metadata richness [S12].

**Flan-PaLM on MedQA** \u2014 using few-shot prompting combined with chain-of-thought and self-consistency, Flan-PaLM (540B) reached **67.6%**, over 17% above the prior state of the art and surpassing PubMedGPT's 50.3% by 17.3 percentage points; on the 5-option variant of MedQA it scores **62.0%** [S1][S3].

Two qualifications from the context:
- The 67.6% is the *combined* few-shot + CoT + self-consistency configuration. Under plain 5-shot prompting the 540B model moves from 58.9 (PaLM) to 60.3 (Flan-PaLM) [S7], and the paper reports that CoT alone did not improve over standard few-shot prompting across MedQA/MedMCQA/PubMedQA [S8].
- Strong MCQ scores did not translate to safe long-form answers: only 61.9% of Flan-PaLM long-form consumer answers were judged aligned with scientific consensus, and 29.7% were rated potentially harmful [S1] \u2014 the gap that motivated instruction prompt tuning and Med-PaLM.
\end{AnswerText}

\textbf{Retrieved evidence index (all returned units; paths and scores) }
\begin{EvidenceText}
[1] score=0.6831 path=(raptor-node-unattributed)
[2] score=0.6676 path=2212.13138
[3] score=0.6545 path=2212.13138
[4] score=0.6483 path=(raptor-node-unattributed)
[5] score=0.6432 path=2212.13138
[6] score=0.639 path=2212.13138
[7] score=0.6368 path=(raptor-node-unattributed)
[8] score=0.6305 path=2212.13138
[9] score=0.6265 path=2212.13138
[10] score=0.6137 path=2212.13138
[11] score=0.613 path=(raptor-node-unattributed)
[12] score=0.6069 path=(raptor-node-unattributed)
\end{EvidenceText}
\subsection{BQ04: Which three independent ontologies subdivide the Gene Ontology?}
\begin{tabular}{@{}ll@{}}
Arm & RAPTOR tree retrieval\\
Generator & omp / deepseek-flash\\
HTTP / success & 200 / True\\
Wall time / cost & 44.5 s / 0.016065000000000003 USD\\
Judge scores & F=5, C=2, G=5, A=5; overall=4.25\\
Verdict & incorrect\\
\end{tabular}

\textbf{Index/retrieval metadata}
\begin{EvidenceText}
{
  "retrieval": "collapse-tree cosine over all nodes",
  "attribution": "prefix containment; boundary chunks may mis-attribute",
  "candidates": 324
}
\end{EvidenceText}

\textbf{Dimension rationale}
\begin{EvidenceText}
{
  "factual": "Its one substantive identification \u2014 Molecular Function is the first of the three \u2014 is correct; its account of its context ([S1] reading '...Molecular Function (' and cutting off mid-list; S2 breaking off before the list; S4/S5 covering only graph structure) is internally consistent and plausible, though unverifiable; it makes no claim that contradicts the gold.",
  "completeness": "Only one of the three names is given; the other two are withheld as a declared guess-avoidance, so the question is only one-third answered.",
  "groundedness": "Exemplary honesty: it pinpoints the exact truncation point of its source and states 'I can't supply the remaining two from this context without guessing' rather than fabricating names it clearly 'knows'.",
  "sources": "Attributions are precise, including the exact mid-list cutoff and what each of [S1][S2][S4][S5] contains."
}
\end{EvidenceText}

\textbf{Judge summary}
\begin{EvidenceText}
Honest truncation-aware refusal that names only Molecular Function, leaving two of the three required ontologies unanswered.
\end{EvidenceText}

\textbf{Complete answer}
\begin{AnswerText}
The context is truncated exactly where the answer would appear. [S1] states the ontology "is subdivided into three distinct, non-redundant ontologies sharing a common identifier space and syntax: Molecular Function (" \u2014 and cuts off mid-list.

So only one of the three is given: **Molecular Function**. The names of the other two are not present in the supplied passages (S2 repeats the abstract but also breaks off before the list; S4/S5 describe graph structure and relationships, not the three sub-ontologies).

I can't supply the remaining two from this context without guessing.
\end{AnswerText}

\textbf{Retrieved evidence index (all returned units; paths and scores) }
\begin{EvidenceText}
[1] score=0.5799 path=1602.01876
[2] score=0.539 path=1602.01876
[3] score=0.4869 path=1602.01876
[4] score=0.4825 path=1602.01876
[5] score=0.4692 path=(raptor-node-unattributed)
[6] score=0.4531 path=1602.01876
[7] score=0.4503 path=(raptor-node-unattributed)
[8] score=0.4422 path=1602.01876
[9] score=0.4386 path=(raptor-node-unattributed)
[10] score=0.4376 path=(raptor-node-unattributed)
[11] score=0.4355 path=1602.01876
[12] score=0.4323 path=1602.01876
\end{EvidenceText}
\subsection{BQ05: Why does historical European economic growth contradict Unified Growth Theory?}
\begin{tabular}{@{}ll@{}}
Arm & RAPTOR tree retrieval\\
Generator & omp / deepseek-flash\\
HTTP / success & 200 / True\\
Wall time / cost & 74.7 s / 0.01962 USD\\
Judge scores & F=5, C=5, G=5, A=5; overall=5.00\\
Verdict & correct\\
\end{tabular}

\textbf{Index/retrieval metadata}
\begin{EvidenceText}
{
  "retrieval": "collapse-tree cosine over all nodes",
  "attribution": "prefix containment; boundary chunks may mis-attribute",
  "candidates": 324
}
\end{EvidenceText}

\textbf{Dimension rationale}
\begin{EvidenceText}
{
  "factual": "Correct throughout: hyperbolic growth over AD 1500-1900 with no stagnation signature; no trap so no 'remarkable'/'stunning' escape; absent Malthusian oscillations with reciprocal values 'fall[ing] on a decreasing straight line'; Industrial Revolution 'had absolutely no effect' on trajectories (gold says 'no impact', a faithful paraphrase) and no boost in Eastern Europe; diversion to a slower trajectory around 1880-1890 with the 1915 (Eastern) and 1929 (30-country) singularities bypassed \u2014 consistent with other systems' quotes.",
  "completeness": "Four numbered reasons plus a conclusion \u2014 covers no-stagnation, no-trap/no-escape, no transition/no-IR-effect, wrong post-1880 behavior, and the same-data paradox and four-point data display \u2014 fully answers the question.",
  "groundedness": "Each point is attributed ([S1][S2][S3][S4][S5][S6][S8][S10][S11]) with quotes consistent with the gold and document; the one soft spot is the '~300 years' display window, which is plausible but not in the shown gold.",
  "sources": "Clear per-claim attributions including which dataset each fit applies to."
}
\end{EvidenceText}

\textbf{Judge summary}
\begin{EvidenceText}
Complete and accurate answer tying the theory's failed predictions to the hyperbolic-growth evidence with per-claim sourcing.
\end{EvidenceText}

\textbf{Complete answer}
\begin{AnswerText}
**Because the data show continuous hyperbolic growth where the theory requires distinct regimes.**

According to the paper's analysis of Maddison's GDP data for Western and Eastern Europe [S1][S2]:

1. **No stagnation epoch.** The theory's "Malthusian regime of stagnation" did not exist; growth was steadily increasing along a hyperbolic trajectory from at least AD 1500 to 1900 \u2014 precisely the range where the stagnation signature should be clearly displayed for ~300 years [S5][S6]. Europe should have been the strongest confirmation case, since it was the centre of the Industrial Revolution [S1][S2].

2. **No Malthusian trap, so no "escape".** The alleged "remarkable"/"stunning" escape from the trap did not occur because there was no trap [S2][S8]. If stagnation existed, random Malthusian oscillations should appear in both the data and their reciprocals; they are absent, while reciprocal values fall on a decreasing straight line \u2014 the signature of hyperbolic growth [S10].

3. **No transition from stagnation to growth, and no effect of the Industrial Revolution.** The Industrial Revolution did not boost growth in Eastern Europe [S3] and had "absolutely no effect" on growth trajectories [S4]; there was no transition because there was no stagnation [S4][S8].

4. **Wrong post-1880 behaviour.** Actually observed growth was diverted around 1880\u20131890 from a hyperbolic path to a slower (though still fast-increasing) one \u2014 e.g., a 1915 singularity bypassed by ~25 years, or a 1929 singularity for the 30-country set \u2014 the opposite of the theory's predicted escape to sustained growth [S3][S5].

Conclusion: Galor's three regimes of growth "have no relevance to the description, let alone to the explanation, of the mechanism of the economic growth" [S3][S8]; the theory is contradicted even by the same data used during its development, which the paper argues were displayed only as four selected points joined by straight lines [S11].
\end{AnswerText}

\textbf{Retrieved evidence index (all returned units; paths and scores) }
\begin{EvidenceText}
[1] score=0.766 path=(raptor-node-unattributed)
[2] score=0.7286 path=(raptor-node-unattributed)
[3] score=0.712 path=(raptor-node-unattributed)
[4] score=0.709 path=1512.08067
[5] score=0.6983 path=(raptor-node-unattributed)
[6] score=0.6845 path=1512.08067
[7] score=0.6766 path=(raptor-node-unattributed)
[8] score=0.6538 path=1512.08067
[9] score=0.646 path=(raptor-node-unattributed)
[10] score=0.6441 path=1512.08067
[11] score=0.6183 path=1512.08067
[12] score=0.4356 path=(raptor-node-unattributed)
\end{EvidenceText}
\subsection{BQ06: Which physical factor primarily controls the response of precipitation extremes to climate change?}
\begin{tabular}{@{}ll@{}}
Arm & RAPTOR tree retrieval\\
Generator & omp / deepseek-flash\\
HTTP / success & 200 / True\\
Wall time / cost & 67.2 s / 0.0123144 USD\\
Judge scores & F=5, C=5, G=5, A=5; overall=5.00\\
Verdict & correct\\
\end{tabular}

\textbf{Index/retrieval metadata}
\begin{EvidenceText}
{
  "retrieval": "collapse-tree cosine over all nodes",
  "attribution": "prefix containment; boundary chunks may mis-attribute",
  "candidates": 324
}
\end{EvidenceText}

\textbf{Dimension rationale}
\begin{EvidenceText}
{
  "factual": "Correct: the thermodynamic contribution (saturation vapor pressure increase, Clausius-Clapeyron) is the robust, well-understood factor; microphysical and dynamical understanding still developing; sensitivity uncertain where convection matters; possibly higher tropical sensitivity; orographic and snowfall extremes differ \u2014 all gold-verbatim or consistent with the paper.",
  "completeness": "Concise but complete: names the factor as the best-supported control, states why (robust/understood vs under development), and includes the gold's main caveats.",
  "groundedness": "Claims are tagged to [S1][S9] and match the gold's abstract language; nothing unsupported.",
  "sources": "Attributions to the abstract and the extratropics/dynamics passage, clear and accurate."
}
\end{EvidenceText}

\textbf{Judge summary}
\begin{EvidenceText}
Compact but complete: correct thermodynamic answer with the paper's robustness ranking and key caveats, all attributed.
\end{EvidenceText}

\textbf{Complete answer}
\begin{AnswerText}
The context does not state an explicit ranking of a single \u201cprimary\u201d control, but it does identify the **thermodynamic contribution \u2014 the increase in saturation vapor pressure with warming (the Clausius\u2013Clapeyron relationship)** \u2014 as the one governing factor that is robust and well understood. Microphysical and dynamical contributions are described as still under development, so the thermodynamic contribution is the best-supported control on the response of precipitation extremes to warming [S1][S9].

Caveats in the context: sensitivity is uncertain where convection is important and may be higher in the tropics than the extratropics, and orographic and snowfall extremes respond differently from other precipitation extremes [S1][S9].
\end{AnswerText}

\textbf{Retrieved evidence index (all returned units; paths and scores) }
\begin{EvidenceText}
[1] score=0.7298 path=1503.07557
[2] score=0.6997 path=1503.07557
[3] score=0.697 path=1503.07557
[4] score=0.6968 path=1503.07557
[5] score=0.6962 path=1503.07557
[6] score=0.6556 path=(raptor-node-unattributed)
[7] score=0.6548 path=1801.00862
[8] score=0.6545 path=1503.07557
[9] score=0.6545 path=(raptor-node-unattributed)
[10] score=0.652 path=1503.07557
[11] score=0.6466 path=(raptor-node-unattributed)
[12] score=0.6445 path=1503.07557
\end{EvidenceText}
\subsection{BQ07: What is MameLoshnLM and what parameter scale does it have?}
\begin{tabular}{@{}ll@{}}
Arm & RAPTOR tree retrieval\\
Generator & omp / deepseek-flash\\
HTTP / success & 200 / True\\
Wall time / cost & 78.5 s / 0.0103722 USD\\
Judge scores & F=5, C=5, G=5, A=5; overall=5.00\\
Verdict & correct\\
\end{tabular}

\textbf{Index/retrieval metadata}
\begin{EvidenceText}
{
  "retrieval": "collapse-tree cosine over all nodes",
  "attribution": "prefix containment; boundary chunks may mis-attribute",
  "candidates": 324
}
\end{EvidenceText}

\textbf{Dimension rationale}
\begin{EvidenceText}
{
  "factual": "All main claims correct: first open-source LM for Yiddish; the four institutions; continued pretraining of Llama 3.1 8B on Oytser (~915M words) with evaluation on Kashes; 8B scale (gold-verbatim); best overall score 62.6 ahead of Gemma-2 9B, Llama 3.1 8B, Qwen3 8B \u2014 consistent with independent answers. Its '14-task Kashes benchmark' count and '9 sources' for Oytser conflict with another answer's '9 tasks' enumeration and cannot be adjudicated from the shown gold.",
  "completeness": "Both parts answered directly, plus creation method, corpus size, benchmark results, and baselines \u2014 fully answered.",
  "groundedness": "Every claim is tagged ([S2][S3][S4][S5][S6]) and matches the gold or the paper; no fabrication.",
  "sources": "Per-claim attributions, with the scale stated to [S2][S3][S4]."
}
\end{EvidenceText}

\textbf{Judge summary}
\begin{EvidenceText}
Complete, accurate and well-attributed answer; only the unverifiable task/source counts differ from a peer answer.
\end{EvidenceText}

\textbf{Complete answer}
\begin{AnswerText}
**MameLoshnLM** is the first open-source language model built specifically for Yiddish, created by researchers at Bar-Ilan University, University of Cambridge, University of Washington, and the Allen Institute for AI [S2]. It was obtained by **continued pretraining of Llama 3.1 8B** on *Oytser*, a high-quality Yiddish pretraining corpus (~915M words from 9 online and literary sources), and evaluated on *Kashes*, a multi-task Yiddish benchmark [S3][S4][S6].

**Parameter scale: 8B** [S2][S3][S4].

It achieves the best overall score on the 14-task Kashes benchmark (62.6 avg, 5-shot), ahead of Gemma-2 9B, Llama 3.1 8B, and Qwen3 8B [S5].
\end{AnswerText}

\textbf{Retrieved evidence index (all returned units; paths and scores) }
\begin{EvidenceText}
[1] score=0.5928 path=2608.05850
[2] score=0.5676 path=(raptor-node-unattributed)
[3] score=0.5638 path=(raptor-node-unattributed)
[4] score=0.5543 path=2608.05850
[5] score=0.5495 path=(raptor-node-unattributed)
[6] score=0.5394 path=(raptor-node-unattributed)
[7] score=0.5281 path=2608.05850
[8] score=0.5115 path=2608.05850
[9] score=0.5087 path=(raptor-node-unattributed)
[10] score=0.5053 path=2608.05850
[11] score=0.5042 path=(raptor-node-unattributed)
[12] score=0.4891 path=(raptor-node-unattributed)
\end{EvidenceText}
\subsection{BQ08: How does oxygen redox contribute to the capacity of battery cathodes?}
\begin{tabular}{@{}ll@{}}
Arm & RAPTOR tree retrieval\\
Generator & omp / deepseek-flash\\
HTTP / success & 200 / True\\
Wall time / cost & 142.1 s / 0.02145 USD\\
Judge scores & F=5, C=5, G=5, A=5; overall=5.00\\
Verdict & correct\\
\end{tabular}

\textbf{Index/retrieval metadata}
\begin{EvidenceText}
{
  "retrieval": "collapse-tree cosine over all nodes",
  "attribution": "prefix containment; boundary chunks may mis-attribute",
  "candidates": 324
}
\end{EvidenceText}

\textbf{Dimension rationale}
\begin{EvidenceText}
{
  "factual": "All correct: the conventional LiTMO2 baseline with 200-300 mAh/g rarely reached due to phase changes, interphase barriers, and oxygen release (gold); Li2MnO3 as the canonical case with Mn staying +4 and ~1.4 Li per formula unit (gold); orphaned O 2p states from fewer-than-three TM coordination and ilmenite under-coordination; stronger O redox in 4d oxides via higher 4d-2p covalency; the Li4Mn2O5 charge-compensation equation with x as oxygen participation; TM diffusion and surface O2 release with post-first-cycle structural decay.",
  "completeness": "Baseline ceiling, where O redox takes over, why those states are redox-active, quantified participation, and the cost side \u2014 fully answers the question.",
  "groundedness": "Claims tagged ([S2][S3][S4][S5][S6][S8][S9][S11]) with honest disclosures: [S1] identified as an unrelated immune-model paper contributing nothing, and the [S11] truncation point ('tendency towards bo...') stated rather than papered over.",
  "sources": "Per-claim attributions plus explicit handling of an irrelevant and a truncated source."
}
\end{EvidenceText}

\textbf{Judge summary}
\begin{EvidenceText}
Complete, accurately sourced answer from the baseline capacity ceiling through the orbital mechanism to quantified participation and degradation costs.
\end{EvidenceText}

\textbf{Complete answer}
\begin{AnswerText}
**Oxygen redox raises cathode capacity above the transition-metal limit by letting lattice O\xb2\u207b supply electrons in place of the metals.**

- **Baseline mechanism and its ceiling.** In conventional layered LiTMO\u2082 cathodes (O3/P2/P3, Ni/Mn/Co), charge compensation occurs via transition-metal oxidation; theoretical capacities of ~200\u2013300 mAh/g are rarely reached because of phase changes, interphase diffusion barriers, and oxygen release [S5].
- **Where O redox takes over.** Li\u2082MnO\u2083 is the canonical case: Mn stays +4 (its t\u2082g frontier level is too low to oxidize) yet Li de-intercalates up to ~1.4 Li per formula unit, so the charge must be compensated by oxygen [S2][S6]. Li-excess layered hexagonal Li\u2082TMO\u2083 puts Li in the TM layer, creating "orphaned" O states [S5][S6].
- **Why those states are redox-active.** Orphaned O 2p states arise from O ions with fewer than three coordinating TM ions; in ilmenite-type oxides under-coordination by both alkali and TM (with vacancies) creates them. They sit closer to the Fermi level than fully hybridized O states and are therefore more labile for redox [S11][S3]. Oxygen redox is stronger in 4d (fifth-row) oxides because higher 4d\u20132p covalency pushes O 2p population up toward E_Fermi [S9][S11].
- **Quantified participation.** For Li\u2084Mn\u2082O\u2085 the combined cationic/anionic process is written with x as the extent of oxygen participation: Li\u2084\u207aMn\u2082\xb3\u207aO\u2085\xb2\u207b \u2192 4Li\u207a + 4e\u207b + Mn_x\u2074\u207aMn\u2082\u208b\u2093\u2075\u207aO\u2085\u208b\u2093\xb2\u207bO\u2093\xb9\u207b [S8].
- **Cost side.** The extra capacity is not free: reduced O ions have a strong tendency toward... [S11, truncated in context]; oxygen oxidation in oxides lacking sufficient metal\u2013oxygen covalency can drive TM diffusion and O\u2082 release near particle surfaces, and post-first-cycle irreversible reactions (including spinel transformation) decay structure and capacity \u2014 the known failure mode of oxygen-redox cathodes in Li- and Na-ion systems [S4].

Caveat: [S1] is an unrelated immune-model bifurcation paper and contributes nothing here. The exact sentence completing "reduced O ions have a strong tendency towards bo\u2026" is truncated in [S11], so the specific degradation product is not stated in the provided context beyond the O\u2082-evolution/diffusion account in [S4].
\end{AnswerText}

\textbf{Retrieved evidence index (all returned units; paths and scores) }
\begin{EvidenceText}
[1] score=0.7221 path=2604.01385
[2] score=0.7142 path=2408.09942
[3] score=0.7131 path=(raptor-node-unattributed)
[4] score=0.6971 path=2408.09942
[5] score=0.6897 path=(raptor-node-unattributed)
[6] score=0.6862 path=(raptor-node-unattributed)
[7] score=0.6861 path=2408.09942
[8] score=0.6853 path=2408.09942
[9] score=0.6451 path=2408.09942
[10] score=0.6385 path=(raptor-node-unattributed)
[11] score=0.6319 path=2408.09942
[12] score=0.6009 path=2408.09942
\end{EvidenceText}
\subsection{BQ09: What problem of visual SLAM in dynamic environments does MLP-SLAM address and how?}
\begin{tabular}{@{}ll@{}}
Arm & RAPTOR tree retrieval\\
Generator & omp / deepseek-flash\\
HTTP / success & 200 / True\\
Wall time / cost & 94.4 s / 0.0192738 USD\\
Judge scores & F=5, C=5, G=5, A=5; overall=5.00\\
Verdict & correct\\
\end{tabular}

\textbf{Index/retrieval metadata}
\begin{EvidenceText}
{
  "retrieval": "collapse-tree cosine over all nodes",
  "attribution": "prefix containment; boundary chunks may mis-attribute",
  "candidates": 324
}
\end{EvidenceText}

\textbf{Dimension rationale}
\begin{EvidenceText}
{
  "factual": "All correct: degradation with pedestrians/cars/buses (gold); the prior defects (manual parameter establishment reducing generality, incomplete geometric information, information loss \u2014 gold-aligned); the MLP consuming the three geometric errors as complete information; the first-ML-discriminator claim (gold-verbatim); ORB-SLAM2-based real-time stereo architecture; the DFM rule; YOLOv7 + Deep Sort replacing CFP-SLAM's EKF + Hungarian association; MLP inference 0.18 ms and tracking 62.81 ms vs 45.31 ms; highest precision and fastest speed on KITTI Odometry and KITTI360; sequence 00 (parked cars) showing real-dynamic-object recognition; ablation; the 50,000+ point 70/20/10 1:1 dataset.",
  "completeness": "Problem with three enumerated defects, full method description, and reported outcomes including latency, precision, ablation, and dataset \u2014 fully answered.",
  "groundedness": "Every claim tagged ([S1][S2][S3][S4][S5][S6][S7][S8][S10][S11]) and consistent with gold and paper.",
  "sources": "Dense per-claim attributions, including specific sources for the DBSCAN/epipolar prior-work characterization."
}
\end{EvidenceText}

\textbf{Judge summary}
\begin{EvidenceText}
Complete, precise and fully sourced answer covering the problem's three defects, the complete-geometric-information MLP solution, and quantitative outcomes.
\end{EvidenceText}

\textbf{Complete answer}
\begin{AnswerText}
**Problem addressed.** V-SLAM is already precise when few dynamic objects are present, but accuracy degrades in scenes crowded with movable objects \u2014 pedestrians, cars, buses \u2014 typical of outdoor driving [S2]. Prior dynamic-SLAM approaches attack this with hand-crafted mathematical/geometric pipelines (e.g., DBSCAN clustering over bounding boxes plus re-projection and epipolar error checks [S5], [S6]) or semantics-based masking. Those suffer three concrete defects: manual parameter establishment and tuning that restricts applicable scenarios and reduces generality, incomplete exploitation of geometric information, and information loss [S2], [S3], [S5], [S6].

**How MLP-SLAM addresses it** [S3], [S4]:

- **Learned dynamic/static feature discrimination.** A Multilayer Perceptron consumes *complete geometric information* per feature point \u2014 image intensity error $e_I$, re-projection error $e_{Re}$, and epipolar error $e_D$ \u2014 and outputs the dynamic/static label, replacing tuned geometric heuristics and semantics-only classification [S3], [S7]. The authors claim it as the first machine-learning method for dynamic/static feature discrimination in SLAM [S1], [S3].
- **System architecture** (real-time stereo, built on ORB-SLAM2, C++): pre-processing/feature extraction/object detection/tracking \u2192 coarse estimation \u2192 MLP classifier \u2192 fine estimation [S4], [S7].
- **Depth Filter Module (DFM)** compensates for loose bounding boxes from object detection (vs. segmentation masks): given mean depth $\mu_d$ and std $\delta_d$ of in-box feature points, points outside $[\mu_d-\eta\delta_d,\ \mu_d+\eta\delta_d]$ are dropped, so background points inside a box are treated as static during coarse estimation [S4].
- **Tracking front end:** YOLOv7 detection on the left image + Deep Sort tracking, replacing CFP-SLAM's EKF + Hungarian association [S4].

**Reported outcome.** MLP inference costs only 0.18 ms, total tracking 62.81 ms vs. 45.31 ms for ORB-SLAM2 \u2014 i.e., higher precision for 17.50 ms/frame extra [S10]. Highest average precision (lowest average ATE) and fastest speed among compared methods on KITTI Odometry and KITTI360 [S1], [S3]; notably best on sequence 00 (many parked cars), showing the MLP identifies *real* dynamic objects rather than all movable ones [S11]. Ablation: MLP-only beats the object-tracking-only variant, and combining both is best [S8]. They also release a 50,000+ feature-point KITTI Odometry set (70/20/10 train/test/val, 1:1 static:dynamic) for *direct* discriminator evaluation, since prior evaluation was only indirect via SLAM metrics and unreliable due to RANSAC randomness [S1], [S3], [S7].
\end{AnswerText}

\textbf{Retrieved evidence index (all returned units; paths and scores) }
\begin{EvidenceText}
[1] score=0.7065 path=(raptor-node-unattributed)
[2] score=0.7036 path=2410.10669
[3] score=0.7036 path=(raptor-node-unattributed)
[4] score=0.6802 path=(raptor-node-unattributed)
[5] score=0.6723 path=2410.10669
[6] score=0.6704 path=(raptor-node-unattributed)
[7] score=0.6546 path=(raptor-node-unattributed)
[8] score=0.6499 path=2410.10669
[9] score=0.6456 path=(raptor-node-unattributed)
[10] score=0.632 path=(raptor-node-unattributed)
[11] score=0.6241 path=2410.10669
[12] score=0.6225 path=2410.10669
\end{EvidenceText}
\subsection{BQ10: What strategies for tumor control emerge under immune evasion and chemotherapy resistance?}
\begin{tabular}{@{}ll@{}}
Arm & RAPTOR tree retrieval\\
Generator & omp / deepseek-flash\\
HTTP / success & 200 / True\\
Wall time / cost & 76.5 s / 0.020196 USD\\
Judge scores & F=5, C=5, G=5, A=5; overall=5.00\\
Verdict & correct\\
\end{tabular}

\textbf{Index/retrieval metadata}
\begin{EvidenceText}
{
  "retrieval": "collapse-tree cosine over all nodes",
  "attribution": "prefix containment; boundary chunks may mis-attribute",
  "candidates": 324
}
\end{EvidenceText}

\textbf{Dimension rationale}
\begin{EvidenceText}
{
  "factual": "All correct: phenotype-matched targets under both mechanisms (checkpoint evasion: down mu -> inhibitors, up sigma/down delta -> adoptive/cytotoxic; reduced antigen presentation: up rho, down mu -> CAR-T/vaccines, with effector-supply increases unhelpful in that regime); the threshold around phi of about 0.003 above which no immune-parameter modulation controls the T1-only equilibrium except altering theta1 (cross-consistent with a peer answer); the combination-expansion and nonlinearity claims (gold-aligned); the bistability landscape for theta1/mu/rho and the theta2 = 0.5 vs 0.8 contrast; narrow coexistence with sizable tumors; resistance economics (theta2 scaling the kill rate, a 30%-of-alpha growth cost, chi at 50% of alpha). One garbled clause ('T1 fully susceptible at zeta2 = 0, theta2 = 1') reads like a typo for T2 but is not resolvable against the shown gold.",
  "completeness": "Both evasion mechanisms, combination therapy, the parameter landscape, resistance economics, and an explicit gap note \u2014 fully answers which strategies emerge.",
  "groundedness": "Claims tagged ([S1][S2][S3][S4][S6][S7][S10][S11]) with quotes and numeric details consistent with the paper and peers.",
  "sources": "Per-claim attributions including the honest gap statement about the missing chemo-resistant-plus-reduced-presentation row."
}
\end{EvidenceText}

\textbf{Judge summary}
\begin{EvidenceText}
Complete and precise answer distinguishing both evasion mechanisms with quantitative bifurcation detail, one minor garbled clause notwithstanding.
\end{EvidenceText}

\textbf{Complete answer}
\begin{AnswerText}
**No single lever controls both phenotypes \u2014 the model points to phenotype-matched targets, with chemotherapy and immunotherapy acting synergistically rather than additively.**

**Immune checkpoint evasion** [S2][S4]
- **T1-only (immune-evasive dominant)**: best target is decreasing \u03bc (tumor-induced effector suppression) \u2192 checkpoint inhibitors [S2].
- **T2-only (immune-sensitive dominant)**: best targets are increasing \u03c3 and decreasing \u03b4 \u2192 adoptive T-cell transfer, cytotoxic therapy [S2][S4].
- Caveat reported: a threshold \u03c6 \u2248 0.003 separates outcomes \u2014 below it the same therapies work for both equilibria; above it (e.g., \u03c6 = \u03bc = 0.00311) **no feasible modulation of immune parameters controls the T1-only equilibrium except by altering \u03b8\u2081 itself** [S10].

**Reduced antigen presentation** [S2]
- T1-only: best targets are **increasing \u03c1 and decreasing \u03bc** \u2192 CAR-T therapy, vaccines [S2][S4].
- Notably, increasing effector-cell presentation is *not* significantly helpful in this regime [S7].

**Combination therapy** [S1]
- Chemotherapy reduces tumor burden and targets sensitive populations; immunotherapy enhances immune-mediated suppression; given together they **expand the parameter region exhibiting tumor-free or low-population equilibria** [S1].
- Nonlinearity warning: modest parameter changes can induce qualitative transitions between immune control and tumor escape [S1].

**Why evasion parameters dominate** [S6]
- Very low \u03b8\u2081 permits stable large-*T*\u2081 equilibria; increasing \u03bc can create bistability near baseline when \u03b8\u2082 = 0.5; decreasing \u03c1 likewise creates bistability, and further decrease collapses to a single stable high-population equilibrium. **No bistability is observed for \u03b8\u2082 = 0.8**, and *T*\u2081/*T*\u2082 coexistence occupies a narrow range with still-sizeable tumors [S6].

**Resistance economics** [S3]
- Chemo-resistance scales kill rate by \u03b8\u2082 \u2208 [0,1] (T\u2081 fully susceptible at \u03b6\u2082 = 0, \u03b8\u2082 = 1); resistance carries a 30%-of-\u03b1 growth cost, and \u03c7 is set at 50% of \u03b1 precisely because chemo is treated as complementary to immunotherapy [S3].

**Gap**: the context gives explicit target tables only for the two evasion mechanisms; it does not quantify an optimal dosing schedule or report therapeutic focus for the combined chemo-resistant-plus-reduced-presentation case [S2][S11].
\end{AnswerText}

\textbf{Retrieved evidence index (all returned units; paths and scores) }
\begin{EvidenceText}
[1] score=0.7176 path=2604.01385
[2] score=0.6839 path=2604.01385
[3] score=0.6773 path=(raptor-node-unattributed)
[4] score=0.6764 path=(raptor-node-unattributed)
[5] score=0.6756 path=2604.01385
[6] score=0.6657 path=2604.01385
[7] score=0.6622 path=2604.01385
[8] score=0.6532 path=2604.01385
[9] score=0.6517 path=2604.01385
[10] score=0.6495 path=(raptor-node-unattributed)
[11] score=0.6494 path=(raptor-node-unattributed)
[12] score=0.6436 path=(raptor-node-unattributed)
\end{EvidenceText}
\section{LightRAG mixed retrieval}
\subsection{BQ01: How does the Transformer compute attention, and why does it replace recurrence and convolution?}
\begin{tabular}{@{}ll@{}}
Arm & LightRAG mixed retrieval\\
Generator & omp / deepseek-flash\\
HTTP / success & 200 / True\\
Wall time / cost & 162.0 s / 0.027018 USD\\
Judge scores & F=5, C=2, G=5, A=5; overall=4.25\\
Verdict & partial\\
\end{tabular}

\textbf{Index/retrieval metadata}
\begin{EvidenceText}
{
  "retrieval": "lightrag aquery_data mode=mix (chunks + entity descriptions)",
  "n_entities": 75,
  "n_relationships": 94
}
\end{EvidenceText}

\textbf{Dimension rationale}
\begin{EvidenceText}
{
  "factual": "All stated facts are correct: attention 'maps a query and a set of key-value pairs to an output' and models dependencies regardless of distance; the self-attention/intra-attention definition; multi-headed self-attention replacing recurrent layers; sinusoidal positional encodings added to embeddings because the authors hypothesized it lets the model easily learn to attend by relative positions (matches paper section 3.5). Nothing contradicts gold.",
  "completeness": "The 'how' half stays at definition level \u2014 it never says the output is computed as scaled dot products, softmax normalization, or a weighted value sum, and itself flags the formula as absent. The 'why' half mentions only the constant-sequential-operations advantage and three general advantages, omitting the ConvS2S-linear/ByteNet-logarithmic contrast, the reduced-resolution trade-off, and any empirical evidence (which it notes are not in the excerpts). Both halves are thin.",
  "groundedness": "Each claim is tagged to a source ([S12][S11][S10][S7][S5][S8][S1]) and matches gold or the real paper; it explicitly states quantitative BLEU/complexity comparisons 'are not included in the excerpts here' rather than inventing numbers.",
  "sources": "Clear per-claim attributions, plus the bibliographic note that the architectural claim is attributed to Vaswani et al., NeurIPS 2017 via [S8]."
}
\end{EvidenceText}

\textbf{Judge summary}
\begin{EvidenceText}
Error-free and clearly attributed, but both halves \u2014 especially how attention is actually computed \u2014 remain at a shallow definitional level.
\end{EvidenceText}

\textbf{Complete answer}
\begin{AnswerText}
## How attention is computed (per the provided context)

- Attention maps a **query** and a set of **key\u2013value pairs** to an output, and it models dependencies "without regard to their distance in input or output sequences" [S12].
- Cross-attention heads transform decoder and encoder states into **query, key, and value vectors** [S11].
- Applied within a single sequence, this is called **self-attention** (intra-attention): it relates different positions of one sequence to compute a representation of that sequence [S10].
- The Transformer uses **multi-headed self-attention**, the mechanism that replaces recurrent layers in the encoder\u2013decoder [S7].
- Position information is not inherent to attention, so it is injected via **sinusoidal positional encodings** added to the embeddings, chosen because the authors hypothesized this lets the model easily learn to attend by relative positions [S5].

**Gap in the supplied context:** the explicit scaled dot-product equation (softmax of $QK^{\top}/\sqrt{d_k}$ applied to $V$) does not appear in the excerpts provided. The above is the mechanism as described, not the formula.

## Why it replaces recurrence and convolution

- The Transformer "dispenses entirely with recurrence and convolutions, relying solely on attention mechanisms," and is based on self-attention [S8].
- Self-attention **connects all positions with a constant number of sequentially executed operations**, giving advantages in **computational complexity, parallelization, and path length for long-range dependencies** [S10]. Constant-depth connectivity is the core argument against recurrence, whose sequential nature lengthens both the operation chain and the dependency path.
- The paper's framing is that dominant sequence-transduction models had been built on complex recurrent or convolutional networks [S1]; the Transformer replaces them with attention alone [S7, S8].

Note: the context attributes the architectural claim (attention replaces recurrence/convolution) to Vaswani et al., "Attention Is All You Need," NeurIPS 2017 [S8]; quantitative BLEU/complexity comparisons in that paper are not included in the excerpts here.
\end{AnswerText}

\textbf{Retrieved evidence index (all returned units; paths and scores) }
\begin{EvidenceText}
[1] score=None path=artificial-intelligence__1706.03762__attention-is-all-you-need.md (part 1 of 2).md
[2] score=None path=artificial-intelligence__2504.16021__navigating-the-state-of-cognitive-flow-conte.md.md
[3] score=None path=artificial-intelligence__1706.03762__attention-is-all-you-need.md (part 2 of 2).md
[4] score=None path=genomics__2401.06823__interpretable-deep-learning-in-single-cell-o (part 1 of 2).md
[5] score=None path=artificial-intelligence__1706.03762__attention-is-all-you-need.md (part 1 of 2).md
[6] score=None path=Lost in the Middle (part 2 of 3).md
[7] score=None path=artificial-intelligence__1706.03762__attention-is-all-you-need.md (part 2 of 2).md
[8] score=None path=Lost in the Middle (part 2 of 3).md<SEP>Lost in the Middle (part 1 of 3).md<SEP>nlp__1810.04805__bert-pre-training-of-deep-bidirectional-transfor.md (part 1 of 3).md<SEP>large-language-models__2005.14165__language-models-are-few-shot-learners.md (part 6 of 13).md<SEP>nlp__1810.04805__bert-pre-training-of-deep-bidirectional-transfor.md (part 2 of 3).md<SEP>artificial-intelligence__2510.13343__aoad-mat-transformer-based-multi-agent-deep (part 1 of 2).md<SEP>large-language-models__2005.14165__language-models-are-few-shot-learners.md (part 13 of 13).md<SEP>nlp__2503.21676__how-do-language-models-learn-facts-dynamics.md (part 2 of 4).md<SEP>nlp__2110.07844__modeling-endorsement-for-multi-document-abst (part 1 of 3).md<SEP>robotics__2410.10669__mlp-slam-multilayer-perceptron-based-simultaneou.md (part 2 of 2).md<SEP>nlp__2503.21676__how-do-language-models-learn-facts-dynamics.md (part 1 of 4).md<SEP>artificial-intelligence__1706.03762__attention-is-all-you-need.md (part 1 of 2).md<SEP>artificial-intelligence__1706.03762__attention-is-all-you-need.md (part 2 of 2).md<SEP>artificial-intelligence__2510.13343__aoad-mat-transformer-based-multi-agent-deep (part 2 of 2).md<SEP>speech__2109.14357__comparison-of-self-supervised-speech-pre-tra.md (part 1 of 2).md<SEP>large-language-models__2005.14165__language-models-are-few-shot-learners.md (part 1 of 13).md<SEP>clinical-medicine__2212.13138__large-language-models-encode-clinical-knowledge.md (part 1 of 8).md<SEP>genomics__2401.06823__interpretable-deep-learning-in-single-cell-o (part 2 of 2).md
[9] score=None path=nlp__2503.21676__how-do-language-models-learn-facts-dynamics.md (part 2 of 4).md
[10] score=None path=Lost in the Middle (part 2 of 3).md<SEP>nlp__2110.07844__modeling-endorsement-for-multi-document-abst (part 1 of 3).md<SEP>artificial-intelligence__1706.03762__attention-is-all-you-need.md (part 1 of 2).md
[11] score=None path=nlp__2110.07844__modeling-endorsement-for-multi-document-abst (part 1 of 3).md
[12] score=None path=artificial-intelligence__1706.03762__attention-is-all-you-need.md (part 1 of 2).md<SEP>artificial-intelligence__1706.03762__attention-is-all-you-need.md (part 2 of 2).md
\end{EvidenceText}
\subsection{BQ02: What defines NISQ technology and what is its central limitation?}
\begin{tabular}{@{}ll@{}}
Arm & LightRAG mixed retrieval\\
Generator & omp / deepseek-flash\\
HTTP / success & 200 / True\\
Wall time / cost & 100.3 s / 0.017931000000000002 USD\\
Judge scores & F=5, C=5, G=5, A=5; overall=5.00\\
Verdict & correct\\
\end{tabular}

\textbf{Index/retrieval metadata}
\begin{EvidenceText}
{
  "retrieval": "lightrag aquery_data mode=mix (chunks + entity descriptions)",
  "n_entities": 80,
  "n_relationships": 65
}
\end{EvidenceText}

\textbf{Dimension rationale}
\begin{EvidenceText}
{
  "factual": "All correct: 50 to a few hundred qubits, noisy and 'not protected by quantum error correction', near future; 50-100 qubits may surpass classical digital computers (gold abstract); noise limits reliably executable circuit size (gold); QEC's high overhead keeps NISQ away from correction near term, making mitigation the relevant strategy; imperfect performance complicates validating outputs and limits computational power.",
  "completeness": "Definition (including the no-QEC clause) and central limitation (gate noise capping reliable circuit depth, with the QEC-overhead reason and the validation consequence) both fully answered; the closing note honestly states there is no single 'official' sentence beyond this recurring one.",
  "groundedness": "Each sentence carries a source tag ([S2][S3][S8][S11]) and the claims are consistent with gold and the paper; the mitigation-vs-correction point is attributed to the benchmarking source rather than invented.",
  "sources": "Clear per-claim attributions plus a meta-note identifying which sources jointly support the recurring limitation."
}
\end{EvidenceText}

\textbf{Judge summary}
\begin{EvidenceText}
Concise but complete: correct definition including the error-correction-free aspect, the noise limitation with its QEC-overhead explanation, all attributed.
\end{EvidenceText}

\textbf{Complete answer}
\begin{AnswerText}
**Definition.** NISQ stands for "Noisy Intermediate-Scale Quantum" technology: quantum computers with roughly 50 to a few hundred qubits whose gates are noisy and not protected by quantum error correction, expected in the near future [S2][S8]. Devices of 50\u2013100 qubits may perform tasks beyond today's classical digital computers [S2].

**Central limitation.** Noise in quantum gates limits the size of quantum circuits that can be executed reliably [S2]. Because quantum error correction carries a high overhead cost, NISQ devices will not use it in the near term, so noise mitigation rather than correction is the relevant strategy [S3]. That imperfect performance also makes it hard to validate that the computer produced the right answer and severely limits its computational power [S11].

Note the context gives no single-sentence "official" statement of a central limitation beyond these \u2014 the recurring one across [S2], [S3], and [S11] is gate noise capping reliable circuit depth.
\end{AnswerText}

\textbf{Retrieved evidence index (all returned units; paths and scores) }
\begin{EvidenceText}
[1] score=None path=quantum-physics__1801.00862__quantum-computing-in-the-nisq-era-and-beyond.md (part 2 of 3).md
[2] score=None path=quantum-physics__1801.00862__quantum-computing-in-the-nisq-era-and-beyond.md (part 1 of 3).md
[3] score=None path=quantum-physics__1801.00862__quantum-computing-in-the-nisq-era-and-beyond.md (part 2 of 3).md
[4] score=None path=quantum-physics__2303.00445__benchmarking-noisy-intermediate-scale-quantu (part 1 of 4).md
[5] score=None path=quantum-physics__1801.00862__quantum-computing-in-the-nisq-era-and-beyond.md (part 1 of 3).md
[6] score=None path=quantum-physics__1801.00862__quantum-computing-in-the-nisq-era-and-beyond.md (part 3 of 3).md
[7] score=None path=quantum-physics__1801.00862__quantum-computing-in-the-nisq-era-and-beyond.md (part 2 of 3).md
[8] score=None path=quantum-physics__1801.00862__quantum-computing-in-the-nisq-era-and-beyond.md (part 1 of 3).md
[9] score=None path=quantum-physics__1311.2485__continuous-time-quantum-error-correction.md (part 1 of 4).md<SEP>quantum-physics__1610.04013__entanglement-assisted-quantum-error-correcti.md (part 2 of 3).md<SEP>quantum-physics__1801.00862__quantum-computing-in-the-nisq-era-and-beyond.md (part 1 of 3).md<SEP>quantum-physics__1801.00862__quantum-computing-in-the-nisq-era-and-beyond.md (part 2 of 3).md
[10] score=None path=quantum-physics__2311.00487__echo-evolution-data-generation-for-quantum-e (part 2 of 3).md
[11] score=None path=quantum-physics__1801.00862__quantum-computing-in-the-nisq-era-and-beyond.md (part 1 of 3).md<SEP>quantum-physics__1801.00862__quantum-computing-in-the-nisq-era-and-beyond.md (part 2 of 3).md
[12] score=None path=quantum-physics__2303.00445__benchmarking-noisy-intermediate-scale-quantu (part 3 of 4).md
\end{EvidenceText}
\subsection{BQ03: What is MultiMedQA and how did Flan-PaLM perform on the MedQA benchmark?}
\begin{tabular}{@{}ll@{}}
Arm & LightRAG mixed retrieval\\
Generator & omp / deepseek-flash\\
HTTP / success & 200 / True\\
Wall time / cost & 93.5 s / 0.022287 USD\\
Judge scores & F=5, C=4, G=5, A=5; overall=4.75\\
Verdict & partial\\
\end{tabular}

\textbf{Index/retrieval metadata}
\begin{EvidenceText}
{
  "retrieval": "lightrag aquery_data mode=mix (chunks + entity descriptions)",
  "n_entities": 46,
  "n_relationships": 141
}
\end{EvidenceText}

\textbf{Dimension rationale}
\begin{EvidenceText}
{
  "factual": "All recoverable claims are correct: six existing datasets named plus the new HealthSearchQA; MedQA as the USMLE-style subset with 4/5 options and 11,450 dev/1,273 test; the gold-verbatim headline that Flan-PaLM 'exceeds the previous state of the art on MedQA (USMLE) by over 17%'. Its refusal on the absolute 67.6% figure is honest and correctly explained ('the supplied excerpt of [S1] cuts off mid-sentence').",
  "completeness": "The definition half is thorough (composition, spans, MedQA's structure and role). The performance half gives the relative margin (>17%) and prompting comparisons qualitatively, but the key absolute number (67.6%) is missing because of its truncated context, and no quantitative breakdown (few-shot/CoT/self-consistency values) is given \u2014 acknowledged as a limit.",
  "groundedness": "Every claim is tagged ([S6][S8][S9][S1][S3]) and supported; the truncation is reported rather than papered over, so no fabricated numbers appear.",
  "sources": "Clear per-claim attributions, including the specific note of where [S1] cuts off mid-sentence ('On the MedQA dataset con\u2026')."
}
\end{EvidenceText}

\textbf{Judge summary}
\begin{EvidenceText}
Accurate and honest, with a complete benchmark definition, but the headline 67.6% accuracy is absent (declared unrecoverable from its truncated excerpts).
\end{EvidenceText}

\textbf{Complete answer}
\begin{AnswerText}
**MultiMedQA** \u2014 a benchmark for medical question answering that combines six existing open QA datasets \u2014 MedQA, MedMCQA, PubMedQA, LiveQA, MedicationQA, and MMLU clinical topics \u2014 plus the newly introduced **HealthSearchQA** dataset. It spans professional medical exams, medical research, and consumer health queries, and it was the evaluation suite used for both Flan-PaLM and Med-PaLM [S6][S8]. Within it, **MedQA** is the USMLE-style multiple-choice subset (four or five options; 11,450 dev / 1,273 test questions) used to measure multiple-choice accuracy and scaling behavior [S9].

**Flan-PaLM on MedQA** \u2014 the paper's headline result is that Flan-PaLM **exceeds the previous state of the art on MedQA (USMLE) by over 17%** [S1]. Supporting detail: scaling plots show PaLM and Flan-PaLM under few-shot prompting on MedQA/MedMCQA with performance rising steeply with model size, and Flan-PaLM is also compared under few-shot vs. self-consistency prompting [S3].

Limit: the supplied excerpt of [S1] cuts off mid-sentence ("On the MedQA dataset con\u2026"), so the **absolute accuracy percentage** for Flan-PaLM on MedQA is not recoverable from this context \u2014 only the "over 17%" margin over prior SOTA [S1].
\end{AnswerText}

\textbf{Retrieved evidence index (all returned units; paths and scores) }
\begin{EvidenceText}
[1] score=None path=clinical-medicine__2212.13138__large-language-models-encode-clinical-knowledge.md (part 2 of 8).md
[2] score=None path=large-language-models__2502.00306__riddle-me-this-stealthy-membership-inference (part 4 of 4).md
[3] score=None path=clinical-medicine__2212.13138__large-language-models-encode-clinical-knowledge.md (part 5 of 8).md
[4] score=None path=Lost in the Middle (part 1 of 3).md
[5] score=None path=clinical-medicine__2212.13138__large-language-models-encode-clinical-knowledge.md (part 1 of 8).md
[6] score=None path=clinical-medicine__2212.13138__large-language-models-encode-clinical-knowledge.md (part 2 of 8).md<SEP>clinical-medicine__2212.13138__large-language-models-encode-clinical-knowledge.md (part 3 of 8).md<SEP>clinical-medicine__2212.13138__large-language-models-encode-clinical-knowledge.md (part 1 of 8).md
[7] score=None path=large-language-models__2502.00306__riddle-me-this-stealthy-membership-inference (part 3 of 4).md<SEP>Chain-of-Thought Prompting (part 2 of 3).md<SEP>nlp__2503.21676__how-do-language-models-learn-facts-dynamics.md (part 2 of 4).md<SEP>large-language-models__2402.12317__evor-evolving-retrieval-for-code-generation (part 2 of 3).md<SEP>large-language-models__2402.12317__evor-evolving-retrieval-for-code-generation (part 1 of 3).md<SEP>nlp__2503.21676__how-do-language-models-learn-facts-dynamics.md (part 1 of 4).md<SEP>speech__2509.05359__an-empirical-analysis-of-discrete-unit-repre (part 1 of 2).md<SEP>large-language-models__2005.14165__language-models-are-few-shot-learners.md (part 9 of 13).md<SEP>large-language-models__2502.00306__riddle-me-this-stealthy-membership-inference (part 1 of 4).md<SEP>large-language-models__2312.10793__demystifying-instruction-mixing-for-fine-tun.md (part 1 of 2).md<SEP>Chain-of-Thought Prompting (part 1 of 3).md<SEP>clinical-medicine__2212.13138__large-language-models-encode-clinical-knowledge.md (part 4 of 8).md<SEP>clinical-medicine__2212.13138__large-language-models-encode-clinical-knowledge.md (part 3 of 8).md<SEP>clinical-medicine__2212.13138__large-language-models-encode-clinical-knowledge.md (part 1 of 8).md<SEP>clinical-medicine__2212.13138__large-language-models-encode-clinical-knowledge.md (part 5 of 8).md
[8] score=None path=clinical-medicine__2212.13138__large-language-models-encode-clinical-knowledge.md (part 2 of 8).md
[9] score=None path=clinical-medicine__2212.13138__large-language-models-encode-clinical-knowledge.md (part 2 of 8).md<SEP>clinical-medicine__2212.13138__large-language-models-encode-clinical-knowledge.md (part 3 of 8).md<SEP>clinical-medicine__2212.13138__large-language-models-encode-clinical-knowledge.md (part 1 of 8).md<SEP>clinical-medicine__2212.13138__large-language-models-encode-clinical-knowledge.md (part 5 of 8).md
[10] score=None path=clinical-medicine__2212.13138__large-language-models-encode-clinical-knowledge.md (part 2 of 8).md<SEP>clinical-medicine__2212.13138__large-language-models-encode-clinical-knowledge.md (part 3 of 8).md<SEP>clinical-medicine__2212.13138__large-language-models-encode-clinical-knowledge.md (part 6 of 8).md
[11] score=None path=clinical-medicine__2212.13138__large-language-models-encode-clinical-knowledge.md (part 2 of 8).md
[12] score=None path=large-language-models__2005.14165__language-models-are-few-shot-learners.md (part 6 of 13).md<SEP>nlp__2310.10981__instructive-dialogue-summarization-with-quer (part 2 of 4).md<SEP>large-language-models__2005.14165__language-models-are-few-shot-learners.md (part 12 of 13).md<SEP>large-language-models__2005.14165__language-models-are-few-shot-learners.md (part 9 of 13).md<SEP>large-language-models__2005.14165__language-models-are-few-shot-learners.md (part 1 of 13).md<SEP>large-language-models__2005.14165__language-models-are-few-shot-learners.md (part 11 of 13).md
\end{EvidenceText}
\subsection{BQ04: Which three independent ontologies subdivide the Gene Ontology?}
\begin{tabular}{@{}ll@{}}
Arm & LightRAG mixed retrieval\\
Generator & omp / deepseek-flash\\
HTTP / success & 200 / True\\
Wall time / cost & 38.5 s / 0.014175 USD\\
Judge scores & F=5, C=5, G=5, A=5; overall=5.00\\
Verdict & correct\\
\end{tabular}

\textbf{Index/retrieval metadata}
\begin{EvidenceText}
{
  "retrieval": "lightrag aquery_data mode=mix (chunks + entity descriptions)",
  "n_entities": 84,
  "n_relationships": 45
}
\end{EvidenceText}

\textbf{Dimension rationale}
\begin{EvidenceText}
{
  "factual": "The three names are correct; the glosses ('molecular activities of gene products', 'biological processes', 'locations within the cell') match the primer's actual definitions of each ontology; the quoted characterization 'subdivided into three ontologies: Molecular Function, Biological Process, and Cellular Component' closely tracks the gold sentence.",
  "completeness": "All three ontologies named with brief definitions \u2014 fully answers the question.",
  "groundedness": "Every line carries a source tag ([S12][S8][S10][S9]) and the content matches the gold and the primer; no unsupported claims.",
  "sources": "Per-item attributions for each ontology plus a quoted sentence attributed to [S9] \u2014 clear and specific."
}
\end{EvidenceText}

\textbf{Judge summary}
\begin{EvidenceText}
Correct and complete answer with per-item sourcing and accurate ontology definitions.
\end{EvidenceText}

\textbf{Complete answer}
\begin{AnswerText}
The Gene Ontology is subdivided into three independent ontologies:

1. **Molecular Function** \u2014 molecular activities of gene products [S12]
2. **Biological Process** \u2014 biological processes [S8]
3. **Cellular Component** \u2014 locations within the cell [S10]

The Gene Ontology is described as a controlled vocabulary "subdivided into three ontologies: Molecular Function, Biological Process, and Cellular Component" [S9].
\end{AnswerText}

\textbf{Retrieved evidence index (all returned units; paths and scores) }
\begin{EvidenceText}
[1] score=None path=genomics__1602.01876__primer-on-the-gene-ontology.md (part 2 of 2).md
[2] score=None path=genomics__1602.01876__primer-on-the-gene-ontology.md (part 1 of 2).md
[3] score=None path=genomics__2508.18304__sci2cl-effectively-integrating-single-cell-m (part 1 of 4).md
[4] score=None path=genomics__2401.06823__interpretable-deep-learning-in-single-cell-o (part 2 of 2).md
[5] score=None path=databases__2009.10669__there-is-no-such-thing-as-an-index-or-the-ne (part 3 of 4).md
[6] score=None path=genomics__1602.01876__primer-on-the-gene-ontology.md (part 1 of 2).md
[7] score=None path=genomics__2401.06823__interpretable-deep-learning-in-single-cell-o (part 1 of 2).md
[8] score=None path=genomics__1602.01876__primer-on-the-gene-ontology.md (part 1 of 2).md
[9] score=None path=genomics__1602.01876__primer-on-the-gene-ontology.md (part 1 of 2).md<SEP>genomics__2401.06823__interpretable-deep-learning-in-single-cell-o (part 1 of 2).md<SEP>genomics__1602.01876__primer-on-the-gene-ontology.md (part 2 of 2).md
[10] score=None path=genomics__1602.01876__primer-on-the-gene-ontology.md (part 1 of 2).md
[11] score=None path=genomics__2110.07787__analysis-and-visualization-of-spatial-transc.md (part 3 of 3).md
[12] score=None path=genomics__1602.01876__primer-on-the-gene-ontology.md (part 1 of 2).md
\end{EvidenceText}
\subsection{BQ05: Why does historical European economic growth contradict Unified Growth Theory?}
\begin{tabular}{@{}ll@{}}
Arm & LightRAG mixed retrieval\\
Generator & omp / deepseek-flash\\
HTTP / success & 200 / True\\
Wall time / cost & 55.6 s / 0.020783999999999997 USD\\
Judge scores & F=5, C=4, G=5, A=5; overall=4.75\\
Verdict & correct\\
\end{tabular}

\textbf{Index/retrieval metadata}
\begin{EvidenceText}
{
  "retrieval": "lightrag aquery_data mode=mix (chunks + entity descriptions)",
  "n_entities": 75,
  "n_relationships": 80
}
\end{EvidenceText}

\textbf{Dimension rationale}
\begin{EvidenceText}
{
  "factual": "All claims correct: 'should have produced the best and the most convincing confirmation' and 'no connection, whatever' are gold quotes; regimes 'have no relevance to the description, let alone to the explanation'; 'there was never a transition from stagnation'; Industrial Revolution 'did not boost the economic growth in Eastern Europe'; the Eastern Europe fit (AD 1 overshoot 51%, singularity 1915, diversion ~1890, 25-year bypass) is consistent with other answers citing the paper.",
  "completeness": "Covers the test-case framing, regime irrelevance, absent stagnation epoch, Industrial-Revolution failure, and hyperbolic growth \u2014 but omits the explicit no-trap/no-escape argument and the paper's data-misuse (four-selected-points) critique that several other answers include.",
  "groundedness": "Every bullet carries source tags with quotes matching the gold or the paper; the closing note honestly lists what its context lacks (Western Europe figures, exact fitting method).",
  "sources": "Per-claim attributions [S1][S3][S7][S9][S10][S12] plus an explicit statement of the context's gaps."
}
\end{EvidenceText}

\textbf{Judge summary}
\begin{EvidenceText}
Accurate and well-sourced on all main contradictions, slightly less complete than peers on the trap-escape and data-handling arguments.
\end{EvidenceText}

\textbf{Complete answer}
\begin{AnswerText}
The sources argue the contradiction is empirical, at the level of the growth trajectory itself rather than a detail:

- **Wrong test case fails.** Western and Eastern Europe "should have produced the best and the most convincing confirmation" of Unified Growth Theory, since Western Europe was the centre of the Industrial Revolution \u2014 which Galor treats as the prime engine of growth. Instead the data "show a remarkable disagreement" with the theory, with "no connection, whatever" between data and theory [S3][S10].
- **The three regimes don't map onto the data.** Galor's Malthusian stagnation, post-Malthusian, and sustained-growth regimes "have no relevance to the description, let alone to the explanation, of the mechanism of the economic growth" [S1][S9].
- **The Malthusian stagnation epoch is absent.** The "alleged Malthusian regime of stagnation did not exist," and the data show "there was never a transition from stagnation" [S1][S3].
- **The Industrial Revolution didn't change the trajectory.** It "did not boost the economic growth in Eastern Europe" [S1][S12].
- **Growth is hyperbolic, not regime-switching.** The fitted curve overshoots at AD 1 by 51%, with a singularity in 1915; actual growth was diverted to a slower trajectory around 1890, bypassing the singularity by 25 years \u2014 i.e., growth follows a single hyperbolic pattern rather than stage-wise transitions [S1][S7].

So the contradiction is that the theory's predicted sequence of regimes and its stagnation-to-sustained-growth transition are not observable in the European series, and the supposed engine of the escape (the Industrial Revolution) left no visible break in the growth path [S1][S3][S12].

The context does not supply the paper's Western Europe figures, its exact fitting method, or the full "escape from Malthusian trap" discussion beyond the quoted fragments.
\end{AnswerText}

\textbf{Retrieved evidence index (all returned units; paths and scores) }
\begin{EvidenceText}
[1] score=None path=economics__1512.08067__unified-growth-theory-contradicted-by-the-econom.md.md
[2] score=None path=databases__2009.10669__there-is-no-such-thing-as-an-index-or-the-ne (part 3 of 4).md
[3] score=None path=economics__1512.08067__unified-growth-theory-contradicted-by-the-econom.md.md
[4] score=None path=economics__1807.09475__cap-and-monetary-policy.md
[5] score=None path=economics__2011.14424__on-the-effectiveness-of-the-european-central (part 2 of 5).md
[6] score=None path=economics__1512.08067__unified-growth-theory-contradicted-by-the-econom.md.md
[7] score=None path=economics__1512.08067__unified-growth-theory-contradicted-by-the-econom.md.md
[8] score=None path=economics__1512.08067__unified-growth-theory-contradicted-by-the-econom.md.md
[9] score=None path=economics__1512.08067__unified-growth-theory-contradicted-by-the-econom.md.md
[10] score=None path=climate-science__1503.07557__precipitation-extremes-under-climate-change.md (part 2 of 3).md<SEP>economics__1512.08067__unified-growth-theory-contradicted-by-the-econom.md.md
[11] score=None path=economics__1512.08067__unified-growth-theory-contradicted-by-the-econom.md.md
[12] score=None path=economics__1512.08067__unified-growth-theory-contradicted-by-the-econom.md.md
\end{EvidenceText}
\subsection{BQ06: Which physical factor primarily controls the response of precipitation extremes to climate change?}
\begin{tabular}{@{}ll@{}}
Arm & LightRAG mixed retrieval\\
Generator & omp / deepseek-flash\\
HTTP / success & 200 / True\\
Wall time / cost & 136.1 s / 0.025785000000000002 USD\\
Judge scores & F=5, C=5, G=5, A=5; overall=5.00\\
Verdict & correct\\
\end{tabular}

\textbf{Index/retrieval metadata}
\begin{EvidenceText}
{
  "retrieval": "lightrag aquery_data mode=mix (chunks + entity descriptions)",
  "n_entities": 78,
  "n_relationships": 50
}
\end{EvidenceText}

\textbf{Dimension rationale}
\begin{EvidenceText}
{
  "factual": "All correct: thermodynamic contribution described as robust/well-understood changes in saturation specific humidity; tracked changes 'close to what is expected from scaling with surface specific humidity, implying a lower rate of increase than scaling with column water vapor' (a real distinction from the paper); Eq. (1) with fixed epsilon reproducing CMIP3 fractional changes; microphysical contribution defined via precipitation efficiency ('ratio of surface precipitation to column-integrated net condensation', fall-speed related); abstract's 'several physical contributions govern' framing.",
  "completeness": "Names the thermodynamic contribution as the strongest, robustly attested factor, explains why (it reproduces simulated responses when epsilon fixed), lists the modulators, and adds the convection/tropics caveat \u2014 fully answers with honest calibration.",
  "groundedness": "Every statement attributed to [S2][S4][S7][S8][S11][S12]; the truncation of the abstract sentence is reported rather than papered over; no unsupported claims.",
  "sources": "Per-claim attributions with an explicit caveat about the fragmentary context."
}
\end{EvidenceText}

\textbf{Judge summary}
\begin{EvidenceText}
Well-calibrated answer identifying the thermodynamic/humidity-scaling control with the paper's supporting evidence and honestly noting the absence of an explicit 'primary' statement.
\end{EvidenceText}

\textbf{Complete answer}
\begin{AnswerText}
**Short answer:** the excerpts do not contain an explicit statement naming a single factor as *the* primary control \u2014 the abstract says only that "several physical contributions govern" the response [S2]. The strongest and most robustly attested factor in the context is the **thermodynamic contribution: scaling with surface specific humidity** (Clausius\u2013Clapeyron-type scaling).

Evidence:

- The thermodynamic component is described as "the robust, well-understood component of the precipitation extreme response arising from changes in **saturation specific humidity with temperature**" [S11], and tracked changes are "close to what is expected from scaling with surface specific humidity, implying a lower rate of increase than scaling with column water vapor" [S11].
- Eq. (1) with precipitation efficiency \u03b5 held fixed **reproduces** fractional changes in precipitation extremes in CMIP3 simulations [S4] \u2014 i.e., the humidity-scaling part accounts for the bulk of the response, with other terms being corrections.
- Deviations from this scaling are attributed to the **microphysical contribution**, defined as the part "arising from changes in precipitation efficiency" [S7], i.e., \u03b5 is "the ratio of surface precipitation to column-integrated net condensation" and "a factor whose changes with warming contribute to deviations from Clausius\u2013Clapeyron scaling, related partly to hydrometeor fall speed" [S12].
- Additional factors named in the context include vertical velocities (in the global-warming/precipitation-extremes analysis framing [S8]) and, for extratropical extremes, cyclone/frontal association [S4].

Caveat: the context is fragmentary (the abstract sentence is truncated at "Several physical contributions govern the response of precipitation extr\u2026" [S2]), so a definitive "primarily controls" claim cannot be sourced from these passages. Based only on what is present, the surface-specific-humidity (thermodynamic) scaling is the indicated dominant, well-understood term, with precipitation efficiency/microphysics and dynamical (vertical velocity, cyclones/fronts) effects as modulators \u2014 the sensitivity remaining uncertain where convection is important, and possibly higher in the tropics than extratropics [S2].
\end{AnswerText}

\textbf{Retrieved evidence index (all returned units; paths and scores) }
\begin{EvidenceText}
[1] score=None path=climate-science__1503.07557__precipitation-extremes-under-climate-change.md (part 3 of 3).md
[2] score=None path=climate-science__1503.07557__precipitation-extremes-under-climate-change.md (part 1 of 3).md
[3] score=None path=climate-science__1503.07557__precipitation-extremes-under-climate-change.md (part 3 of 3).md
[4] score=None path=climate-science__1503.07557__precipitation-extremes-under-climate-change.md (part 1 of 3).md
[5] score=None path=climate-science__1503.07557__precipitation-extremes-under-climate-change.md (part 2 of 3).md
[6] score=None path=climate-science__2409.09639__gpc-m-global-precipitation-climatology-by-ma (part 2 of 2).md
[7] score=None path=climate-science__1503.07557__precipitation-extremes-under-climate-change.md (part 1 of 3).md
[8] score=None path=climate-science__2409.13934__on-the-extrapolation-of-generative-adversari (part 2 of 3).md<SEP>climate-science__1503.07557__precipitation-extremes-under-climate-change.md (part 1 of 3).md<SEP>power-systems__2207.08146__mapping-disruption-sources-in-the-power-grid.md (part 1 of 2).md<SEP>climate-science__1503.07557__precipitation-extremes-under-climate-change.md (part 3 of 3).md
[9] score=None path=climate-science__1503.07557__precipitation-extremes-under-climate-change.md (part 1 of 3).md
[10] score=None path=climate-science__2007.10574__impacts-of-the-madden-julian-oscillation-on (part 2 of 2).md<SEP>climate-science__2409.13934__on-the-extrapolation-of-generative-adversari (part 2 of 3).md<SEP>climate-science__2409.13934__on-the-extrapolation-of-generative-adversari (part 3 of 3).md<SEP>climate-science__1503.07557__precipitation-extremes-under-climate-change.md (part 1 of 3).md<SEP>climate-science__1503.07557__precipitation-extremes-under-climate-change.md (part 2 of 3).md<SEP>climate-science__1503.07557__precipitation-extremes-under-climate-change.md (part 3 of 3).md
[11] score=None path=climate-science__1503.07557__precipitation-extremes-under-climate-change.md (part 1 of 3).md
[12] score=None path=climate-science__1503.07557__precipitation-extremes-under-climate-change.md (part 1 of 3).md<SEP>climate-science__1503.07557__precipitation-extremes-under-climate-change.md (part 2 of 3).md
\end{EvidenceText}
\subsection{BQ07: What is MameLoshnLM and what parameter scale does it have?}
\begin{tabular}{@{}ll@{}}
Arm & LightRAG mixed retrieval\\
Generator & omp / deepseek-flash\\
HTTP / success & 200 / True\\
Wall time / cost & 47.9 s / 0.014565 USD\\
Judge scores & F=5, C=4, G=5, A=5; overall=4.75\\
Verdict & correct\\
\end{tabular}

\textbf{Index/retrieval metadata}
\begin{EvidenceText}
{
  "retrieval": "lightrag aquery_data mode=mix (chunks + entity descriptions)",
  "n_entities": 52,
  "n_relationships": 129
}
\end{EvidenceText}

\textbf{Dimension rationale}
\begin{EvidenceText}
{
  "factual": "Correct on everything stated: first open-source LM built specifically for Yiddish; 8 billion parameters; the four institutions match the gold author list; the note that other parts compare it against Llama 3.1 8B 'consistent with the 8B scale' is accurate and appropriately hedged ('no other parameter sizes are reported').",
  "completeness": "Both parts answered (identity and 8B scale), but thinly: it omits the continued-pretraining recipe from Llama 3.1 8B, does not name Oytser or Kashes, and gives no benchmark or training detail beyond the title-based benchmark mention.",
  "groundedness": "Every statement is attributed to [S1] with cross-references to [S3][S6][S12]; the content matches the gold abstract; nothing unsupported.",
  "sources": "Clear attributions with an explicit note about which sources say what and what is absent."
}
\end{EvidenceText}

\textbf{Judge summary}
\begin{EvidenceText}
Correct and honestly sourced answer to both parts, though with minimal elaboration beyond the abstract's definition and scale.
\end{EvidenceText}

\textbf{Complete answer}
\begin{AnswerText}
MameLoshnLM is the first open-source language model built specifically for Yiddish, introduced alongside an evaluation benchmark (as the paper title indicates, "Yiddish Language Model and Evaluation Benchmark") [S1]. It has **8 billion parameters** \u2014 an 8B scale model [S1].

Context notes: the abstract in [S1] attributes it to Bar-Ilan University, University of Cambridge, University of Washington, and Allen Institute for AI. Other parts of the same paper [S3][S6][S12] discuss analysis and comparisons against Llama 3.1 8B, consistent with the 8B scale, but no other parameter sizes are reported in the provided context.
\end{AnswerText}

\textbf{Retrieved evidence index (all returned units; paths and scores) }
\begin{EvidenceText}
[1] score=None path=nlp__2608.05850__mameloshnlm-yiddish-language-model-and-evalu.md (part 1 of 4).md
[2] score=None path=nlp__2310.10981__instructive-dialogue-summarization-with-quer (part 1 of 4).md
[3] score=None path=nlp__2608.05850__mameloshnlm-yiddish-language-model-and-evalu.md (part 2 of 4).md
[4] score=None path=nlp__1810.04805__bert-pre-training-of-deep-bidirectional-transfor.md (part 3 of 3).md
[5] score=None path=robotics__2410.10669__mlp-slam-multilayer-perceptron-based-simultaneou.md (part 2 of 2).md
[6] score=None path=nlp__2608.05850__mameloshnlm-yiddish-language-model-and-evalu.md (part 4 of 4).md
[7] score=None path=large-language-models__2402.11690__vision-flan-scaling-human-labeled-tasks-in-v.md (part 1 of 9).md
[8] score=None path=nlp__2503.21676__how-do-language-models-learn-facts-dynamics.md (part 4 of 4).md
[9] score=None path=large-language-models__2402.11690__vision-flan-scaling-human-labeled-tasks-in-v.md (part 2 of 9).md<SEP>large-language-models__2402.11690__vision-flan-scaling-human-labeled-tasks-in-v.md (part 1 of 9).md
[10] score=None path=statistics__2604.22320__nonparametric-estimation-of-isotropic-covari (part 7 of 9).md
[11] score=None path=economics__1911.06206__bayesian-state-space-modeling-for-analyzing (part 4 of 5).md
[12] score=None path=Chain-of-Thought Prompting (part 2 of 3).md<SEP>nlp__2608.05850__mameloshnlm-yiddish-language-model-and-evalu.md (part 4 of 4).md<SEP>Chain-of-Thought Prompting (part 1 of 3).md
\end{EvidenceText}
\subsection{BQ08: How does oxygen redox contribute to the capacity of battery cathodes?}
\begin{tabular}{@{}ll@{}}
Arm & LightRAG mixed retrieval\\
Generator & omp / deepseek-flash\\
HTTP / success & 200 / True\\
Wall time / cost & 102.7 s / 0.022293 USD\\
Judge scores & F=5, C=4, G=5, A=5; overall=4.75\\
Verdict & correct\\
\end{tabular}

\textbf{Index/retrieval metadata}
\begin{EvidenceText}
{
  "retrieval": "lightrag aquery_data mode=mix (chunks + entity descriptions)",
  "n_entities": 85,
  "n_relationships": 57
}
\end{EvidenceText}

\textbf{Dimension rationale}
\begin{EvidenceText}
{
  "factual": "All claims correct: anionic redox supplying capacity in addition to TM redox; 'useful capacity' (gold-verbatim); compensation for absent oxidizable TMs; Li2MnO3 as the prototypical material; Koga/Sathiya 2013 first conclusive demonstrations; orphaned oxygen states in Li-excess disordered spinels; ilmenite-type NaTMO3 under-coordination; Na2-xRuO3 delivering ~180 mAh/g (~1.4 electrons per formula unit) \u2014 consistent with the paper and peer answers; the debate over proportioning capacity between cationic and anionic effects.",
  "completeness": "Covers the addition mechanism, material origins, magnitude, and the attribution caveat \u2014 but omits the reversibility/cost side (irreversible O2 release, first-cycle capacity loss, fluorination), which the gold itself frames as central ('reversible reactions... offer pathways towards the reversible capacities'; 'inevitable diminishment of capacity... due to irreversible or poorly reversible redox phenomena').",
  "groundedness": "Every claim tagged ([S1][S3][S4][S6][S7][S9][S10][S11][S12]) with quotes; nothing unsupported.",
  "sources": "Clear per-claim attributions including the debate caveat attributed to [S4]."
}
\end{EvidenceText}

\textbf{Judge summary}
\begin{EvidenceText}
Accurate, well-sourced answer on how oxygen redox adds capacity and its magnitude, though it omits the reversibility and degradation caveats the gold emphasizes.
\end{EvidenceText}

\textbf{Complete answer}
\begin{AnswerText}
Oxygen (anionic) redox supplies capacity **in addition to** transition-metal redox: oxide ions themselves donate electrons during charge, so the measured capacity can exceed the limit set by cation oxidation alone [S4][S6]. In Li- and Na-ion oxide cathodes this participation is described as governing the "useful capacity" of the cell [S1][S6].

Mechanistic and material-specific points from the context:

- **Origin of the extra capacity.** Because anionic redox can compensate for the absence of oxidizable transition metals, it enables excess capacity that cation redox cannot reach [S7]. Li\u2082MnO\u2083 is the prototypical material for this behaviour [S7]; reversible anionic redox in TM-oxide Li-ion cathodes was first conclusively demonstrated in 2013 (Koga and Sathiya) [S7][S9].
- **Structural origin of the active oxygen.** In Li-excess disordered spinel cathodes, oxygen redox was first found to contribute significantly to overall capacity; it is attributed to "orphaned" oxygen states on oxygen ions coordinated by fewer than three transition metals [S10]. By contrast, in ilmenite-type Na\u2093TMO\u2083 the effect arises not from Na over-coordinating oxygen but from oxygen being **under-coordinated** [S3].
- **Magnitude.** Na\u2082\u208b\u2093RuO\u2083 (0 < x < 1.5) delivered ~180 mAh g\u207b\xb9, corresponding to ~1.4 electrons per formula unit [S11]. Anionic redox also underlies high-capacity designs such as Li\u2088ZrO\u2086 and Y-doped Na\u2082ZrO\u2083 [S7][S12].

Caveat stated in the source: attributing a measured capacity quantitatively to cationic versus anionic redox remains debated \u2014 experimental confirmation of oxygen redox is "not at all trivial," and proportioning capacity between the two effects is still under discussion [S4].
\end{AnswerText}

\textbf{Retrieved evidence index (all returned units; paths and scores) }
\begin{EvidenceText}
[1] score=None path=materials__2408.09942__oxygen-redox-in-battery-cathodes-a-brief-overvie.md (part 1 of 2).md
[2] score=None path=chemistry__1911.02792__machine-learning-for-molecular-simulation (part 2 of 4).md
[3] score=None path=materials__2408.09942__oxygen-redox-in-battery-cathodes-a-brief-overvie.md (part 2 of 2).md
[4] score=None path=materials__2408.09942__oxygen-redox-in-battery-cathodes-a-brief-overvie.md (part 1 of 2).md
[5] score=None path=materials__2211.04977__design-of-battery-materials-via-defects-and.md (part 3 of 3).md
[6] score=None path=materials__2408.09942__oxygen-redox-in-battery-cathodes-a-brief-overvie.md (part 2 of 2).md<SEP>materials__2408.09942__oxygen-redox-in-battery-cathodes-a-brief-overvie.md (part 1 of 2).md
[7] score=None path=materials__2211.04977__design-of-battery-materials-via-defects-and.md (part 3 of 3).md<SEP>materials__2408.09942__oxygen-redox-in-battery-cathodes-a-brief-overvie.md (part 1 of 2).md<SEP>materials__2211.04977__design-of-battery-materials-via-defects-and.md (part 2 of 3).md
[8] score=None path=materials__2408.09942__oxygen-redox-in-battery-cathodes-a-brief-overvie.md (part 2 of 2).md
[9] score=None path=materials__2408.09942__oxygen-redox-in-battery-cathodes-a-brief-overvie.md (part 1 of 2).md
[10] score=None path=materials__2408.09942__oxygen-redox-in-battery-cathodes-a-brief-overvie.md (part 2 of 2).md
[11] score=None path=materials__2408.09942__oxygen-redox-in-battery-cathodes-a-brief-overvie.md (part 2 of 2).md
[12] score=None path=materials__2211.04977__design-of-battery-materials-via-defects-and.md (part 3 of 3).md<SEP>materials__2211.04977__design-of-battery-materials-via-defects-and.md (part 2 of 3).md
\end{EvidenceText}
\subsection{BQ09: What problem of visual SLAM in dynamic environments does MLP-SLAM address and how?}
\begin{tabular}{@{}ll@{}}
Arm & LightRAG mixed retrieval\\
Generator & omp / deepseek-flash\\
HTTP / success & 200 / True\\
Wall time / cost & 143.3 s / 0.0219072 USD\\
Judge scores & F=5, C=4, G=5, A=5; overall=4.75\\
Verdict & correct\\
\end{tabular}

\textbf{Index/retrieval metadata}
\begin{EvidenceText}
{
  "retrieval": "lightrag aquery_data mode=mix (chunks + entity descriptions)",
  "n_entities": 79,
  "n_relationships": 94
}
\end{EvidenceText}

\textbf{Dimension rationale}
\begin{EvidenceText}
{
  "factual": "All claims correct: degradation in crowded outdoor scenes with the 'artificial mathematical methods' framing of limited generality and information loss (gold-verbatim); the stated target of addressing SLAM degradation in dynamic environments; a real-time stereo MLP system using 'complete geometric information'; ORB-SLAM2 base with discriminator module and a preprocessing/feature-extraction/detection/tracking stage; the DFM mean-depth/standard-deviation interval rule reclassifying out-of-interval points as static.",
  "completeness": "Both halves addressed at core level (problem; MLP discriminator + complete geometry + DFM), but it omits the >50,000-point dataset contribution, any quantitative results, and the MLP's input specifics \u2014 the latter honestly declared outside its context.",
  "groundedness": "Claims tagged ([S3][S4][S8]) and matching gold; the 'Not in the provided context' list is accurate to its excerpt limits.",
  "sources": "Clear per-claim attributions with an explicit gaps list."
}
\end{EvidenceText}

\textbf{Judge summary}
\begin{EvidenceText}
Correctly answers the problem and the core mechanism with honest context-gap disclosure, though without the dataset contribution or results.
\end{EvidenceText}

\textbf{Complete answer}
\begin{AnswerText}
**Problem.** V-SLAM reaches high precision when few objects move, but its accuracy degrades in environments crowded with movable objects (pedestrians, cars, buses) \u2014 typical outdoor scenes. Prior attempts use "artificial mathematical methods" to handle this, which the authors say suffer from limited generality and information loss [S3]. MLP-SLAM's stated target is therefore SLAM degradation in dynamic environments [S8].

**How.**
- Core mechanism: an MLP-based real-time **stereo** system that distinguishes dynamic from static feature points, using "complete geometric information" rather than hand-crafted mathematical rules [S8].
- Architecture: built on **ORB-SLAM2** with an added dynamic/static object discriminator module; four main parts, with a preprocessing/feature-extraction/object-detection/object-tracking stage [S4], [S8].
- Concretely in the coarse estimation stage: it computes the mean depth and standard deviation $\delta_d$ of candidate points and discards any feature point outside $[\mu_d - \eta\delta_d,\ \mu_d + \eta\delta_d]$ from the potential-dynamic set, treating it as static [S4].

**Not in the provided context:** the details of the MLP's inputs/training, how the discriminator is integrated into tracking/local mapping, and the full abstract sentence following "we propose a\u2026" [S3]. So the exact network design cannot be stated from these excerpts alone.
\end{AnswerText}

\textbf{Retrieved evidence index (all returned units; paths and scores) }
\begin{EvidenceText}
[1] score=None path=robotics__2410.10669__mlp-slam-multilayer-perceptron-based-simultaneou.md (part 2 of 2).md
[2] score=None path=robotics__2410.10669__mlp-slam-multilayer-perceptron-based-simultaneou.md (part 2 of 2).md
[3] score=None path=robotics__2410.10669__mlp-slam-multilayer-perceptron-based-simultaneou.md (part 1 of 2).md
[4] score=None path=robotics__2410.10669__mlp-slam-multilayer-perceptron-based-simultaneou.md (part 1 of 2).md
[5] score=None path=robotics__2509.14191__mcgs-slam-a-multi-camera-slam-framework-usin.md (part 2 of 2).md
[6] score=None path=robotics__2410.10669__mlp-slam-multilayer-perceptron-based-simultaneou.md (part 1 of 2).md
[7] score=None path=robotics__2410.10669__mlp-slam-multilayer-perceptron-based-simultaneou.md (part 2 of 2).md<SEP>robotics__2410.10669__mlp-slam-multilayer-perceptron-based-simultaneou.md (part 1 of 2).md
[8] score=None path=robotics__2410.10669__mlp-slam-multilayer-perceptron-based-simultaneou.md (part 2 of 2).md<SEP>robotics__2410.10669__mlp-slam-multilayer-perceptron-based-simultaneou.md (part 1 of 2).md
[9] score=None path=robotics__2410.10669__mlp-slam-multilayer-perceptron-based-simultaneou.md (part 2 of 2).md
[10] score=None path=robotics__2410.10669__mlp-slam-multilayer-perceptron-based-simultaneou.md (part 1 of 2).md
[11] score=None path=robotics__2509.14191__mcgs-slam-a-multi-camera-slam-framework-usin.md (part 2 of 2).md
[12] score=None path=robotics__2410.10669__mlp-slam-multilayer-perceptron-based-simultaneou.md (part 1 of 2).md
\end{EvidenceText}
\subsection{BQ10: What strategies for tumor control emerge under immune evasion and chemotherapy resistance?}
\begin{tabular}{@{}ll@{}}
Arm & LightRAG mixed retrieval\\
Generator & omp / deepseek-flash\\
HTTP / success & 200 / True\\
Wall time / cost & 101.1 s / 0.028377 USD\\
Judge scores & F=5, C=3, G=5, A=5; overall=4.50\\
Verdict & partial\\
\end{tabular}

\textbf{Index/retrieval metadata}
\begin{EvidenceText}
{
  "retrieval": "lightrag aquery_data mode=mix (chunks + entity descriptions)",
  "n_entities": 63,
  "n_relationships": 75
}
\end{EvidenceText}

\textbf{Dimension rationale}
\begin{EvidenceText}
{
  "factual": "Everything stated is correct: combination therapy quotes (gold-verbatim); 'consistently expands the parameter region corresponding to tumor control and delays resistance emergence'; the checkpoint-evasion table rows (T1-only down mu -> checkpoint inhibitors; T2-only up sigma, down delta -> adoptive T-cell transfer and cytotoxic therapy, with the paper's reference numbers); the PD-L1/PD-1 mechanism with T1 cells gaining PD-L1 by mutation; exhaustion-rate sensitivity; the four equilibria with the tumor-free state's sigma/(delta+omega*chi) form; the T1-mutant resurgence caveat quoted; future-work items.",
  "completeness": "Covers combination therapy and only the checkpoint-evasion branch of phenotype matching \u2014 the reduced-antigen-presentation mechanism (with its up rho/CAR-T/vaccines targets) is entirely absent, as are the paper's sequential composition-guided schedules; the two-row table and honest truncation notes reflect its narrower excerpt set.",
  "groundedness": "Every claim tagged ([S1][S3][S4][S5][S6][S7][S9][S12]) with quotes; the declared context limits (truncated S1/S3/S7, no dosing numbers) are honest.",
  "sources": "Clear per-claim attributions, including reference-numbered table entries."
}
\end{EvidenceText}

\textbf{Judge summary}
\begin{EvidenceText}
Accurate and well-sourced on combination and checkpoint-evasion strategies, but omits the reduced-antigen-presentation branch and the sequential scheduling strategies.
\end{EvidenceText}

\textbf{Complete answer}
\begin{AnswerText}
**Combination therapy is the central strategy.** Immune dynamics alone "are generally insufficient to fully eliminate the tumor," so the paper analyzes chemotherapy plus immunotherapy together: chemotherapy "significantly enlarges the parameter region corresponding to tumor control or tumor-free equilibria," reduces tumor burden, and directly targets resistant subpopulations, while immunotherapy enhances immune-mediated tumor suppression [S1]. Across the bifurcation analysis, combination therapy "consistently expands the parameter region corresponding to tumor control and delays resistance emergence," and optimal choices depend on both tumor phenotype and immune-evasion mechanism [S4].

**Phenotype-specific targeting** (from the checkpoint-evasion table) [S3]:

| Tumor phenotype | Best target | Therapeutic focus |
|---|---|---|
| T1-only (immune-evasive dominant) | \u2193\u03bc | Checkpoint inhibitors [34] |
| T2-only (immune-sensitive dominant) | \u2191\u03c3, \u2193\u03b4 | Adoptive T-cell transfer [13], cytotoxic therapy [31] |

The mechanistic basis for the first row: PD-L1 binding to PD-1 inactivates effector cells, and T1 cells can gain the PD-L1 receptor by mutation, so blocking checkpoint regulation (lowering \u03bc) restores immune killing [S9], [S12]. Control under checkpoint evasion is noted as "highly sensitive to immune exhaustion rate" [S12].

**Stability framing.** The analysis derives four equilibria \u2014 tumor-free $\mathcal{E}_0=(0,0,\sigma/(\delta+\omega\chi))$, T1-only $\mathcal{E}_1$, T2-only, and a coexistence state \u2014 with stability conditions depending on immune-evasion mechanism and on whether T2 is chemo-resistant [S6], [S7]. Control strategies therefore aim to land the system in $\mathcal{E}_0$ or a controlled state rather than relying on immune suppression of growth.

**Key caveat / limitation.** "If the tumor is only controlled, the emergence of $T_1$ mutants could result in resurgence of the tumor" [S5] \u2014 control is not cure. The paper flags as future work optimal-control scheduling to minimize burden while reducing toxicity, other evasion strategies, microenvironment feedback (breast/prostate), and oncolytic viral therapy [S5].

Not in the provided context: concrete dosing schedules, parameter values, or quantitative survival/control probabilities \u2014 the excerpts are truncated (S1, S3, S7 cut mid-sentence), so specific numbers cannot be cited.
\end{AnswerText}

\textbf{Retrieved evidence index (all returned units; paths and scores) }
\begin{EvidenceText}
[1] score=None path=immunology__2604.01385__strategies-for-tumor-elimination-and-control.md (part 3 of 3).md
[2] score=None path=immunology__2604.01385__strategies-for-tumor-elimination-and-control.md (part 1 of 3).md
[3] score=None path=immunology__2604.01385__strategies-for-tumor-elimination-and-control.md (part 2 of 3).md
[4] score=None path=immunology__2604.01385__strategies-for-tumor-elimination-and-control.md (part 3 of 3).md
[5] score=None path=immunology__2604.01385__strategies-for-tumor-elimination-and-control.md (part 3 of 3).md
[6] score=None path=immunology__2604.01385__strategies-for-tumor-elimination-and-control.md (part 2 of 3).md
[7] score=None path=immunology__2604.01385__strategies-for-tumor-elimination-and-control.md (part 2 of 3).md
[8] score=None path=immunology__2604.01385__strategies-for-tumor-elimination-and-control.md (part 2 of 3).md<SEP>immunology__2604.01385__strategies-for-tumor-elimination-and-control.md (part 1 of 3).md<SEP>immunology__2604.01385__strategies-for-tumor-elimination-and-control.md (part 3 of 3).md
[9] score=None path=immunology__2604.01385__strategies-for-tumor-elimination-and-control.md (part 2 of 3).md<SEP>immunology__2604.01385__strategies-for-tumor-elimination-and-control.md (part 1 of 3).md<SEP>immunology__2604.01385__strategies-for-tumor-elimination-and-control.md (part 3 of 3).md
[10] score=None path=seismology__1308.1846__developing-and-testing-the-automated-post-ev.md (part 1 of 4).md<SEP>immunology__2604.01385__strategies-for-tumor-elimination-and-control.md (part 2 of 3).md<SEP>immunology__2604.01385__strategies-for-tumor-elimination-and-control.md (part 3 of 3).md
[11] score=None path=immunology__2604.01385__strategies-for-tumor-elimination-and-control.md (part 1 of 3).md<SEP>immunology__2604.01385__strategies-for-tumor-elimination-and-control.md (part 3 of 3).md
[12] score=None path=immunology__2604.01385__strategies-for-tumor-elimination-and-control.md (part 1 of 3).md<SEP>immunology__2604.01385__strategies-for-tumor-elimination-and-control.md (part 3 of 3).md
\end{EvidenceText}
\section{DeepRead native pipeline}
\subsection{BQ01: How does the Transformer compute attention, and why does it replace recurrence and convolution?}
\begin{tabular}{@{}ll@{}}
Arm & DeepRead native pipeline\\
Generator & deepread-agent / deepseek-flash\\
HTTP / success & 200 / True\\
Wall time / cost & 40.6 s / not recorded\\
Judge scores & F=5, C=5, G=5, A=5; overall=5.00\\
Verdict & correct\\
\end{tabular}

\textbf{Index/retrieval metadata}
\begin{EvidenceText}
{
  "retrieval": "deepread semantic retrieval + agentic read",
  "generator_disclosure": "native DeepRead agent pipeline (disclosed)"
}
\end{EvidenceText}

\textbf{Dimension rationale}
\begin{EvidenceText}
{
  "factual": "Mechanism claims all match the actual paper: Attention(Q,K,V)=softmax(QK^T/sqrt(d_k))V, additive vs dot-product attention, 'much faster and more space-efficient in practice', scaling rationale (large d_k pushes softmax into regions with extremely small gradients), h=8 with d_k=d_v=64, and the three uses with decoder masking by setting illegal inputs to -infinity. Why-claims match gold and section 4: recurrent O(n) sequential ops vs constant, O(n/k) stacked or O(log_k n) dilated conv layers, separable convolution O(k*n*d + n*d^2) equal to self-attention + FFN at k=n, and the gold's 28.4/41.8 BLEU numbers. No errors found.",
  "completeness": "Both parts answered exhaustively: attention function definition, scaled dot-product procedure and matrix form, full multi-head equations and three roles, then three desiderata with per-architecture comparisons, shorter-path argument, BLEU evidence, and even the interpretability side benefit.",
  "groundedness": "Every claim is anchored to a cited section (section 3.2.1/3.2.2/3.2.3, section 4) and the quoted phrases (e.g. 'jointly attend to information from different representation subspaces') match the paper's actual wording; I could not identify any unsupported assertion.",
  "sources": "Attributions are section-level ('doc 17, section 3.2', 'section 4 Why Self-Attention') and accurate throughout."
}
\end{EvidenceText}

\textbf{Judge summary}
\begin{EvidenceText}
Gives the full scaled dot-product and multi-head equations plus the three-desiderata motivation; both halves of the question are answered completely and correctly.
\end{EvidenceText}

\textbf{Complete answer}
\begin{AnswerText}
## How the Transformer computes attention

**The attention function in general** (doc 17, \xa73.2): an attention function maps a **query** and a set of **key\u2013value pairs** to an output, all of which are vectors. The output is a **weighted sum of the values**, where the weight assigned to each value comes from a **compatibility function of the query with the corresponding key**.

**Scaled Dot-Product Attention** (\xa73.2.1): inputs are queries and keys of dimension $d_k$, and values of dimension $d_v$. The procedure is:
1. Compute the dot products of the query with all keys;
2. Divide each by $\sqrt{d_k}$;
3. Apply a softmax to obtain the weights on the values.

In practice all queries are packed into a matrix $Q$ (keys into $K$, values into $V$), giving

$$\mathrm{Attention}(Q,K,V)=\mathrm{softmax}\!\left(\frac{QK^{T}}{\sqrt{d_k}}\right)V$$

The paper notes the two common attention functions are **additive attention** (compatibility computed by a single-hidden-layer feed-forward network) and **dot-product (multiplicative) attention**; dot-product attention is identical to their algorithm apart from the $1/\sqrt{d_k}$ scaling. Theoretically similar in complexity, dot-product attention is **much faster and more space-efficient in practice** because it can use highly optimized matrix-multiplication code. The scaling is needed because for large $d_k$ the dot products grow large in magnitude, pushing softmax into regions with extremely small gradients; dividing by $\sqrt{d_k}$ counteracts this.

**Multi-Head Attention** (\xa73.2.2): rather than one attention function with $d_{model}$-dimensional keys/values/queries, the model linearly projects Q, K, V $h$ times with different learned projections to $d_k$, $d_k$, $d_v$ dimensions, runs attention in parallel on each projected version, concatenates the $d_v$-dimensional outputs, and projects once more:

$$\mathrm{MultiHead}(Q,K,V)=\mathrm{Concat}(\mathrm{head}_1,\dots,\mathrm{head}_h)W^O,\quad \mathrm{head}_i=\mathrm{Attention}(QW_i^Q,KW_i^K,VW_i^V)$$

This lets the model **jointly attend to information from different representation subspaces at different positions** \u2014 with a single head, averaging inhibits this. The paper uses $h=8$ heads with $d_k=d_v=d_{model}/h=64$, so the total computational cost is similar to single-head attention at full dimensionality.

**Three uses of multi-head attention** (\xa73.2.3):
- **Encoder\u2013decoder attention**: queries from the previous decoder layer, keys/values from the encoder output, so every decoder position can attend over all input positions.
- **Encoder self-attention**: Q, K, V all come from the previous encoder layer; each position attends to all positions of the previous layer.
- **Decoder self-attention**: each position attends to all decoder positions up to and including itself; leftward information flow is prevented to preserve the auto-regressive property, implemented by masking (setting to $-\infty$) the illegal softmax inputs.

## Why it replaces recurrence and convolution

The motivation (\xa74, "Why Self-Attention") is framed as three desiderata: (1) total computational complexity per layer, (2) the amount of computation that can be parallelized (minimum number of sequential operations), and (3) the **path length between long-range dependencies**.

- **Against recurrence**: recurrent layers factor computation along symbol positions, generating hidden state $h_t$ from $h_{t-1}$ and the input at $t$. This **inherently sequential nature precludes parallelization within training examples**, which becomes critical at longer sequence lengths since memory constraints limit batching. A recurrent layer requires $O(n)$ sequential operations, whereas **a self-attention layer connects all positions with a constant number of sequentially executed operations**. In complexity terms, self-attention layers are faster than recurrent layers when the sequence length $n$ is smaller than the representation dimensionality $d$ \u2014 most often the case for sentence representations such as word-piece and byte-pair encodings.
- **Against convolution**: a single convolutional layer with kernel width $k<n$ does not connect all pairs of input and output positions; doing so requires a stack of $O(n/k)$ layers for contiguous kernels, or $O(\log_k(n))$ for dilated convolutions, which **increases the length of the longest paths** between any two positions. Convolutional layers are also generally more expensive than recurrent layers by a factor of $k$ (separable convolutions reduce this to $O(k\cdot n\cdot d + n\cdot d^2)$), and even with $k=n$ a separable convolution's complexity equals the combination of a self-attention layer plus a point-wise feed-forward layer \u2014 the approach taken in this model.
- **Shorter dependency paths**: learning long-range dependencies depends on how long the paths forward and backward signals must traverse; the shorter the paths between any input and output positions, the easier it is to learn such dependencies. Self-attention gives the shortest possible path (constant) between any two positions.
- **Overall claim** (Abstract/Introduction): the Transformer is based **solely on attention mechanisms, dispensing with recurrence and convolutions entirely**, allowing significantly more parallelization and better translation quality \u2014 28.4 BLEU on WMT 2014 English\u2013German (over 2 BLEU above the previous best, including ensembles) and 41.8 BLEU single-model state of the art on English\u2013French after 3.5 days on eight GPUs, a small fraction of the training cost of the best prior models.
- **Side benefit**: self-attention may yield more interpretable models \u2014 individual heads clearly learn different tasks, and many exhibit behavior related to syntactic and semantic structure of sentences.
\end{AnswerText}

\textbf{Retrieved evidence index (all returned units; paths and scores) }
\begin{EvidenceText}
[1] score=None path=1706.03762
[2] score=None path=1706.03762
[3] score=None path=1706.03762
[4] score=None path=1706.03762
[5] score=None path=1706.03762
[6] score=None path=1706.03762
[7] score=None path=1706.03762
[8] score=None path=1706.03762
[9] score=None path=1706.03762
[10] score=None path=1706.03762
[11] score=None path=1706.03762
[12] score=None path=1706.03762
\end{EvidenceText}
\subsection{BQ02: What defines NISQ technology and what is its central limitation?}
\begin{tabular}{@{}ll@{}}
Arm & DeepRead native pipeline\\
Generator & deepread-agent / deepseek-flash\\
HTTP / success & 200 / True\\
Wall time / cost & 28.4 s / not recorded\\
Judge scores & F=5, C=5, G=5, A=5; overall=5.00\\
Verdict & correct\\
\end{tabular}

\textbf{Index/retrieval metadata}
\begin{EvidenceText}
{
  "retrieval": "deepread semantic retrieval + agentic read",
  "generator_disclosure": "native DeepRead agent pipeline (disclosed)"
}
\end{EvidenceText}

\textbf{Dimension rationale}
\begin{EvidenceText}
{
  "factual": "Quotes check out against the paper: 'I needed a name to describe this impending new era, so I made up a word: NISQ'; intermediate scale as '50 to a few hundred' qubits with 50 beyond brute-force simulation; noisy = 'imperfect control'; 'noisy gates unprotected by quantum error correction'; the abstract circuit-size sentence; two-qubit gate error 'above 0.1% (often much worse)'; the verbatim ~1000-gate passage and ceiling sentence; QEC 'heavy overhead cost'. All correct.",
  "completeness": "Exhaustive: coinage story, both halves of the name, the error-correction-unprotected clause, the central limitation with its quantitative chain (gate error -> 1000-gate ceiling -> compute ceiling), the QEC escape route with overhead, and the corollary that NISQ 'is not likely to change the world all by itself'.",
  "groundedness": "Every quotation is attributed to a specific section/node (abstract node 2, section 4.1 node 9, section 4.2 node 10, section 7 node 25, section 8 node 26) and matches the paper's actual text; nothing unsupported.",
  "sources": "Precise document-and-node attributions (doc 18, abstract, 4.1, 4.2, 7, 8) quoted directly \u2014 the clearest attribution scheme in this packet."
}
\end{EvidenceText}

\textbf{Judge summary}
\begin{EvidenceText}
Fully answers definition and central limitation with verbatim section-level quotes covering every key fact, including the 1000-gate ceiling and QEC overhead.
\end{EvidenceText}

\textbf{Complete answer}
\begin{AnswerText}
## NISQ: Definition

In *Quantum Computing in the NISQ era and beyond* (doc 18), John Preskill explains that he coined the term for his talk because he "needed a name to describe this impending new era":

> "NISQ. This stands for **Noisy Intermediate-Scale Quantum**." (\xa74.1, node 9)

The two halves of the name are defined as follows:

- **"Intermediate scale"** \u2014 "refers to the size of quantum computers which will be available in the next few years, with a number of qubits ranging from **50 to a few hundred**." Fifty qubits is a significant milestone "because that's beyond what can be simulated by brute force using the most powerful existing digital supercomputers." (The abstract states the range as 50\u2013100 qubits.)
- **"Noisy"** \u2014 "emphasizes that we'll have **imperfect control over those qubits**; the noise will place serious limitations on what quantum devices can achieve in the near term."

Crucially, the NISQ era is defined as the regime of **noisy gates unprotected by quantum error correction**: "When I speak of the NISQ era, I'm imagining quantum computers with noisy gates unprotected by quantum error correction" (\xa74.2, node 10).

## The Central Limitation

The defining constraint is that **noise limits the size of the quantum circuit that can be run reliably**, which in turn caps the computational power of the whole technology:

> "noise in quantum gates will **limit the size of quantum circuits that can be executed reliably**." (Abstract, node 2)

Preskill quantifies this in \xa74.2 (node 10): with the best current hardware (trapped ions, superconducting circuits), the two-qubit gate error rate is **above 0.1%** (often much worse), and it is unknown whether such rates can be maintained in larger devices. Consequently:

> "with these noisy devices we don't expect to be able to execute a circuit that contains many more than about **1000 gates** \u2014 that is, 1000 fundamental two-qubit operations \u2014 because the noise will overwhelm the signal in a circuit much larger than that. **That limitation on circuit size imposes a ceiling on the computational power of NISQ technology.**"

The escape route \u2014 quantum error correction \u2014 is itself the reason the limitation persists: "quantum error correction imposes a heavy overhead cost in number of qubits and number of gates," so scaling up that way "is a more distant goal" (and reaching thousands of protected qubits would require millions of physical qubits, \xa77, node 25).

## Corollary: Modest Expectations

Because of this ceiling, Preskill stresses that "NISQ is **not likely to change the world all by itself**"; it "should be regarded as a step toward more powerful quantum technologies we'll develop in the future," and near-term algorithms should be designed "with noise resilience in mind" to extend the power of NISQ devices (\xa78, node 26).
\end{AnswerText}

\textbf{Retrieved evidence index (all returned units; paths and scores) }
\begin{EvidenceText}
[1] score=None path=1801.00862
[2] score=None path=1801.00862
[3] score=None path=1801.00862
[4] score=None path=1801.00862
[5] score=None path=1801.00862
[6] score=None path=1801.00862
[7] score=None path=1801.00862
[8] score=None path=1801.00862
[9] score=None path=1801.00862
[10] score=None path=1801.00862
[11] score=None path=1801.00862
[12] score=None path=1801.00862
\end{EvidenceText}
\subsection{BQ03: What is MultiMedQA and how did Flan-PaLM perform on the MedQA benchmark?}
\begin{tabular}{@{}ll@{}}
Arm & DeepRead native pipeline\\
Generator & deepread-agent / deepseek-flash\\
HTTP / success & 200 / True\\
Wall time / cost & 32.5 s / not recorded\\
Judge scores & F=5, C=5, G=5, A=2; overall=4.25\\
Verdict & correct\\
\end{tabular}

\textbf{Index/retrieval metadata}
\begin{EvidenceText}
{
  "retrieval": "deepread semantic retrieval + agentic read",
  "generator_disclosure": "native DeepRead agent pipeline (disclosed)"
}
\end{EvidenceText}

\textbf{Dimension rationale}
\begin{EvidenceText}
{
  "factual": "Everything checks out: seven datasets correctly enumerated including HealthSearchQA; the five axes of dataset variation (format, capabilities, domain, source, labels) match the paper's dataset description; MedQA as USMLE-style with 11,450 dev/1,273 test; 67.6% on 4-option MedQA, DRAGON +20.1%, concurrent PubMedGPT 50.3% exceeded by 17.3%; 62.0% on 5-option; Table 5 numbers (58.9 vs 60.3, 25.7 vs 35.4, 40.9 vs 46.1) are the paper's actual scaling results.",
  "completeness": "Both parts answered completely: definition (composition, spans, variation axes, motivation) and performance (headline result, comparison baselines, 5-option variant, scale table, SOTA-on-every-MCQ-dataset claim, and the human-eval gap leading to Med-PaLM).",
  "groundedness": "No unsupported claims \u2014 every statement matches the gold abstract or the paper's actual content \u2014 but support comes from the material itself rather than from citations.",
  "sources": "Almost no attributions: only one incidental internal reference ('Under plain few-shot prompting (Table 5)'); no document or passage IDs are given anywhere, so statements are not traceable to specific sources."
}
\end{EvidenceText}

\textbf{Judge summary}
\begin{EvidenceText}
Complete and factually flawless answer to both halves, but it cites essentially no identifiable sources beyond a single '(Table 5)' mention.
\end{EvidenceText}

\textbf{Complete answer}
\begin{AnswerText}
## What is MultiMedQA?

MultiMedQA is a benchmark for medical question answering curated by the paper's authors to evaluate how well large language models encode clinical knowledge. It comprises **seven medical question answering datasets**:

- **Six existing datasets:** MedQA, MedMCQA, PubMedQA, LiveQA, MedicationQA, and MMLU clinical topics
- **One newly introduced dataset:** **HealthSearchQA**, consisting of commonly searched health questions

The benchmark spans **medical exam questions, medical research questions, and consumer medical questions**, and the datasets vary along several axes: format (multiple-choice vs. long-form answers), capabilities tested (recall of medical facts in isolation vs. medical reasoning plus recall), domain (open vs. closed), question source (professional exams, medical research, or consumers), and labels/metadata (presence of labels or explanations and their sources).

The motivation was that existing medical QA benchmarks were often limited to classification accuracy or automated NLG metrics (e.g., BLEU) and did not enable the detailed analysis needed for real-world clinical applications \u2014 assessing factuality, use of expert knowledge in reasoning, helpfulness, precision, health equity, and potential harm.

## Flan-PaLM's performance on MedQA

MedQA consists of USMLE-style questions with a choice of 4 or 5 possible answers (development set: 11,450 questions; test set: 1,273 questions).

- Using a combination of **few-shot, chain-of-thought (CoT), and self-consistency prompting**, **Flan-PaLM 540B achieved 67.6% MCQ accuracy on the 4-option MedQA set**, surpassing the DRAGON model by 20.1% and exceeding the concurrent PubMedGPT (50.3%) \u2014 the prior state of the art \u2014 **by 17.3%**.
- On the **more difficult 5-option set, the model scored 62.0%**.
- Under plain few-shot prompting (Table 5), Flan-PaLM 540B scored **60.3%** on MedQA (4 options), versus 58.9% for the baseline PaLM 540B; smaller variants scored 35.4% (8B) and 46.1% (62B) versus 25.7% and 40.9% for the corresponding PaLM models \u2014 showing consistent gains from instruction tuning.

Overall, Flan-PaLM set a new state of the art on MedQA (USMLE), MedMCQA, PubMedQA, and MMLU clinical topics, though its answers to consumer medical questions still revealed key gaps \u2014 which motivated instruction prompt tuning to produce Med-PaLM.
\end{AnswerText}

\textbf{Retrieved evidence index (all returned units; paths and scores) }
\begin{EvidenceText}
[1] score=None path=2212.13138
[2] score=None path=2212.13138
[3] score=None path=2212.13138
[4] score=None path=2212.13138
[5] score=None path=2212.13138
[6] score=None path=2212.13138
[7] score=None path=2212.13138
[8] score=None path=2212.13138
[9] score=None path=2212.13138
[10] score=None path=2212.13138
[11] score=None path=2212.13138
[12] score=None path=2212.13138
\end{EvidenceText}
\subsection{BQ04: Which three independent ontologies subdivide the Gene Ontology?}
\begin{tabular}{@{}ll@{}}
Arm & DeepRead native pipeline\\
Generator & deepread-agent / deepseek-flash\\
HTTP / success & 200 / True\\
Wall time / cost & 6.8 s / not recorded\\
Judge scores & F=5, C=5, G=5, A=4; overall=4.75\\
Verdict & correct\\
\end{tabular}

\textbf{Index/retrieval metadata}
\begin{EvidenceText}
{
  "retrieval": "deepread semantic retrieval + agentic read",
  "generator_disclosure": "native DeepRead agent pipeline (disclosed)"
}
\end{EvidenceText}

\textbf{Dimension rationale}
\begin{EvidenceText}
{
  "factual": "The three names are correct and match gold; the added glosses (MF as 'activities of a gene product at the molecular level', BP as 'biological objectives or processes', CC as 'where in the cell the gene product is located/active') are the GO paper's actual definitions; the note that they 'share a common space of identifiers and a well-specified syntax' is gold-verbatim.",
  "completeness": "All three ontologies named with explanatory glosses \u2014 fully answers the question.",
  "groundedness": "Names and the shared-identifier claim are supported by the cited section 2 of the gold document; the glosses match the primer's own definitions elsewhere, so nothing is wrong or invented.",
  "sources": "One attribution at the end ('doc 16, \"2. What is the Gene Ontology?\"') covers the subdivision claim, but the three per-ontology glosses carry no individual attribution."
}
\end{EvidenceText}

\textbf{Judge summary}
\begin{EvidenceText}
Names all three ontologies correctly with accurate definitions, attributed to the right section, though only via a single trailing citation.
\end{EvidenceText}

\textbf{Complete answer}
\begin{AnswerText}
The Gene Ontology is subdivided into three distinct, non-redundant ontologies representing different biological aspects:

1. **Molecular Function (MF)** \u2014 the activities of a gene product at the molecular level
2. **Biological Process (BP)** \u2014 the biological objectives or processes the gene product contributes to
3. **Cellular Component (CC)** \u2014 where in the cell the gene product is located/active

These three share a common space of identifiers and a well-specified syntax (doc 16, "2. What is the Gene Ontology?").
\end{AnswerText}

\textbf{Retrieved evidence index (all returned units; paths and scores) }
\begin{EvidenceText}
[No source list recorded for this arm.]
\end{EvidenceText}
\subsection{BQ05: Why does historical European economic growth contradict Unified Growth Theory?}
\begin{tabular}{@{}ll@{}}
Arm & DeepRead native pipeline\\
Generator & deepread-agent / deepseek-flash\\
HTTP / success & 200 / True\\
Wall time / cost & 34.8 s / not recorded\\
Judge scores & F=5, C=5, G=4, A=5; overall=4.75\\
Verdict & correct\\
\end{tabular}

\textbf{Index/retrieval metadata}
\begin{EvidenceText}
{
  "retrieval": "deepread semantic retrieval + agentic read",
  "generator_disclosure": "native DeepRead agent pipeline (disclosed)"
}
\end{EvidenceText}

\textbf{Dimension rationale}
\begin{EvidenceText}
{
  "factual": "Correct on all gold-verifiable content: the 100,000 BC-1750 Malthusian regime, 1750-1870 post-Malthusian, 1870 sustained-growth dates; no stagnation/trap/transition; Industrial Revolution with 'absolutely no impact'; Maddison data shown as four points joined by straight lines; the paradox of being contradicted by the same data used to develop the theory. The fit parameters for 12 countries (1923 singularity, 23-year bypass, AD 1 +27%, AD 1000 -54%), 30 countries (1929, 29y) and Eastern Europe (1915, 25y, AD 1 +51%) are consistent across independent answers citing this paper.",
  "completeness": "Most comprehensive answer in the packet: test-case logic, a claim-vs-data table, quantitative fits, the reciprocal-coordinate method, methodological criticism, and the conclusion \u2014 all parts of the question plus context.",
  "groundedness": "Nearly everything is attributed and cross-consistent, but one paragraph is unique to this answer and unverifiable from the gold: the claim that UGT contains '\u5371\u9669\u7684\u8bef\u5bfc\u6027\u4fe1\u606f' (dangerous misleading information) about unconstrained future growth pointing to an 'unsafe future' appears in no other system's answer and has no counterpart in the shown gold excerpts.",
  "sources": "Strong sourcing: paper title, arXiv id 1512.08067, KB name and id, P0 evidence tier, plus honest retrieval notes (other hits like 'leisure class' rejected as irrelevant; absence of any Galor rebuttal in the KB disclosed)."
}
\end{EvidenceText}

\textbf{Judge summary}
\begin{EvidenceText}
Thorough and accurate on both the conceptual and quantitative contradictions, with only one unique, unverifiable rhetorical claim beyond the evidence shown.
\end{EvidenceText}

\textbf{Complete answer}
\begin{AnswerText}
## \u6838\u5fc3\u7ed3\u8bba

\u6839\u636e Nielsen\uff082015\uff09\u5bf9 Maddison \u6570\u636e\u7684\u5206\u6790\uff0c**\u6b27\u6d32\uff08\u5c24\u5176\u662f\u897f\u6b27\uff09\u672c\u5e94\u662f\u7edf\u4e00\u589e\u957f\u7406\u8bba\uff08Unified Growth Theory, Galor\uff09\u6700\u6709\u529b\u7684\u9a8c\u8bc1\u573a\uff0c\u4f46\u6570\u636e\u5374\u4e0e\u8be5\u7406\u8bba"\u60ca\u4eba\u5730\u4e0d\u4e00\u81f4"**\u2014\u2014\u4e24\u8005\u4e4b\u95f4"\u6ca1\u6709\u4efb\u4f55\u8054\u7cfb"\u3002

## \u4e00\u3001\u4e3a\u4ec0\u4e48\u6b27\u6d32\u662f\u6700\u5173\u952e\u7684\u68c0\u9a8c\u573a

- \u897f\u6b27\u662f\u5de5\u4e1a\u9769\u547d\u7684\u4e2d\u5fc3\uff0c\u800c\u5de5\u4e1a\u9769\u547d\u88ab Galor \u79f0\u4e3a"\u7ecf\u6d4e\u589e\u957f\u7684\u9996\u8981\u5f15\u64ce"\uff08Galor, 2005a, p. 212\uff09\u3002\u56e0\u6b64\u8be5\u7406\u8bba\u6240\u8ff0\u7684\u5404\u9879\u673a\u5236\u82e5\u6709\u6548\uff0c\u7406\u5e94\u5728\u6b27\u6d32\u6570\u636e\u4e2d\u8868\u73b0\u5f97\u6700\u4e3a\u7a81\u51fa\u3002
- \u4f5c\u8005\u5206\u522b\u5206\u6790\u4e86\u897f\u6b27\u4e24\u7ec4\u6570\u636e\uff1a12 \u4e2a\u6700\u5148\u8fdb\u56fd\u5bb6\uff08\u5360 30 \u56fd\u603b GDP \u7684 85%\uff09\u4e0e 30 \u56fd\u603b\u91cf\uff0c\u4ee5\u53ca\u4e1c\u6b27\u6570\u636e\u300212 \u56fd\u7ec4"\u672c\u5e94\u7ed9\u51fa\u7406\u8bba\u4e0e\u6570\u636e\u7684\u6700\u4f73\u543b\u5408"\u3002

## \u4e8c\u3001\u6570\u636e\u4e0e\u7406\u8bba\u7684\u5177\u4f53\u77db\u76fe\uff08\u9010\u6761\u5bf9\u7167\uff09

| \u7edf\u4e00\u589e\u957f\u7406\u8bba\u7684\u4e3b\u5f20 | \u6570\u636e\u5b9e\u9645\u663e\u793a |
|---|---|
| \u6f2b\u957f\u7684\u9a6c\u5c14\u8428\u65af\u505c\u6ede\u671f\uff08\u81ea\u516c\u5143\u524d 10 \u4e07\u5e74\u76f4\u81f3 1750 \u5e74\uff09 | \u7ecf\u6d4e\u6cbf\u53cc\u66f2\u7ebf\u8f68\u8ff9\u6301\u7eed\u589e\u957f\uff0c**\u4ece\u672a\u6709\u8fc7\u505c\u6ede**\uff08\u516c\u5143 1500\u20131900 \u5e74\u62df\u5408\u6781\u4f73\uff09 |
| \u5b58\u5728"\u9a6c\u5c14\u8428\u65af\u9677\u9631" | **\u6ca1\u6709\u9677\u9631** |
| "\u975e\u51e1\u7684""\u60ca\u4eba\u7684"\u9003\u79bb\u9a6c\u5c14\u8428\u65af\u9677\u9631\uff08\u8d77\u98de\uff09 | \u589e\u957f\u672a\u53d7\u6270\u52a8\u5730\u6cbf\u53cc\u66f2\u7ebf\u4e0a\u5347\uff0c**\u6ca1\u6709\u9003\u79bb\uff0c\u56e0\u4e3a\u6ca1\u6709\u9677\u9631** |
| \u4ece\u505c\u6ede\u5230\u589e\u957f\u7684\u620f\u5267\u6027\u8f6c\u53d8 | **\u6ca1\u6709\u8f6c\u53d8**\uff0c\u56e0\u4e3a\u672c\u65e0\u505c\u6ede |
| \u5de5\u4e1a\u9769\u547d\u662f\u589e\u957f\u7684\u9996\u8981\u5f15\u64ce | \u5de5\u4e1a\u9769\u547d**\u5bf9\u589e\u957f\u8f68\u8ff9\u6beb\u65e0\u5f71\u54cd** |
| 1870 \u5e74\u540e\u8fdb\u5165"\u6301\u7eed\u589e\u957f"\u65b0\u4f53\u5236 | \u589e\u957f**\u4e00\u76f4\u662f\u6301\u7eed\u7684**\uff1b\u53cd\u800c\u5728\u6240\u8c13"\u6301\u7eed\u589e\u957f\u4f53\u5236"\u5f00\u59cb\u524d\u540e\uff0c\u589e\u957f\u88ab\u8f6c\u5411**\u66f4\u6162**\u7684\u8f68\u8ff9 |

## \u4e09\u3001\u5173\u952e\u5b9a\u91cf\u8bc1\u636e

- **\u897f\u6b27 12 \u56fd**\uff1aa = 1.147\xd710\u207b\xb9\uff0ck = 5.961\xd710\u207b\u2075\uff1b\u5947\u70b9\u5e94\u5728 1923 \u5e74\uff0c\u4f46\u589e\u957f\u7ea6\u5728 1900 \u5e74\u8f6c\u5411\u8f83\u6162\u8f68\u8ff9\uff0c\u7ed5\u8fc7\u5947\u70b9\u7ea6 23 \u5e74\u3002\u516c\u5143 1 \u5e74\u6570\u636e\u70b9\u6bd4\u62df\u5408\u503c\u9ad8 27%\uff0c\u516c\u5143 1000 \u5e74\u4f4e 54%\u3002
- **\u897f\u6b27 30 \u56fd**\uff1aa = 9.859\xd710\u207b\xb2\uff0ck = 5.112\xd710\u207b\u2075\uff1b\u5947\u70b9 1929 \u5e74\uff0c\u7ea6 1900 \u5e74\u8f6c\u5411\uff0c\u7ed5\u8fc7\u7ea6 29 \u5e74\u3002
- **\u4e1c\u6b27**\uff1aa = 7.749\xd710\u207b\xb9\uff0ck = 4.048\xd710\u207b\u2074\uff1b\u5947\u70b9 1915 \u5e74\uff0c\u7ea6 1890 \u5e74\u8f6c\u5411\uff0c\u7ed5\u8fc7\u7ea6 25 \u5e74\uff1b\u516c\u5143 1 \u5e74\u6570\u636e\u70b9\u6bd4\u66f2\u7ebf\u9ad8 51%\u3002

\u65b9\u6cd5\u4e0a\uff0c\u4f5c\u8005\u7528\u4e24\u79cd\u5c55\u793a\u65b9\u5f0f\uff1aGDP \u7684\u534a\u5bf9\u6570\u56fe\uff0c\u4ee5\u53ca\u5012\u6570 1/GDP \u56fe\u3002\u53cc\u66f2\u7ebf\u589e\u957f S(t) = (a \u2212 kt)\u207b\xb9 \u7684\u5012\u6570 1/S(t) = a \u2212 kt \u662f\u4e00\u6761\u76f4\u7ebf\uff1b\u6570\u636e\u5012\u6570\u5448\u4e0b\u964d\u76f4\u7ebf\uff0c\u5373\u8bc1\u660e\u589e\u957f\u662f\u53cc\u66f2\u7ebf\u5f0f\u7684\u800c\u975e\u505c\u6ede\u7684\u3002\u4f5c\u8005\u5f3a\u8c03\uff1a\u628a\u4e00\u6761\u76f4\u7ebf\u4efb\u610f\u5207\u6210\u4e24\u4e09\u6bb5\u5e76\u5ba3\u79f0\u5b58\u5728\u4e0d\u540c\u589e\u957f\u4f53\u5236\uff0c\u662f\u6beb\u65e0\u610f\u4e49\u7684\u3002

## \u56db\u3001\u65b9\u6cd5\u8bba\u6279\u8bc4

- Maddison \u7684\u6570\u636e"\u5c61\u906d\u80a2\u89e3"\u2014\u2014\u5e38\u53ea\u53d6\u56db\u4e2a\u70b9\u7528\u76f4\u7ebf\u8fde\u63a5\uff08Ashraf, 2009\uff1bGalor \u591a\u7bc7\uff1bGalor & Moav, 2002\uff1bSnowdon & Galor, 2008\uff09\uff0c\u8fd9\u79cd\u5904\u7406"\u4e0d\u79d1\u5b66"\u3002
- \u53cc\u66f2\u7ebf\u5206\u5e03\u672c\u8eab\u6613\u8bef\u5bfc\uff0c\u5fc5\u987b\u7528\u5012\u6570\u503c\u5206\u6790\uff1b\u4ece\u8fd9\u7c7b\u5206\u5e03\u4e2d\u6458\u53d6\u5b64\u7acb\u6570\u5b57\u6216\u4f5c\u6697\u793a\u6027\u4f5c\u56fe\uff0c\u96be\u4ee5\u5f97\u51fa\u53ef\u9760\u7ed3\u8bba\u3002\u7edf\u4e00\u589e\u957f\u7406\u8bba"\u5f3a\u70c8\u5730\u3001\u751a\u81f3\u5b8c\u5168\u5730"\u5efa\u7acb\u5728\u8fd9\u79cd\u4e0d\u79d1\u5b66\u7684\u7528\u6570\u636e\u65b9\u5f0f\u4e0a\u3002
- \u6096\u8bba\u5728\u4e8e\uff1a\u8be5\u7406\u8bba\u88ab**\u5b83\u81ea\u5df1\u53d1\u5c55\u65f6\u6240\u4f7f\u7528\u7684\u540c\u4e00\u6279\u6570\u636e**\u6240\u5426\u5b9a\uff0c\u539f\u56e0\u6b63\u662f\u8fd9\u4e9b\u6570\u636e\u4ece\u672a\u88ab\u6070\u5f53\u5206\u6790\u3002

## \u4e94\u3001\u4f5c\u8005\u7684\u7ed3\u8bba\u4e0e\u8b66\u793a

- \u8be5\u7406\u8bba"\u4e0e\u73b0\u5b9e\u8131\u8282"\uff0c\u5176\u589e\u957f\u673a\u5236\u5728\u6570\u636e\u4e2d\u6beb\u65e0\u53cd\u6620\uff1b\u5b83\u8bb2\u8ff0\u7684\u662f"\u795e\u8bdd\u5f0f\u7279\u5f81"\u3002
- \u5b83\u8fd8\u5305\u542b"\u5371\u9669\u7684\u8bef\u5bfc\u6027\u4fe1\u606f"\uff1a\u5ba3\u79f0\u4eba\u7c7b\u5df2\u8d70\u51fa\u505c\u6ede\u3001\u8fdb\u5165\u65e0\u7ea6\u675f\u7684\u6301\u7eed\u589e\u957f\u3001\u8bb8\u8bfa\u65e5\u76ca\u7e41\u8363\u3002\u6570\u636e\u5374\u663e\u793a\u76f8\u53cd\u56fe\u666f\u2014\u2014\u8fc7\u53bb\u7684\u589e\u957f\u672c\u5c31\u662f\u6301\u7eed\u7684\uff0c\u800c\u5f53\u524d\u589e\u957f\u867d\u5df2\u8f6c\u5411\u8f83\u6162\u8f68\u8ff9\uff0c\u4ecd\u589e\u957f\u8fc7\u5feb\uff0c\u6307\u5411**\u4e0d\u5b89\u5168\u7684\u672a\u6765**\u3002
- \u56e0\u6b64\u8be5\u7406\u8bba"\u9700\u8981\u5f7b\u5e95\u4fee\u8ba2\uff0c\u4f46\u6700\u53ef\u80fd\u5e94\u88ab\u53d6\u4ee3"\uff0c\u4ee3\u4e4b\u4ee5\u57fa\u4e8e\u7ecf\u9a8c\u8bc1\u636e\uff08\u5c24\u5176 Maddison \u6570\u636e\uff09\u79d1\u5b66\u5206\u6790\u7684\u65b0\u7406\u8bba\u3002
\end{AnswerText}

\textbf{Retrieved evidence index (all returned units; paths and scores) }
\begin{EvidenceText}
[1] score=None path=1512.08067
[2] score=None path=1512.08067
[3] score=None path=1512.08067
[4] score=None path=1512.08067
[5] score=None path=1512.08067
[6] score=None path=1512.08067
[7] score=None path=1512.08067
[8] score=None path=1512.08067
[9] score=None path=1512.08067
[10] score=None path=1512.08067
[11] score=None path=1512.08067
[12] score=None path=1512.08067
\end{EvidenceText}
\subsection{BQ06: Which physical factor primarily controls the response of precipitation extremes to climate change?}
\begin{tabular}{@{}ll@{}}
Arm & DeepRead native pipeline\\
Generator & deepread-agent / deepseek-flash\\
HTTP / success & 200 / True\\
Wall time / cost & 71.4 s / not recorded\\
Judge scores & F=5, C=5, G=5, A=5; overall=5.00\\
Verdict & correct\\
\end{tabular}

\textbf{Index/retrieval metadata}
\begin{EvidenceText}
{
  "retrieval": "deepread semantic retrieval + agentic read",
  "generator_disclosure": "native DeepRead agent pipeline (disclosed)"
}
\end{EvidenceText}

\textbf{Dimension rationale}
\begin{EvidenceText}
{
  "factual": "The decomposition P = -eps*omega*S with dynamical/thermodynamic/microphysical contributions matches the paper; the abstract quote ('thermodynamic contribution is robust and well understood') is gold-verbatim; 'the thermodynamic contribution is the easiest to understand and always gives an intensification with warming', the extratropical closeness to the thermodynamic contribution, surface-specific-humidity scaling (lower than column water vapor), reduction to Clausius-Clapeyron ~6-7%/K, and the below-295K deviations all match the paper's content.",
  "completeness": "Answers directly ('the thermodynamic contribution') and then supplies the governing equation, the robustness ranking, the extratropical evidence, the scaling magnitude, and the paper's caveats (convection, tropical vs extratropical sensitivity) \u2014 fully answered.",
  "groundedness": "Each claim is tied to a node of the document with quotes; the ~9-11%/K tropical vs ~4-6%/K extratropical sensitivities are specific but attributed and plausible for this paper; nothing contradicts gold.",
  "sources": "Precise node-level attributions (doc 14, nodes 2, 4, 5, 7, 8, 12, 13)."
}
\end{EvidenceText}

\textbf{Judge summary}
\begin{EvidenceText}
Decisive, correct answer (thermodynamic/Clausius-Clapeyron control) backed by equation-level decomposition, robustness quotes, and quantitative caveats.
\end{EvidenceText}

\textbf{Complete answer}
\begin{AnswerText}
**The thermodynamic contribution \u2014 i.e., the increase in the saturation vapor pressure (water-vapor content) of the atmosphere with warming, which gives roughly Clausius\u2013Clapeyron scaling of about 6\u20137 % K\u207b\xb9.**

How the paper frames it:

- The starting point is the approximation of extreme precipitation rate, *P \u2248 \u2212\u03b5{\u03c9 S(T,p)}*, where \u03b5 is precipitation efficiency, \u03c9 the vertical (pressure) velocity, and S the derivative of saturation specific humidity with respect to pressure. Changes in *P* therefore decompose into a **dynamical** contribution (changes in \u03c9), a **thermodynamic** contribution (changes in S, which depends only on temperature and pressure), and a **microphysical** contribution (changes in \u03b5) (doc 14, node 5).
- Of these, "the thermodynamic contribution is robust and well understood, but theoretical understanding of the microphysical and dynamical contributions is still being developed" (Abstract, node 2). The Conclusions likewise state that "the thermodynamic contribution is the easiest to understand and always gives an intensification with warming" (node 13).
- In the extratropics, "the simulated rate of increase of precipitation extremes is close to the thermodynamic contribution at all latitudes, and there is little dynamical contribution from changes in vertical velocities" (node 7). Equation (1) with \u03b5 held fixed reproduces the fractional changes in precipitation extremes in CMIP3 simulations.
- The thermodynamic contribution is close to scaling with surface specific humidity (a lower rate than scaling with column water vapor), and for moist-adiabatic stratification with near-surface convergence it reduces to Clausius\u2013Clapeyron scaling, ~6\u20137 % K\u207b\xb9 for typical surface temperatures (node 5).

Caveats noted in the paper: the sensitivity is less certain where convection matters (tropics, summer midlatitude convective systems), where dynamical and microphysical (precipitation-efficiency) contributions can be substantial \u2014 e.g., deviations from Clausius\u2013Clapeyron scaling at surface temperatures below 295 K, and higher tropical sensitivities (~9\u201311 % K\u207b\xb9) than extratropical ones (~4\u20136 % K\u207b\xb9) (nodes 4, 7, 8, 12).
\end{AnswerText}

\textbf{Retrieved evidence index (all returned units; paths and scores) }
\begin{EvidenceText}
[1] score=None path=1503.07557
[2] score=None path=1503.07557
[3] score=None path=1503.07557
[4] score=None path=1503.07557
[5] score=None path=1503.07557
[6] score=None path=1503.07557
[7] score=None path=1503.07557
[8] score=None path=1503.07557
[9] score=None path=1503.07557
[10] score=None path=1503.07557
[11] score=None path=1503.07557
[12] score=None path=1503.07557
\end{EvidenceText}
\subsection{BQ07: What is MameLoshnLM and what parameter scale does it have?}
\begin{tabular}{@{}ll@{}}
Arm & DeepRead native pipeline\\
Generator & deepread-agent / deepseek-flash\\
HTTP / success & 200 / True\\
Wall time / cost & 15.6 s / not recorded\\
Judge scores & F=5, C=5, G=5, A=2; overall=4.25\\
Verdict & correct\\
\end{tabular}

\textbf{Index/retrieval metadata}
\begin{EvidenceText}
{
  "retrieval": "deepread semantic retrieval + agentic read",
  "generator_disclosure": "native DeepRead agent pipeline (disclosed)"
}
\end{EvidenceText}

\textbf{Dimension rationale}
\begin{EvidenceText}
{
  "factual": "Everything matches: first open-source LLM built specifically for Yiddish; 8B parameters; continued pretraining of Llama-3.1-8B on the Oytser corpus; the mame-loshn ('mother tongue') etymology is genuine; training details (~5.7B tokens, 36,663 steps, single H200, ~207 GPU-hours, Yiddish 72% words/92% tokens) and the Kashes result (62.6 vs 56.8 for Llama 3.1 8B, with Gemma-2 9B, EuroLLM 9B, BLOOMZ 7B, Qwen3 8B baselines) are consistent with the paper.",
  "completeness": "Both parts answered completely, plus name origin, training configuration, and benchmark results.",
  "groundedness": "All claims are supported by the gold abstract and the paper's content \u2014 no fabrication \u2014 but the support comes from the material itself rather than any citation.",
  "sources": "No source attributions at all: no document, passage, or section ids anywhere in the answer."
}
\end{EvidenceText}

\textbf{Judge summary}
\begin{EvidenceText}
Complete and accurate answer to both parts, but entirely uncited \u2014 nothing indicates which sources back which claims.
\end{EvidenceText}

\textbf{Complete answer}
\begin{AnswerText}
**MameLoshnLM (MAMELOSHNLM)** is described in the paper as **the first open-source large language model built specifically for Yiddish** (an 8B-parameter model).

Key points:

- **Name/origin:** "Mame-Loshn" (\u05de\u05d0\u05b7\u05de\u05e2\u05be\u05dc\u05e9\u05d5\u05df, "mother tongue") is the traditional Yiddish term evoking the language's intimate connection to home and family.
- **Parameter scale:** **8 billion parameters** \u2014 it is created by **continued pretraining of Llama-3.1-8B** on a Yiddish corpus (the "Oytser" corpus) using a causal language modeling objective. Training used ~5.7 billion tokens over 36,663 optimization steps on a single NVIDIA H200 GPU (~207 GPU-hours), with a Yiddish-dominant mixture (Yiddish 72% of words / 92% of tokens, English 28% / 8%).
- **Evaluation:** It is assessed on **Kashes**, a new multi-task Yiddish benchmark (translation, linguistic analysis, information extraction, language understanding), where it outperforms open baselines of similar scale (e.g., Llama 3.1 8B, Qwen3 8B, Gemma-2 9B, EuroLLM 9B, BLOOMZ 7B), with an average score of 62.6 vs. 56.8 for Llama 3.1 8B.
\end{AnswerText}

\textbf{Retrieved evidence index (all returned units; paths and scores) }
\begin{EvidenceText}
[1] score=None path=2608.05850
[2] score=None path=2608.05850
[3] score=None path=2608.05850
\end{EvidenceText}
\subsection{BQ08: How does oxygen redox contribute to the capacity of battery cathodes?}
\begin{tabular}{@{}ll@{}}
Arm & DeepRead native pipeline\\
Generator & deepread-agent / deepseek-flash\\
HTTP / success & 200 / True\\
Wall time / cost & 50.6 s / not recorded\\
Judge scores & F=5, C=5, G=5, A=5; overall=5.00\\
Verdict & correct\\
\end{tabular}

\textbf{Index/retrieval metadata}
\begin{EvidenceText}
{
  "retrieval": "deepread semantic retrieval + agentic read",
  "generator_disclosure": "native DeepRead agent pipeline (disclosed)"
}
\end{EvidenceText}

\textbf{Dimension rationale}
\begin{EvidenceText}
{
  "factual": "Everything checks: the gold's Eq. (1) charge-compensation reaction, 200-300 mAh/g theoretical capacities rarely reached, 'achieved and even exceeded' the TM oxidation limit, Li2MnO3 with 1.4 Li extracted despite too-low t2g; the peroxo O2-(n-) / electron-hole representations, orphaned O 2p states from Li-O-Li environments, ilmenite-type NaTMO3 under-coordination, the Na2-xRuO3 ~180 mAh/g with voltage spike 2.2 V to 3.4 V, the 4d-oxide DFT survey numbers, and the fluorination mitigation are all consistent with the paper and cross-corroborated by other answers.",
  "completeness": "Six sections cover the capacity problem, the anionic mechanism, the orbital origin, evidence and magnitude, capacity fade costs, and how the effect is probed (DFT, hybrid functionals/GGA+U) \u2014 the most complete treatment in the packet.",
  "groundedness": "No unsupported claims; every assertion is consistent with the gold excerpt or the paper's content.",
  "sources": "Consistent section-level attributions to the named review ('doc 20') \u2014 Abstract, \xa71, \xa72, \xa73.1, \xa73.2, \xa74, \xa75, Fig. 2, Tables 1-2 \u2014 one reference per claim."
}
\end{EvidenceText}

\textbf{Judge summary}
\begin{EvidenceText}
Exhaustive, section-referenced explanation from the capacity gap through orbital mechanism, quantified evidence, degradation costs, and DFT methodology.
\end{EvidenceText}

\textbf{Complete answer}
\begin{AnswerText}
## How oxygen redox contributes to cathode capacity

Based on the review *"Oxygen redox in battery cathodes: A brief overview"* (doc 20), the contribution works as follows.

### 1. The basic capacity problem it solves
In a conventional cathode, extracting one Li\u207a (or Na\u207a) is charge-compensated by oxidizing a transition metal (TM):

LiTM\xb3\u207aO\u2082 \u2192 Li\u208d\u2081\u208b\u2093\u208eTM\xb3\u207a\u208d\u2081\u208b\u2093\u208eTM\u2074\u207a\u2093O\u2082 + xLi\u207a + xe\u207b

Theoretical capacities of ~200\u2013300 mAh/g are rarely reached because of phase changes, interphase diffusion barriers, and O\u2082 release. Yet real cathodes have **achieved and even exceeded the transition-metal oxidation limit** \u2014 an "overcapacity" that requires an additional charge-storage mechanism (Abstract; \xa71).

### 2. The anionic mechanism itself
Oxygen redox means that **oxide anions O\xb2\u207b donate electrons during charging** \u2014 i.e., they act as reducing agents \u2014 thereby compensating for the absence of oxidizable transition metals. This "may allow greater energy storage capacity in battery electrodes than would otherwise be theoretically predicted" (\xa72). Reversibly, it is summarized as formation of peroxo-type species:

O\xb2\u207b \u2194 (O\u2082)\u207f\u207b

or alternatively as the formation of **electron holes on oxygen atoms**, which explains anion-driven excess capacity in alkali-rich cathodes *without* detrimental O\u2082 gas evolution (\xa72).

### 3. Why it happens: orbital/structural origin
- In normal LiTMO\u2082, TM s and p orbitals hybridize with the six O 2p orbitals, pushing the occupied, bonding p-like states ~10 eV below the Fermi level. The non-hybridized TM t\u2082g states sit closest to E_F and are therefore the electrochemically labile ones \u2192 TM-borne redox (\xa72).
- For O redox, **p-character states must be lifted toward the Fermi level**. This happens when oxygen's coordination changes: in **Li-excess / Na-excess** materials (e.g., Li\u2082MnO\u2083-type, where the TM layer also contains alkali), some O ions are coordinated by fewer than three TM ions (Li\u2013O\u2013Li / Na\u2013O\u2013Na). These **"orphaned," non-bonding O 2p states** lie near the TM t\u2082g states, so both O and TM can donate/accept electrons. The O-vs-TM split depends on M\u2013O covalency, TM d population, and depth of delithiation/desodiation (\xa72, Fig. 2).
- In **ilmenite-type NaTMO\u2083**, O redox arises not from Na over-coordination but from O being **under-coordinated by both Na and TM** (2 TM + 2 Na), creating non-bonding O 2p states just below E_F (\xa74).

### 4. Evidence and magnitude
- **Li\u2082MnO\u2083**: Mn remains +4 (its t\u2082g level is too low to oxidize), so the observed capacity is attributed to oxygen redox (\xa71, \xa72).
- **Li-excess disordered spinel cathodes**: O redox was first found to contribute significantly to overall capacity, thanks to orphaned O states (\xa75).
- **Na\u2082\u208b\u2093RuO\u2083**: ~180 mAh/g, i.e. ~1.4 e\u207b per formula unit; the extra ~0.4 e\u207b was assigned to reversible O redox after desodiation to ilmenite Na\u2081RuO\u2083, accompanied by a voltage spike from 2.2 V to 3.4 V \u2014 indicating O redox can give a *higher* operating potential than TM redox alone (\xa74).
- DFT surveys of 3d and 4d NaTMO\u2083 show O 2p states dominating the region below E_F, with **4d oxides generally having a larger labile O 2p population** (higher 4d\u20132p covalency), e.g. NaNbO\u2083 2.58 e\u207b, NaAgO\u2083 2.51 e\u207b, NaPdO\u2083 2.10 e\u207b, versus NaVO\u2083 0.17 e\u207b (\xa74, Tables 1\u20132). O redox is thus more dominant in 4d (fifth-row) oxides (\xa75).

### 5. The cost: capacity fade and O\u2082 release
Although anionic redox raises capacity, reduced O ions strongly tend to bond to one another and **leave the lattice as O\u2082 gas**. Besides being a fire/ignition hazard, O\u2082 release causes **rapid and permanent capacity fade**, voltage fade, and often a layered\u2192spinel transformation; this is a key factor limiting real-world use of Li-rich layered Mn oxides. Mitigation strategies such as **surface fluorination** strengthen TM\u2013O bonds, raise the barrier to oxygen-vacancy formation, and reduce Jahn\u2013Teller distortion \u2014 at the cost of some initial capacity (\xa73.1, \xa75).

### 6. How it is probed
Direct probing of anionic redox is not possible; **DFT density-of-states analysis** is the main route. Because cathodes are strongly correlated materials, plain GGA is inadequate \u2014 Hartree\u2013Fock hybrid functionals (e.g., \u03b1 = 0.2) or GGA+U (U_eff \u2248 2.0 eV) are needed, though choosing these phenomenological parameters remains a challenge (Abstract; \xa73.2, \xa74, \xa75).
\end{AnswerText}

\textbf{Retrieved evidence index (all returned units; paths and scores) }
\begin{EvidenceText}
[1] score=None path=2408.09942
[2] score=None path=2408.09942
[3] score=None path=2408.09942
[4] score=None path=2408.09942
[5] score=None path=2408.09942
[6] score=None path=2408.09942
[7] score=None path=2408.09942
[8] score=None path=2408.09942
[9] score=None path=2408.09942
[10] score=None path=2408.09942
[11] score=None path=2408.09942
[12] score=None path=2408.09942
\end{EvidenceText}
\subsection{BQ09: What problem of visual SLAM in dynamic environments does MLP-SLAM address and how?}
\begin{tabular}{@{}ll@{}}
Arm & DeepRead native pipeline\\
Generator & deepread-agent / deepseek-flash\\
HTTP / success & 200 / True\\
Wall time / cost & 29.7 s / not recorded\\
Judge scores & F=5, C=5, G=5, A=2; overall=4.25\\
Verdict & correct\\
\end{tabular}

\textbf{Index/retrieval metadata}
\begin{EvidenceText}
{
  "retrieval": "deepread semantic retrieval + agentic read",
  "generator_disclosure": "native DeepRead agent pipeline (disclosed)"
}
\end{EvidenceText}

\textbf{Dimension rationale}
\begin{EvidenceText}
{
  "factual": "Everything matches gold and paper: the static-world assumption and dynamic-scene degradation; the three prior categories with their exact defects (Type II information loss on parked cars and highway guardrails \u2014 the gold's own Figure 1 examples; Type III partial geometric information and manual tuning); the full pipeline (ORB-SLAM2 base, YOLOv7 + Deep Sort, the DFM reclassifying out-of-interval points as static background using mean depth and standard deviation, the 3-D input vector of intensity/epipolar/reprojection errors, zero-mean normalization, coarse then fine pose estimation); the >50,000-point dataset; and the ablation ATE progression 2.67 to 1.97 to 1.48, cross-consistent with other answers.",
  "completeness": "The most detailed answer in the packet: the problem, all three prior-approach failure modes, a five-step pipeline, the dataset contribution, and results \u2014 fully answered.",
  "groundedness": "All content is supported by the gold and the paper (numbers cross-check against independent answers); no fabrication.",
  "sources": "No source citations at all \u2014 no document, section, or passage ids anywhere in the answer."
}
\end{EvidenceText}

\textbf{Judge summary}
\begin{EvidenceText}
Exhaustive and accurate answer to both halves, but entirely uncited despite rich verifiable detail.
\end{EvidenceText}

\textbf{Complete answer}
\begin{AnswerText}
## The problem MLP-SLAM addresses

Most existing V-SLAM systems (ORB-SLAM2/3, LSD-SLAM, DSO, etc.) **assume a static world**. They work with high precision when scenes are mostly static, but their accuracy **degrades in dynamic environments** \u2014 outdoor scenes with pedestrians, cars, buses, and other movable objects \u2014 which are common in real-world robotics, autonomous driving, and space exploration.

Prior attempts to fix this fall into three categories, each with shortcomings:

- **Type I (geometry only):** classifies feature points using only geometric cues (epipolar error, reprojection error, intensity/depth, optical flow).
- **Type II (semantics only, e.g., DynaSLAM):** discards all features inside masks/bounding boxes of "dynamic" classes. This causes **information loss** \u2014 e.g., features on parked cars or highway guardrails are wrongly thrown away even though they are static and useful for pose estimation. It is also **high-latency** (MaskRCNN).
- **Type III (semantic + geometric, e.g., DynaSLAM II, CFP-SLAM, VDO-SLAM):** the mainstream approach, but it uses only **partial geometric information** (e.g., only epipolar error or only reprojection error), requires **manual parameter tuning**, **lacks generality**, and is laborious to fine-tune. Object-pose-based extensions also suffer high latency and poor object-pose accuracy.

A further gap: there was **no public dataset** for directly evaluating dynamic/static feature classification, so evaluation was only indirect via SLAM metrics, which is unreliable because of RANSAC randomness in pose estimation.

## How MLP-SLAM addresses it

MLP-SLAM is an **MLP-based real-time stereo SLAM system built on ORB-SLAM2**, with an added dynamic/static feature discriminator module. Its pipeline:

1. **Pre-processing:** ORB feature extraction on stereo images; object detection and tracking on the left image using **YOLOv7 + Deep Sort** (detection chosen over semantic segmentation to reduce latency; Deep Sort replaces CFP-SLAM's EKF + Hungarian matching).
2. **Depth Filter Module (DFM):** because detection boxes are less precise than segmentation masks, the DFM computes the mean depth \u03bc_d and standard deviation \u03b4_d of features in each box and reclassifies points outside [\u03bc_d \u2212 \u03b7\u03b4_d, \u03bc_d + \u03b7\u03b4_d] as **static background**, recovering information that mask-based methods would discard.
3. **Coarse estimation:** camera pose is estimated from features outside the boxes plus the DFM-recovered static background features.
4. **MLP classification:** each remaining in-box feature x_i is turned into a 3-D input vector v_i = (image intensity error e\u2081, epipolar error e_D, reprojection error e_Re) \u2014 i.e., **complete geometric information**, unlike prior work that used only some of these. After zero-mean normalization, the MLP labels each point static or dynamic. This is claimed to be the **first machine-learning-based dynamic/static feature discriminator in SLAM**, and it needs **no manual parameter configuration**.
5. **Fine estimation:** the pose is refined by minimizing reprojection error over all static points (outside-box + DFM-filtered + MLP-classified static).

**Supporting contribution:** a public dataset of **>50,000 feature points** from KITTI Odometry for training/testing the discriminator, enabling *direct* evaluation instead of indirect SLAM-metric evaluation.

**Results:** the discriminator outperforms other methods on that dataset, and the full system achieves the **highest average precision and fastest speed** on KITTI Odometry and KITTI360 compared with other dynamic SLAM systems (ablation average ATE: 2.67 \u2192 1.97 with MLP only \u2192 1.48 with both modules).
\end{AnswerText}

\textbf{Retrieved evidence index (all returned units; paths and scores) }
\begin{EvidenceText}
[1] score=None path=2410.10669
[2] score=None path=2410.10669
[3] score=None path=2410.10669
[4] score=None path=2410.10669
[5] score=None path=2410.10669
[6] score=None path=2410.10669
[7] score=None path=2410.10669
[8] score=None path=2410.10669
[9] score=None path=2410.10669
[10] score=None path=2410.10669
[11] score=None path=2410.10669
[12] score=None path=2410.10669
\end{EvidenceText}
\subsection{BQ10: What strategies for tumor control emerge under immune evasion and chemotherapy resistance?}
\begin{tabular}{@{}ll@{}}
Arm & DeepRead native pipeline\\
Generator & deepread-agent / deepseek-flash\\
HTTP / success & 200 / True\\
Wall time / cost & 55.5 s / not recorded\\
Judge scores & F=5, C=5, G=5, A=4; overall=4.75\\
Verdict & correct\\
\end{tabular}

\textbf{Index/retrieval metadata}
\begin{EvidenceText}
{
  "retrieval": "deepread semantic retrieval + agentic read",
  "generator_disclosure": "native DeepRead agent pipeline (disclosed)"
}
\end{EvidenceText}

\textbf{Dimension rationale}
\begin{EvidenceText}
{
  "factual": "All correct: the T1/T2 subpopulation model; the phenotype-guided table (checkpoint evasion -> down mu -> checkpoint inhibitors; reduced antigen presentation -> up rho, down mu -> CAR-T/vaccines; T2-dominant -> up sigma, down delta -> adoptive T-cell transfer/cytotoxic therapy); 'checkpoint inhibitors combined with moderate chemotherapy provide the most effective control strategy'; the three sequential protocols including the T1-mutant resurgence caveat; the explicit tumor-free-stability inequality sigma/(delta+omega*chi) > max{alpha2 - theta2*chi, (alpha1 - chi)/theta1}, consistent with the equilibrium form quoted by another answer; adaptive therapy and biomarker guidance; and the mechanism-specific caveats (antigen-presentation robustness to weak chemo; high-phi checkpoint evasion uncontrollable except via theta1), cross-consistent with peer answers.",
  "completeness": "Six numbered strategies spanning mechanism-matched therapy, combination, sequencing, thresholds, adaptive/precision medicine, and caveats \u2014 the most comprehensive treatment in the packet.",
  "groundedness": "No unsupported claims; content is consistent with the gold and the paper throughout.",
  "sources": "Attributes the answer to the named document (doc 22) with section references (Table 2, section 3.2, 3.3, 4) in the headings, but individual claims inside sections carry no citation tags."
}
\end{EvidenceText}

\textbf{Judge summary}
\begin{EvidenceText}
Most comprehensive and accurate account of the emerging strategies, though attributions are heading-level rather than per-claim.
\end{EvidenceText}

\textbf{Complete answer}
\begin{AnswerText}
## Strategies for tumor control under immune evasion and chemotherapy resistance

The paper (doc 22, *Strategies for tumor elimination and control under immune evasion and chemotherapy resistance*) builds an ODE model of effector cells plus two tumor subpopulations \u2014 an immune-resistant type **T\u2081** and an immune-sensitive (wild-type) type **T\u2082**, which may also be chemo-resistant \u2014 and derives control strategies from stability/bifurcation analysis. The main strategies that emerge are:

### 1. Phenotype-guided (mechanism-specific) immunotherapy
Because the two immune-evasion mechanisms are dynamically distinct, therapy must target the dominant mechanism (Table 2, \xa73.2):

| Dominant phenotype | Best parameter target | Therapeutic focus |
|---|---|---|
| T\u2081-only, **immune checkpoint evasion** | \u2193 \u03bc (immune exhaustion) | Checkpoint inhibitors |
| T\u2081-only, **reduced antigen presentation** | \u2191 \u03c1 (effector activation), \u2193 \u03bc | CAR-T therapy, cancer vaccines, checkpoint inhibitors |
| T\u2082-only (immune-sensitive dominant) | \u2191 \u03c3 (effector recruitment), \u2193 \u03b4 (effector death) | Adoptive T-cell transfer, cytotoxic therapy |

- Under checkpoint evasion, control is highly sensitive to the immune exhaustion rate \u03bc, so checkpoint inhibitors are the most effective single option.
- Under reduced antigen presentation, stability depends more on effector activation/exhaustion, so strategies that restore immune *recognition* (CAR-T, vaccines) are favored.
- When the immune-sensitive phenotype dominates, boosting effector recruitment and persistence is critical.

### 2. Combination immuno-chemotherapy
Immune dynamics alone are generally insufficient to eliminate the tumor. Chemotherapy reduces tumor burden and directly targets the sensitive subpopulation, while immunotherapy enhances immune-mediated suppression; together they substantially enlarge the parameter region with tumor-free or low-population equilibria (\xa73.3). Notably, effector recruitment (\u03c3) and effector death (\u03b4) modulation is *not* very helpful against T\u2081 without chemotherapy, but becomes effective when combined with it. When T\u2081 dominates via checkpoint evasion, "checkpoint inhibitors combined with moderate chemotherapy provide the most effective control strategy."

### 3. Sequential, composition-guided treatment protocols (\xa74)
- **T\u2081-dominant start:** chemotherapy first to reduce burden, then immunotherapy to maintain control; alternatively immunotherapy first to shrink the tumor, then chemotherapy to clear residual resistant cells.
- **T\u2082-dominant start:** chemotherapy reduces burden but risks resistant escape, so immunotherapy is applied afterward to stabilize the system.
- **Predicted coexistence of both resistances:** give chemotherapy first to eliminate T\u2081 cells, leaving a low T\u2082-only equilibrium, then immunotherapy to potentially eliminate the tumor \u2014 with the caveat that if only controlled, emergence of T\u2081 mutants could cause resurgence.

### 4. Threshold-based immune strengthening
The tumor-free equilibrium is stable iff the effective immune strength exceeds the maximal effective tumor growth rate:
\u03c3/(\u03b4 + \u03c9\u03c7) > max{\u03b1\u2082 \u2212 \u03b8\u2082\u03c7, (\u03b1\u2081 \u2212 \u03c7)/\u03b8\u2081}.
Practically this means increasing effector recruitment \u03c3, decreasing effector death \u03b4, and limiting chemotherapy-induced suppression of effectors (\u03c9). Once this threshold is crossed, the system transitions to the tumor-free equilibrium.

### 5. Adaptive therapy and precision medicine
Consistent with the adaptive-therapy paradigm, the aim is often to *control* rather than eliminate the tumor, exploiting competition between sensitive and resistant subpopulations and adjusting treatment dynamically rather than using maximum tolerated dose. Because stability thresholds depend on \u03c3/\u03b4 and max{\u03b1, \u03b1\u2081/\u03b8\u2081}, small changes in immune competence or tumor immunogenicity can flip the outcome \u2014 supporting therapy selection guided by measurable biomarkers (PD-L1 expression, T-cell infiltration, exhaustion markers, antigen presentation profiles).

### 6. Mechanism-specific caveats
Reduced antigen presentation is more robust against relatively weak chemotherapy (larger coexistence regions, stable T\u2081-only states even at low \u03c3), so such tumors require stronger immune-strengthening or recognition-restoring interventions. Also, when the immune-resistant subpopulation dominates under checkpoint evasion with high \u03c6, no feasible modulation of immune parameters controls the T\u2081-only equilibrium except altering \u03b8\u2081 \u2014 i.e., uniform, non-personalized treatment often fails.
\end{AnswerText}

\textbf{Retrieved evidence index (all returned units; paths and scores) }
\begin{EvidenceText}
[1] score=None path=2604.01385
[2] score=None path=2604.01385
[3] score=None path=2604.01385
[4] score=None path=2604.01385
[5] score=None path=2604.01385
[6] score=None path=2604.01385
[7] score=None path=2604.01385
[8] score=None path=2604.01385
[9] score=None path=2604.01385
[10] score=None path=2604.01385
[11] score=None path=2604.01385
[12] score=None path=2604.01385
\end{EvidenceText}
\section{GraphRAG local search}
\subsection{BQ01: How does the Transformer compute attention, and why does it replace recurrence and convolution?}
\begin{tabular}{@{}ll@{}}
Arm & GraphRAG local search\\
Generator & omp / deepseek-flash\\
HTTP / success & 200 / True\\
Wall time / cost & 154.9 s / 0.032133 USD\\
Judge scores & F=5, C=4, G=5, A=5; overall=4.75\\
Verdict & partial\\
\end{tabular}

\textbf{Index/retrieval metadata}
\begin{EvidenceText}
{
  "retrieval": "graphrag local-search style entity-centric assembly (fast NLP entities; shim embeddings; LanceDB absent)",
  "top_entities": [
    "ATTENTION MECHANISM",
    "ATTENTION MECHANISMS",
    "ATTENTION FUNCTION",
    "TRANSFORMER",
    "ATTENTION NETWORKS",
    "MULTI-HEAD ATTENTION",
    "ATTENTION VISUALIZATIONS",
    "SINGLE-HEAD ATTENTION",
    "MAGNITUDE CHANGES",
    "ATTENTION HEADS",
    "TRANSFORMATIVE QUANTUM TECHNOLOGIES",
    "RANDOM PERFORMANCE"
  ],
  "n_entities_total": 1816
}
\end{EvidenceText}

\textbf{Dimension rationale}
\begin{EvidenceText}
{
  "factual": "Matches gold and the real paper: encoder maps (x1..xn) to z and decodes auto-regressively (gold section 3 opening); attention scaled by 1/sqrt(d_k); multi-head h-way projection to d_k,d_k,d_v then concat and re-projection; single-head is 0.9 BLEU worse than the best head setting and quality degrades with too many heads (paper Table 3 discussion); Figure 3 layer-5 head attending a distant dependency of 'making'; constant sequential operations vs recurrence \u2014 all correct, none contradicted by gold.",
  "completeness": "'How' covers scaling, Q/K/V projections, multi-head mechanics and empirical head ablations, but explicitly states the closed-form softmax(QK^T/sqrt(d_k))V and Table 1's quantitative path-length numbers are truncated/absent. 'Why' is fully answered (sequential-operations and path-length criteria, first attention-only transduction model, faster training, new SOTA). Near-complete with the central formula honestly missing.",
  "groundedness": "All claims carry source tags ([S2][S3][S8][S11][S4][S5][S9]) and the answer adds a dedicated 'Gaps in the provided context' section refusing to treat the truncated formula and Table 1 numbers as answered \u2014 no overreach.",
  "sources": "Attributions are specific, including where [S3] truncates and that [S4] contains only the constant-sequential-operations advantage; descriptions of each source's coverage are accurate."
}
\end{EvidenceText}

\textbf{Judge summary}
\begin{EvidenceText}
Answers both halves substantively with correct facts, but the explicit attention equation and Table 1 numbers are acknowledged missing from its excerpts.
\end{EvidenceText}

\textbf{Complete answer}
\begin{AnswerText}
**How attention is computed** (as far as the context states)

- The architecture is encoder\u2013decoder: the encoder maps $(x_1,\dots,x_n)\to \mathbf{z}=(z_1,\dots,z_n)$, and the decoder auto-regressively generates $(y_1,\dots,y_m)$ one symbol at a time, consuming previously generated symbols [S2].
- Attention is **scaled dot-product** attention: dot products are scaled by $\frac{1}{\sqrt{d_k}}$ [S3]. The mechanism is described in terms of queries, keys and values, with the attention function applied to projected versions of them [S3].
- **Multi-head attention**: instead of one attention function over $d_{\text{model}}$-dimensional keys/values/queries, Q, K, V are linearly projected $h$ times to $d_k,d_k,d_v$, attention is run in parallel on each projection yielding $d_v$-dimensional outputs, and these are concatenated and projected again to the final values [S3]. The stated benefit: the model can jointly attend to information from different representation subspaces [S3].
- Empirical support: with computation held constant, single-head attention is 0.9 BLEU worse than the best head setting, and quality also degrades with too many heads; shrinking the key size $d_k$ hurts quality, suggesting dot product may be a less-than-ideal compatibility function [S8]. Attention heads do follow long-distance dependencies (Figure 3 shows a layer-5 encoder head attending to a distant dependency of "making") [S11].

**Why it replaces recurrence and convolution** (per the context)

- The comparison criterion is the **minimum number of sequential operations** ("computation that can be parallelized") and the **path length between long-range dependencies**: shorter forward/backward paths between any input/output position pair make long-range dependencies easier to learn [S4].
- A self-attention layer connects all positions with a **constant** number of sequentially executed operations, whereas recurrent layers require a number of sequential steps proportional to sequence length \u2014 hence attention is far more parallelizable and has shorter dependency paths [S4].
- Result: the Transformer is "the first sequence transduction model based entirely on attention, replacing the recurrent layers most commonly used in encoder-decoder architectures with multi-headed self-attention," and it trains significantly faster than recurrent- or convolutional-layer architectures while setting new state of the art on WMT 2014 EN\u2013DE and EN\u2013FR [S9]; Table 2 reports BLEU and FLOPs against ByteNet, ConvS2S, GNMT+RL, MoE, etc. [S5].

**Gaps in the provided context** \u2014 do not treat these as answered:

- The explicit closed-form $\mathrm{softmax}(QK^\top/\sqrt{d_k})V$ and the exact wording of why scaling is needed ("large magnitude \u2192 small gradients") are truncated in [S3]; only the scaling factor survives.
- Table 1's quantitative path-length comparison (self-attention $O(1)$, recurrent $O(n)$, convolutional $O(\log_k n)$) and the argument that convolutional layers need $O(\log_k n)$ layers to relate distant positions are **not present** in the supplied excerpts \u2014 [S4] only states the constant-sequential-operations advantage for self-attention.
\end{AnswerText}

\textbf{Retrieved evidence index (all returned units; paths and scores) }
\begin{EvidenceText}
[1] score=None path=part_0057___artificial-intelligence__1706.03762__attention-is-all-you-need.md_part_1_of_2_.md.txt
[2] score=None path=part_0057___artificial-intelligence__1706.03762__attention-is-all-you-need.md_part_1_of_2_.md.txt
[3] score=None path=part_0057___artificial-intelligence__1706.03762__attention-is-all-you-need.md_part_1_of_2_.md.txt
[4] score=None path=part_0057___artificial-intelligence__1706.03762__attention-is-all-you-need.md_part_1_of_2_.md.txt
[5] score=None path=part_0057___artificial-intelligence__1706.03762__attention-is-all-you-need.md_part_1_of_2_.md.txt
[6] score=None path=part_0057___artificial-intelligence__1706.03762__attention-is-all-you-need.md_part_1_of_2_.md.txt
[7] score=None path=part_0058___artificial-intelligence__1706.03762__attention-is-all-you-need.md_part_2_of_2_.md.txt
[8] score=None path=part_0058___artificial-intelligence__1706.03762__attention-is-all-you-need.md_part_2_of_2_.md.txt
[9] score=None path=part_0058___artificial-intelligence__1706.03762__attention-is-all-you-need.md_part_2_of_2_.md.txt
[10] score=None path=part_0058___artificial-intelligence__1706.03762__attention-is-all-you-need.md_part_2_of_2_.md.txt
[11] score=None path=part_0058___artificial-intelligence__1706.03762__attention-is-all-you-need.md_part_2_of_2_.md.txt
[12] score=None path=part_0177___quantum-physics__1801.00862__quantum-computing-in-the-nisq-era-and-beyond.md_part_2_of_3_.md.txt
\end{EvidenceText}
\subsection{BQ02: What defines NISQ technology and what is its central limitation?}
\begin{tabular}{@{}ll@{}}
Arm & GraphRAG local search\\
Generator & omp / deepseek-flash\\
HTTP / success & 200 / True\\
Wall time / cost & 38.6 s / 0.022130999999999998 USD\\
Judge scores & F=5, C=4, G=5, A=5; overall=4.75\\
Verdict & correct\\
\end{tabular}

\textbf{Index/retrieval metadata}
\begin{EvidenceText}
{
  "retrieval": "graphrag local-search style entity-centric assembly (fast NLP entities; shim embeddings; LanceDB absent)",
  "top_entities": [
    "NISQ",
    "AHRQ",
    "IMPORTANT LIMITATIONS",
    "PIQA",
    "LIVEQA",
    "MEDICATIONQA",
    "LACK GENERALITY",
    "NI",
    "POWERFUL QUANTUM TECHNOLOGIES",
    "INEQUALITY CHARACTERIZES",
    "QUANTUM TECHNOLOGY",
    "NARCCAP"
  ],
  "n_entities_total": 1816
}
\end{EvidenceText}

\textbf{Dimension rationale}
\begin{EvidenceText}
{
  "factual": "No errors: 50-100 qubit devices 'will be available in the near future' that 'may be able to perform tasks which surpass the capabilities of today's classical digital computers' matches the gold abstract; the central limitation 'noise in quantum gates limits the size of quantum circuits that can be executed reliably' is the gold's exact statement; the measurement-disturbance root-cause claim is standard physics attributed to [S12], and 'useful tools for exploring many-body quantum physics' matches the abstract.",
  "completeness": "Both parts answered but at abstract-level depth: the definition is given (acronym expansion, 50-100 qubits, near future) without unpacking 'intermediate-scale' vs 'noisy' or the no-error-correction clause, and the limitation is stated qualitatively without the ~1000-gate quantification; a hardware-difficulty root cause and the many-body-physics consequence are added.",
  "groundedness": "All claims are explicitly sourced to [S11] and [S12]; it transparently states the other nine sources 'concern SLAM and Yiddish NLP and contain nothing about NISQ', avoiding cross-contamination.",
  "sources": "Attributions are clear and honest, with an explicit note on which sources were used and why the rest were ignored."
}
\end{EvidenceText}

\textbf{Judge summary}
\begin{EvidenceText}
Correctly answers both halves from its (narrower) sources with full honesty about source coverage, though with less definitional and quantitative depth than the best answers.
\end{EvidenceText}

\textbf{Complete answer}
\begin{AnswerText}
**Definition (from [S11]):** NISQ stands for *Noisy Intermediate-Scale Quantum* technology \u2014 devices with roughly **50\u2013100 qubits** that "will be available in the near future" and "may be able to perform tasks which surpass the capabilities of today's classical digital computers."

**Central limitation (from [S11]):** **Noise in quantum gates limits the size of quantum circuits that can be executed reliably.** So despite being potentially classically intractable in principle, these devices cannot run long/deep circuits without errors corrupting the computation.

**Root cause (from [S12]):** Noise is intrinsic to the physics \u2014 "we cannot observe a quantum system without producing an uncontrollable disturbance in the system." Reliable storage/processing requires near-perfect isolation from the environment, yet the same qubits must interact strongly with each other, be externally controllable, and be readable out. Satisfying all of these simultaneously is what makes building the hardware "very challenging" [S12].

Consequence noted in [S11]: NISQ devices are therefore positioned as useful tools for **exploring many-body quantum physics**, with other applications still speculative \u2014 not as general-purpose fault-tolerant machines.

*Note: the other provided sources ([S1]\u2013[S10]) concern SLAM and Yiddish NLP and contain nothing about NISQ; the answer draws solely on [S11] and [S12].*
\end{AnswerText}

\textbf{Retrieved evidence index (all returned units; paths and scores) }
\begin{EvidenceText}
[1] score=None path=part_0081___robotics__2410.10669__mlp-slam-multilayer-perceptron-based-simultaneou.md_part_1_of_2_.md.txt
[2] score=None path=part_0081___robotics__2410.10669__mlp-slam-multilayer-perceptron-based-simultaneou.md_part_1_of_2_.md.txt
[3] score=None path=part_0091___nlp__2608.05850__mameloshnlm-yiddish-language-model-and-evalu.md_part_1_of_4_.md.txt
[4] score=None path=part_0091___nlp__2608.05850__mameloshnlm-yiddish-language-model-and-evalu.md_part_1_of_4_.md.txt
[5] score=None path=part_0091___nlp__2608.05850__mameloshnlm-yiddish-language-model-and-evalu.md_part_1_of_4_.md.txt
[6] score=None path=part_0093___nlp__2608.05850__mameloshnlm-yiddish-language-model-and-evalu.md_part_3_of_4_.md.txt
[7] score=None path=part_0093___nlp__2608.05850__mameloshnlm-yiddish-language-model-and-evalu.md_part_3_of_4_.md.txt
[8] score=None path=part_0093___nlp__2608.05850__mameloshnlm-yiddish-language-model-and-evalu.md_part_3_of_4_.md.txt
[9] score=None path=part_0093___nlp__2608.05850__mameloshnlm-yiddish-language-model-and-evalu.md_part_3_of_4_.md.txt
[10] score=None path=part_0094___nlp__2608.05850__mameloshnlm-yiddish-language-model-and-evalu.md_part_4_of_4_.md.txt
[11] score=None path=part_0176___quantum-physics__1801.00862__quantum-computing-in-the-nisq-era-and-beyond.md_part_1_of_3_.md.txt
[12] score=None path=part_0176___quantum-physics__1801.00862__quantum-computing-in-the-nisq-era-and-beyond.md_part_1_of_3_.md.txt
\end{EvidenceText}
\subsection{BQ03: What is MultiMedQA and how did Flan-PaLM perform on the MedQA benchmark?}
\begin{tabular}{@{}ll@{}}
Arm & GraphRAG local search\\
Generator & omp / deepseek-flash\\
HTTP / success & 200 / True\\
Wall time / cost & 129.3 s / 0.0258462 USD\\
Judge scores & F=2, C=1, G=2, A=3; overall=2.00\\
Verdict & incorrect\\
\end{tabular}

\textbf{Index/retrieval metadata}
\begin{EvidenceText}
{
  "retrieval": "graphrag local-search style entity-centric assembly (fast NLP entities; shim embeddings; LanceDB absent)",
  "top_entities": [
    "MULTIMEDQA",
    "MEDMCQA",
    "MEDQA",
    "MEDICATIONQA",
    "PUBMEDQA",
    "MED-PALM",
    "LIVEQA",
    "FLAN-PALM",
    "HEALTHSEARCHQA",
    "MULTIMODEL ANALYSIS",
    "MULTIMODEL MEDIAN",
    "MED-PALM FLAN-PALM"
  ],
  "n_entities_total": 1816
}
\end{EvidenceText}

\textbf{Dimension rationale}
\begin{EvidenceText}
{
  "factual": "Its two central factual assertions are contradicted by the GOLD material: it claims 'The provided context never uses the term MultiMedQA', but the gold excerpt's abstract says 'we present MultiMedQA, a benchmark combining six existing open question answering datasets...'; and it claims the context 'never mentions Flan-PaLM and reports no MedQA accuracy figure', but the gold states 'Flan-PaLM achieves state-of-the-art accuracy... including 67.6% accuracy on MedQA'. Secondary details are correct (MedQA 11,450 dev/1,273 test, the BLEU-metrics motivation, paper identity).",
  "completeness": "Declines both parts of the question ('Insufficient context to answer'), so neither the MultiMedQA definition nor Flan-PaLM's MedQA performance is answered.",
  "groundedness": "This is not an honest refusal grounded in the material \u2014 the absence claims ('never uses the term', 'never mentions Flan-PaLM') are false against the gold, which contains both; only the retained observations (dataset details, motivation) are supported.",
  "sources": "It does make explicit attributions ([S5] for the paper's identity, [S6] for the motivation, [S8][S9][S10] for dataset listings), but its source coverage evidently missed the abstract, producing the wrong absence claims."
}
\end{EvidenceText}

\textbf{Judge summary}
\begin{EvidenceText}
Refuses to answer based on the false premise that MultiMedQA and Flan-PaLM are absent from the material, when the gold excerpt defines the benchmark and reports the 67.6% result directly.
\end{EvidenceText}

\textbf{Complete answer}
\begin{AnswerText}
**Insufficient context to answer.**

- The provided context never uses the term **MultiMedQA**, so it does not define or describe that benchmark. It only lists individual medical QA datasets \u2014 MedQA (USMLE, 4-choice, 11,450 dev / 1,273 test), MedMCQA, PubMedQA, MMLU medical subsets, LiveQA TREC-2017 \u2014 as component tasks [S8][S9][S10].
- The context never mentions **Flan-PaLM** and reports no MedQA accuracy figure for it [S5].

What the context does support: the paper is *Large Language Models Encode Clinical Knowledge* (Singhal et al.) [S5], and its motivation is that existing medical QA benchmarks are limited to classification accuracy or automated NLG metrics (e.g., BLEU) and therefore do not permit detailed analysis of model answers [S6].

Given only the supplied excerpts (parts 1\u20132 of 8), the MultiMedQA definition and Flan-PaLM's MedQA score are not present, so I cannot state them without guessing.
\end{AnswerText}

\textbf{Retrieved evidence index (all returned units; paths and scores) }
\begin{EvidenceText}
[1] score=None path=part_0179___climate-science__1503.07557__precipitation-extremes-under-climate-change.md_part_1_of_3_.md.txt
[2] score=None path=part_0179___climate-science__1503.07557__precipitation-extremes-under-climate-change.md_part_1_of_3_.md.txt
[3] score=None path=part_0181___climate-science__1503.07557__precipitation-extremes-under-climate-change.md_part_3_of_3_.md.txt
[4] score=None path=part_0181___climate-science__1503.07557__precipitation-extremes-under-climate-change.md_part_3_of_3_.md.txt
[5] score=None path=part_0241___clinical-medicine__2212.13138__large-language-models-encode-clinical-knowledge.md_part_1_of_8_.md.txt
[6] score=None path=part_0241___clinical-medicine__2212.13138__large-language-models-encode-clinical-knowledge.md_part_1_of_8_.md.txt
[7] score=None path=part_0241___clinical-medicine__2212.13138__large-language-models-encode-clinical-knowledge.md_part_1_of_8_.md.txt
[8] score=None path=part_0241___clinical-medicine__2212.13138__large-language-models-encode-clinical-knowledge.md_part_1_of_8_.md.txt
[9] score=None path=part_0241___clinical-medicine__2212.13138__large-language-models-encode-clinical-knowledge.md_part_1_of_8_.md.txt
[10] score=None path=part_0241___clinical-medicine__2212.13138__large-language-models-encode-clinical-knowledge.md_part_1_of_8_.md.txt
[11] score=None path=part_0242___clinical-medicine__2212.13138__large-language-models-encode-clinical-knowledge.md_part_2_of_8_.md.txt
[12] score=None path=part_0242___clinical-medicine__2212.13138__large-language-models-encode-clinical-knowledge.md_part_2_of_8_.md.txt
\end{EvidenceText}
\subsection{BQ04: Which three independent ontologies subdivide the Gene Ontology?}
\begin{tabular}{@{}ll@{}}
Arm & GraphRAG local search\\
Generator & omp / deepseek-flash\\
HTTP / success & 200 / True\\
Wall time / cost & 181.4 s / 0.026448 USD\\
Judge scores & F=5, C=1, G=5, A=5; overall=4.00\\
Verdict & incorrect\\
\end{tabular}

\textbf{Index/retrieval metadata}
\begin{EvidenceText}
{
  "retrieval": "graphrag local-search style entity-centric assembly (fast NLP entities; shim embeddings; LanceDB absent)",
  "top_entities": [
    "GENE ONTOLOGY",
    "GENE ONTOLOGY CONSORTIUM",
    "GENE",
    "GENOME",
    "MEDICAL GENETICS",
    "SEPARATE ANALYSIS",
    "REFERENCE GENOME PROJECT",
    "IMMUNOLOGY",
    "GENE PRODUCT",
    "GENE PRODUCTS",
    "DIFFERENT SUBPOPULATIONS",
    "BIOLOGY"
  ],
  "n_entities_total": 1816
}
\end{EvidenceText}

\textbf{Dimension rationale}
\begin{EvidenceText}
{
  "factual": "No incorrect claims: its description of its context (chunks on term structure and obsoletion [S5], GAF 2.1 [S6], evidence codes [S7][S8]) is plausible, and its clarification that the 'three main categories' in [S7] are 'experimental, curated non-experimental, and electronically inferred evidence codes' \u2014 not ontology subdivisions \u2014 is true of the primer. The refusal itself is a coverage failure, not a factual error.",
  "completeness": "It does not name any of the three ontologies ('The provided context does not contain the answer'), so the question is unanswered; this is an honest refusal case, so groundedness is preserved but completeness sits at the floor.",
  "groundedness": "It explicitly refuses to supply names from memory 'since that claim is not grounded in the supplied sources' \u2014 a model of honest non-fabrication; every observation it offers concerns what its chunks do contain.",
  "sources": "Attributions are specific about which chunk covers what, and it correctly disambiguates the evidence-code taxonomy from the ontology subdivision."
}
\end{EvidenceText}

\textbf{Judge summary}
\begin{EvidenceText}
Honest, non-fabricating refusal that correctly parses its sources, but its chunks lacked the subdivision sentence so the question goes unanswered.
\end{EvidenceText}

\textbf{Complete answer}
\begin{AnswerText}
The provided context does not contain the answer.

No excerpt states which three ontologies subdivide the Gene Ontology. The primer's chunks here cover GO term structure, revisions/obsoletion [S5], the GAF 2.1 annotation format [S6], and evidence codes [S7][S8] \u2014 the "three main categories" in [S7] refers to experimental, curated non-experimental, and electronically inferred *evidence codes*, not ontology subdivisions. I will not supply the three ontology names from memory, since that claim is not grounded in the supplied sources.
\end{AnswerText}

\textbf{Retrieved evidence index (all returned units; paths and scores) }
\begin{EvidenceText}
[1] score=None path=part_0241___clinical-medicine__2212.13138__large-language-models-encode-clinical-knowledge.md_part_1_of_8_.md.txt
[2] score=None path=part_0241___clinical-medicine__2212.13138__large-language-models-encode-clinical-knowledge.md_part_1_of_8_.md.txt
[3] score=None path=part_0241___clinical-medicine__2212.13138__large-language-models-encode-clinical-knowledge.md_part_1_of_8_.md.txt
[4] score=None path=part_0249___genomics__1602.01876__primer-on-the-gene-ontology.md_part_1_of_2_.md.txt
[5] score=None path=part_0249___genomics__1602.01876__primer-on-the-gene-ontology.md_part_1_of_2_.md.txt
[6] score=None path=part_0249___genomics__1602.01876__primer-on-the-gene-ontology.md_part_1_of_2_.md.txt
[7] score=None path=part_0249___genomics__1602.01876__primer-on-the-gene-ontology.md_part_1_of_2_.md.txt
[8] score=None path=part_0249___genomics__1602.01876__primer-on-the-gene-ontology.md_part_1_of_2_.md.txt
[9] score=None path=part_0249___genomics__1602.01876__primer-on-the-gene-ontology.md_part_1_of_2_.md.txt
[10] score=None path=part_0250___genomics__1602.01876__primer-on-the-gene-ontology.md_part_2_of_2_.md.txt
[11] score=None path=part_0250___genomics__1602.01876__primer-on-the-gene-ontology.md_part_2_of_2_.md.txt
[12] score=None path=part_0263___immunology__2604.01385__strategies-for-tumor-elimination-and-control.md_part_1_of_3_.md.txt
\end{EvidenceText}
\subsection{BQ05: Why does historical European economic growth contradict Unified Growth Theory?}
\begin{tabular}{@{}ll@{}}
Arm & GraphRAG local search\\
Generator & omp / deepseek-flash\\
HTTP / success & 200 / True\\
Wall time / cost & 132.5 s / 0.026118000000000002 USD\\
Judge scores & F=5, C=5, G=5, A=5; overall=5.00\\
Verdict & correct\\
\end{tabular}

\textbf{Index/retrieval metadata}
\begin{EvidenceText}
{
  "retrieval": "graphrag local-search style entity-centric assembly (fast NLP entities; shim embeddings; LanceDB absent)",
  "top_entities": [
    "THEORY CONTRADICTED",
    "ECONOMIC GROWTH TRAJECTORIES",
    "ECONOMIC GROWTH DATA",
    "GROWTH REGIME",
    "ECONOMIC GROWTH",
    "GROWTH RATE",
    "ECONOMIC DEVELOPMENT",
    "GDP",
    "DATA PROBLEM",
    "DIFFERENT DATASETS",
    "HYPERBOLIC GROWTH",
    "DIFFERENT CLIMATES"
  ],
  "n_entities_total": 1816
}
\end{EvidenceText}

\textbf{Dimension rationale}
\begin{EvidenceText}
{
  "factual": "Quotes check against gold and the paper: 'There was no transition from stagnation to growth at any time'; escape 'did not happen because there was no trap'; the post-Malthusian regime overlapping the Industrial Revolution 1760-1840; 'no connection, whatever, between the data and the Unified Growth Theory'; growth 'started to be diverted to a slower trajectory' at the start of the 'mythical regime of sustained-growth'. All correct.",
  "completeness": "Covers every core reason: no transition, no trap/escape, Industrial Revolution failure to boost growth in the theory's best test region, the perverse deceleration at the alleged regime launch, and the overall verdict \u2014 the question is fully answered at the argumentative level its excerpts support.",
  "groundedness": "Every bullet is anchored to [S9][S10][S11][S12] with direct quotes; the caveat honestly states the quantitative/hyperbolic model details are absent from its excerpts, so no numbers are fabricated.",
  "sources": "Clear per-claim attributions with quotes, plus an explicit caveat about what its sources do not contain."
}
\end{EvidenceText}

\textbf{Judge summary}
\begin{EvidenceText}
Complete qualitative answer with verbatim quotes for each contradiction and honest disclosure that its excerpts lack the quantitative fits.
\end{EvidenceText}

\textbf{Complete answer}
\begin{AnswerText}
**The contradiction, per the paper [S9]:** European growth data show no regime transition at all \u2014 so the three-stage Galor story (Malthusian stagnation \u2192 post-Malthusian \u2192 sustained growth) fails to describe, let alone explain, the data.

Key points from the excerpts:

- **No transition from stagnation to growth ever occurred.** For the 12 leading Western European countries, "There was no transition from stagnation to growth at any time" [S11]. Growth was hyperbolic both before and after the alleged transition, so the claimed escape from the Malthusian trap "did not happen because there was no trap" [S11].
- **The Industrial Revolution did not boost growth.** The paper notes the alleged post-Malthusian regime overlaps the Industrial Revolution (1760\u20131840) [S10], yet the data show the Industrial Revolution "did not boost the economic growth" [S11]. Since Western Europe \u2014 the center of the Industrial Revolution, which Galor called "the prime engine of economic growth" \u2014 is where the theory should fit best, its failure there is especially damaging [S9][S10].
- **Growth decelerated, not accelerated, at the supposed launch point.** At "the beginning of the mythical regime of sustained-growth, when the economic growth was supposed to have been launched from stagnation to a fast-increasing trajectory, the growth started to be diverted to a slower trajectory" [S12].
- **Overall verdict:** "There is no connection, whatever, between the data and the Unified Growth Theory" [S9]; "Galor's three regimes of growth have no relevance to the description, let alone to the explanation, of the mechanism of the economic growth" [S11].

So the contradiction is empirical: the persistent hyperbolic growth pattern and absent stagnation-to-growth breakpoint are incompatible with the theory's staged regime structure, precisely in the region where the theory should be most convincing [S9][S11][S12].

*Caveat:* the supplied excerpts do not include the paper's quantitative/hyperbolic-growth model details beyond these claims, so I can only report the argument at this level.
\end{AnswerText}

\textbf{Retrieved evidence index (all returned units; paths and scores) }
\begin{EvidenceText}
[1] score=None path=part_0179___climate-science__1503.07557__precipitation-extremes-under-climate-change.md_part_1_of_3_.md.txt
[2] score=None path=part_0181___climate-science__1503.07557__precipitation-extremes-under-climate-change.md_part_3_of_3_.md.txt
[3] score=None path=part_0241___clinical-medicine__2212.13138__large-language-models-encode-clinical-knowledge.md_part_1_of_8_.md.txt
[4] score=None path=part_0241___clinical-medicine__2212.13138__large-language-models-encode-clinical-knowledge.md_part_1_of_8_.md.txt
[5] score=None path=part_0245___clinical-medicine__2212.13138__large-language-models-encode-clinical-knowledge.md_part_5_of_8_.md.txt
[6] score=None path=part_0245___clinical-medicine__2212.13138__large-language-models-encode-clinical-knowledge.md_part_5_of_8_.md.txt
[7] score=None path=part_0263___immunology__2604.01385__strategies-for-tumor-elimination-and-control.md_part_1_of_3_.md.txt
[8] score=None path=part_0263___immunology__2604.01385__strategies-for-tumor-elimination-and-control.md_part_1_of_3_.md.txt
[9] score=None path=part_0320___economics__1512.08067__unified-growth-theory-contradicted-by-the-econom.md.md.txt
[10] score=None path=part_0320___economics__1512.08067__unified-growth-theory-contradicted-by-the-econom.md.md.txt
[11] score=None path=part_0320___economics__1512.08067__unified-growth-theory-contradicted-by-the-econom.md.md.txt
[12] score=None path=part_0320___economics__1512.08067__unified-growth-theory-contradicted-by-the-econom.md.md.txt
\end{EvidenceText}
\subsection{BQ06: Which physical factor primarily controls the response of precipitation extremes to climate change?}
\begin{tabular}{@{}ll@{}}
Arm & GraphRAG local search\\
Generator & omp / deepseek-flash\\
HTTP / success & 200 / True\\
Wall time / cost & 72.2 s / 0.022884 USD\\
Judge scores & F=5, C=3, G=5, A=5; overall=4.50\\
Verdict & partial\\
\end{tabular}

\textbf{Index/retrieval metadata}
\begin{EvidenceText}
{
  "retrieval": "graphrag local-search style entity-centric assembly (fast NLP entities; shim embeddings; LanceDB absent)",
  "top_entities": [
    "TEMPERATURE CHANGES",
    "PHYSICAL FACTORS",
    "RAINFALL EXTREMES",
    "CURRENT CLIMATE",
    "SNOWFALL EXTREMES BEHAVE",
    "DIFFERENT CLIMATES",
    "NATURE PHYSICS",
    "CLIMATE",
    "CLIMATE MODELS",
    "TEMPERATURE ANOMALIES",
    "CLIMATE-CHANGE SIMULATIONS",
    "GLOBAL CLIMATE MODELS"
  ],
  "n_entities_total": 1816
}
\end{EvidenceText}

\textbf{Dimension rationale}
\begin{EvidenceText}
{
  "factual": "Everything it states is true against gold: 'several physical contributions govern the response' is gold-verbatim; the decomposition into dynamical (omega), thermodynamic, and precipitation-efficiency contributions matches the paper; the short-duration factors list (relative humidity, large-scale dynamics, temperature gradients, stratiform-convective transitions) is plausible paper content. No errors.",
  "completeness": "It reports only that 'no single factor is named as the primary control' in its retrieved excerpts and stops there, never identifying the thermodynamic contribution as the robust, well-understood one \u2014 the paper's own text (abstract) does support exactly that ranking, so the core answer is withheld.",
  "groundedness": "An honest refusal: it explicitly notes the passage truncation ('the truncated passage does not complete the list') and attributes each item to [S3][S5][S10]; no fabrication.",
  "sources": "Clear attributions for each listed contribution and the caveat that its excerpts' ranking information is absent."
}
\end{EvidenceText}

\textbf{Judge summary}
\begin{EvidenceText}
Factually clean and honest, but it fails to give the thermodynamic answer the question seeks, offering only the multi-factor decomposition.
\end{EvidenceText}

\textbf{Complete answer}
\begin{AnswerText}
The provided context does not explicitly identify a single physical factor as the primary control. It only states that "several physical contributions govern the response" [S3] and decomposes changes in extreme precipitation rates into:

- a **dynamical** contribution from changes in vertical velocity $\omega$ [S5],
- a **thermodynamic** contribution from temperature/moisture changes [S5] (the truncated passage does not complete the list),
- plus **precipitation efficiency** effects [S5].

For short-duration scaling with local temperature, it further lists relative humidity, large-scale dynamics and temperature gradients, and stratiform-to-convective transitions as "thought to be important" [S10].

So: per this context, no single factor is named as the primary control \u2014 the paper frames the response as governed by multiple contributions (dynamics/$\omega$, thermodynamics, precipitation efficiency), with the ranking not stated in the retrieved excerpts.
\end{AnswerText}

\textbf{Retrieved evidence index (all returned units; paths and scores) }
\begin{EvidenceText}
[1] score=None path=part_0178___quantum-physics__1801.00862__quantum-computing-in-the-nisq-era-and-beyond.md_part_3_of_3_.md.txt
[2] score=None path=part_0178___quantum-physics__1801.00862__quantum-computing-in-the-nisq-era-and-beyond.md_part_3_of_3_.md.txt
[3] score=None path=part_0179___climate-science__1503.07557__precipitation-extremes-under-climate-change.md_part_1_of_3_.md.txt
[4] score=None path=part_0179___climate-science__1503.07557__precipitation-extremes-under-climate-change.md_part_1_of_3_.md.txt
[5] score=None path=part_0179___climate-science__1503.07557__precipitation-extremes-under-climate-change.md_part_1_of_3_.md.txt
[6] score=None path=part_0179___climate-science__1503.07557__precipitation-extremes-under-climate-change.md_part_1_of_3_.md.txt
[7] score=None path=part_0179___climate-science__1503.07557__precipitation-extremes-under-climate-change.md_part_1_of_3_.md.txt
[8] score=None path=part_0179___climate-science__1503.07557__precipitation-extremes-under-climate-change.md_part_1_of_3_.md.txt
[9] score=None path=part_0180___climate-science__1503.07557__precipitation-extremes-under-climate-change.md_part_2_of_3_.md.txt
[10] score=None path=part_0180___climate-science__1503.07557__precipitation-extremes-under-climate-change.md_part_2_of_3_.md.txt
[11] score=None path=part_0180___climate-science__1503.07557__precipitation-extremes-under-climate-change.md_part_2_of_3_.md.txt
[12] score=None path=part_0181___climate-science__1503.07557__precipitation-extremes-under-climate-change.md_part_3_of_3_.md.txt
\end{EvidenceText}
\subsection{BQ07: What is MameLoshnLM and what parameter scale does it have?}
\begin{tabular}{@{}ll@{}}
Arm & GraphRAG local search\\
Generator & omp / deepseek-flash\\
HTTP / success & 200 / True\\
Wall time / cost & 31.1 s / 0.016961999999999998 USD\\
Judge scores & F=5, C=1, G=5, A=4; overall=3.75\\
Verdict & incorrect\\
\end{tabular}

\textbf{Index/retrieval metadata}
\begin{EvidenceText}
{
  "retrieval": "graphrag local-search style entity-centric assembly (fast NLP entities; shim embeddings; LanceDB absent)",
  "top_entities": [
    "MAMELOSHNLM",
    "LEARNABLE PARAMETERS",
    "EUROLLM",
    "BASELINE PARAMETER VALUES",
    "PARAMETER RANGE",
    "MODEL PARAMETERS",
    "MLP-BASED",
    "MLP",
    "LLM",
    "USMLE",
    "MINH-THANG LUONG",
    "MLP-BASED SLAM"
  ],
  "n_entities_total": 1816
}
\end{EvidenceText}

\textbf{Dimension rationale}
\begin{EvidenceText}
{
  "factual": "Makes no incorrect claims about content: it truthfully reports that nothing in its sources [S1]-[S12] mentions the name, and describes its sources' topics (Transformer for [S1,S2], MLP-SLAM for [S3]-[S12]) plausibly and specifically; the gold material does contain the answer, but Q's retrieved context evidently does not.",
  "completeness": "Refuses both parts ('I cannot answer from this context') \u2014 the question goes unanswered.",
  "groundedness": "An honest refusal, not fabrication: it declines to guess the definition or scale from off-topic sources.",
  "sources": "Its description of what each source covers is specific and plausible, though it provides no content-bearing attribution for the question itself."
}
\end{EvidenceText}

\textbf{Judge summary}
\begin{EvidenceText}
Honest refusal after apparently retrieving off-topic sources, leaving both the definition and the parameter scale unanswered.
\end{EvidenceText}

\textbf{Complete answer}
\begin{AnswerText}
The provided context does not contain any information about "MameLoshnLM." Nothing in [S1]\u2013[S12] mentions that name or its parameter scale \u2014 the sources cover the Transformer [S1, S2] and MLP-SLAM [S3]\u2013[S12]. I cannot answer from this context.
\end{AnswerText}

\textbf{Retrieved evidence index (all returned units; paths and scores) }
\begin{EvidenceText}
[1] score=None path=part_0058___artificial-intelligence__1706.03762__attention-is-all-you-need.md_part_2_of_2_.md.txt
[2] score=None path=part_0058___artificial-intelligence__1706.03762__attention-is-all-you-need.md_part_2_of_2_.md.txt
[3] score=None path=part_0081___robotics__2410.10669__mlp-slam-multilayer-perceptron-based-simultaneou.md_part_1_of_2_.md.txt
[4] score=None path=part_0081___robotics__2410.10669__mlp-slam-multilayer-perceptron-based-simultaneou.md_part_1_of_2_.md.txt
[5] score=None path=part_0081___robotics__2410.10669__mlp-slam-multilayer-perceptron-based-simultaneou.md_part_1_of_2_.md.txt
[6] score=None path=part_0081___robotics__2410.10669__mlp-slam-multilayer-perceptron-based-simultaneou.md_part_1_of_2_.md.txt
[7] score=None path=part_0081___robotics__2410.10669__mlp-slam-multilayer-perceptron-based-simultaneou.md_part_1_of_2_.md.txt
[8] score=None path=part_0081___robotics__2410.10669__mlp-slam-multilayer-perceptron-based-simultaneou.md_part_1_of_2_.md.txt
[9] score=None path=part_0081___robotics__2410.10669__mlp-slam-multilayer-perceptron-based-simultaneou.md_part_1_of_2_.md.txt
[10] score=None path=part_0082___robotics__2410.10669__mlp-slam-multilayer-perceptron-based-simultaneou.md_part_2_of_2_.md.txt
[11] score=None path=part_0082___robotics__2410.10669__mlp-slam-multilayer-perceptron-based-simultaneou.md_part_2_of_2_.md.txt
[12] score=None path=part_0082___robotics__2410.10669__mlp-slam-multilayer-perceptron-based-simultaneou.md_part_2_of_2_.md.txt
\end{EvidenceText}
\subsection{BQ08: How does oxygen redox contribute to the capacity of battery cathodes?}
\begin{tabular}{@{}ll@{}}
Arm & GraphRAG local search\\
Generator & omp / deepseek-flash\\
HTTP / success & 200 / True\\
Wall time / cost & 265.8 s / 0.032568 USD\\
Judge scores & F=5, C=5, G=5, A=5; overall=5.00\\
Verdict & correct\\
\end{tabular}

\textbf{Index/retrieval metadata}
\begin{EvidenceText}
{
  "retrieval": "graphrag local-search style entity-centric assembly (fast NLP entities; shim embeddings; LanceDB absent)",
  "top_entities": [
    "BATTERY CATHODES",
    "O REDOX REACTION",
    "OXYGEN REDOX",
    "REDOX REACTION",
    "REDOX REACTIONS",
    "REDOX ACTIVITY",
    "BATTERY MATERIALS",
    "O REDOX",
    "ANIONIC REDOX ACTIVITY",
    "METAL OXIDES",
    "BATTERY ELECTRODES",
    "OXYGEN ATOMS"
  ],
  "n_entities_total": 1816
}
\end{EvidenceText}

\textbf{Dimension rationale}
\begin{EvidenceText}
{
  "factual": "All correct: 'play[s] a role in governing the useful capacity' is gold-verbatim; bulk reversible O redox as a capacity source; TM t2g states near EF barely hybridizing with O 2p; Li2TMO3 Li-in-TM-layer producing orphaned O 2p states; the Li4Mn2O5 worked example (Mn3+ to Mn4+ first, then O2- to O1- accommodating additional de-lithiation, Mn4+ to Mn5+ only at x<1) cross-corroborated by two other answers; ilmenite-type O under-coordination; the O 2p population tables; O2 evolution as the detrimental branch.",
  "completeness": "Short answer plus mechanism, electronic-structure origin, material cases, a quantified worked example, the cost branch, and honest notes on what is absent (no mAh/g figure for O redox; an OCR-corrupted formula flagged) \u2014 fully answered.",
  "groundedness": "Every claim tagged ([S1]-[S8]) with quotes matching gold or the paper; the two declared gaps are honesty markers, not errors.",
  "sources": "Per-claim attributions including the specific source for each material system and the table reference."
}
\end{EvidenceText}

\textbf{Judge summary}
\begin{EvidenceText}
Complete, precisely sourced answer with the orbital mechanism, material cases, and a quantified worked example, plus transparent notes on context gaps.
\end{EvidenceText}

\textbf{Complete answer}
\begin{AnswerText}
**Short answer:** oxygen (anionic) redox adds capacity *on top of* transition-metal cation redox, but the context gives the mechanism qualitatively rather than as a capacity number.

**Mechanism and where the extra capacity comes from**

- Oxygen participation in redox is stated to "play a role in governing the useful capacity" of Li- and Na-ion cathodes [S1]. Reversible O redox in the **bulk** of the cathode is what "offer[s] pathways towards the reversible capacit[y]" \u2014 i.e., it is a capacity source, not just a degradation channel [S2].
- The electronic-structure origin: TM\u2013O hybridization produces M\u2013O antibonding states; the TM $t_{2g}$ orbitals barely hybridize with O $2p$ and sit as "d-like" states between bonding and antibonding levels, near $E_F$ [S3]. When the cation's frontier level is too low to be oxidized further, charge compensation must come from elsewhere \u2014 ideally from oxygen [S2].
- In **Li-excess layered oxides** ($\mathrm{Li_2TMO_3}$), Li sits in the TM layer, so some O ions are coordinated by more than three Li. This produces "orphaned O states" \u2014 nonbonding O $2p$ states \u2014 which are the states that can be oxidized [S4].
- Worked example: on de-lithiation, $\mathrm{Mn^{3+}\!\to\!Mn^{4+}}$ occurs first; **additional de-lithiation capacity is accommodated by oxygen oxidation, $\mathrm{O^{2-}\!\to\!O^{1-}}$**; only at low Li content ($x<1$) does $\mathrm{Mn^{4+}\!\to\!Mn^{5+}}$ (with octahedral\u2192tetrahedral change) take over [S5]. So O redox fills the capacity gap the cation alone cannot supply.
- In **ilmenite-type $\mathrm{NaTMO_3}$**, O redox arises from a different trigger: **O under-coordination**, not Na over-coordinating oxygen [S7]. The paper tabulates O $2p$ populations below $E_F$ against computed potentials for 3d and 4d members [S8].
- Computationally, anionic redox is pursued precisely to "maximiz[e] capacities beyond what would be expected solely based on cationic redox pathways" [S6].

**Caveat / cost**
The same anionic activity has a detrimental branch: O$_2$ evolution (and electrolyte-solvent decomposition / oxygen loss) impairs reversibility and limits cycling [S2][S5].

**Not in the provided context:** any quantitative capacity figure (e.g., mAh g\u207b\xb9) attributable to O redox, or the fraction of total capacity it supplies. The excerpts also do not resolve the OCR-corrupted formula near "$\mathrm{Li_xMnO_5}$ (x<1)" [S5].
\end{AnswerText}

\textbf{Retrieved evidence index (all returned units; paths and scores) }
\begin{EvidenceText}
[1] score=None path=part_0293___materials__2408.09942__oxygen-redox-in-battery-cathodes-a-brief-overvie.md_part_1_of_2_.md.txt
[2] score=None path=part_0293___materials__2408.09942__oxygen-redox-in-battery-cathodes-a-brief-overvie.md_part_1_of_2_.md.txt
[3] score=None path=part_0293___materials__2408.09942__oxygen-redox-in-battery-cathodes-a-brief-overvie.md_part_1_of_2_.md.txt
[4] score=None path=part_0293___materials__2408.09942__oxygen-redox-in-battery-cathodes-a-brief-overvie.md_part_1_of_2_.md.txt
[5] score=None path=part_0293___materials__2408.09942__oxygen-redox-in-battery-cathodes-a-brief-overvie.md_part_1_of_2_.md.txt
[6] score=None path=part_0293___materials__2408.09942__oxygen-redox-in-battery-cathodes-a-brief-overvie.md_part_1_of_2_.md.txt
[7] score=None path=part_0294___materials__2408.09942__oxygen-redox-in-battery-cathodes-a-brief-overvie.md_part_2_of_2_.md.txt
[8] score=None path=part_0294___materials__2408.09942__oxygen-redox-in-battery-cathodes-a-brief-overvie.md_part_2_of_2_.md.txt
\end{EvidenceText}
\subsection{BQ09: What problem of visual SLAM in dynamic environments does MLP-SLAM address and how?}
\begin{tabular}{@{}ll@{}}
Arm & GraphRAG local search\\
Generator & omp / deepseek-flash\\
HTTP / success & 200 / True\\
Wall time / cost & 122.4 s / 0.0286302 USD\\
Judge scores & F=5, C=5, G=5, A=5; overall=5.00\\
Verdict & correct\\
\end{tabular}

\textbf{Index/retrieval metadata}
\begin{EvidenceText}
{
  "retrieval": "graphrag local-search style entity-centric assembly (fast NLP entities; shim embeddings; LanceDB absent)",
  "top_entities": [
    "MLP-BASED SLAM",
    "MLP-SLAM",
    "MLP-BASED",
    "ORB-SLAM2",
    "ORB-SLAM3",
    "CFP-SLAM",
    "DYNAMIC SCENES",
    "SLAM",
    "MLP",
    "VDO-SLAM",
    "DYNAMIC OBJECTS",
    "MULTI-HEAD ATTENTION"
  ],
  "n_entities_total": 1816
}
\end{EvidenceText}

\textbf{Dimension rationale}
\begin{EvidenceText}
{
  "factual": "All correct: degradation amid movable objects (gold); hand-crafted geometric tests requiring manual tuning, lacking generality, losing information (gold); indirect evaluation unreliable due to RANSAC randomness; the classifier formulation c_i = argmax(y_i) with y_i = M(v_i); inputs as image intensity, re-projection, and epipolar errors; ORB-SLAM2 base with the discriminator filtering dynamic points before tracking; Deep Sort boxes with in-box background treated as static; static-only reprojection-error minimization; the >50,000-point 70/20/10 split with 1:1 class ratio; Accuracy/F1 direct evaluation with SVM second; ATE 1.48 vs 1.99/2.43/1.77, cross-consistent with other answers.",
  "completeness": "Problem (two gaps) and how (classifier, inputs, pipeline, dataset, evaluation, results) fully answered; the truncation caveat is honest.",
  "groundedness": "Every claim tagged ([S4][S5][S6][S7][S8][S9][S10][S11][S12]) and consistent with gold and the paper.",
  "sources": "Per-claim attributions with an explicit note on where its context cuts off."
}
\end{EvidenceText}

\textbf{Judge summary}
\begin{EvidenceText}
Complete, accurate, and well-attributed answer covering the problem, the learned-classifier mechanism, the dataset, and quantitative results.
\end{EvidenceText}

\textbf{Complete answer}
\begin{AnswerText}
**Problem addressed:** V-SLAM is accurate in scenes with few dynamic objects, but degrades when many movable objects (pedestrians, cars, buses) are present \u2014 typical of outdoor scenes [S4]. Prior fixes relied on hand-crafted geometric tests (epipolar error, re-projection error, or combinations of both), which require manual parameter tuning, lack generality, are laborious to fine-tune, and lose information [S4][S5]. A second gap: since no public dataset targets dynamic/static feature classification, such methods were only evaluated indirectly through SLAM metrics, which is unreliable because of RANSAC randomness in pose estimation [S5].

**How MLP-SLAM addresses it:**
- Replaces hand-tuned geometric heuristics with a learned classifier: an MLP module classifies each feature point as dynamic (1) or static (0), i.e. $c_i = \arg\max(y_i)$ with $y_i = M(v_i)$ [S7].
- Inputs are geometric residuals \u2014 image intensity error $e_I$, re-projection error $e_{Re}$, epipolar error $e_D$ \u2014 generated using ground-truth poses [S9].
- Built on ORB-SLAM2 plus an extra dynamic/static discriminator module that filters dynamic feature points before tracking [S7]; object detection (Deep Sort) supplies bounding boxes, and background features inside them are treated as static, combined with all static points outside the boxes for pose estimation [S8].
- Pose is recovered by minimizing re-projection error over static correspondences only, explicitly avoiding dynamic-object influence [S6].
- Training set: >50,000 feature points, shuffled, 70% train / 20% test / 10% validation, with a 1:1 static-to-dynamic ratio to fix class imbalance [S9].
- Direct evaluation of the classifier (Accuracy, F1-Score) closes the indirect-evaluation gap: MLP beats the compared methods, with SVM second [S10]. On KITTI odometry it reports average ATE 1.48 vs 1.99, 2.43, 1.77 for the baselines [S11][S12].

Caveat: the context is truncated mid-sentence at the proposal statement [S4], so the exact wording of the stated contribution and full baseline names are not available.
\end{AnswerText}

\textbf{Retrieved evidence index (all returned units; paths and scores) }
\begin{EvidenceText}
[1] score=None path=part_0057___artificial-intelligence__1706.03762__attention-is-all-you-need.md_part_1_of_2_.md.txt
[2] score=None path=part_0057___artificial-intelligence__1706.03762__attention-is-all-you-need.md_part_1_of_2_.md.txt
[3] score=None path=part_0057___artificial-intelligence__1706.03762__attention-is-all-you-need.md_part_1_of_2_.md.txt
[4] score=None path=part_0081___robotics__2410.10669__mlp-slam-multilayer-perceptron-based-simultaneou.md_part_1_of_2_.md.txt
[5] score=None path=part_0081___robotics__2410.10669__mlp-slam-multilayer-perceptron-based-simultaneou.md_part_1_of_2_.md.txt
[6] score=None path=part_0081___robotics__2410.10669__mlp-slam-multilayer-perceptron-based-simultaneou.md_part_1_of_2_.md.txt
[7] score=None path=part_0081___robotics__2410.10669__mlp-slam-multilayer-perceptron-based-simultaneou.md_part_1_of_2_.md.txt
[8] score=None path=part_0081___robotics__2410.10669__mlp-slam-multilayer-perceptron-based-simultaneou.md_part_1_of_2_.md.txt
[9] score=None path=part_0081___robotics__2410.10669__mlp-slam-multilayer-perceptron-based-simultaneou.md_part_1_of_2_.md.txt
[10] score=None path=part_0081___robotics__2410.10669__mlp-slam-multilayer-perceptron-based-simultaneou.md_part_1_of_2_.md.txt
[11] score=None path=part_0082___robotics__2410.10669__mlp-slam-multilayer-perceptron-based-simultaneou.md_part_2_of_2_.md.txt
[12] score=None path=part_0082___robotics__2410.10669__mlp-slam-multilayer-perceptron-based-simultaneou.md_part_2_of_2_.md.txt
\end{EvidenceText}
\subsection{BQ10: What strategies for tumor control emerge under immune evasion and chemotherapy resistance?}
\begin{tabular}{@{}ll@{}}
Arm & GraphRAG local search\\
Generator & omp / deepseek-flash\\
HTTP / success & 200 / True\\
Wall time / cost & 144.2 s / 0.032079 USD\\
Judge scores & F=4, C=4, G=4, A=4; overall=4.00\\
Verdict & partial\\
\end{tabular}

\textbf{Index/retrieval metadata}
\begin{EvidenceText}
{
  "retrieval": "graphrag local-search style entity-centric assembly (fast NLP entities; shim embeddings; LanceDB absent)",
  "top_entities": [
    "TUMOR SUPPRESSION",
    "SUPPRESS TUMOR GROWTH",
    "IMMUNE EVASION",
    "IMMUNE-EVASION MECHANISM",
    "TUMOR ELIMINATION",
    "TUMOR CONTROL",
    "REDUCES TUMOR BURDEN",
    "LONG-TERM TUMOR CONTROL",
    "TUMOR-IMMUNE DYNAMICS",
    "IMMUNE CHECKPOINT INHIBITION STRATEGY",
    "TUMOR DYNAMICS",
    "TUMOR PERSISTENCE"
  ],
  "n_entities_total": 1816
}
\end{EvidenceText}

\textbf{Dimension rationale}
\begin{EvidenceText}
{
  "factual": "Mostly correct: the two-lever framing; combination quotes matching peers ('significantly enlarges the parameter region...', 'expand the parameter region exhibiting tumor-free or low population equilibria'); checkpoint-evasion control 'primarily by decreasing immune exhaustion (mu) or increasing immune activation (rho and sigma)' and 'checkpoint inhibitors combined with moderate chemotherapy provide the most effective control strategy'; reduced antigen presentation handled by increasing the effector population via sigma/delta with the sigma requirement 'significantly larger' for T1-only; narrow coexistence with sizable tumors. One error: its summary table attributed to [S7] lists the T1-only therapeutic focus as 'Checkpoint inhibitors; adoptive T-cell transfer', conflating the checkpoint-evasion row with the adoptive-transfer focus that peers consistently report for the T2-only / antigen-presentation rows.",
  "completeness": "Combination and mechanism-matched strategies are well covered, but the sequential composition-guided schedules that the paper proposes (reported by four peer answers from section 4) are declared absent, and the answer covers only a two-row table without distinguishing both evasion mechanisms.",
  "groundedness": "Prose claims are supported by tags and cross-consistent quotes, but the conflated table row is presented as the paper's summary table and misstates it, so one attribution does not hold.",
  "sources": "Claims are tagged ([S1][S2][S6][S7][S8][S9][S11]) but the central table attribution is inaccurate."
}
\end{EvidenceText}

\textbf{Judge summary}
\begin{EvidenceText}
Correct on combination and checkpoint-matched strategies with honest context limits, but its summary table misattributes the T1-only therapeutic focus and it omits the paper's sequential strategies.
\end{EvidenceText}

\textbf{Complete answer}
\begin{AnswerText}
Two levers, phenotype-dependent, and both operate by expanding the region where a tumor-free or low-burden equilibrium is stable.

**1. Combine chemotherapy with immunotherapy.** Chemotherapy alone is insufficient in the model: immune dynamics alone rarely eliminate the tumor [S8], but adding chemo "significantly enlarges the parameter region corresponding to tumor control or tumor-free equilibria," reduces tumor burden, and directly targets resistant subpopulations, while immunotherapy enhances immune-mediated suppression [S8]. Jointly applied, the two "expand the parameter region exhibiting tumor-free or low population equilibria" [S11].

**2. Target the dominant phenotype's escape mechanism** (the paper's summary table, [S7]):

| Tumor phenotype | Best target | Therapeutic focus |
|---|---|---|
| T1-only (immune-evasive dominant) | \u2193 \u03bc (lower immune exhaustion) | Checkpoint inhibitors; adoptive T-cell transfer |
| T2-only (immune-sensitive dominant) | \u2191 \u03c3, \u2193 \u03b4 | Increase effector supply / reduce effector death |

- Under **immune checkpoint evasion**, control is achieved "primarily by decreasing immune exhaustion (\u03bc) or increasing immune activation (\u03c1 and \u03c3)"; "checkpoint inhibitors combined with moderate chemotherapy provide the most effective control strategy" [S9].
- Under **reduced antigen presentation**, the primary strategy "is to increase the overall effector cell population," via higher effector birth rate \u03c3 or lower death rate \u03b4; raising \u03c3 can also yield a small controlled tumor in the T1-only case, but the \u03c3 required is "significantly larger" than for T2-only [S6].

**3. Exploit nonlinearity, don't assume monotone response.** "Modest parameter changes can induce qualitative transitions between immune control and tumor escape," and efficacy depends on both the evasion mechanism and which phenotype dominates [S11]. T1/T2 coexistence occurs only over a narrow range and "tumors can still be sizable" even there [S9].

**Not answered by this context:** the specific optimal sequences, combinations, timing, or pause/switch schedules for chemotherapy and immunotherapy. These are posed as open questions the ODE model is built to investigate [S2]; the excerpts supply no dosing schedule or timing result. Also, the abstract only states that resistant subpopulations arise and drive relapse with more aggressive, treatment-resistant disease [S1]. The specific relapse dynamics are outside these excerpts.
\end{AnswerText}

\textbf{Retrieved evidence index (all returned units; paths and scores) }
\begin{EvidenceText}
[1] score=None path=part_0263___immunology__2604.01385__strategies-for-tumor-elimination-and-control.md_part_1_of_3_.md.txt
[2] score=None path=part_0263___immunology__2604.01385__strategies-for-tumor-elimination-and-control.md_part_1_of_3_.md.txt
[3] score=None path=part_0263___immunology__2604.01385__strategies-for-tumor-elimination-and-control.md_part_1_of_3_.md.txt
[4] score=None path=part_0263___immunology__2604.01385__strategies-for-tumor-elimination-and-control.md_part_1_of_3_.md.txt
[5] score=None path=part_0264___immunology__2604.01385__strategies-for-tumor-elimination-and-control.md_part_2_of_3_.md.txt
[6] score=None path=part_0264___immunology__2604.01385__strategies-for-tumor-elimination-and-control.md_part_2_of_3_.md.txt
[7] score=None path=part_0264___immunology__2604.01385__strategies-for-tumor-elimination-and-control.md_part_2_of_3_.md.txt
[8] score=None path=part_0265___immunology__2604.01385__strategies-for-tumor-elimination-and-control.md_part_3_of_3_.md.txt
[9] score=None path=part_0265___immunology__2604.01385__strategies-for-tumor-elimination-and-control.md_part_3_of_3_.md.txt
[10] score=None path=part_0265___immunology__2604.01385__strategies-for-tumor-elimination-and-control.md_part_3_of_3_.md.txt
[11] score=None path=part_0265___immunology__2604.01385__strategies-for-tumor-elimination-and-control.md_part_3_of_3_.md.txt
[12] score=None path=part_0265___immunology__2604.01385__strategies-for-tumor-elimination-and-control.md_part_3_of_3_.md.txt
\end{EvidenceText}
\end{document}